% Generated by build_chapter.py from the GY1964 files. Do not edit by hand; edit the files and rebuild.
\chapter{The academic year 1964--65}\label{ch:1964}
\nocite{*}
\paragraph{Scope and how to read this chapter.}
This chapter covers all \textbf{49 records} in the supplied 1964--65
catalogue, with \textbf{248 mathematical occurrences}.  All 247 supplied
identifiers are retained; the additional support-of-measure observation
is numbered 1964-6.9.  Dates range from
21 September 1964 to 10 May 1965; the February planning entry is included
without inventing mathematical content.  The sequence follows the supplied
MCCME record order, which is not chronological.  Individual page citations
locate the actual, sometimes interleaved, notebook text.  P000090 starts the
following academic year's Kohn material and is excluded from the present
occurrence count.  Russian transcription is the textual authority; scans
decide consequential doubtful readings.  In each entry, \emph{The notes},
\emph{Statement}, and \emph{Reconstruction} distinguish the record, a precise
modern formulation, and the editor's additions.

\paragraph{The year's mathematical lines.}
The seminar alternates sustained exposition with short announcements and
questions.  Blowups, representation theory, Hilbert kernels, distributions,
Fredholm operators, algebraic jets, and sites supply common language.
Statistical mechanics leads from contour arguments to a walk question;
random matrices lead from circular ensembles to gap probabilities and
spectral traces.  Arnold's geometric optics connects caustics to quantization,
while Novikov's exposition places vector fields on spheres in real
\(K\)-theory.  The final Halmos session collects several decades of
invariant-subspace methods and ends with unusually concrete weighted
operators.  These connections are mathematical and textual; the chapter does
not infer influence merely from two related topics appearing together.

\paragraph{Publication and identification.}
Some announcements have close later matches; others were already published
or circulating.  Dobrushin's phase-separation work requires its correction
as well as the original theorem.  The different-rank Weyl-algebra question
(1964-38.1) has an affirmative distinction through dimension growth; it is
not the fixed-rank endomorphism conjecture.  Sz.-Nagy--Foia\c{s}'s contraction
theorem is available before the Halmos talk, and its linked erratum leaves
that theorem intact.  Domar's later weighted-shift construction resolves the
question of an invariant-subspace chain indexed by all integers, after
adjoining zero and the whole space.  Submission dates are reported as submission dates:
they are not silently identified with first public circulation.  A modern
proof of a recorded idea, such as the Slater-determinant route to Sturm's
theorem, is distinguished from a new theorem discovered in the notebook.

\begin{center}
\begin{tabular}{@{}lrr@{}}\toprule
Class&This chapter&Supplied first run\\\midrule
A&24&22\\
B&22&21\\
C&5&7\\
D&0&0\\
E&160&163\\
F&37&34\\
\midrule Total&248&247\\\bottomrule
\end{tabular}
\end{center}
Here A denotes a close subsequent publication by the speaker or Gelfand;
B a later development by others; C a partly settled precise line; D a
searched but undeveloped line; E prior mathematics; and F an identification
the record cannot sustain.  The separate correctness and match fields are
essential: a familiar theorem may be false in the notebook's literal
wording.  E is an expository classification, not a separately audited
priority claim for every assertion: some record-6 entries lack a verified
pre-talk theorem-level publication date, as their fates expressly state.
No absence of a search hit is used as evidence of an orphan idea.

\paragraph{Questions and bounded development.}
The 17 problem entries include both settled and unresolved questions,
with explicit hypotheses.  The major residual operator questions concern
the Hilbert invariant-subspace problem and Bishop operators.  Weight
classification is reported only to the scope of the verified criteria.
The walk question needs its sampling law specified before a
modern asymptotic question is assigned to it.  Sinai's observation on
spectral multiplicity is not converted into a conjecture absent from the
source.  Recent 2026 claims are reported with their actual verification
level, not treated as accepted resolutions.

One development is provided, for Bishop cyclicity.  It proves rational
cyclicity and a finite-orbit continuity bound, supplies exact approximation
certificates checked by two independent implementations, and derives a
terminating sufficient continued-fraction criterion.  These are bounded
certification and exposition results, not a solution for every irrational
rotation and not a claim of historical novelty; compare
\citet{chamizo2020bishop} and \citet{behani2023bishop}.

\paragraph{Continuity with the supplied index.}
Despite its filename, \texttt{previous\_index.md} contains this same year's
247 first-run items, no earlier-year occurrences, and no existing problem
identifiers.  It is therefore a comparison baseline, not evidence for
cross-year links.  All retained identifiers remain stable; 14 class
assignments change after renewed source checking.  Within-year links join
the two invariant-subspace announcements, repeated gap formulas, and the
successive algebraic-jet discussions.  A genuine earlier-year index can be
joined later without fabricating a history that was not supplied.

\GStalk{2}{1964-09-21}{I. M. Gelfand}{P000001-P000004}{Resolution, blowups, and projective spectra}
After a group-theoretic heading, the notes discuss Hironaka's resolution theorem through the universal property of blowing up, graded algebras and two explicit centers. The algebraic examples occupy P000001--P000003; Dobrushin begins on P000003. The distinction between a symmetric algebra and a Rees algebra matters, and the crossing-center example is not itself a resolution.

\begin{GSitem}{1964-2.1}{F}{Cartesian and Boolean products of groups}
\GSfield{Kind; source}{remark; \GSpage{P000001}}
\GSfield{The notes}{Only the heading \ru{Декартово и булево произведение групп} survives before the change to resolution.}
\GSfield{Statement}{No groups, category, construction of a Boolean product, or assertion are specified. The phrase cannot determine one mathematical item more precisely.}
\GSfield{Fate}{Title searches distinguish several uses of Boolean products; none can identify the intended construction from this heading. No publication match is asserted.}
\GSfield{Present status}{As of November 2026: Unidentified heading, rather than evidence of an undeveloped problem.}
\GSfield{Assessment}{literal statement: not determinable; publication match: none found; verification: unresolved lead; development: none; confidence: low}
\end{GSitem}
\begin{GSitem}{1964-2.2}{E}{Hironaka resolution in characteristic zero}
\GSfield{Kind; source}{theorem; \GSpage{P000001}}
\GSfield{The notes}{A singular divisor \(D\) is replaced by a smooth hypersurface over it, under the heading Hironaka. The base field and the distinction between total and strict transforms are not stated.}
\GSfield{Statement}{For a reduced variety over a characteristic-zero field, resolution gives a proper birational morphism from a smooth variety, an isomorphism over a prescribed smooth locus. In embedded resolution of a divisor in a smooth ambient variety, the strict transform is smooth and the total transform has simple normal crossings after suitable blowups; the total inverse image need not be a smooth hypersurface.}
\GSfield{Reconstruction}{The field hypothesis and strict/total-transform distinction are editorial precisions.}
\GSfield{Fate}{The named result had appeared in \citet{hironaka1964resolution}, Part I, introduction and main theorem, in 1964, before this September discussion; the Russian translation followed in 1965. The opened archive record verifies the publication, not every proof in the 95-page paper.}
\GSfield{Present status}{As of November 2026: The characteristic-zero theorem is established. Removing the characteristic-zero hypothesis is a different question and is not posed here.}
\GSfield{Assessment}{literal statement: correct with missing hypotheses; publication match: exact; verification: publisher metadata checked; development: annotated; confidence: high}
\end{GSitem}
\begin{GSitem}{1964-2.3}{E}{Universal blowup and invertibility of the center}
\GSfield{Kind; source}{construction; \GSpage{P000001-P000002}}
\GSfield{The notes}{The notebook seeks \(V'\to V\) making an ideal \(J\) locally principal, with a universal mapping property, and then writes its symmetric tensor algebra.}
\GSfield{Statement}{For a scheme \(X\) and a finite-type quasi-coherent ideal \(I\), the blowup is \(\operatorname{Proj}_X\mathcal R(I)\), where \(\mathcal R(I)=\bigoplus_{n\ge0}I^n\). It is terminal among \(X\)-schemes on which the inverse-image ideal defines an effective Cartier divisor. Locally principal alone is insufficient when a generator is a zero divisor. In general \(\operatorname{Sym}(I)\twoheadrightarrow\mathcal R(I)\) is not an isomorphism.}
\GSfield{Reconstruction}{Replace the symmetric algebra by the Rees algebra in the general assertion; retain the original notation in the notes field.}
\GSfield{Fate}{This is the standard blowup universal property, recorded precisely in \citet[tags 01OF and 0806]{stacks2026blowups}. The modern reference is used to correct the algebra, not to date its invention. The pre-seminar source is \citet[Chapter II, Section 8.1]{grothendieck1961ega2} (1961); its publisher metadata were opened, but the 27 MB scan exceeded the retrieval limit. The exact modern construction was checked in Stacks.}
\GSfield{Present status}{As of November 2026: Established construction. The symmetric algebra works for ideals of linear type, including the two regular-sequence examples below.}
\GSfield{Assessment}{literal statement: false as written; publication match: exact; verification: read the result; development: annotated; confidence: high}
\end{GSitem}
\begin{GSitem}{1964-2.4}{E}{Proj, homogeneous primes, and the local sheaf}
\GSfield{Kind; source}{construction; \GSpage{P000002}}
\GSfield{The notes}{Homogeneous prime ideals and homogeneous fractions of equal degree are discussed; the abbreviation \ru{Выбр.} precedes primes containing all positive-degree elements.}
\GSfield{Statement}{For a positively graded ring \(S\), \(\operatorname{Proj}S\) consists of homogeneous primes not containing \(S_+=\bigoplus_{d>0}S_d\). On \(D_+(f)\), \(\mathcal O\) is the degree-zero part of \(S_f\); the degree-one component gives \(\mathcal O(1)\), with the usual degree conventions.}
\GSfield{Reconstruction}{The intended exclusion and the distinction between degree zero and degree one are explicit modern reconstructions.}
\GSfield{Fate}{The projective-spectrum construction underlying \citet[Section 31.33]{stacks2026blowups} predates the lecture. The scan supports the abbreviation but does not justify expanding it as ``choose'' rather than ``discard. The pre-seminar source is \citet[Chapter II, Section 8.1]{grothendieck1961ega2} (1961); its publisher metadata were opened, but the 27 MB scan exceeded the retrieval limit. The exact modern construction was checked in Stacks.}
\GSfield{Present status}{As of November 2026: Standard scheme construction; reversing the irrelevant-ideal exclusion would destroy the intended Proj.}
\GSfield{Assessment}{literal statement: correct with missing hypotheses; publication match: partial; verification: read the result; development: annotated; confidence: high}
\end{GSitem}
\begin{GSitem}{1964-2.5}{E}{Blowing up the origin of the affine plane}
\GSfield{Kind; source}{construction; \GSpage{P000002-P000003}}
\GSfield{The notes}{For the ideal of the origin in the plane the exceptional fiber is a projective line and local coordinate ratios are introduced.}
\GSfield{Statement}{For \(I=(x,y)\subset k[x,y]\), \(\operatorname{Bl}_0\mathbb A^2\subset\mathbb A^2\times\mathbb P^1\) has equation \(xV=yU\). Its two charts are \(y=xv\) and \(x=yu\), each an affine plane, and the fiber over the origin is \(\mathbb P^1_k\).}
\GSfield{Reconstruction}{The homogeneous-coordinate equation and chart equations spell out the drawing.}
\GSfield{Fate}{This is the affine-chart construction of \citet[Lemma 31.33.2]{stacks2026blowups}; the notebook example is an exposition of it. The pre-seminar source is \citet[Chapter II, Section 8.1]{grothendieck1961ega2} (1961); its publisher metadata were opened, but the 27 MB scan exceeded the retrieval limit. The exact modern construction was checked in Stacks.}
\GSfield{Present status}{As of November 2026: The calculation is valid over any field and gives a smooth blowup.}
\GSfield{Assessment}{literal statement: correct; publication match: exact; verification: read the result; development: annotated; confidence: high}
\end{GSitem}
\begin{GSitem}{1964-2.6}{E}{Blowing up the crossing ideal (z,xy)}
\GSfield{Kind; source}{construction; \GSpage{P000003}}
\GSfield{The notes}{The next center is \(z=0,xy=0\), the union of two lines, with the two generators \(z,xy\).}
\GSfield{Statement}{For \(I=(z,xy)\subset k[x,y,z]\), the Rees presentation is \(k[x,y,z,U,V]/(xyU-zV)\). On \(U\ne0\) the chart is \(xy=zv\); on \(V\ne0\) it is \(z=xyu\). The former has a singular point at the origin; therefore this single blowup is not a general resolution procedure.}
\GSfield{Reconstruction}{The Rees relation, two charts, and singular-point diagnosis are editorial algebraic checks; the notebook does not claim a completed resolution.}
\GSfield{Fate}{The example is an application of the blowup charts in \citet[Lemma 31.33.2]{stacks2026blowups}. No later publication is needed to account for it. The pre-seminar source is \citet[Chapter II, Section 8.1]{grothendieck1961ega2} (1961); its publisher metadata were opened, but the 27 MB scan exceeded the retrieval limit. The exact modern construction was checked in Stacks.}
\GSfield{Present status}{As of November 2026: The calculation remains a useful warning that an arbitrary center can produce a singular blowup.}
\GSfield{Assessment}{literal statement: correct with missing hypotheses; publication match: partial; verification: read the result; development: annotated; confidence: high}
\end{GSitem}

\GStalk{3}{1964-09-21}{R. L. Dobrushin}{P000002-P000012}{Canonical Ising coexistence, contours, and a no-return walk}
Dobrushin develops the canonical lattice gas, a criterion for phase coexistence through macroscopic density fluctuations, and the contour estimates underlying his 1965 phase-transition paper. A short walk question follows. Signs and the density/specific-volume notation vary between pages, so the modern statements fix one convention. The walk question lacks enough information to identify its probability law uniquely.

\begin{GSitem}{1964-3.1}{E}{Canonical lattice gas and Ising energy}
\GSfield{Kind; source}{definition; \GSpage{P000003-P000004}}
\GSfield{The notes}{There are \(N\) particles at distinct sites of a lattice box of volume \(V\), with energy minus the number of occupied nearest-neighbor pairs and weights \(e^{-U/T}\).}
\GSfield{Statement}{Let \(\Lambda\subset\mathbb Z^\nu\) be finite and choose \(N\in\{0,\ldots,|\Lambda|\}\). For \(\eta_x\in\{0,1\}\) with \(\sum_{x\in\Lambda}\eta_x=N\), set \[ H_\Lambda(\eta)=-\sum_{\substack{\{x,y\}\subset\Lambda\\|x-y|_1=1}} \eta_x\eta_y,\qquad Z_{\Lambda,N}=\sum_{\eta:\,\sum_x\eta_x=N}e^{-H_\Lambda(\eta)/T}. \] Here \(T>0\), each unordered nearest-neighbor pair is counted once, and the canonical law is the normalized weight. With \(\eta_x=(1+\sigma_x)/2\) this is an Ising model with the corresponding fixed magnetization and boundary terms.}
\GSfield{Fate}{The model is the starting point of \citet[Section 1]{dobrushin1965phase}; its definition predates the lecture. That paper was received on 30 September 1964 and published in 1965.}
\GSfield{Present status}{As of November 2026: Standard finite-volume probability law; the sum counts configurations once, not ordered particle labels.}
\GSfield{Assessment}{literal statement: correct; publication match: exact; verification: read the result; development: annotated; confidence: high}
\end{GSitem}
\begin{GSitem}{1964-3.2}{E}{Thermodynamic limit and independence of regular shape}
\GSfield{Kind; source}{theorem; \GSpage{P000004-P000006}}
\GSfield{The notes}{The notes introduce \(f(v,T)=\lim N^{-1}\log Z\), \(v=V/N\), assert independence of the expanding region, and call the resulting function convex.}
\GSfield{Statement}{For this finite-range lattice gas along regular expanding boxes, with \(N/|\Lambda|\to\rho\in(0,1)\), the canonical log-partition density \(s_T(\rho)=\lim|\Lambda|^{-1}\log Z_{\Lambda,N}\) exists and is concave. Its perspective \(f(v,T)=v s_T(1/v)\), \(v>1\), is concave. Boundary changes of order \(|\partial\Lambda|\) do not affect the limit for van Hove sequences.}
\GSfield{Reconstruction}{A single positive-log convention is imposed; regularity of the exhausting sequence is supplied.}
\GSfield{Fate}{\citet[Section 1]{dobrushin1965phase} explicitly uses the thermodynamic limit as prior input. The positive-log convention in the notebook forces concavity; its later minus sign cannot be retained simultaneously.}
\GSfield{Present status}{As of November 2026: Established for the stated finite-range interaction and regular limits. Arbitrary irregular regions need not have negligible boundary effects.}
\GSfield{Assessment}{literal statement: false as written; publication match: exact; verification: read the result; development: annotated; confidence: high}
\end{GSitem}
\begin{GSitem}{1964-3.3}{E}{Coexistence plateau and the meaning of one transition}
\GSfield{Kind; source}{claim; \GSpage{P000004-P000008}}
\GSfield{The notes}{A pressure plateau is drawn and the van der Waals picture is compared with the lattice model; one line later says \(\partial f/\partial v=0\).}
\GSfield{Statement}{For \(F=-NT f(v,T)\), pressure is \(p=T\partial_v f\) wherever the derivative exists. A coexistence interval is an affine segment of \(f\), hence a constant-pressure interval; the pressure need not be zero. In density variables, the grand-canonical transition corresponds to distinct equilibrium densities at one chemical potential.}
\GSfield{Reconstruction}{The relationship between canonical and grand-canonical variables is made explicit; ``zero derivative is replaced by ``constant derivative.}
\GSfield{Fate}{The thermodynamic interpretation is the framework of \citet[Section 1 and Theorem 1]{dobrushin1965phase}. The general claim of exactly one transition is not justified by the sketch alone.}
\GSfield{Present status}{As of November 2026: Phase coexistence is rigorous at sufficiently low temperature in the models discussed; the picture is not a proof of a complete phase diagram.}
\GSfield{Assessment}{literal statement: false as written; publication match: partial; verification: read the result; development: annotated; confidence: high}
\end{GSitem}
\begin{GSitem}{1964-3.4}{E}{No positive-temperature transition in one dimension}
\GSfield{Kind; source}{theorem; \GSpage{P000004}}
\GSfield{The notes}{The one-dimensional model is said to have no phase transition at positive temperature.}
\GSfield{Statement}{The nearest-neighbor one-dimensional Ising model has analytic infinite-volume pressure for every finite real field and \(T>0\), because its positive transfer matrix has a simple Perron eigenvalue. Its canonical free energy has no low-temperature coexistence segment of the type sought in dimensions two and three.}
\GSfield{Fate}{This classical contrast is part of the introduction in \citet{dobrushin1965phase}; the transfer-matrix mechanism predates 1964.}
\GSfield{Present status}{As of November 2026: Established for finite-range one-dimensional Ising interactions; long-range interactions are outside the statement.}
\GSfield{Assessment}{literal statement: correct with missing hypotheses; publication match: partial; verification: read the result; development: annotated; confidence: high}
\end{GSitem}
\begin{GSitem}{1964-3.5}{E}{Onsager and Kac--Helfand antecedents}
\GSfield{Kind; source}{remark; \GSpage{P000004}}
\GSfield{The notes}{Onsager's two-dimensional result and work attributed to Kac and Helfand are named as earlier achievements, without a complete bibliographic or parameter specification.}
\GSfield{Statement}{Onsager solved the zero-field nearest-neighbor square-lattice Ising free energy. The Kac--Helfand reference is an antecedent in the discussion of phase transition, but the notebook alone does not identify a particular long-range limiting regime.}
\GSfield{Fate}{The prior literature is discussed in \citet[Section 1 and bibliography]{dobrushin1965phase}. The Onsager identification is secure; a unique Kac--Helfand theorem is not inferred merely from the names.}
\GSfield{Present status}{As of November 2026: The historical comparison is expository. It must not be read as an exact solution of the three-dimensional model.}
\GSfield{Assessment}{literal statement: not determinable; publication match: partial; verification: read the result; development: annotated; confidence: medium}
\end{GSitem}
\begin{GSitem}{1964-3.6}{A}{Macroscopic density deviations and coexistence}
\GSfield{Kind; source}{theorem; \GSpage{P000005-P000006}}
\GSfield{The notes}{The number of particles in a macroscopic subcube is compared with its expected density. Failure of a volume-scale exponential deviation bound is used to detect a linear piece of the limiting free energy.}
\GSfield{Statement}{For a fixed volume fraction \(a\in(0,1)\), compare a density split \(\rho=a\rho_1+(1-a)\rho_2\). Up to subvolume boundary errors the log probability has volume rate \(a s_T(\rho_1)+(1-a)s_T(\rho_2)-s_T(\rho)\). Strict inequality gives exponential suppression; equality along an interval is the canonical signature of coexistence. The criterion requires the allowed density splits and thermodynamic sequences to be specified.}
\GSfield{Reconstruction}{The rate function and the precise role of concavity replace compressed probability notation.}
\GSfield{Fate}{The matching criterion appears in \citet[Theorem 2 and the reduction following it]{dobrushin1965phase}, received nine days after the seminar and published in 1965. The shorter \citet{dobrushin1965announcement} was received and presented by Kolmogorov on 15 September 1964, then published in 1965. Receipt and presentation are recorded separately from publication; no earlier public distribution date is inferred.}
\GSfield{Present status}{As of November 2026: The criterion is established in this lattice-gas setting. A failure of an unspecified law of large numbers in an arbitrary model is not an equivalent theorem.}
\GSfield{Assessment}{literal statement: correct with missing hypotheses; publication match: exact; verification: read the result; development: annotated; confidence: high}
\end{GSitem}
\begin{GSitem}{1964-3.7}{A}{Boundary-energy factorization of partition functions}
\GSfield{Kind; source}{method; \GSpage{P000005-P000006}}
\GSfield{The notes}{The partition sum is split into a subbox and its complement, with a factor \(\exp\{\theta N^{2/3}\}\) accounting for crossing bonds.}
\GSfield{Statement}{For fixed particle allocations \(k,N-k\), removing all bonds crossing an interface changes each Gibbs weight by factors between \(e^{-C|\partial H|/T}\) and \(e^{C|\partial H|/T}\). The restricted partition sum is therefore comparable to \(Z_{H,k}Z_{\Lambda\setminus H,N-k}\). For geometrically similar boxes in \(\nu\) dimensions the exponent is \(O(N^{(\nu-1)/\nu})\), and \(2/3\) is the three-dimensional case.}
\GSfield{Reconstruction}{The dependence on dimension and temperature is restored.}
\GSfield{Fate}{This is the finite-range comparison used in \citet[proof of Theorem 2]{dobrushin1965phase}; the 1965 publication supplies the boundary bookkeeping. The shorter \citet{dobrushin1965announcement} was received and presented by Kolmogorov on 15 September 1964, then published in 1965. Receipt and presentation are recorded separately from publication; no earlier public distribution date is inferred.}
\GSfield{Present status}{As of November 2026: A standard and valid surface-error estimate. It is negligible only at volume scale, not at the surface scale of interface fluctuations.}
\GSfield{Assessment}{literal statement: correct with missing hypotheses; publication match: exact; verification: read the result; development: annotated; confidence: high}
\end{GSitem}
\begin{GSitem}{1964-3.8}{A}{Large and small contours force spatial separation}
\GSfield{Kind; source}{method; \GSpage{P000006-P000010}}
\GSfield{The notes}{Contours are divided at a cutoff \(M\); a large contour has surface of order \(l^{\nu-1}\), while the total enclosed volume of small contours is controlled by an isoperimetric sum. Earlier exponents differ from the later restatement.}
\GSfield{Statement}{For closed interfaces of surface sizes \(|\Gamma_i|\), isoperimetry bounds enclosed volumes by \(S_\nu|\Gamma_i|^{\nu/(\nu-1)}\), with geometry-dependent lattice constants. If the small-contour contribution cannot accommodate the minority phase, a macroscopic interface or separated dense and dilute regions must occur. A surface in a box of side \(l\) has scale \(l^{\nu-1}\), not \(l^{(\nu-1)/\nu}\).}
\GSfield{Reconstruction}{Only the geometric implication common to the displayed estimates is asserted; overwritten summation limits are not silently completed.}
\GSfield{Fate}{The contour decomposition and geometric estimates belong to \citet[Sections 2--4]{dobrushin1965phase}. The later notebook restatement is used to diagnose the earlier exponent, not to conceal it. The shorter \citet{dobrushin1965announcement} was received and presented by Kolmogorov on 15 September 1964, then published in 1965. Receipt and presentation are recorded separately from publication; no earlier public distribution date is inferred. \citet[paragraphs 1--2]{dobrushin1966addendum} repairs an error in Lemma 4 by grouping adjacent contours into generalized cycles and also corrects inequality (8.8); the original coexistence conclusion and its \(Q(T)\) are preserved.}
\GSfield{Present status}{As of November 2026: Established proof strategy for low-temperature coexistence. Precise constants depend on the contour convention and cannot be transferred unchanged between lattice surfaces and Euclidean spheres.}
\GSfield{Assessment}{literal statement: false as written; publication match: exact; verification: read the result; development: annotated; confidence: high}
\end{GSitem}
\begin{GSitem}{1964-3.9}{A}{Contour entropy against Boltzmann cost}
\GSfield{Kind; source}{method; \GSpage{P000009-P000011}}
\GSfield{The notes}{The proof counts contours by exponential bounds, including \(3^k\), and compares their multiplicity with the energetic penalty.}
\GSfield{Statement}{A rooted planar nearest-neighbor contour with a prescribed first edge has at most \(3^{k-1}\) nonbacktracking continuations of length \(k\); self-avoidance only decreases that number. In higher dimensions a separate connected-surface enumeration supplies \(C^k\). If the energy cost per unit interface exceeds \(T\log C\), the contour sum converges.}
\GSfield{Fate}{This is the entropy-energy estimate in \citet[Sections 3--4]{dobrushin1965phase}. The seminar's planar \(3^k\) cannot by itself enumerate arbitrary three-dimensional surfaces. The shorter \citet{dobrushin1965announcement} was received and presented by Kolmogorov on 15 September 1964, then published in 1965. Receipt and presentation are recorded separately from publication; no earlier public distribution date is inferred. \citet[paragraphs 1--2]{dobrushin1966addendum} repairs an error in Lemma 4 by grouping adjacent contours into generalized cycles and also corrects inequality (8.8); the original coexistence conclusion and its \(Q(T)\) are preserved.}
\GSfield{Present status}{As of November 2026: Rigorous with a specified root and contour class; it is the low-temperature Peierls mechanism.}
\GSfield{Assessment}{literal statement: correct with missing hypotheses; publication match: exact; verification: read the result; development: annotated; confidence: high}
\end{GSitem}
\begin{GSitem}{1964-3.10}{A}{An explicit low-temperature coexistence window}
\GSfield{Kind; source}{theorem; \GSpage{P000008-P000010}}
\GSfield{The notes}{P000009 gives \(Q(T)<1/v<1-Q(T)\), with a displayed series transcribed as \[Q(T)=\frac{S_\nu e(1-e^{-1/(2T)})}{3(1+e^{-\nu/T})}\sum_{k=2\nu}^{\infty}3\exp\{-1/(2T)\}^{k}k^{\nu/(\nu-1)}.\] The factor resembling \(e\), the upper limit and especially whether the exponent \(k\) includes the factor 3 remain ambiguous. This display is a source reading, not an endorsed numerical bound.}
\GSfield{Statement}{For the nearest-neighbor attractive lattice gas in dimension \(\nu=2\) or \(3\), and each density \(\rho\in(0,1)\), there is \(T_0(\rho)>0\) such that for \(0<T<T_0(\rho)\) the canonical limiting free energy has an affine segment containing the specific volume \(v=1/\rho\). The original publication gives an explicit sufficient density window by its equations (1.6)--(1.7). Its formula could not be visually recovered at sufficient confidence in this run; the existence statement here is deliberately qualitative but has complete quantifiers.}
\GSfield{Reconstruction}{The qualitative theorem is separated from the partly ambiguous quantitative series. No unverified numerical threshold or repaired placement of the power in \(Q\) is asserted.}
\GSfield{Fate}{The direct match is \citet[Theorem 1, equations (1.6)--(1.7)]{dobrushin1965phase}. Its receipt date is 30 September 1964 and publication year 1965. The shorter \citet{dobrushin1965announcement} was received and presented by Kolmogorov on 15 September 1964, then published in 1965. Receipt and presentation are recorded separately from publication; no earlier public distribution date is inferred. \citet[paragraphs 1--2]{dobrushin1966addendum} repairs an error in Lemma 4 by grouping adjacent contours into generalized cycles and also corrects inequality (8.8); the original coexistence conclusion and its \(Q(T)\) are preserved.}
\GSfield{Present status}{As of November 2026: low-temperature coexistence and the repaired published argument are established. The source display should not be used for a numerical estimate until the original formula is visually verified.}
\GSfield{Assessment}{literal statement: correct with missing hypotheses; publication match: partial; verification: read the result; development: annotated; confidence: high}
\end{GSitem}
\begin{GSitem}{1964-3.11}{C}{Expected displacement of a walk with returns forbidden}
\GSfield{Kind; source}{problem; \GSpage{P000011}}
\GSfield{The notes}{The notebook labels an unsolved question: turn with probability \(1/3\), returning forbidden, and estimate expected distance after \(N\) steps. It compares the unrestricted walk with \(\sqrt N\).}
\GSfield{Statement}{The written rule does not determine a probability measure: forbidding only immediate reversals, forbidding every revisit during growth, and choosing uniformly among all length-\(N\) self-avoiding walks give different laws. A definite modern problem associated with the passage is to determine the displacement exponent of uniform nearest-neighbor self-avoiding walks on \(\mathbb Z^2\). That association does not identify the original law uniquely.}
\GSfield{Reconstruction}{Uniform SAW and the exponent conjecture are modern comparisons, explicitly not words or a theorem in the notebook.}
\GSfield{Fate}{For the definite uniform-SAW comparison, \citet[Theorem 1.1 and Corollary 1.2]{duminil2013subballistic}, preprint 2 May 2012 and publication 13 October 2013, proves exponential suppression of linear displacement: for every \(v>0\), some \(c(v)>0\) gives \(\mathbb P(\max_{k\le n}\|\gamma_k\|\ge vn)\le e^{-c(v)n}\) for all sufficiently large \(n\). In particular \(\mathbb E\|\gamma_n\|^2=o(n^2)\). This is a partial development of a natural interpretation of the question, not an identification of the notebook law.}
\GSfield{Present status}{As of November 2026: for the square-lattice uniform fixed-length SAW, the prediction \(\mathbb E\|\gamma_n\|^2=n^{3/2+o(1)}\) and the expected endpoint-distance exponent \(3/4\) are not established by the sources reviewed. The mean square and mean distance are distinct observables, so even a theorem for one requires a moment argument before being transferred to the other. The notebook states no \(3/4\) exponent.}
\GSfield{Assessment}{literal statement: not determinable; publication match: partial; verification: read the result; development: annotated; confidence: medium; problem \hyperref[prob:GP-1964-01]{GP-1964-01}}
\GSfield{Worth developing}{The uniform-SAW problem remains significant, but a bounded companion contribution should clarify the law and known estimates rather than promise the planar exponent.}
\end{GSitem}
\begin{GSitem}{1964-3.12}{B}{The conjectured exponential count e\textasciicircum{}n}
\GSfield{Kind; source}{conjecture; \GSpage{P000011}}
\GSfield{The notes}{The following line gives rough bases \(2.6\) and \(2.8\), followed by a conjecture \(e^n\); the phrase about allowing returns conflicts with this count.}
\GSfield{Statement}{If \(c_n\) counts nearest-neighbor self-avoiding walks of length \(n\) from the origin in \(\mathbb Z^2\), submultiplicativity defines \(\mu=\lim c_n^{1/n}\). The conjecture \(\mu=e\) is false by rigorous upper bounds below \(e\). If immediate reversals alone are prohibited, the count is exactly \(4\cdot3^{n-1}\), also incompatible with \(e^n\).}
\GSfield{Reconstruction}{The connective-constant formulation is conditional on the self-avoiding-walk interpretation.}
\GSfield{Fate}{Under the square-lattice SAW interpretation, \citet[abstract and upper-bound result]{ponitz2000bounds}, published 13 April 2000, gives \(\mu\le2.679192495<e\), decisively refuting \(\mu=e\). The numerical bound and publication metadata were read on the journal page; its computational certificate was not rerun. The exact square-lattice connective constant is a separate question.}
\GSfield{Present status}{As of November 2026: The base \(e\) is excluded for uniform square-lattice SAW counting, while the exact constant remains unknown in the reviewed literature. The later published bound is \(\mu\le2.662343\), from \citet{couronne2025bound} (preprint first posted 29 November 2022; journal publication 31 October 2025). An author-posted preprint of 21 October 2026, \citet{tomar2026connective}, claims \(\mu\le2.6586513\); its certificate was not independently rerun here. Neither value is an exact determination of the square-lattice constant.}
\GSfield{Assessment}{literal statement: not determinable; publication match: partial; verification: read the result; development: annotated; confidence: medium; problem \hyperref[prob:GP-1964-01]{GP-1964-01}}
\end{GSitem}

\GStalk{4}{1964-09-28}{Irving E. Segal}{P000010-P000014}{Nonlinear wave evolution and geometry of solutions}
A preliminary note explicitly dated Thursday 24 September appears on P000007--P000008; the main 28 September exposition starts on P000011. Both develop the nonlinear Klein--Gordon equation as an evolution in a Hilbert space, then its manifold of solutions, tangent equation, symplectic form and canonical quantization. These are linked to Segal's 1964--65 Leray seminar; they are not an unrestricted existence theorem for interacting quantum fields.

\begin{GSitem}{1964-4.1}{F}{A quotient algebra for nonlinear equations}
\GSfield{Kind; source}{construction; \GSpage{P000007-P000008}}
\GSfield{The notes}{The preliminary note encodes the field equation by functionals \(F_f\), forms a quotient \(\Phi/\Omega\), and interprets maximal ideals as generalized solutions.}
\GSfield{Statement}{For a specified commutative algebra \(\Phi\) of functionals on a classical configuration space, the ideal generated by weak equation functionals vanishes at every classical solution; evaluation gives a scalar-valued character of \(\Phi/\Omega\). The converse requires a specified character space and continuity or algebraic hypotheses. Infinite-dimensional maximal ideals need not be evaluations at ordinary fields.}
\GSfield{Reconstruction}{Only the implication from classical solutions to characters is retained unconditionally.}
\GSfield{Fate}{The broad algebraic quantization program matches \citet[Parts I--III]{segal1965variete}, but the notebook does not define \(\Phi\), its topology or the allowed characters well enough to identify this maximal-ideal assertion as a theorem.}
\GSfield{Present status}{As of November 2026: The quotient construction is meaningful once its algebra is fixed. The asserted exhaustive solution interpretation is not established by the surviving text.}
\GSfield{Assessment}{literal statement: not determinable; publication match: thematic only; verification: read the result; development: annotated; confidence: medium}
\end{GSitem}
\begin{GSitem}{1964-4.2}{E}{Semilinear evolution and Duhamel formula}
\GSfield{Kind; source}{method; \GSpage{P000008; P000011}}
\GSfield{The notes}{The field equation becomes \(u'=Au+K(u)\), with \(A=-A^*\), and an integral equation with \(e^{tA}\).}
\GSfield{Statement}{Let \(A\) generate a strongly continuous group on a Banach space and let \(K\) be locally Lipschitz. A mild solution satisfies \(u(t)=e^{(t-t_0)A}u_0+\int_{t_0}^t e^{(t-s)A}K(u(s))\,ds\), with local uniqueness. A strong solution needs the appropriate domain regularity; global existence needs an a priori bound preventing finite-time escape.}
\GSfield{Reconstruction}{Mild/strong solutions and local/global existence are distinguished.}
\GSfield{Fate}{\citet[pp. 130--131]{segal1963global} explicitly formulates this abstract equation with skew-adjoint \(A\) and locally Lipschitz \(K\), citing his already circulating nonlinear-semigroup work.}
\GSfield{Present status}{As of November 2026: Established local theory. The nonlinear field map is not locally Lipschitz on every energy space in every dimension or for every power.}
\GSfield{Assessment}{literal statement: correct with missing hypotheses; publication match: exact; verification: read the result; development: annotated; confidence: high}
\end{GSitem}
\begin{GSitem}{1964-4.3}{E}{The Klein--Gordon energy space and its generator}
\GSfield{Kind; source}{construction; \GSpage{P000011-P000012}}
\GSfield{The notes}{With \(B=(m^2-\Delta)^{1/2}\), the notes write \(E=D(B)\oplus L^2\), an energy expression, and call the matrix generator self-adjoint.}
\GSfield{Statement}{For \(m>0\), put \(\|(\phi,\pi)\|_E^2=\|B\phi\|_2^2+\|\pi\|_2^2\). On \(D(B^2)\oplus D(B)\), \(A(\phi,\pi)=(\pi,-B^2\phi)\) is skew-adjoint in this energy inner product. Thus \(e^{tA}\) is unitary; \(iA\), after complexification, is self-adjoint.}
\GSfield{Reconstruction}{The missing norm square, operator domain and skew-adjointness are supplied.}
\GSfield{Fate}{The Hilbert-space formulation appears in \citet[pp. 130--132]{segal1963global}. Its sign agrees with P000011 and contradicts the later word ``self-adjoint.}
\GSfield{Present status}{As of November 2026: Standard free Klein--Gordon evolution. For \(m=0\), the homogeneous energy completion requires a separate definition.}
\GSfield{Assessment}{literal statement: false as written; publication match: exact; verification: read the result; development: annotated; confidence: high}
\end{GSitem}
\begin{GSitem}{1964-4.4}{A}{A smooth manifold of global solutions}
\GSfield{Kind; source}{theorem; \GSpage{P000012}}
\GSfield{The notes}{The totality \(M\) of solutions is a manifold, with evaluation at any time giving coordinates in \(E\).}
\GSfield{Statement}{If a semilinear evolution has a globally defined \(C^r\) flow \(\Phi_t:E\to E\), with \(\Phi_t\) a \(C^r\) diffeomorphism for every real \(t\), then the set of entire solutions inherits a \(C^r\) Banach-manifold structure via \(u\mapsto u(t_0)\). Local flows give only the corresponding local solution manifold.}
\GSfield{Reconstruction}{The required flow hypotheses are explicit rather than attributed to the abbreviated notes.}
\GSfield{Fate}{The solution-manifold program is developed in \citet[Part II]{segal1965variete}, a 1964--65 lecture text, and in the related 1965 journal paper cited there. The formal assertion is conditional on the analytic hypotheses.}
\GSfield{Present status}{As of November 2026: Established under smooth global well-posedness; neither global existence nor a smooth flow follows merely from writing the differential equation.}
\GSfield{Assessment}{literal statement: correct with missing hypotheses; publication match: exact; verification: read the result; development: annotated; confidence: high}
\end{GSitem}
\begin{GSitem}{1964-4.5}{A}{The variational equation identifies tangent vectors}
\GSfield{Kind; source}{theorem; \GSpage{P000012}}
\GSfield{The notes}{The tangent equation is written schematically as \(v'=Av+K'_u\), omitting application of the derivative to \(v\).}
\GSfield{Statement}{For a differentiable flow of \(u'=Au+K(u)\), the derivative in initial data solves \(v'=Av+DK(u(t))v\). Tangent vectors to the solution manifold at \(u\) are these linearized solutions, in the regularity class in which differentiation of the flow is justified.}
\GSfield{Reconstruction}{Restore the factor \(v\) and the derivative-of-flow hypothesis.}
\GSfield{Fate}{This geometric linearization is part of \citet[Part II]{segal1965variete}. Its omission in the notebook is a missing factor, not an alternative evolution equation.}
\GSfield{Present status}{As of November 2026: A standard variational equation under differentiable well-posedness.}
\GSfield{Assessment}{literal statement: false as written; publication match: exact; verification: read the result; development: annotated; confidence: high}
\end{GSitem}
\begin{GSitem}{1964-4.6}{A}{Symplectic form, Poisson brackets and vector fields}
\GSfield{Kind; source}{construction; \GSpage{P000012-P000013}}
\GSfield{The notes}{A skew form on tangent solutions is asserted independent of time and used to form Poisson brackets and associated vector fields.}
\GSfield{Statement}{For a Hamiltonian scalar-wave evolution whose linearized potential is symmetric, \(\Omega(v,w)=\int(v_1w_2-v_2w_1)\,dx\) is conserved. On an infinite-dimensional energy manifold it may be weakly symplectic: a functional has a Hamiltonian vector field only if its derivative lies in the range of \(v\mapsto\Omega(v,\cdot)\).}
\GSfield{Reconstruction}{The Hamiltonian and weak-symplectic qualifications are supplied.}
\GSfield{Fate}{\citet[Part II, discussion of the form and Poisson brackets]{segal1965variete} develops this construction. Arbitrary \(K\) in the preceding abstract evolution need not preserve a symplectic form.}
\GSfield{Present status}{As of November 2026: Established Hamiltonian construction with domain and range conditions; it does not make every differentiable functional Hamiltonian.}
\GSfield{Assessment}{literal statement: correct with missing hypotheses; publication match: exact; verification: read the result; development: annotated; confidence: high}
\end{GSitem}
\begin{GSitem}{1964-4.7}{A}{Canonical commutation relations from the solution geometry}
\GSfield{Kind; source}{construction; \GSpage{P000013}}
\GSfield{The notes}{The classical form and associated differential operators lead to canonical commutation relations, with an explicit operator construction sketched.}
\GSfield{Statement}{On a real symplectic space, the Weyl relations are \(W(f)W(g)=e^{-i\Omega(f,g)/2}W(f+g)\). Infinitesimal field commutators require common invariant domains. A representation and a state are extra choices; infinite-dimensional CCR do not have the finite-dimensional uniqueness property.}
\GSfield{Reconstruction}{The Weyl form is a modern precise comparison; it is not a claim that the notebook contains all domain arguments.}
\GSfield{Fate}{The concrete field construction is treated in \citet[Part III, theorem on printed p. 45]{segal1965variete}, circulated as 1964--65 lectures. The source's finite-dimensional analogy does not establish equivalence of all field representations. The accompanying \citet{segal1965erratum} revises the theorem on p. 45 because of inconsistent d’Alembertian signs and introduces a second renormalization constant besides mass; this limits any literal identification with a one-constant field equation.}
\GSfield{Present status}{As of November 2026: The CCR algebra and its representations are established, but this alone does not construct a renormalized interacting field theory.}
\GSfield{Assessment}{literal statement: correct with missing hypotheses; publication match: partial; verification: read the result; development: annotated; confidence: high}
\end{GSitem}

\GStalk{5}{1964-10-05}{Irving E. Segal}{P000012-P000014}{Scattering and canonical fields}
The short continuation on P000013 mentions renormalization, scattering, generalized probabilities, particle interpretation and equal-time commutators. Several headings state intentions without enough hypotheses for a theorem. The classical wave-operator construction can be identified more precisely than the asserted quantum existence and uniqueness.

\begin{GSitem}{1964-5.1}{F}{Existence and uniqueness for a renormalized cubic field}
\GSfield{Kind; source}{claim; \GSpage{P000013}}
\GSfield{The notes}{Existence, uniqueness and renormalization of a nonlinear field are listed, with no cutoff scheme, dimension, operator domain or theorem statement.}
\GSfield{Statement}{No single existence theorem can be recovered: classical \(\phi^3\) evolution, a formal operator-valued equation, a regularized Hamiltonian, and a renormalized quantum field are distinct objects.}
\GSfield{Fate}{\citet[Parts III--V]{segal1965variete} provides a close contemporary program, but the headings do not identify an exact renormalization result. The accompanying \citet{segal1965erratum} revises the theorem on p. 45 because of inconsistent d’Alembertian signs and introduces a second renormalization constant besides mass; this limits any literal identification with a one-constant field equation.}
\GSfield{Present status}{As of November 2026: No unconditional quantum-field existence or uniqueness is inferred from this passage.}
\GSfield{Assessment}{literal statement: not determinable; publication match: thematic only; verification: read the result; development: none; confidence: low}
\end{GSitem}
\begin{GSitem}{1964-5.2}{A}{The classical scattering transformation}
\GSfield{Kind; source}{construction; \GSpage{P000013}}
\GSfield{The notes}{A scattering map \(S\) compares asymptotic free fields in the distant past and future.}
\GSfield{Statement}{If the nonlinear flow admits incoming and outgoing wave operators \(W_-\) and \(W_+\), with the needed injectivity and range agreement, then the scattering transformation on their common asymptotic data space is \(S=W_+^{-1}W_-\). Existence and completeness depend on the interaction, dimension, sign, size and decay of the data.}
\GSfield{Reconstruction}{The common domain of the two wave operators and invertibility requirements are stated explicitly.}
\GSfield{Fate}{\citet[Part IV, printed pp. 51--60]{segal1965variete} treats asymptotic solutions; \citet{segal1968dispersion} is a later systematic dispersion paper. The latter's metadata were checked, not its complete proof.}
\GSfield{Present status}{As of November 2026: The definition is standard; the unspecified notebook interaction does not support unrestricted asymptotic completeness.}
\GSfield{Assessment}{literal statement: correct with missing hypotheses; publication match: partial; verification: read the result; development: annotated; confidence: high}
\end{GSitem}
\begin{GSitem}{1964-5.3}{F}{Generalized probabilities and particle interpretation}
\GSfield{Kind; source}{remark; \GSpage{P000013}}
\GSfield{The notes}{Generalized probabilities and a particle interpretation appear as headings in the continuation.}
\GSfield{Statement}{The notes do not specify a probability algebra, a state, asymptotic particle observables, or a proposed uniqueness criterion. They identify a physical program, not one determinate mathematical statement.}
\GSfield{Fate}{The discussion is thematically connected to the states and quantization in \citet[Parts I and V]{segal1965variete}; an exact result match is not possible.}
\GSfield{Present status}{As of November 2026: The passage remains identifiable only at the level of research program.}
\GSfield{Assessment}{literal statement: not determinable; publication match: thematic only; verification: read the result; development: none; confidence: low}
\end{GSitem}
\begin{GSitem}{1964-5.4}{E}{Equal-time canonical commutators}
\GSfield{Kind; source}{definition; \GSpage{P000013}}
\GSfield{The notes}{Three equal-time commutators are written; the mixed line has an unclear mark on the second field and an abbreviated delta distribution.}
\GSfield{Statement}{For smeared scalar fields and conjugate momenta on a common invariant domain, \([\phi(f),\phi(g)]=[\pi(f),\pi(g)]=0\) and \([\phi(f),\pi(g)]=i\hbar\int fg\,dx\,I\). The pointwise notation \([\phi(x),\pi(y)]=i\hbar\delta(x-y)I\) is distributional.}
\GSfield{Reconstruction}{The displayed smeared formulation is an editorial normalization, not a diplomatic emendation of the unclear mark.}
\GSfield{Fate}{These are the canonical relations used, rather than invented, in \citet[Parts I and III]{segal1965variete}. The scan does not warrant silently replacing every ambiguous mark in the mixed line.}
\GSfield{Present status}{As of November 2026: Established formalism with smearing and domain conditions; field products at coincident points require additional analysis.}
\GSfield{Assessment}{literal statement: correct with missing hypotheses; publication match: exact; verification: read the result; development: annotated; confidence: high}
\end{GSitem}

\GStalk{6}{1964-10-05}{Laurent Schwartz}{P000012-P000018}{Radon measures beyond local compactness}
Starting on P000013, Schwartz compares abstract measures with Radon measures, pushforwards and completions, compactifications, Radon spaces and stochastic path spaces. The argument separates a cylindrical process law from a Radon law on a suitable function space. Several expansive shorthand claims need restrictions on the topology or a countable family of compact sets.

\begin{GSitem}{1964-6.1}{E}{Abstract measures and Radon measures}
\GSfield{Kind; source}{definition; \GSpage{P000013-P000014}}
\GSfield{The notes}{A countably additive set measure is compared with integration of compactly supported continuous functions and with compact approximation.}
\GSfield{Statement}{A finite Radon measure on a Hausdorff space is a finite Borel measure inner regular by compact sets. On a locally compact Hausdorff space, positive linear functionals on \(C_c(X)\) give Radon measures with the customary local finiteness. For Banach-valued Bochner integration, strong measurability and integrability of the norm are additional conditions.}
\GSfield{Reconstruction}{A finite-measure convention is fixed to make the subsequent statements unambiguous.}
\GSfield{Fate}{The classical foundation is already used by \citet[Sections 1 and 4]{courrege1963mesures}. The later synthesis \citet{schwartz1973radon} was bibliographically verified; its full text was not available in this run.}
\GSfield{Present status}{As of November 2026: Standard measure theory. Positivity, topology, finiteness and the chosen vector-valued integral must be specified.}
\GSfield{Assessment}{literal statement: correct with missing hypotheses; publication match: partial; verification: read the result; development: annotated; confidence: high}
\end{GSitem}
\begin{GSitem}{1964-6.2}{E}{Pushforward and completion of a Radon measure}
\GSfield{Kind; source}{theorem; \GSpage{P000014}}
\GSfield{The notes}{For a map \(f:X\to Y\), the notes compare measurability of \(A\subset Y\) with measurability of \(f^{-1}A\), emphasizing completion.}
\GSfield{Statement}{Let \(\mu\) be a finite Radon measure on a Hausdorff space \(X\), let \(f:X\to Y\) be continuous with \(Y\) Hausdorff, and put \(\nu=f_*\mu\). Then \(A\) is measurable in the completion of \(\nu\) if and only if \(f^{-1}A\) is measurable in the completion of \(\mu\). In the reverse direction compact approximations to the inverse image and its complement give disjoint compact images whose masses exhaust \(\nu(Y)\).}
\GSfield{Reconstruction}{Continuity, Hausdorff separation, finiteness and completed sigma algebras are explicit editorial hypotheses.}
\GSfield{Fate}{This is part of the Radon-measure framework developed in \citet{schwartz1973radon}; only metadata of that book were checked. The compact approximation argument supplied here verifies the stated continuous-map form directly.}
\GSfield{Present status}{As of November 2026: Valid with the hypotheses displayed. The unrestricted abstract-measure version is false, as the next example demonstrates.}
\GSfield{Assessment}{literal statement: correct with missing hypotheses; publication match: partial; verification: indirect; development: annotated; confidence: high}
\end{GSitem}
\begin{GSitem}{1964-6.3}{E}{A non-Radon image loses measurability information}
\GSfield{Kind; source}{construction; \GSpage{P000014-P000015}}
\GSfield{The notes}{A nonmeasurable subset of an interval, of full outer measure, carries the restriction of Lebesgue measure; its inclusion is used to show that inverse-image measurability need not imply image measurability.}
\GSfield{Statement}{Choose a non-Lebesgue-measurable \(X\subset[0,1]\) of outer measure one and define \(\mu(X\cap B)=m(B)\) for Borel \(B\). This is well defined because the complement of \(X\) has inner measure zero. For inclusion \(i\), \(i_*\mu=m\), but \(i^{-1}([0,1]\setminus X)=\varnothing\), although \([0,1]\setminus X\) is not Lebesgue measurable. The source measure is not Radon.}
\GSfield{Reconstruction}{The complement itself supplies the required nonmeasurable set, avoiding the false implication non-Borel implies nonmeasurable.}
\GSfield{Fate}{The example illustrates the precise issue in \citet{schwartz1973radon}; the later book is a thematic source with checked metadata, while the set-theoretic verification here is explicit.}
\GSfield{Present status}{As of November 2026: A valid counterexample once non-Lebesgue-measurability is required; being non-Borel alone would not suffice.}
\GSfield{Assessment}{literal statement: correct with missing hypotheses; publication match: thematic only; verification: indirect; development: annotated; confidence: high}
\end{GSitem}
\begin{GSitem}{1964-6.4}{E}{Compact regularity and compactification}
\GSfield{Kind; source}{theorem; \GSpage{P000015-P000016}}
\GSfield{The notes}{A space is embedded in a compact space by continuous functions, and Radon measures are described by compact approximation. One displayed quantifier ranges over compact sets where Borel sets are needed.}
\GSfield{Statement}{For a completely regular Hausdorff space, the Stone--Čech compactification extends every bounded continuous real function. A finite Borel measure is Radon exactly when every Borel set is approximated from inside by compact sets. When a measure is represented on a compactification, concentration on the original space must be expressed using outer measure if that space is not measurable.}
\GSfield{Reconstruction}{Boundedness and the Borel-set quantifier repair the compressed statements.}
\GSfield{Fate}{\citet[Theorem 1.1 and Section 4]{courrege1963mesures}, a lecture of 6 June 1963, supplies the relevant compactification method. It explicitly uses bounded functions.}
\GSfield{Present status}{As of November 2026: Established. Compact approximation of compact sets alone is tautological and cannot define Radon regularity.}
\GSfield{Assessment}{literal statement: false as written; publication match: partial; verification: read the result; development: annotated; confidence: high}
\end{GSitem}
\begin{GSitem}{1964-6.5}{E}{Radon spaces and universal measurability}
\GSfield{Kind; source}{definition; \GSpage{P000016}}
\GSfield{The notes}{A Radon space is one on which every abstract measure is Radon. A subset of a Radon space is Radon precisely when universally measurable; analytic and Suslin examples are mentioned. A brief analogy compares this definition with nuclear spaces, described through coincidence of two tensor-product constructions.}
\GSfield{Statement}{Under the Hausdorff, finite-Borel-measure convention, a Radon space is one on which every finite Borel measure is inner regular by compact sets. Universal measurability means measurability in every completed finite Radon measure. Suslin subspaces of Hausdorff spaces are universally measurable for Radon measures. The equivalence for subspaces needs the stated ambient Radon and subspace-topology assumptions.}
\GSfield{Reconstruction}{The sigma algebra and meaning of universal measurability are specified. The nuclear-space sentence is an analogy rather than a new measure-theoretic assertion; a literal tensor criterion needs the tensor topologies and the quantifier over partner spaces specified.}
\GSfield{Fate}{These are the Radon/Suslin principles synthesized in \citet{schwartz1973radon}. The notebook is an early exposition, but exact theorem-by-theorem priority was not established from the later metadata.}
\GSfield{Present status}{As of November 2026: Established framework; ``every abstract measure cannot mean every measure on every unrelated sigma algebra.}
\GSfield{Assessment}{literal statement: correct with missing hypotheses; publication match: partial; verification: indirect; development: annotated; confidence: medium}
\end{GSitem}
\begin{GSitem}{1964-6.6}{E}{Function spaces and spaces of operators as Radon spaces}
\GSfield{Kind; source}{claim; \GSpage{P000016}}
\GSfield{The notes}{Separable Fréchet spaces, countable inductive limits, duals and \(\operatorname{Hom}(E,F)\) are grouped together; weak, strong and compact-convergence topologies are mentioned.}
\GSfield{Statement}{Separable Fréchet spaces are Polish and hence Radon. The standard test-function and distribution spaces fit the Suslin/Radon theory with their customary locally convex topologies. More general assertions about inductive limits and operator spaces require specification of regularity, separability and topology; \(\operatorname{Hom}(E,F)\) need not itself be a Fréchet space.}
\GSfield{Reconstruction}{The modern statement intentionally separates a verified elementary subclass from the broader insufficiently specified assertion.}
\GSfield{Fate}{The thematic source is \citet{schwartz1973radon}; its metadata confirm the scope, but no precise universal theorem matching all the shorthand topologies was read. The safe Polish-space case is isolated here.}
\GSfield{Present status}{As of November 2026: The named standard examples are accounted for; an arbitrary countable inductive limit or arbitrary operator topology is not covered by the literal statement.}
\GSfield{Assessment}{literal statement: correct with missing hypotheses; publication match: partial; verification: indirect; development: annotated; confidence: medium}
\end{GSitem}
\begin{GSitem}{1964-6.7}{E}{Measurability of the continuous-path event}
\GSfield{Kind; source}{theorem; \GSpage{P000016-P000017}}
\GSfield{The notes}{The space of continuous functions is viewed inside \(\mathbb R^{\mathbb R}\), and the probability of continuous sample paths is said always to exist.}
\GSfield{Statement}{For compact metric \(K\), the inclusion \(C(K)\to\mathbb R^K\), from the uniform topology to the product topology, is continuous; its image is Suslin and therefore measurable for the completion of any Radon Borel measure on the product. A merely cylindrical law from finite-dimensional distributions does not automatically give this event a probability in its cylinder sigma algebra.}
\GSfield{Reconstruction}{Compact metric time sets and a completed Radon product law make the assertion precise; continuous paths on the real line are handled on successive compact intervals.}
\GSfield{Fate}{This is an application of the Suslin/Radon method described by \citet{schwartz1973radon}. The source identification is thematic rather than a checked numbered theorem.}
\GSfield{Present status}{As of November 2026: Valid for the specified Radon-law setting. The statement must distinguish a law on continuous paths from an arbitrary process on a cylindrical measurable space.}
\GSfield{Assessment}{literal statement: correct with missing hypotheses; publication match: partial; verification: indirect; development: annotated; confidence: high}
\end{GSitem}
\begin{GSitem}{1964-6.8}{E}{Tightness, compact layers and a common modulus of continuity}
\GSfield{Kind; source}{claim; \GSpage{P000017}}
\GSfield{The notes}{A continuous-path law is concentrated on a countable union of compact sets, followed by the assertion that almost all paths have a common modulus.}
\GSfield{Statement}{A Borel probability on \(C(K)\), for compact metric \(K\), is tight and is concentrated on \(\bigcup_j K_j\) with compact \(K_j\). Each \(K_j\) is bounded and equicontinuous. A single deterministic modulus with coefficient one need not hold almost surely: \(X_t=Zt\) with unbounded \(Z\) is a counterexample. One can use compact-dependent moduli, or a fixed vanishing modulus with an almost surely finite random multiplier.}
\GSfield{Reconstruction}{Random multipliers or compact-dependent constants are editorial repairs, not a transcription change.}
\GSfield{Fate}{The compact-layer method is the Radon interpretation in \citet{schwartz1973radon}; the displayed counterexample checks the overstrong notebook conclusion without attributing it to Schwartz as a published theorem.}
\GSfield{Present status}{As of November 2026: The countable compact exhaustion is correct; the unqualified coefficient-one common-modulus assertion is false.}
\GSfield{Assessment}{literal statement: false as written; publication match: partial; verification: indirect; development: annotated; confidence: high}
\end{GSitem}
\begin{GSitem}{1964-6.9}{E}{Support of a finite Radon measure}
\GSfield{Kind; source}{claim; \GSpage{P000015}}
\GSfield{The notes}{The closing sentence says that one cannot define the support of arbitrary Radon measures after dropping local compactness.}
\GSfield{Statement}{For a finite positive Radon Borel measure \(\mu\) on a Hausdorff space \(X\), with inner compact regularity on every Borel set, define \(\operatorname{supp}\mu=X\setminus\bigcup\{U\subset X:U\text{ open},\mu(U)=0\}\). This is closed and has full measure. Every compact subset of its complement has a finite cover by null open sets and hence measure zero; inner compact regularity gives \(\mu(X\setminus\operatorname{supp}\mu)=0\). Local compactness is not needed for this conclusion.}
\GSfield{Reconstruction}{The support definition and compact-cover proof are standard editorial clarification of an otherwise omitted source assertion.}
\GSfield{Fate}{The finite Hausdorff Radon framework is the one used in \citet{courrege1963mesures}. The compact-cover argument above is an editorial check of this sentence, not a claimed quotation from a numbered theorem. The notebook’s exception does not hold under the hypotheses fixed in this record.}
\GSfield{Present status}{As of November 2026: Support exists and is a closed set of full measure in the stated setting. A different, unstated convention for generalized or nonfinite measures cannot be inferred from this sentence.}
\GSfield{Assessment}{literal statement: false as written; publication match: partial; verification: indirect; development: annotated; confidence: high}
\end{GSitem}

\GStalk{7}{1964-10-12}{Laurent Schwartz}{P000016-P000018}{Hilbert subspaces and positive kernels}
The actual mathematical exposition begins on P000018 and continues through P000022, beyond the catalogue span. Hilbert spaces continuously injected into a locally convex space form an ordered cone, identified with positive kernels. Distribution and Sobolev examples explain how this language recovers weak existence theorems. The correspondence was already presented in Schwartz’s 1962 Bourbaki lecture.

\begin{GSitem}{1964-7.1}{E}{Hilbert subspaces of a locally convex space}
\GSfield{Kind; source}{definition; \GSpage{P000018}}
\GSfield{The notes}{\(\operatorname{Hilb}(E)\) consists of Hilbert spaces contained in \(E\), with inclusion continuous; the inclusion need not be a topological embedding onto its image.}
\GSfield{Statement}{For a Hausdorff locally convex space \(E\), a Hilbert subspace is a Hilbert space \(H\) with a continuous injective linear map \(i:H\to E\), identified with its image but carrying its Hilbert norm. Equality includes equality of that norm, not only of the underlying set.}
\GSfield{Fate}{\citet[Section 1]{schwartz1962hilbert} presented this in May 1962, before the Moscow lecture.}
\GSfield{Present status}{As of November 2026: Established definition; an unbounded inverse to the inclusion is permitted.}
\GSfield{Assessment}{literal statement: correct; publication match: exact; verification: read the result; development: annotated; confidence: high}
\end{GSitem}
\begin{GSitem}{1964-7.2}{E}{Scalar multiplication and the Hilbert quotient sum}
\GSfield{Kind; source}{construction; \GSpage{P000018-P000019}}
\GSfield{The notes}{Multiplication rescales the inner product by \(1/\lambda\); the sum norm is the infimum of the sum of squared norms, but the square on the left is absent.}
\GSfield{Statement}{For \(a>0\), \(aH\) has \(\|h\|_{aH}^2=a^{-1}\|h\|_H^2\). For two continuously injected Hilbert spaces, \(H_1+H_2\) is the quotient of \(H_1\oplus H_2\) by the closed kernel of addition, with \(\|h\|^2=\inf_{h=h_1+h_2}(\|h_1\|_1^2+\|h_2\|_2^2)\). Set \(0H=\{0\}\).}
\GSfield{Reconstruction}{The norm square and the zero scalar convention are made explicit.}
\GSfield{Fate}{These operations are given in \citet[Section 1]{schwartz1962hilbert}. The missing square is confirmed by the norm homogeneity.}
\GSfield{Present status}{As of November 2026: Well-defined Hilbert-cone operations; the infimum is attained by orthogonal projection in the direct sum.}
\GSfield{Assessment}{literal statement: false as written; publication match: exact; verification: read the result; development: annotated; confidence: high}
\end{GSitem}
\begin{GSitem}{1964-7.3}{E}{Contractive order and the cancellative cone}
\GSfield{Kind; source}{theorem; \GSpage{P000019-P000020}}
\GSfield{The notes}{The order is inclusion with norm at most one. The cone is called strictly convex, and subtraction is described through a unique positive remainder.}
\GSfield{Statement}{Define \(H_1\le H_2\) when the inclusion is contractive. Under the kernel correspondence this is positive-operator order, and \(H_2=H_1+H_3\) for a unique Hilbert subspace \(H_3\) exactly when \(H_1\le H_2\). The cone is pointed and cancellative. It is not strictly convex in the usual geometric sense: positive semidefinite matrices already have nontrivial faces.}
\GSfield{Reconstruction}{Pointedness replaces the potentially misleading modern reading of strict convexity.}
\GSfield{Fate}{\citet[Sections 1--2, Theorem 1]{schwartz1962hilbert} uses a pointed convex cone. This supplies the meaning behind the notebook terminology.}
\GSfield{Present status}{As of November 2026: Established order structure. The scalar comparison \(aH\ge H\iff a\ge1\) needs \(H\ne0\).}
\GSfield{Assessment}{literal statement: correct with missing hypotheses; publication match: exact; verification: read the result; development: annotated; confidence: high}
\end{GSitem}
\begin{GSitem}{1964-7.4}{E}{The Hilbert-subspace/positive-kernel correspondence}
\GSfield{Kind; source}{theorem; \GSpage{P000020-P000021}}
\GSfield{The notes}{The inclusion \(i\) is composed with its adjoint to form a positive operator from the antidual, and conversely a positive kernel is completed to recover a Hilbert space.}
\GSfield{Statement}{For a quasi-complete Hausdorff locally convex \(E\), a Hilbert subspace \(i:H\to E\) determines the positive Hermitian kernel \(K=ii^*:E^{\prime}_{\!a}\to E\), with the compatible weak topologies. Conversely such a continuous positive kernel determines a unique Hilbert subspace, obtained by completing its range modulo its null space with \((K\phi,K\psi)_H=\langle K\phi,\psi\rangle\).}
\GSfield{Reconstruction}{The topology and completeness assumptions in the published theorem are restored.}
\GSfield{Fate}{The exact match is \citet[Theorem 1, printed p. 257]{schwartz1962hilbert}, announced in May 1962. Quasi-completeness is used to realize the completion inside \(E\).}
\GSfield{Present status}{As of November 2026: Established and widely applicable; the appropriate antidual convention avoids a hidden complex conjugation error.}
\GSfield{Assessment}{literal statement: correct with missing hypotheses; publication match: exact; verification: read the result; development: annotated; confidence: high}
\end{GSitem}
\begin{GSitem}{1964-7.5}{E}{Weak boundedness and completion of a Hilbert kernel}
\GSfield{Kind; source}{theorem; \GSpage{P000021-P000022}}
\GSfield{The notes}{A theorem attributed to Mackey says weakly bounded sets are bounded; the Hilbert unit ball is shown weakly bounded during the completion argument.}
\GSfield{Statement}{In a Hausdorff locally convex space, a subset bounded for the weak topology \(\sigma(E,E')\) is bounded for the original topology. Applied to the kernel unit ball, this gives continuity estimates; quasi-completeness of \(E\) then prevents the Hilbert completion from leaving the ambient space.}
\GSfield{Fate}{This functional-analytic step is used in \citet[proof of Theorem 1]{schwartz1962hilbert}; it is an older boundedness theorem, not a new claim of the seminar.}
\GSfield{Present status}{As of November 2026: Established. Weak boundedness alone is not weak compactness, and completion still requires its separate hypothesis.}
\GSfield{Assessment}{literal statement: correct; publication match: exact; verification: read the result; development: annotated; confidence: high}
\end{GSitem}
\begin{GSitem}{1964-7.6}{E}{Positive distribution kernels and Sobolev examples}
\GSfield{Kind; source}{construction; \GSpage{P000022}}
\GSfield{The notes}{For \(E=\mathcal D'\), positive kernels are distributions in two variables; \(L^2\), \(H^1\), \(H^{-1}\) are examples. The notebook assigns \(-\Delta+\lambda\) to \(H^1\).}
\GSfield{Statement}{A positive distribution kernel \(K\in\mathcal D'(\mathbb R^n\times\mathbb R^n)\) satisfies \(K(\phi\otimes\overline\phi)\ge0\) and gives the corresponding Hilbert subspace of distributions. For \(H^1\) with norm \(\int(\lambda|u|^2+|\nabla u|^2)\), \(\lambda>0\), the associated operator is \((\lambda-\Delta)^{-1}\). The operator \(\lambda-\Delta\) instead corresponds to the reciprocal Sobolev weight.}
\GSfield{Reconstruction}{The inverse operator is an emendation of a faithfully transcribed formula.}
\GSfield{Fate}{The distribution-kernel examples belong to \citet[Section 3]{schwartz1962hilbert}. The inverse is also checked directly from the Fourier multiplier.}
\GSfield{Present status}{As of November 2026: The positive-kernel correspondence is established. Omitting \(\lambda\) on all of Euclidean space changes to a homogeneous space and needs separate low-frequency treatment.}
\GSfield{Assessment}{literal statement: false as written; publication match: exact; verification: read the result; development: annotated; confidence: high}
\end{GSitem}
\begin{GSitem}{1964-7.7}{E}{Hilbert representation of weak elliptic solutions}
\GSfield{Kind; source}{method; \GSpage{P000022}}
\GSfield{The notes}{The talk ends by saying the Hilbert-subspace construction yields existence theorems for generalized solutions, mentioning domain Sobolev spaces.}
\GSfield{Statement}{A precise representative case is: on a bounded domain \(\Omega\), a continuous coercive sesquilinear form \(a\) on \(H_0^1(\Omega)\) gives, for each \(f\in H^{-1}(\Omega)\), a unique \(u\in H_0^1(\Omega)\) with \(a(u,v)=\langle f,v\rangle\) for every \(v\). For Hermitian positive \(a\), the inverse is a positive kernel in the Hilbert-subspace framework.}
\GSfield{Reconstruction}{The displayed variational problem is an explicitly chosen representative example, not a recovered verbatim theorem.}
\GSfield{Fate}{\citet[applications following Theorem 1]{schwartz1962hilbert} gives the Hilbert-kernel context. The notebook has no particular coefficients or boundary problem to match more exactly.}
\GSfield{Present status}{As of November 2026: Established coercive weak-solution method. Noncoercive equations need additional Fredholm or spectral hypotheses.}
\GSfield{Assessment}{literal statement: correct with missing hypotheses; publication match: partial; verification: read the result; development: annotated; confidence: high}
\end{GSitem}

\GStalk{8}{1964-10-12}{S. G. Gindikin}{P000016-P000024}{Weil rigidity and group cohomology}
The actual note is on P000017, between the Radon discussion and the Hilbert-kernel lecture. Gindikin reports Weil’s local rigidity argument, its cohomological criterion, and a dimension formula for surface groups. The brief nondiscrete-subgroup aside lacks a usable hypothesis list.

\begin{GSitem}{1964-8.1}{E}{Local rigidity of cocompact semisimple lattices}
\GSfield{Kind; source}{theorem; \GSpage{P000017}}
\GSfield{The notes}{A semisimple Lie group without a real two-by-two matrix factor and a discrete \(\Gamma\) with compact \(G/\Gamma\) are introduced. All small deformations are said to be conjugate.}
\GSfield{Statement}{For a connected semisimple real Lie group with finite center and no compact factors or simple factor locally isomorphic to \(\mathrm{SL}_2(\mathbb R)\), a uniform lattice is locally rigid: sufficiently nearby embeddings in \(G\) are conjugate by an element of \(G\). The topology is pointwise convergence on a finite generating set.}
\GSfield{Reconstruction}{Finite center, absence of compact factors and the precise local-isomorphism exception give a safe standard formulation of the shorthand.}
\GSfield{Fate}{\citet[Section 3, discussion after its local-rigidity theorem]{weil1964cohomology} relates the cohomological argument to his earlier rigidity papers. The July 1964 paper preceded the October report.}
\GSfield{Present status}{As of November 2026: Established local rigidity in this scope. The excluded surface-group case has a nontrivial deformation space.}
\GSfield{Assessment}{literal statement: correct with missing hypotheses; publication match: exact; verification: read the result; development: annotated; confidence: high}
\end{GSitem}
\begin{GSitem}{1964-8.2}{E}{Vanishing of H1 implies local rigidity}
\GSfield{Kind; source}{theorem; \GSpage{P000017}}
\GSfield{The notes}{The sufficient condition is \(H^1(\Gamma,\mathfrak g)=0\), with \(\Gamma\) acting through the adjoint representation.}
\GSfield{Statement}{For a finitely generated group \(\Gamma\), a real Lie group \(G\), and a homomorphism \(\rho:\Gamma\to G\), if \(H^1(\Gamma,\mathfrak g_{\operatorname{Ad}\rho})=0\), every sufficiently close homomorphism is conjugate to \(\rho\). Cocycles describe the infinitesimal relation equations; coboundaries describe infinitesimal conjugation.}
\GSfield{Fate}{This is the theorem at the end of Section 3 of \citet{weil1964cohomology}, published July 1964.}
\GSfield{Present status}{As of November 2026: Established sufficient criterion. Nonzero \(H^1\) alone does not guarantee integrability of every infinitesimal deformation.}
\GSfield{Assessment}{literal statement: correct with missing hypotheses; publication match: exact; verification: read the result; development: annotated; confidence: high}
\end{GSitem}
\begin{GSitem}{1964-8.3}{F}{The nondiscrete-subgroup aside}
\GSfield{Kind; source}{remark; \GSpage{P000017}}
\GSfield{The notes}{The note says that for nondiscrete \(\Gamma\) the theorem can hold although \(H^1\ne0\). No topology or subgroup class is added.}
\GSfield{Statement}{No single theorem is determined: continuous cohomology of a topological subgroup and abstract group cohomology are different, as are deformations of embeddings and of homomorphisms.}
\GSfield{Fate}{\citet[Section 4]{weil1964cohomology} discusses cases with nonzero cohomology, but the source does not justify identifying those cases with the notebook word ``nondiscrete.}
\GSfield{Present status}{As of November 2026: The aside remains unresolved at the level of identification; it is not a new open rigidity problem.}
\GSfield{Assessment}{literal statement: not determinable; publication match: thematic only; verification: read the result; development: annotated; confidence: low}
\end{GSitem}
\begin{GSitem}{1964-8.4}{E}{Dimension of first cohomology of a surface group}
\GSfield{Kind; source}{theorem; \GSpage{P000017}}
\GSfield{The notes}{For a \(d\)-dimensional representation a formula \((2g-2)d+i+i'\) is given for the first cohomology dimension.}
\GSfield{Statement}{Let \(\Gamma=\pi_1(\Sigma_g)\), \(g\ge2\), and let \(V\) be a finite-dimensional characteristic-zero representation. Then \(\dim H^1(\Gamma,V)=(2g-2)\dim V+\dim V^\Gamma+\dim(V^*)^\Gamma\), by Euler characteristic and Poincaré duality with local coefficients.}
\GSfield{Reconstruction}{The two invariant dimensions, compactness and the absence of orbifold generators are supplied.}
\GSfield{Fate}{The torsion-free compact-surface case of \citet[Sections 4--6]{weil1964cohomology} gives this formula; the paper also treats punctures and elliptic generators, whose extra terms are absent here.}
\GSfield{Present status}{As of November 2026: Established. The two invariant-space terms are essential and need not vanish.}
\GSfield{Assessment}{literal statement: correct with missing hypotheses; publication match: exact; verification: read the result; development: annotated; confidence: high}
\end{GSitem}

\GStalk{9}{1964-10-19}{A. Ya. Povzner}{P000022-P000024}{Thermodynamic potentials and Wigner distributions}
After a brief Kostant bibliographic remark, the notes on P000024--P000025 compare canonical and microcanonical thermodynamics and then introduce the phase-space expression now called the Wigner distribution. Temperature factors and concavity signs require care. The Wigner marginals do not make it a positive classical joint probability.

\begin{GSitem}{1964-9.1}{E}{The Kostant multiplicity reference}
\GSfield{Kind; source}{remark; \GSpage{P000024}}
\GSfield{The notes}{An unattributed preliminary line points to Kostant, volume 93, page 53, 1959, concerning representations of Lie algebras.}
\GSfield{Statement}{The bibliographic coordinates identify Kostant’s weight-multiplicity paper: for a finite-dimensional irreducible highest-weight representation of a complex semisimple Lie algebra, multiplicities are an alternating Weyl-group sum of the positive-root partition function.}
\GSfield{Reconstruction}{The mathematical description comes from the identified title and classical result, not an extra notebook formula.}
\GSfield{Fate}{\citet{kostant1959multiplicity} is exactly Transactions of the AMS 93(1), 53--73; the opened Math-Net record verifies these coordinates. The seminar note contains no displayed multiplicity formula.}
\GSfield{Present status}{As of November 2026: Secure identification of the reference; the speaker of the preliminary line is not established.}
\GSfield{Assessment}{literal statement: correct; publication match: exact; verification: publisher metadata checked; development: annotated; confidence: high}
\end{GSitem}
\begin{GSitem}{1964-9.2}{E}{Canonical free energy and thermodynamic derivatives}
\GSfield{Kind; source}{definition; \GSpage{P000024}}
\GSfield{The notes}{The canonical partition integral is written with \(1/N!\), followed by a per-particle minus log and formulas for pressure and entropy.}
\GSfield{Statement}{For a stable finite-volume classical Hamiltonian \(H_N\), \[Z_N=(N!h^{dN})^{-1}\int e^{-H_N/T}\,dq\,dp.\] Assume \(T>0\) and \(0<Z_N<\infty\). Physical Helmholtz free energy is \(F_N=-T\log Z_N\), with \(S_N=-\partial_TF_N\) and \(p_N=-\partial_VF_N\) when derivatives exist. The reduced quantity \(-\log Z_N/N\) is \(F_N/(NT)\), not \(F_N/N\).}
\GSfield{Reconstruction}{Boltzmann’s constant is set to one and the physical versus reduced free energies are separated.}
\GSfield{Fate}{These are standard canonical definitions used in the setting of \citet[Section 1]{dobrushin1965phase}; the continuum integral here is not a claim that the lattice paper proves every continuum thermodynamic limit.}
\GSfield{Present status}{As of November 2026: Correct with a consistent temperature convention. A thermodynamic-limit assertion for general pair potentials additionally needs stability and short-range/temperedness assumptions.}
\GSfield{Assessment}{literal statement: correct with missing hypotheses; publication match: thematic only; verification: read the result; development: annotated; confidence: high}
\end{GSitem}
\begin{GSitem}{1964-9.3}{E}{Microcanonical entropy and phase volume}
\GSfield{Kind; source}{definition; \GSpage{P000024-P000025}}
\GSfield{The notes}{The logarithm of the phase volume below an energy is introduced as entropy; convexity and a kinetic \(\log T\) term are mentioned.}
\GSfield{Statement}{For \(\Omega_N(E,V)=(N!h^{dN})^{-1}\int_{H_N\le E}dq\,dp\), the microcanonical entropy is \(S_N=\log\Omega_N\) (up to convention-dependent constants and surface-volume choices). In a regular thermodynamic limit for short-range stable systems the entropy density is concave in the extensive densities; canonical free energy is its Legendre--Fenchel transform. Integrating quadratic momenta contributes \((dN/2)\log T\) to \(\log Z_N\).}
\GSfield{Reconstruction}{Concavity and ensemble conventions are corrected explicitly rather than merging the two potentials.}
\GSfield{Fate}{The concavity/phase-coexistence framework is the same thermodynamic setting used in \citet{dobrushin1965phase}. No specific new equivalence-of-ensembles theorem is stated in the notebook.}
\GSfield{Present status}{As of November 2026: Established with the required limit hypotheses; the sign of a kinetic term changes when passing from \(\log Z\) to a minus-log potential.}
\GSfield{Assessment}{literal statement: correct with missing hypotheses; publication match: thematic only; verification: read the result; development: annotated; confidence: high}
\end{GSitem}
\begin{GSitem}{1964-9.4}{E}{Wigner quasidistribution and its marginals}
\GSfield{Kind; source}{construction; \GSpage{P000025}}
\GSfield{The notes}{The notebook gives the position density \(|\psi(x)|^2\) and momentum density \(|\widetilde\psi(p)|^2\), asks for a phase-space expression with these marginals, and says it is constructed in Wigner’s 1932 work. It does not display the bilinear Fourier formula supplied below.}
\GSfield{Statement}{For a normalized Schwartz function \(\psi\), \(W_\psi(x,p)=(2\pi\hbar)^{-d}\int e^{-ip\cdot y/\hbar}\psi(x+y/2)\overline{\psi(x-y/2)}\,dy\). Its marginals are \(|\psi(x)|^2\) and the momentum density in the matching Fourier normalization. It is real but may be negative; the two marginals alone do not uniquely characterize it.}
\GSfield{Reconstruction}{The full Wigner formula, Fourier normalization, regularity assumptions and nonuniqueness from two marginals are editorial supplementation, not a formula transcribed from this page.}
\GSfield{Fate}{\citet{wigner1932quantum}, received 14 March and published 1 June 1932, introduced this phase-space construction. Publisher metadata and abstract were checked; the full article was behind an access restriction.}
\GSfield{Present status}{As of November 2026: An established quasiprobability representation, not a positive classical joint law of position and momentum.}
\GSfield{Assessment}{literal statement: correct with missing hypotheses; publication match: exact; verification: publisher metadata checked; development: annotated; confidence: high}
\end{GSitem}

\GStalk{10}{1964-10-12}{V. I. Arnold}{P000022-P000026}{Kushnirenko’s entropy bound}
P000023--P000024 reports Kushnirenko’s estimate of measure entropy by hypersurface expansion, with a partition-boundary argument. His short paper was received on 1 October and presented by Kolmogorov on 7 October 1964, before this report; publication followed in 1965. The notebook’s sharper invertible estimate is distinguished from the theorem explicitly proved in that announcement.

\begin{GSitem}{1964-10.1}{E}{Partition entropy and Kolmogorov--Sinai entropy}
\GSfield{Kind; source}{definition; \GSpage{P000023}}
\GSfield{The notes}{For a partition \(\xi\), \(-\sum p_i\log p_i\) is followed by the entropy growth of successive refinements and the supremum over partitions.}
\GSfield{Statement}{For a probability-preserving transformation \(T\) and a finite measurable partition \(\xi\), \(h_\mu(T,\xi)=\lim_{n\to\infty}n^{-1}H_\mu(\bigvee_{j=0}^{n-1}T^{-j}\xi)\). The limit exists by subadditivity and \(h_\mu(T)=\sup_\xi h_\mu(T,\xi)\). Fixing the logarithm base fixes the entropy unit.}
\GSfield{Fate}{\citet[opening definitions and cited Rokhlin survey]{kushnirenko1965entropy} uses this already established definition. His announcement uses logarithms to base two.}
\GSfield{Present status}{As of November 2026: Standard definition; forward images of partitions require invertibility, whereas inverse-image refinements work for endomorphisms.}
\GSfield{Assessment}{literal statement: correct with missing hypotheses; publication match: exact; verification: read the result; development: annotated; confidence: high}
\end{GSitem}
\begin{GSitem}{1964-10.2}{B}{Finite entropy bounded by hypersurface expansion}
\GSfield{Kind; source}{theorem; \GSpage{P000023-P000024}}
\GSfield{The notes}{For a smooth transformation preserving Riemannian volume, \(h(T)\le n\log\lambda\), where \(\lambda\) bounds \((n-1)\)-area expansion.}
\GSfield{Statement}{Let \(T\) be a volume-preserving smooth diffeomorphism of a compact \(n\)-dimensional Riemannian manifold, with normalized volume. If every smooth hypersurface piece satisfies \(S(TA)\le\lambda S(A)\), then \(h_\mu(T)\le n\log_2\lambda\) in bits. For endomorphisms the sufficient bound concerns the full inverse image \(T^{-1}A\), not only the forward image.}
\GSfield{Reconstruction}{The forward/preimage distinction and logarithm base are restored.}
\GSfield{Fate}{Exact match: \citet[Theorem 1, pp. 37--38]{kushnirenko1965entropy}; received 1 October 1964, presented 7 October, published 1965. Arnold is thanked for posing the problem.}
\GSfield{Present status}{As of November 2026: Established upper bound. The constant is an area-expansion bound, not automatically the ordinary Lipschitz constant.}
\GSfield{Assessment}{literal statement: correct with missing hypotheses; publication match: exact; verification: read the result; development: annotated; confidence: high}
\end{GSitem}
\begin{GSitem}{1964-10.3}{B}{Boundary proof and the sharper invertible estimate}
\GSfield{Kind; source}{claim; \GSpage{P000024}}
\GSfield{The notes}{The boundary of a joined partition is estimated, and for invertible maps the coefficient \(n-1\) is proposed in place of \(n\).}
\GSfield{Statement}{Let \(T\) be a volume-preserving \(C^1\) diffeomorphism of a compact \(n\)-dimensional Riemannian manifold, with normalized Riemannian volume and \(n\ge2\). If \(\|\wedge^{n-1}DT_x\|\le\lambda\) for all \(x\), then \(\|DT^{-1}_{T(x)}\|\le\lambda\). Lyapunov exponents sum to zero; each negative exponent has magnitude at most \(\log\lambda\), and there are at most \(n-1\) negative exponents unless all vanish. Ruelle’s inequality gives \(h_\mu(T)\le(n-1)\log\lambda\) in nats. In dimension one, preservation of Riemannian volume forces derivative norm one and entropy zero.}
\GSfield{Reconstruction}{The modern proof route through exponents is editorial; the announced 1965 theorem is kept separate.}
\GSfield{Fate}{\citet[p. 38]{kushnirenko1965entropy} proves the coefficient \(n\) and mentions examples attaining \((n-1)\log\lambda\), but does not state the general improved theorem. The later differential-entropy inequality of \citet{ruelle1978inequality} gives the displayed refinement; this is a mathematical implication, not a claim of direct historical dependence.}
\GSfield{Present status}{As of November 2026: The refined bound is established by the stated later route. The notebook does not contain the needed Lyapunov-exponent argument.}
\GSfield{Assessment}{literal statement: correct with missing hypotheses; publication match: partial; verification: read the result; development: annotated; confidence: high}
\end{GSitem}
\begin{GSitem}{1964-10.4}{F}{Entropy as forgetting and the strength of mixing}
\GSfield{Kind; source}{remark; \GSpage{P000023}}
\GSfield{The notes}{Entropy is glossed as the rate at which the initial state is forgotten; a line places zero near weak mixing and infinity near strong mixing.}
\GSfield{Statement}{Metric entropy is an asymptotic information-production rate of finite observations. Strong mixing is a separate asymptotic decorrelation property; neither an entropy value nor a linear scale of entropy classifies mixing strength.}
\GSfield{Fate}{\citet{kushnirenko1965entropy} proves a numerical entropy bound and does not supply the suggested universal hierarchy of mixing. No precise theorem is determined by the drawing.}
\GSfield{Present status}{As of November 2026: The mnemonic should remain a mnemonic; it is not a quantitative equivalence.}
\GSfield{Assessment}{literal statement: not determinable; publication match: none found; verification: read the result; development: annotated; confidence: medium}
\end{GSitem}

\GStalk{11}{1964-10-19}{A. A. Kirillov}{P000024-P000026}{Announcement of Thoma’s character theory}
A short paragraph on P000025 points to Thoma’s 1964 work on countable groups, including the rational affine group. Kirillov develops the topic at length on 2 November, record 13. The normalized tracial notion of character must be distinguished from an arbitrary irreducible unitary representation.

\begin{GSitem}{1964-11.1}{E}{Thoma characters and the rational affine group}
\GSfield{Kind; source}{remark; \GSpage{P000025}}
\GSfield{The notes}{Thoma, Mathematische Zeitschrift volume 84, and the group \(x\mapsto ax+b\) with rational coefficients are mentioned.}
\GSfield{Statement}{For a countable discrete group, normalized positive-definite functions invariant under conjugation form a compact convex set. Its extreme points correspond via the GNS construction to finite-factor representations with a normalized trace. For finite-dimensional factors these are normalized irreducible characters; the general finite-factor case includes type \(\mathrm{II}_1\).}
\GSfield{Reconstruction}{The finite-factor interpretation resolves the compressed phrase ``irreducible characters.}
\GSfield{Fate}{The coordinates identify \citet{thoma1964classes}, Math. Z. 84, 389--402. Its metadata were opened; the modern definition and factor distinction were read in \citet[Definition 1.1 and following discussion]{petersonthom2016characters}. This notebook citation itself establishes knowledge by 19 October 1964. The publisher’s December issue label is not used as proof that printed copies were already available before the talk.}
\GSfield{Present status}{As of November 2026: Established character theory, developed in record 13; it does not classify the entire unitary dual of a non-type-I group.}
\GSfield{Assessment}{literal statement: correct with missing hypotheses; publication match: exact; verification: read the result; development: annotated; confidence: high}
\end{GSitem}

\GStalk{12}{1964-10-19}{Laurent Schwartz}{P000024-P000030}{Smooth vectors, invariant distributions and fundamental solutions}
On P000025--P000028 Schwartz passes from smoothing in representations to Lorentz-invariant distributions and differential equations on groups. The notes distinguish the solvable Lorentz case from the general invariant-solution question and then discuss global fundamental solutions, with a torus obstruction. The latter has both a constant-mode obstruction and a separate small-divisor obstruction.

\begin{GSitem}{1964-12.1}{E}{Smooth vectors and smoothing by test functions}
\GSfield{Kind; source}{theorem; \GSpage{P000025-P000026}}
\GSfield{The notes}{A group representation is averaged against a smooth compactly supported function to produce differentiable vectors; Haar measure is called two-sided.}
\GSfield{Statement}{For a strongly continuous representation \(\pi\) of a Lie group on a Banach space, \(\pi(f)v=\int f(g)\pi(g)v\,dg\), \(f\in C_c^\infty(G)\), is a smooth vector, and an approximate identity proves density of smooth vectors. A left Haar measure suffices; it is bi-invariant only for a unimodular group.}
\GSfield{Reconstruction}{Strong continuity and a left Haar convention replace the unnecessary two-sided assumption.}
\GSfield{Fate}{This is the smoothing mechanism of \citet{garding1947continuous}, communicated in August and published in November 1947. The opened primary archive verifies its original publication; the PDF endpoint did not return full searchable text.}
\GSfield{Present status}{As of November 2026: Established Gårding-density construction; bi-invariance is not available for every Lie group.}
\GSfield{Assessment}{literal statement: correct with missing hypotheses; publication match: exact; verification: publisher metadata checked; development: annotated; confidence: high}
\end{GSitem}
\begin{GSitem}{1964-12.2}{E}{Distribution vectors and the rigged representation space}
\GSfield{Kind; source}{construction; \GSpage{P000026-P000027}}
\GSfield{The notes}{Smooth vectors, their duals and delta-like vectors are used to enlarge the representation and differentiate the group action.}
\GSfield{Statement}{For a continuous unitary representation on \(H\), give \(H^\infty\) its seminorms \(v\mapsto\|d\pi(D)v\|\), \(D\) in the enveloping algebra. Its continuous antidual \(H^{-\infty}\), formed using the contragredient representation, gives \(H^\infty\subset H\subset H^{-\infty}\). Lie-algebra operators act continuously in this distribution-vector sense.}
\GSfield{Reconstruction}{The contragredient antidual and seminorm topology are editorial specifications.}
\GSfield{Fate}{The smoothing basis goes back to \citet{garding1947continuous}; the notebook adds the familiar distributional extension. The exact strong-topology assertion about all orbit maps is not used here.}
\GSfield{Present status}{As of November 2026: Standard rigged-space construction with the stated topology and dual convention.}
\GSfield{Assessment}{literal statement: correct with missing hypotheses; publication match: partial; verification: indirect; development: annotated; confidence: high}
\end{GSitem}
\begin{GSitem}{1964-12.3}{E}{Lorentz-invariant inhomogeneous wave equations}
\GSfield{Kind; source}{theorem; \GSpage{P000027}}
\GSfield{The notes}{For Lorentz-invariant \(A\), the notes assert a Lorentz-invariant solution of \(\Box X=A\), reducing away from the origin to distributions in a quadratic invariant.}
\GSfield{Statement}{For the proper orthochronous Lorentz group, the identity component, on Minkowski space of dimension at least three, invariant distribution theory distinguishes spacelike and timelike orbits and the two time orientations. The full proper group can interchange future and past; the notebook does not specify its group component. Reduction to a single unrestricted expression \(T(q)\) needs separate analysis on the light cone and at the origin. A precise invariant-solution theorem requires this orbit-distribution framework.}
\GSfield{Reconstruction}{The identity-component convention is an explicit editorial choice; the notebook says only Lorentz group. The modern statement describes the necessary framework and verification limit without inventing a surjectivity theorem or theorem number.}
\GSfield{Fate}{\citet{methee1954lorentz}, received October 1953 and published December 1954, studies exactly Lorentz-invariant distributions and wave equations. Its metadata and introductory scope were checked; the complete theorem supporting every inhomogeneous datum was not retrieved.}
\GSfield{Present status}{As of November 2026: The topic is identified as prior invariant-distribution theory; this entry does not assert an unread blanket surjectivity theorem.}
\GSfield{Assessment}{literal statement: correct with missing hypotheses; publication match: partial; verification: indirect; development: annotated; confidence: medium}
\end{GSitem}
\begin{GSitem}{1964-12.4}{E}{Extension across the origin and a cohomological obstruction}
\GSfield{Kind; source}{method; \GSpage{P000027-P000028}}
\GSfield{The notes}{The invariant distribution is extended across the exceptional point, and a finite-dimensional semisimple-group representation is said to have a fixed point.}
\GSfield{Statement}{Suppose a \(G\)-invariant distribution on the punctured space, for an action fixing the origin, has an extension of bounded order. The difference between its transforms and itself is then a continuous 1-cocycle with values in a finite-dimensional space of delta derivatives, once a common finite-order bound is established. For a connected semisimple group in characteristic zero, the vanishing of this finite-dimensional first cohomology permits an invariant extension. The fixed-point assertion concerns an affine action; every linear representation already fixes zero.}
\GSfield{Reconstruction}{The affine cocycle, finite jet space and extension hypothesis make the proof mechanism explicit. Invariance of the punctured datum and an action fixing the origin are required to localize transform differences in the delta-jet space.}
\GSfield{Fate}{The extension problem is part of the setting of \citet{methee1954lorentz}. The finite-dimensional cohomological correction is a reconstruction of the notebook argument; exact theorem identification remains indirect.}
\GSfield{Present status}{As of November 2026: Valid conditional mechanism. Establishing extension and a uniform jet-order bound is an essential analytic step, not a consequence of the word ``distribution alone.}
\GSfield{Assessment}{literal statement: correct with missing hypotheses; publication match: partial; verification: indirect; development: annotated; confidence: medium}
\end{GSitem}
\begin{GSitem}{1964-12.5}{E}{Does invariant data admit an invariant solution?}
\GSfield{Kind; source}{problem; \GSpage{P000027}}
\GSfield{The notes}{After the Lorentz example the notes ask for the analogous assertion for an arbitrary group and invariant differential operator, noting the lack of compact averaging.}
\GSfield{Statement}{If a compact group acts continuously and \(P\) commutes with that action, averaging any solution of \(Pu=f\) with invariant \(f\) gives an invariant solution. For noncompact groups this is false in general: let translations of \(\mathbb R\) act on distributions, take \(P=d/dx\) and \(f=1\). The equation has solutions \(x+C\), while translation-invariant distributions are constant and are killed by \(P\).}
\GSfield{Reconstruction}{The explicit translation counterexample is an editorial check of the quantifiers, not a claimed later theorem by Schwartz.}
\GSfield{Fate}{The compact averaging principle and elementary translation obstruction settle the unrestricted question; \citet{methee1954lorentz} concerns a special noncompact action where additional structure matters. No universal positive theorem is inferred from that example.}
\GSfield{Present status}{As of November 2026: The unrestricted assertion is false; useful versions require assumptions on the action, operator and representation.}
\GSfield{Assessment}{literal statement: correct with missing hypotheses; publication match: partial; verification: indirect; development: annotated; confidence: high; problem \hyperref[prob:GP-1964-02]{GP-1964-02}}
\end{GSitem}
\begin{GSitem}{1964-12.6}{B}{Fundamental solutions for invariant differential operators}
\GSfield{Kind; source}{problem; \GSpage{P000028}}
\GSfield{The notes}{For a polynomial \(P(X_1,\ldots,X_n)\) in invariant vector fields on a Lie group, the problem is to find \(E\) with \(PE=\delta_e\). A positive conjecture for noncompact groups is suggested.}
\GSfield{Statement}{Global existence is not automatic. On a compact group, a nontrivial kernel of the transpose already obstructs \(PE=\delta_e\). On noncompact groups, complex invariant operators can fail even local solvability, as in the Heisenberg-group form of the Lewy operator. Constant-coefficient operators on \(\mathbb R^n\) form the positive Malgrange--Ehrenpreis case, not the general case.}
\GSfield{Reconstruction}{Local and global solvability, real and complex coefficients, and the compact obstruction are separated.}
\GSfield{Fate}{\citet[Introduction and Theorem 1 in Section 4]{cerezo1969fundamental} gives a precise compact-group Fourier criterion; \citet[Section 1]{lerner2000solvability} records the Lewy obstruction and its distributional strengthening. The latter was known by 1960, before this broad conjectural suggestion.}
\GSfield{Present status}{As of November 2026: The universal noncompact conjecture is false. The compact case and many specified operator classes have precise criteria; there is no single unconditional group-wide existence theorem.}
\GSfield{Assessment}{literal statement: false as written; publication match: partial; verification: read the result; development: annotated; confidence: high; problem \hyperref[prob:GP-1964-03]{GP-1964-03}}
\end{GSitem}
\begin{GSitem}{1964-12.7}{B}{Torus fundamental solutions and small divisors}
\GSfield{Kind; source}{problem; \GSpage{P000028}}
\GSfield{The notes}{An irrational first-order vector field on a torus is said to lack an elementary solution; the note contrasts compact groups with noncompact ones.}
\GSfield{Statement}{On \(\mathbb T^2\), \(L=\partial_x-\alpha\partial_y\) never has \(LE=\delta_0\), since the left side annihilates constants. For irrational \(\alpha\), the reduced equation \(LE=\delta_0-1\), with normalized Haar measure, has a distribution solution exactly when \(|m-\alpha n|\ge C(1+|m|+|n|)^{-N}\) for all nonzero integer pairs and some \(C,N>0\). Thus the reduced problem distinguishes non-Liouville from Liouville slopes.}
\GSfield{Reconstruction}{The reduced right-hand side is an explicit modern refinement, not what the notebook literally writes.}
\GSfield{Fate}{\citet[Theorem and Proposition, pp. 112--114]{greenfield1972global} gives the polynomial-symbol criterion and the irrational-slope classification; received November 1970, revised March 1971, published January 1972. The full author-deposited article was read.}
\GSfield{Present status}{As of November 2026: The precise Fourier problem is settled. The notebook objection to a delta fundamental solution holds for every slope because of the constant mode, independently of small divisors.}
\GSfield{Assessment}{literal statement: correct with missing hypotheses; publication match: exact; verification: read the result; development: annotated; confidence: high; problem \hyperref[prob:GP-1964-03]{GP-1964-03}}
\end{GSitem}

\GStalk{13}{1964-11-02}{A. A. Kirillov}{P000028-P000036}{Thoma characters and representations over different fields}
Kirillov reports three 1964 papers of Elmar Thoma, using the affine group over a finite field and over the rationals as examples. The infinite symmetric group and the type-I criterion follow. A question about rational two-by-two matrices leads into a comparison of rank-one character formulae over complex, real, non-Archimedean and finite fields, ending with adeles. The mathematical text begins on P000029 and ends at the top of P000035; the wider catalogue span also contains neighbouring discussions.

\begin{GSitem}{1964-13.1}{E}{Normalized positive-definite class functions and finite factors}
\GSfield{Kind; source}{definition; \GSpage{P000029}}
\GSfield{The notes}{The notes define positive definiteness by the matrices \((f(g_i g_j^{-1}))\), impose conjugacy invariance, and call extreme points characters of irreducible representations. The finite-dimensional example is \(n^{-1}\operatorname{Tr}T(g)\).}
\GSfield{Statement}{For a countable discrete group \(G\), let \(\operatorname{Char}(G)\) be the convex set of conjugation-invariant positive-definite functions \(\tau:G\to\mathbb C\) with \(\tau(e)=1\). An element is extreme exactly when its tracial GNS representation generates a finite factor. If that factor is \(M_n(\mathbb C)\), it is the normalized character of an irreducible finite-dimensional representation; a type-\(\mathrm{II}_1\) factor supplies an infinite-dimensional example without an ordinary operator trace.}
\GSfield{Reconstruction}{Normalization and the finite-factor interpretation are made explicit; positivity means positive semidefinite, not necessarily strictly positive.}
\GSfield{Fate}{This is the framework of \citet{thoma1964classes}; the original metadata were verified. The precise GNS and factor statement is read in \citet[Definition 1.1 and the following discussion]{petersonthom2016characters}.}
\GSfield{Present status}{As of November 2026: Established. Extreme tracial characters classify the relevant finite-factor representations up to their tracial equivalence, not every irreducible unitary representation of a non-type-I group.}
\GSfield{Assessment}{literal statement: correct with missing hypotheses; publication match: exact; verification: read the result; development: annotated; confidence: high}
\end{GSitem}
\begin{GSitem}{1964-13.2}{E}{Irreducible representations of the finite affine group}
\GSfield{Kind; source}{theorem; \GSpage{P000029-P000030}}
\GSfield{The notes}{For \(x\mapsto ax+b\) over \(\mathbb F_p\), the notes count \(p-1\) one-dimensional representations and conjecture one more of dimension \(p-1\), then construct it from nonzero additive characters.}
\GSfield{Statement}{For prime \(p\), the complex irreducible representations of \(G=\mathbb F_p\rtimes\mathbb F_p^\times\) are the \(p-1\) characters \((a,b)\mapsto\chi(a)\), and one representation of dimension \(p-1\) induced from a nontrivial character of the additive normal subgroup. Multiplication acts transitively and freely on its nonzero dual characters. For a generator \(u\) of \(\mathbb F_p^\times\), \(\chi_j(u^r)=e^{2\pi i jr/(p-1)}\); the written \(e^{2\pi i a/p}\) is not a multiplicative character.}
\GSfield{Reconstruction}{The multiplicative-character formula and a consistent conjugation convention replace the incompatible additive exponential and eigenvector indices.}
\GSfield{Fate}{The notebook itself completes the finite-group construction. It is the elementary normal-subgroup induction method proved in \citet[Section 1.6]{prasad2007representations}, used here as a check of the classical argument, not as its first publication.}
\GSfield{Present status}{As of November 2026: Complete classification. The sum of squared degrees is \((p-1)+(p-1)^2=p(p-1)\), so the displayed list exhausts the irreducibles.}
\GSfield{Assessment}{literal statement: false as written; publication match: partial; verification: read the result; development: annotated; confidence: high}
\end{GSitem}
\begin{GSitem}{1964-13.3}{E}{Permutation realization and explicit affine characters}
\GSfield{Kind; source}{construction; \GSpage{P000031-P000032}}
\GSfield{The notes}{Functions on the finite affine line give a permutation representation; constants are removed, and fixed-point counting produces traces \(0,-1,p-1\). The normalized trace is then sent formally to a limit as \(p\to\infty\).}
\GSfield{Statement}{Let \(G=\operatorname{Aff}(\mathbb F_p)\) act on functions by \((\pi(g)f)(x)=f(g^{-1}x)\). Its zero-sum subspace is the irreducible representation of dimension \(p-1\), with character \(0\) when \(a\ne1\), \(-1\) when \(a=1,b\ne0\), and \(p-1\) at the identity. Division by \(p-1\) gives the normalized tracial character. The unreduced permutation character counts the one, zero, or \(p\) fixed points respectively.}
\GSfield{Reconstruction}{The inverse in the action makes it a left representation; the handwritten convention without the inverse is an antirepresentation. The identity row is reconstructed mathematically while its ambiguous source glyph remains logged.}
\GSfield{Fate}{The formulas follow from the explicit representation in the notebook and the finite-group irreducibility criterion of \citet[Section 1.6]{prasad2007representations}.}
\GSfield{Present status}{As of November 2026: Established finite calculation. A pointwise limit across changing groups requires chosen identifications of the elements; the notebook comparison is a suggestive formula comparison.}
\GSfield{Assessment}{literal statement: correct with missing hypotheses; publication match: partial; verification: read the result; development: annotated; confidence: high}
\end{GSitem}
\begin{GSitem}{1964-13.4}{E}{The rational affine group and the limiting regular character}
\GSfield{Kind; source}{theorem; \GSpage{P000032}}
\GSfield{The notes}{After the finite table, the notes ask what remains for rational coefficients and display a limiting column equal to one at the identity and zero elsewhere.}
\GSfield{Statement}{For the discrete group \(G=\mathbb Q\rtimes\mathbb Q^\times\), the extreme normalized tracial characters are the one-dimensional functions \(\tau(a,b)=\chi(a)\), with \(\chi\) a unitary character of the discrete multiplicative group, and the regular character \(\delta_e\). The nonzero translations form one conjugacy class. For an extreme trace, the projection onto translation-fixed vectors is central in the generated factor, hence is zero or one; in the first case orthogonality of translations forces vanishing away from the identity, and in the second the representation factors through \(\mathbb Q^\times\).}
\GSfield{Reconstruction}{The extremality argument is an editorial verification. The finite-field table does not by itself prove a limiting theorem for rational groups.}
\GSfield{Fate}{The report identifies \citet{thoma1964classes} as the source of this example; record 11 already names its Math. Z. volume on 19 October 1964. Publisher metadata give December 1964 for volume 84 and February 1964 for the related volume-85 paper, an ordering inconsistent with a straightforward reading of their mutual bibliography. Therefore the notebook establishes knowledge by October, while an exact prior public-distribution date is not asserted. The brief factor-projection argument supplied here checks the formula without claiming that the unread original uses the same proof.}
\GSfield{Present status}{As of November 2026: The tracial-character classification is settled. It does not imply that the entire unitary dual of the rational affine group admits a comparable finite list.}
\GSfield{Assessment}{literal statement: correct with missing hypotheses; publication match: exact; verification: indirect; development: annotated; confidence: high}
\end{GSitem}
\begin{GSitem}{1964-13.5}{E}{Spectral decomposition and dynamical systems for affine representations}
\GSfield{Kind; source}{method; \GSpage{P000031}}
\GSfield{The notes}{The translation operators are simultaneously diagonalized; multiplication acts on their characters, and ergodic measures are said to give representations. Vector-valued fibers are explicitly indicated.}
\GSfield{Statement}{For a unitary representation of a locally compact abelian translation group \(A\), the spectral theorem decomposes the representation over \(\widehat A\), with possible multiplicity spaces. A compatible action of dilations transports this spectral measure and its fibers. A concrete realization therefore includes a quasi-invariant measure class, the associated Radon--Nikodym factors and a measurable unitary cocycle; an ergodic measure alone need not determine the representation.}
\GSfield{Reconstruction}{The measure-class and multiplicity qualifications make the spectral construction precise. The uncertain Haar-measure abbreviation is not expanded into a theorem.}
\GSfield{Fate}{This is the semidirect-product spectral method underlying the finite construction in \citet[Section 1.6]{prasad2007representations} and the character discussion of \citet{thoma1964classes}. These references confirm the setting; an exact later theorem matching every abbreviated infinite-dimensional assertion was not read.}
\GSfield{Present status}{As of November 2026: A standard representation-theoretic framework. The topology on \(\mathbb R\) versus discrete \(\mathbb Q\), multiplicity and cocycle data must be fixed before a classification claim is meaningful.}
\GSfield{Assessment}{literal statement: correct with missing hypotheses; publication match: partial; verification: indirect; development: annotated; confidence: medium}
\end{GSitem}
\begin{GSitem}{1964-13.6}{F}{Restriction to a normal subgroup and extreme points}
\GSfield{Kind; source}{claim; \GSpage{P000032}}
\GSfield{The notes}{A normal subgroup \(G_0\) is introduced, followed by a compressed statement that certain extreme points coincide with those for the original group. The exact invariant convex sets and the final abbreviation are unclear.}
\GSfield{Statement}{Restriction of an extreme normalized character to a normal subgroup need not be extreme among all normalized characters of that subgroup: the normalized two-dimensional irreducible character of \(S_3\) restricts to the average of the two nontrivial characters of \(A_3\). Extremality in a smaller convex set invariant under conjugation by the ambient group is a different statement.}
\GSfield{Reconstruction}{The counterexample is an editorial test of one possible reading; it is not attributed to Thoma.}
\GSfield{Fate}{The surrounding reference is \citet{thoma1964classes}, but its original proof was not retrieved and the source does not specify the convex set needed to identify a theorem. The finite counterexample excludes the most literal unrestricted reading.}
\GSfield{Present status}{As of November 2026: The intended lemma is unresolved as a textual identification, not an open theorem about all normal subgroups.}
\GSfield{Assessment}{literal statement: not determinable; publication match: none found; verification: unresolved lead; development: annotated; confidence: low}
\end{GSitem}
\begin{GSitem}{1964-13.7}{E}{Thoma characters of the infinite symmetric group}
\GSfield{Kind; source}{theorem; \GSpage{P000032-P000033}}
\GSfield{The notes}{For finite-support permutations, the cycle counts \(n_k\) determine conjugacy, and every character is written as \(\prod a_k^{n_k}\). The restrictions on the real numbers \(a_k\) are omitted.}
\GSfield{Statement}{Let \(S_\infty\) be the group of finite-support permutations of \(\mathbb N\). Its extreme normalized characters are indexed by nonincreasing sequences \(\alpha_i,\beta_i\ge0\) with \(\sum_i\alpha_i+\sum_i\beta_i\le1\), and satisfy \(\chi(\sigma)=\prod_{k\ge2}(\sum_i\alpha_i^k+(-1)^{k-1}\sum_i\beta_i^k)^{n_k(\sigma)}\). Only finitely many factors differ from one; this formula does not describe arbitrary unconstrained real sequences.}
\GSfield{Reconstruction}{The parameter simplex and the restriction to extreme characters supply the essential missing hypotheses.}
\GSfield{Fate}{The exact source is \citet{thoma1964symmetric}, Math. Z. 85, 40--61 (1964). Its metadata were opened. The full parameter formula and an operator-algebraic proof are read in \citet[Theorem 0.1 and Sections 6--8]{gohm2010characters}, circulated in May 2010.}
\GSfield{Present status}{As of November 2026: Thoma’s classification is established. General normalized characters are mixtures of these extreme characters, so they need not be multiplicative over disjoint cycles.}
\GSfield{Assessment}{literal statement: correct with missing hypotheses; publication match: exact; verification: read the result; development: annotated; confidence: high}
\end{GSitem}
\begin{GSitem}{1964-13.8}{E}{Type I countable discrete groups}
\GSfield{Kind; source}{theorem; \GSpage{P000033}}
\GSfield{The notes}{The notes give a normal subgroup \(G_0\) with finite quotient and insert \ru{изом. G}; they call this necessary and sufficient for type I.}
\GSfield{Statement}{A countable discrete group is type I exactly when it is virtually abelian. Equivalently it has an abelian normal subgroup of finite index, since one can take the core of any abelian finite-index subgroup. Merely requiring a finite-index subgroup isomorphic to the original group is not a type-I criterion.}
\GSfield{Reconstruction}{Replace the asserted criterion by virtual abelianness as a mathematical emendation, not a transcription repair.}
\GSfield{Fate}{This is \citet{thoma1964unitary}, Math. Ann. 153, 111--138, published April 1964 after receipt on 18 January 1963. The metadata and the statement as cited in \citet[Introduction]{petersonthom2016characters} were checked; the full original proof was not retrieved.}
\GSfield{Present status}{As of November 2026: Established classification. The omitted word is mathematically consequential, and the scan does not justify pretending that the manuscript actually says abelian.}
\GSfield{Assessment}{literal statement: false as written; publication match: exact; verification: indirect; development: annotated; confidence: high}
\end{GSitem}
\begin{GSitem}{1964-13.9}{B}{Characters of rational two-by-two matrix groups}
\GSfield{Kind; source}{problem; \GSpage{P000033}}
\GSfield{The notes}{The notes label characters of two-by-two matrices over the rationals as not done. The determinant is not specified in that line, but a later matrix in the same discussion has determinant one.}
\GSfield{Statement}{For the definite interpretation \(G=\mathrm{SL}_2(\mathbb Q)\) with the discrete topology, classify extreme normalized tracial characters. There are exactly the trivial character and the two functions supported on \(Z(G)=\{\pm I\}\) which take value one at \(I\) and value \(+1\) or \(-1\) at \(-I\). For \(\mathrm{PSL}_2(\mathbb Q)\) the list is the trivial and regular characters.}
\GSfield{Reconstruction}{The determinant-one group is an explicit interpretation supported by the later matrix, and the central signs are the specialization of the theorem to \(\mathbb Q\).}
\GSfield{Fate}{\citet[Theorem 2.4]{petersonthom2016characters} proves that every nontrivial extreme character of \(\mathrm{SL}_2(k)\), for any infinite field \(k\), is induced from its center. The preprint appeared 16 March 2013, the article online 11 March 2014, and the print issue on 1 July 2016. The theorem and its proof were read.}
\GSfield{Present status}{As of November 2026: The stated special-linear interpretation is solved. The unqualified source line does not license silently replacing the group by \(\mathrm{GL}_2(\mathbb Q)\) and claiming that the same theorem is its full classification.}
\GSfield{Assessment}{literal statement: correct with missing hypotheses; publication match: partial; verification: read the result; development: annotated; confidence: high; problem \hyperref[prob:GP-1964-04]{GP-1964-04}}
\end{GSitem}
\begin{GSitem}{1964-13.10}{F}{A field-uniform functor of characters}
\GSfield{Kind; source}{question; \GSpage{P000030-P000033}}
\GSfield{The notes}{Representations and characters are proposed as functors of an unknown field; identities valid over several fields are suggested to hold universally.}
\GSfield{Statement}{An algebraic group can be a functor on fields, but its complex unitary or tracial character theory additionally depends on the topology of its field of points, the category of representations, parameters and normalization. No category, transition maps or transfer principle for the proposed character functor is specified.}
\GSfield{Reconstruction}{No modern transfer theorem or Langlands statement is substituted for the unformalized proposal.}
\GSfield{Fate}{The local formulas compared below and the adelic construction in \citet{gelfand1964adeles} account for the motivation. They do not give a theorem that equality over several fields forces equality over every field.}
\GSfield{Present status}{As of November 2026: A recognisable research programme with insufficient data for a definite mathematical claim or an open-problem status.}
\GSfield{Assessment}{literal statement: not determinable; publication match: thematic only; verification: unresolved lead; development: none; confidence: medium}
\end{GSitem}
\begin{GSitem}{1964-13.11}{E}{Complex rank-one distribution characters}
\GSfield{Kind; source}{theorem; \GSpage{P000034}}
\GSfield{The notes}{For complex matrices of determinant one the notes write a numerator \(\pi(\lambda)+\pi(\lambda^{-1})\) divided by \(|\lambda-\lambda^{-1}|^2\), attributed to Gelfand and Naimark, 1947.}
\GSfield{Statement}{For the normalized principal series of \(\mathrm{SL}_2(\mathbb C)\) induced by a unitary character \(\chi\) of \(\mathbb C^\times\), its distribution character on a regular semisimple element with eigenvalues \(\lambda,\lambda^{-1}\) is \((\chi(\lambda)+\chi(\lambda^{-1}))/|\lambda-\lambda^{-1}|^2\), with ordinary complex modulus in the denominator. The trace is defined after integration against a test function; it is not an ordinary trace of the individual infinite-dimensional operator.}
\GSfield{Reconstruction}{Normalized induction, the distributional meaning and the regular-semisimple domain are supplied; the 1947 original was not separately opened.}
\GSfield{Fate}{The notebook’s 1947 attribution is consistent with the historical references in \citet{gelfandgraev1963representations}. Its Section 5.1, equation (5), gives the principal-series formula with the local-field absolute value, which for \(\mathbb C\) is the square of ordinary modulus.}
\GSfield{Present status}{As of November 2026: Established regular-semisimple character formula. Values at the identity or repeated eigenvalues cannot be obtained by substituting directly into this quotient.}
\GSfield{Assessment}{literal statement: correct with missing hypotheses; publication match: exact; verification: read the result; development: annotated; confidence: high}
\end{GSitem}
\begin{GSitem}{1964-13.12}{E}{Real rank-one principal and discrete series}
\GSfield{Kind; source}{theorem; \GSpage{P000034}}
\GSfield{The notes}{The real principal-series expression vanishes for nonreal eigenvalues. A second expression, apparently for a discrete-series family, has an undefined minimum eigenvalue and a row without its condition; the 1947 marginal surname remains unclear.}
\GSfield{Statement}{For a normalized unitary principal series of \(\mathrm{SL}_2(\mathbb R)\) induced from a character \(\chi\) of \(\mathbb R^\times\), the regular character is \((\chi(\lambda)+\chi(\lambda^{-1}))/|\lambda-\lambda^{-1}|\) on split semisimple elements and zero on elliptic regular elements. The discrete series are additional representations. For a holomorphic discrete series with lowest rotation weight \(k\ge2\), one compatible elliptic boundary-value convention gives \(e^{ik\theta}/(1-e^{2i\theta})\) for \(\theta\notin\pi\mathbb Z\); changing the rotation convention conjugates or reverses this formula.}
\GSfield{Reconstruction}{The distinction between the two families and the displayed rotation-weight convention are editorial. The uncertain surname is retained rather than silently replaced with Bargmann.}
\GSfield{Fate}{The principal-series formula is the real specialization of \citet[Section 5.1, equation (5)]{gelfandgraev1963representations}. The discrete-series expression is an explicit modern highest/lowest-weight reconstruction from the geometric series of rotation weights, not a claimed decipherment of the incomplete second brace.}
\GSfield{Present status}{As of November 2026: The representation families and distribution characters are established. The source’s second formula and marginal author cannot be uniquely identified at the same confidence as the principal-series line.}
\GSfield{Assessment}{literal statement: not determinable; publication match: partial; verification: read the result; development: annotated; confidence: medium}
\end{GSitem}
\begin{GSitem}{1964-13.13}{E}{Non-Archimedean principal-series characters}
\GSfield{Kind; source}{theorem; \GSpage{P000034}}
\GSfield{The notes}{The same quotient is written for a p-adic field, zero when the eigenvalue is not in the base field, and attributed to Gelfand and Graev, 1963.}
\GSfield{Statement}{Let \(K\) be a non-Archimedean local field of characteristic zero and odd residue characteristic, and let \(\chi\) be a unitary character of \(K^\times\). The normalized induced principal-series character on a regular split semisimple element of \(\mathrm{SL}_2(K)\), with eigenvalues \(\lambda,\lambda^{-1}\), is \((\chi(\lambda)+\chi(\lambda^{-1}))/|\lambda-\lambda^{-1}|_K\); it is zero on regular elliptic elements. At reducibility parameters this is the character of the entire induced representation, not automatically of either irreducible constituent.}
\GSfield{Reconstruction}{The normalized local absolute value and the distinction between induced and irreducible characters remove hidden conventions; the chosen field range stays within the opened source.}
\GSfield{Fate}{Exact source match: \citet[Section 5.1, equation (5)]{gelfandgraev1963representations}, received 8 April 1963 and published in the 1963 survey. The formula and its regularity qualifications were read.}
\GSfield{Present status}{As of November 2026: Established principal-series formula. Supercuspidal and other discrete-series characters are not exhausted by this one line.}
\GSfield{Assessment}{literal statement: correct with missing hypotheses; publication match: exact; verification: read the result; development: annotated; confidence: high}
\end{GSitem}
\begin{GSitem}{1964-13.14}{E}{Finite-field split and nonsplit character formulae}
\GSfield{Kind; source}{theorem; \GSpage{P000034}}
\GSfield{The notes}{For a finite field, the split formula loses its denominator, and another family has \(-\pi(\lambda)-\pi(\lambda^{-1})\) on nonsplit elements. The margin says Frobenius, 1895.}
\GSfield{Statement}{For odd \(q\), a generic principal-series character of \(\mathrm{SL}_2(\mathbb F_q)\), induced from \(\chi:\mathbb F_q^\times\to\mathbb C^\times\) with \(\chi^2\ne1\), has value \(\chi(\lambda)+\chi(\lambda^{-1})\) on regular split elements and zero on regular nonsplit elements. A generic cuspidal character associated with a character \(\theta\) of the norm-one torus in \(\mathbb F_{q^2}^\times\), \(\theta^2\ne1\), is zero on regular split elements and \(-\theta(\lambda)-\theta(\lambda^{-1})\) on regular nonsplit elements. Central, unipotent and exceptional parameters require further rows.}
\GSfield{Reconstruction}{The two parameter groups differ: the nonsplit parameter belongs to the norm-one torus, not to the multiplicative group of the base field.}
\GSfield{Fate}{The formula is read in \citet[Sections 2.5 and 3.6, especially Exercise 3.44 and Lemma 3.45]{prasad2007representations}, whose construction supplies a modern verification. Its introduction dates the complete finite-group classifications to Jordan and Schur in 1907, so the notebook’s bare Frobenius 1895 label is not used as a verified exact publication citation.}
\GSfield{Present status}{As of November 2026: The finite character theory is established. The two displayed families are not a complete character table without their exceptional and nonregular cases.}
\GSfield{Assessment}{literal statement: correct with missing hypotheses; publication match: exact; verification: read the result; development: annotated; confidence: high}
\end{GSitem}
\begin{GSitem}{1964-13.15}{E}{Trace as a proposed conjugacy coordinate}
\GSfield{Kind; source}{claim; \GSpage{P000033}}
\GSfield{The notes}{For a determinant-one matrix, the notes say the conjugacy class is determined by \(\alpha+\delta\).}
\GSfield{Statement}{Over an algebraically closed field, a regular semisimple \(\mathrm{SL}_2\) conjugacy class is determined by its trace. Without regularity, the identity and a nontrivial unipotent matrix both have trace two but are not conjugate. Over a nonclosed field, a stable conjugacy class can split into several rational conjugacy classes, so trace alone is not a universal coordinate.}
\GSfield{Reconstruction}{The regularity and field restrictions are mathematical corrections, not changes to the visible handwriting.}
\GSfield{Fate}{The explicit split, elliptic, central and unipotent distinctions in \citet[Section 2.4]{prasad2007representations} and \citet[Section 5]{gelfandgraev1963representations} show the missing domain of the notebook assertion.}
\GSfield{Present status}{As of November 2026: The unqualified assertion is false; it is a valid coordinate on the specified regular semisimple algebraically closed setting.}
\GSfield{Assessment}{literal statement: false as written; publication match: partial; verification: read the result; development: annotated; confidence: high}
\end{GSitem}
\begin{GSitem}{1964-13.16}{E}{Adeles as a framework for local-field analogies}
\GSfield{Kind; source}{construction; \GSpage{P000035}}
\GSfield{The notes}{Adeles are described as a direct sum of all p-adic groups, with the rational numbers as the diagonal.}
\GSfield{Statement}{The adele ring of \(\mathbb Q\) is \(\mathbb A_{\mathbb Q}=\mathbb R\times\prod_p'\mathbb Q_p\), restricted relative to \(\mathbb Z_p\): all but finitely many finite coordinates are integral, not necessarily zero. The diagonal \(\mathbb Q\) is discrete and cocompact in this full additive adele group. Omitting the real place gives finite adeles, in which \(\mathbb Q\) is dense rather than discrete.}
\GSfield{Reconstruction}{Restricted product and the Archimedean place replace the shorthand direct-sum description; additive cocompactness is the standard specialization to \(\mathbb G_a\).}
\GSfield{Fate}{\citet[Section 1, p. 487]{gelfand1964adeles}, received 8 January 1964 and published before the talk, explicitly uses all inequivalent valuations and the restricted-product condition for adelic matrix groups. The definition and diagonal-subgroup statement were read.}
\GSfield{Present status}{As of November 2026: Established local-to-global framework. The literal algebraic direct sum cannot contain the proposed nonzero rational diagonal.}
\GSfield{Assessment}{literal statement: false as written; publication match: exact; verification: read the result; development: annotated; confidence: high}
\end{GSitem}

\GStalk{14}{1964-11-02}{I. I. Pyatetskii-Shapiro}{P000034-P000036}{Linnik’s lattice points on spheres}
P000036 reports Linnik’s use of arithmetic dynamics to distribute integer solutions of \(x^2+y^2+z^2=m\) on the sphere. The fixed splitting prime is visible; the source also warns that some radii have few points. The catalogue assigns the report to 2 November, while the preceding page separately announces Linnik for the 16th. The handwritten parenthesis Ilya Iosifovich resolves the speaker’s otherwise merged initials.

\begin{GSitem}{1964-14.1}{E}{Admissible integers represented by three squares}
\GSfield{Kind; source}{theorem; \GSpage{P000036}}
\GSfield{The notes}{The problem asks which \(m\) are sums of three squares and notes the obstruction \(m=8a+7\).}
\GSfield{Statement}{A positive integer \(m\) is represented by three integral squares exactly when \(m\ne4^a(8b+7)\) for nonnegative integers \(a,b\). The visible mod-eight obstruction is necessary but is only one stage of the complete criterion, since divisibility by four forces all three coordinates even.}
\GSfield{Reconstruction}{The full obstruction and the powers-of-four countersequence supplement the source’s single congruence example.}
\GSfield{Fate}{This classical criterion predates the seminar and provides the arithmetic background to \citet[Section 1]{linnik1957sphere}; a checked modern discussion appears in \citet[Section 1]{ellenberg2010linnik}.}
\GSfield{Present status}{As of November 2026: Established existence criterion. Existence alone does not imply equidistribution as the radius grows: multiplying a fixed radius squared by powers of four preserves the finite set of directions.}
\GSfield{Assessment}{literal statement: correct; publication match: partial; verification: read the result; development: annotated; confidence: high}
\end{GSitem}
\begin{GSitem}{1964-14.2}{E}{Linnik angular equidistribution with a fixed splitting prime}
\GSfield{Kind; source}{theorem; \GSpage{P000036}}
\GSfield{The notes}{The proportion of lattice points projecting into a region is said to approach its area when \(\xi^2+m\equiv0\pmod q\) is solvable for a fixed prime \(q\).}
\GSfield{Statement}{Fix a prime \(q>3\). Let \(m\to\infty\) through integers congruent to \(1,2,3,5,6\pmod8\) and satisfying the Legendre-symbol condition \(\bigl(\frac{-m}{q}\bigr)=1\). For a fixed region \(\Omega\subset S^2\) with piecewise smooth boundary, the proportion of points \(x\in\mathbb Z^3\), \(|x|^2=m\), for which \(x/\sqrt m\in\Omega\), tends to \(\operatorname{area}(\Omega)/(4\pi)\). The analogous statement holds for primitive points. Requiring Legendre symbol one also excludes \(q\mid m\), so mere solvability with the zero residue is weaker.}
\GSfield{Reconstruction}{The admissible congruence classes, fixed prime, nonzero square condition, normalized area and limit are supplied from the original theorem.}
\GSfield{Fate}{Exact prior match: \citet[Theorem 1, Section 2, p. 258]{linnik1957sphere}, received 11 July 1956 and published 1957. The paper distinguishes all and primitive points in equations (2.1)--(2.2). \citet[pp. 73--74]{duke1988hyperbolic} later removes the fixed splitting-prime restriction by modular-form methods.}
\GSfield{Present status}{As of November 2026: The stated theorem is established. In particular, for squarefree admissible \(m\), Duke’s 1988 result gives equidistribution without a fixed splitting prime. No unrestricted assertion along \(m=4^a m_0\) follows.}
\GSfield{Assessment}{literal statement: correct with missing hypotheses; publication match: exact; verification: read the result; development: annotated; confidence: high}
\end{GSitem}
\begin{GSitem}{1964-14.3}{E}{Arithmetic dynamics from rational orthogonal matrices}
\GSfield{Kind; source}{method; \GSpage{P000036}}
\GSfield{The notes}{Matrices with entries \(a/q\) are used to move integral sphere points; the note says each point has its own 48 orthogonal matrices, then calls the evolution a Markov chain.}
\GSfield{Statement}{Under the splitting condition, Linnik constructs arithmetic trajectories using rational orthogonal transformations with prime-power denominators, after passing to a fundamental region for the finite integral symmetry group. Choosing a root of \(c^2+m\equiv0\pmod{q^k}\) selects one of two directions in the reduced dynamics. Equidistribution requires arithmetic control of long trajectories and their correlations; it does not follow solely from calling a finite evolution a Markov chain.}
\GSfield{Reconstruction}{The distinction between a deterministic reduced trajectory, integral symmetries and a Markov presentation prevents an unsupported identification of the number 48.}
\GSfield{Fate}{\citet[Section 3, pp. 258--260, Theorems 3--4]{linnik1957sphere} gives the deterministic reduced flow and its ergodic/mixing statements. \citet[Sections 2, 8 and 9]{ellenberg2010linnik} develops a graph and random-walk formulation. Neither read source identifies the notebook’s unexplained 48 as a universal transition degree for each prime.}
\GSfield{Present status}{As of November 2026: The arithmetic method is established and has later graph-theoretic refinements. The written 48 is retained as an unresolved combinatorial shorthand: the full integral orthogonal symmetry group has order 48, but this fact does not establish that the manuscript means precisely that group.}
\GSfield{Assessment}{literal statement: not determinable; publication match: partial; verification: read the result; development: annotated; confidence: medium}
\end{GSitem}

\GStalk{15}{1964-11-02}{M. I. Vishik}{P000034-P000037}{Fredholm index, nonlinear equations, and families}
P000035 moves from Fredholm perturbation invariance to nonlinear equations and parameter families. The homotopy claim is meaningful only after specifying the operator space. Catalogue envelopes overlap adjacent records; the item locators identify the actual passage.

\begin{GSitem}{1964-15.1}{E}{Stability of the Fredholm index}
\GSfield{Kind; source}{theorem; \GSpage{P000035}}
\GSfield{The notes}{The notes set \(k=\dim\ker T\), \(\ell=\dim\operatorname{coker}K\), \(\text{\ae}=k-\ell\), unchanged by small or completely continuous perturbations.}
\GSfield{Statement}{For a bounded Fredholm \(T:X\to Y\) between Banach spaces, the index is locally constant in operator norm and unchanged by addition of a compact operator.}
\GSfield{Reconstruction}{Banach spaces, boundedness and norm topology are supplied; the isolated K is normalized to T.}
\GSfield{Fate}{The background is explicitly recalled by \citet[p.~861, Theorem~1.1]{smale1965sard}; the elliptic setting already appears in \citet[p.~423]{atiyahsinger1963index}.}
\GSfield{Present status}{As of November 2026, both invariance assertions hold with Fredholmness, closed range, and a fixed topology.}
\GSfield{Assessment}{literal statement: correct with missing hypotheses; publication match: exact; verification: read the result; development: annotated; confidence: high}
\end{GSitem}
\begin{GSitem}{1964-15.2}{B}{What the integer index classifies}
\GSfield{Kind; source}{claim; \GSpage{P000035}}
\GSfield{The notes}{Vishik considers index the only homotopy invariant.}
\GSfield{Statement}{For infinite-dimensional separable complex Hilbert \(H\), the path components of the norm-topological \(\operatorname{Fred}(H)\) are classified by the integer index. For compact Hausdorff \(X\), families instead satisfy \([X,\operatorname{Fred}(H)]\cong K^0(X)\).}
\GSfield{Reconstruction}{Ambient Hilbert space and topology supplied.}
\GSfield{Fate}{\citet[pp.~129--142]{janich1965fredholm}, published April 1965, gives the precise family formulation. The notebook does not specify the intended ambient space.}
\GSfield{Present status}{As of November 2026, the component statement is true in the full Fredholm space; it does not classify restricted elliptic symbols or all families.}
\GSfield{Assessment}{literal statement: correct with missing hypotheses; publication match: partial; verification: publisher metadata checked; development: annotated; confidence: medium}
\end{GSitem}
\begin{GSitem}{1964-15.3}{B}{Index for nonlinear elliptic maps}
\GSfield{Kind; source}{question; \GSpage{P000035}}
\GSfield{The notes}{What are kernel, cokernel and index for a nonlinear equation? Linearize at every point.}
\GSfield{Statement}{A \(C^1\) map \(F:M\to N\) of Banach manifolds is Fredholm if every \(DF_x\) is Fredholm. Its derivative index is locally constant, hence constant on connected \(M\). This does not define vector-space kernels or cokernels for \(F\) itself.}
\GSfield{Reconstruction}{The Banach-manifold and Fredholm-derivative assumptions are supplied.}
\GSfield{Fate}{\citet[Theorem~1.3 and Section~2]{smale1965sard} gives the Fredholm-map framework and elliptic applications. Received 2 November 1964, the day of this seminar, and published October 1965; a causal connection is not established.}
\GSfield{Present status}{As of November 2026, Fredholm-map theory supplies the precise answer in this framework.}
\GSfield{Assessment}{literal statement: correct with missing hypotheses; publication match: exact; verification: read the result; development: annotated; confidence: high; problem \hyperref[prob:GP-1964-05]{GP-1964-05}}
\end{GSitem}
\begin{GSitem}{1964-15.4}{B}{Stable invariant of a family}
\GSfield{Kind; source}{question; \GSpage{P000035}}
\GSfield{The notes}{For \(T_x\), \(x\in X\), an invariant of the whole family is sought.}
\GSfield{Statement}{For norm-continuous \(T:X\to\operatorname{Fred}(H)\), with compact Hausdorff \(X\), choose finite-dimensional \(V\) with \(T_xH+V=H\). Then \(E_x=T_x^{-1}V\) is a bundle and \([E]-[X\times V]\in K^0(X)\) is independent of \(V\).}
\GSfield{Reconstruction}{Compactness, continuity and stabilization are supplied.}
\GSfield{Fate}{\citet[pp.~129--142]{janich1965fredholm} identifies the Fredholm classifying space in 1965, giving the family invariant. No transmission from this seminar is claimed.}
\GSfield{Present status}{As of November 2026, the index bundle answers the stable question; its rank is the numerical index.}
\GSfield{Assessment}{literal statement: correct with missing hypotheses; publication match: exact; verification: publisher metadata checked; development: annotated; confidence: high; problem \hyperref[prob:GP-1964-06]{GP-1964-06}}
\end{GSitem}

\GStalk{16}{1964-11-02}{A. Ya. Povzner}{P000035-P000038}{Statistical-mechanics programme and particle-energy fragment}
P000037 proposes Ruelle, Yang--Lee and Kac readings, then sketches a pair-potential minimization claim. The handwritten given-name attribution conflicts with the catalogue; neither a uniform confinement theorem nor its assumptions can be recovered.

\begin{GSitem}{1964-16.1}{E}{Statistical-mechanics reading programme}
\GSfield{Kind; source}{remark; \GSpage{P000037}}
\GSfield{The notes}{Two or three Ruelle works, one Yang--Lee work, and Kac models are proposed.}
\GSfield{Statement}{A reading programme in rigorous equilibrium statistical mechanics, rather than a theorem.}
\GSfield{Fate}{Ruelle author metadata confirms \citet[pp.~183--197]{ruelle1963classical} and adjacent 1963--64 papers as plausible reading. The unspecified list cannot be uniquely reconstructed.}
\GSfield{Present status}{As of November 2026, the subject is identifiable but the exact proposed bibliography is not.}
\GSfield{Assessment}{literal statement: correct; publication match: thematic only; verification: publisher metadata checked; development: annotated; confidence: medium}
\end{GSitem}
\begin{GSitem}{1964-16.2}{F}{Confinement of pair-potential minimizers}
\GSfield{Kind; source}{claim; \GSpage{P000037}}
\GSfield{The notes}{A repulsive-core/attractive-well graph accompanies a double interaction sum; a minimizing configuration is said to lie in a sphere of radius \(C\). Nearby powers of \(N\) are unclear.}
\GSfield{Statement}{No unique mathematical statement can be extracted: potential hypotheses, treatment of self-interactions, existence, dimension and dependence of \(C\) on particle number are absent.}
\GSfield{Fate}{Povzner/Ruelle, ground-state, confinement and pair-potential searches did not locate this exact statement. Thermodynamic stability alone does not supply it.}
\GSfield{Present status}{As of November 2026, the fragment remains unidentifiable from the supplied evidence; this does not establish that a mathematical problem is open.}
\GSfield{Assessment}{literal statement: not determinable; publication match: none found; verification: unresolved lead; development: annotated; confidence: low}
\end{GSitem}

\GStalk{17}{1964-11-16}{I. M. Gelfand}{P000036-P000038}{Invariant-subspace announcement and Schur positivity}
The enthusiastic invariant-subspace announcement on P000037 is followed by Schur positivity. The first records a subsequently retracted claim, not a solution.

\begin{GSitem}{1964-17.1}{C}{Retracted invariant-subspace announcement}
\GSfield{Kind; source}{claim; \GSpage{P000037}}
\GSfield{The notes}{\ru{Решено!!!} accompanies a reference to a 1964 Bulletin AMS de Branges paper and a second author, with Livshits mentioned.}
\GSfield{Statement}{The intended question is whether every bounded operator on a separable infinite-dimensional complex Hilbert space has a closed invariant subspace other than \(0,H\). The de Branges--Rovnyak affirmative announcement was retracted.}
\GSfield{Reconstruction}{Boundedness, complex scalars, separability and closed nontrivial subspace are supplied.}
\GSfield{Fate}{\citet[pp.~718--721]{debrangesrovnyak1964existence} was retracted in \citet[p.~396]{debrangesrovnyak1965correction}. Both publication records and the retraction were checked in the authors’ institutional bibliographies; the original proof and correction were not read. The later preprint \citet[Introduction, p.~2]{marcoux2026generators}, first posted 21 August 2026, explicitly still calls the exact Hilbert-space problem unsolved.}
\GSfield{Present status}{As of 12 November 2026, no generally accepted solution was verified in the literature checked. The unrestricted Hilbert-space problem remains distinct from its solved special classes and unverified claims.}
\GSfield{Assessment}{literal statement: false as written; publication match: exact; verification: indirect; development: annotated; confidence: high; problem \hyperref[prob:GP-1964-07]{GP-1964-07}}
\GSfield{Worth developing}{Worth investigating in precise special classes; the unrestricted problem exceeds a bounded development.}
\end{GSitem}
\begin{GSitem}{1964-17.2}{E}{Schur positivity of generalized determinants}
\GSfield{Kind; source}{theorem; \GSpage{P000037}}
\GSfield{The notes}{Generalized determinants of a positive matrix are positive, by Schur.}
\GSfield{Statement}{For Hermitian positive semidefinite \(A=(a_{ij})\) and an irreducible character \(\chi\) of \(S_n\), the immanant \(\sum_{\sigma}\chi(\sigma)\prod_i a_{i,\sigma(i)}\) is nonnegative; it is positive for positive definite \(A\).}
\GSfield{Reconstruction}{Hermitian positivity, irreducible characters and nonnegative in the semidefinite case are supplied.}
\GSfield{Fate}{The antecedent is \citet[pp.~184--207]{schur1918formen}, published June 1918. The particular generalized determinant is not defined in the note.}
\GSfield{Present status}{As of November 2026, the character-immanant reading is established; entrywise positivity alone is insufficient.}
\GSfield{Assessment}{literal statement: correct with missing hypotheses; publication match: partial; verification: publisher metadata checked; development: annotated; confidence: medium}
\end{GSitem}

\GStalk{18}{1964-11-16}{V. I. Arnold}{P000036-P000039}{Semiclassical transport and caustic phase}
The action/Jacobian description is introduced before the caustic phase problem. This foreshadows the 1967 characteristic-class paper; fuller formulas appear on 23 November.

\begin{GSitem}{1964-18.1}{E}{WKB transport before a caustic}
\GSfield{Kind; source}{method; \GSpage{P000037-P000038}}
\GSfield{The notes}{The wave has an action phase and square-root Jacobian amplitude.}
\GSfield{Statement}{For \(ih\psi_t=(-h^2\Delta/2+V)\psi\), \(\psi=ae^{iS/h}\) gives \(S_t+|\nabla S|^2/2+V=0\) and \(a_t+\nabla S\cdot\nabla a+(\Delta S)a/2=0\) at leading orders. The amplitude is a transported half-density.}
\GSfield{Reconstruction}{A consistent normalization and formal leading-order interpretation are supplied.}
\GSfield{Fate}{The local equations follow by direct substitution of the ansatz. The Introduction of \citet{arnold1967characteristic} supplies the later Maslov-asymptotic context, not a proof of these transport formulas.}
\GSfield{Present status}{As of November 2026, the formula applies where the characteristic projection is nonsingular.}
\GSfield{Assessment}{literal statement: correct with missing hypotheses; publication match: partial; verification: read the result; development: annotated; confidence: high}
\end{GSitem}
\begin{GSitem}{1964-18.2}{A}{Characteristic class of caustic phase}
\GSfield{Kind; source}{claim; \GSpage{P000038}}
\GSfield{The notes}{Caustic crossings carry a topological phase correction.}
\GSfield{Statement}{A smooth Lagrangian Gauss map pulls back the integral degree-one Maslov class of the Lagrangian Grassmannian. In a standard convention its loop phase is \(\exp(-i\pi\mu/2)\).}
\GSfield{Reconstruction}{The Grassmannian construction is the later precise identification.}
\GSfield{Fate}{\citet[Sections~1--3]{arnold1967characteristic}, received 28 September 1966 and published 1967, gives the characteristic class. These notes witness November 1964 oral circulation.}
\GSfield{Present status}{As of November 2026, the Maslov-class formulation is established, with convention-dependent phase sign.}
\GSfield{Assessment}{literal statement: correct with missing hypotheses; publication match: exact; verification: read the result; development: annotated; confidence: high}
\end{GSitem}

\GStalk{19}{1964-11-16}{Yu. V. Linnik}{P000037-P000041}{Tests, nuisance parameters and analytic ideals}
The record passes from sufficiency and randomized tests to similarity, nuisance variances, Bayes-weighted cotests and analytic ideal methods. The closely matching paper was submitted in December 1964.

\begin{GSitem}{1964-19.1}{E}{Experiments and sufficient statistics}
\GSfield{Kind; source}{definition; \GSpage{P000038}}
\GSfield{The notes}{Distribution families and statistics are introduced, with a Gaussian sample-mean example.}
\GSfield{Statement}{In a dominated experiment \(\{P_\theta\}\), sufficiency means that the conditional law of the sample given a statistic is independent of \(\theta\). For independent \(N(\theta,\sigma^2)\) observations with known positive variance, the sample mean is sufficient.}
\GSfield{Reconstruction}{Domination, independence and known variance are supplied.}
\GSfield{Fate}{These are classical preliminaries used in \citet[p.~520]{linnik1965cotest}.}
\GSfield{Present status}{As of November 2026, the Gaussian example needs known variance; the mean alone is insufficient when variance is also unknown.}
\GSfield{Assessment}{literal statement: correct with missing hypotheses; publication match: partial; verification: read the result; development: annotated; confidence: high}
\end{GSitem}
\begin{GSitem}{1964-19.2}{E}{Randomized and similar tests}
\GSfield{Kind; source}{definition; \GSpage{P000038-P000039}}
\GSfield{The notes}{The source defines randomized tests, then explicitly says \ru{Тесты подобные}: the null rejection probability is \(\alpha\) at every null parameter. A separate later line abbreviates invariant tests.}
\GSfield{Statement}{A randomized test is measurable \(\phi:\Omega\to[0,1]\). Level \(\alpha\) requires \(\sup_{\Theta_0}E_\theta\phi\leq\alpha\); similarity requires equality for every null parameter. Group invariance is a different condition.}
\GSfield{Reconstruction}{The terminology is corrected, not an invariance theorem asserted.}
\GSfield{Fate}{\citet[p.~520]{linnik1965cotest} uses the same similarity definition. The P000039 scan and both supplied transcriptions already distinguish the definition from the later abbreviated invariance remark.}
\GSfield{Present status}{As of November 2026, level, similarity and group invariance must be distinguished.}
\GSfield{Assessment}{literal statement: correct; publication match: exact; verification: read the result; development: annotated; confidence: high}
\end{GSitem}
\begin{GSitem}{1964-19.3}{E}{Normal means with unknown unequal variances}
\GSfield{Kind; source}{problem; \GSpage{P000039}}
\GSfield{The notes}{Comparison of two Gaussian means is the nuisance-parameter example.}
\GSfield{Statement}{Given independent samples \(X_i\sim N(\mu_1,\sigma_1^2)\), \(Y_j\sim N(\mu_2,\sigma_2^2)\), with unknown positive variances, test \(\mu_1=\mu_2\). Optimality additionally requires a level, sidedness, admissible test class and alternative power criterion.}
\GSfield{Reconstruction}{The model is normalized; no universal optimum is inferred.}
\GSfield{Fate}{This is the classical Behrens--Fisher setting. \citet[pp.~520--522]{linnik1965cotest} supplies a general weighted optimization framework, not an unrestricted uniformly most powerful solution.}
\GSfield{Present status}{As of November 2026, the model is definite, but no single remaining optimality conjecture is specified.}
\GSfield{Assessment}{literal statement: correct with missing hypotheses; publication match: thematic only; verification: read the result; development: annotated; confidence: medium}
\end{GSitem}
\begin{GSitem}{1964-19.4}{A}{Bayes-weighted cotest optimization}
\GSfield{Kind; source}{method; \GSpage{P000039-P000040}}
\GSfield{The notes}{An alternative measure and a cotest ideal encode the testing constraints.}
\GSfield{Statement}{For a similar test set \(l=\phi-\alpha\). Then \(E_\theta l=0\) on \(\Theta_0\), \(-\alpha\leq l\leq1-\alpha\). A fixed alternative measure \(B\) gives objective \(\int E_\theta l\,dB(\theta)\). Under analytic-transform hypotheses, null-zero constraints become ideal membership.}
\GSfield{Reconstruction}{The bounds and fixed alternative weighting make the problem precise.}
\GSfield{Fate}{\citet[pp.~520--522]{linnik1965cotest} gives this construction. Received 7 December 1964, published 1965, following the seminar by three weeks.}
\GSfield{Present status}{As of November 2026, the framework is identifiable; its conclusions depend on the transform range and weighting.}
\GSfield{Assessment}{literal statement: correct with missing hypotheses; publication match: exact; verification: read the result; development: annotated; confidence: high}
\end{GSitem}
\begin{GSitem}{1964-19.5}{A}{Analytic ideals in statistical testing}
\GSfield{Kind; source}{method; \GSpage{P000040}}
\GSfield{The notes}{Cartan, Oka and division organize cotest constraints.}
\GSfield{Statement}{A coherent null-constraint ideal on a Stein parameter domain admits global patching of local division by Cartan B. This does not give boundary-continuous division or guarantee that every holomorphic function is a statistical transform.}
\GSfield{Reconstruction}{The Stein/coherent formulation explains the method; not all these hypotheses appear in the note.}
\GSfield{Fate}{\citet[p.~521]{linnik1965cotest} applies analytic-ideal generation under specified conditions on the null set; the application is new, the sheaf machinery older.}
\GSfield{Present status}{As of November 2026, the analytic reduction is established with those hypotheses.}
\GSfield{Assessment}{literal statement: correct with missing hypotheses; publication match: partial; verification: read the result; development: annotated; confidence: high}
\end{GSitem}
\begin{GSitem}{1964-19.6}{F}{Unspecified primary ideals on a ball}
\GSfield{Kind; source}{question; \GSpage{P000040}}
\GSfield{The notes}{Primary ideals for functions on a ball are mentioned.}
\GSfield{Statement}{No unique mathematical statement can be extracted: the algebra could be holomorphic, bounded holomorphic or boundary-continuous, and no support or ideal is specified.}
\GSfield{Fate}{Linnik/Palamodov/analytic-division searches identify the surrounding programme, but no unique problem; Gorin next states a precise maximal-ideal question.}
\GSfield{Present status}{As of November 2026, the fragment remains unidentifiable from the supplied evidence; this does not establish that a mathematical problem is open.}
\GSfield{Assessment}{literal statement: not determinable; publication match: none found; verification: unresolved lead; development: annotated; confidence: low}
\end{GSitem}

\GStalk{20}{1964-11-16}{E. A. Gorin}{P000039-P000042}{Gleason problem for the ball algebra}
The precise question is coordinate generation of the maximal ideal at the origin, preserving continuity on the closed ball. Comparisons with the disk, polydisc and unrestricted holomorphic division isolate the boundary issue. The ball case is solved.

\begin{GSitem}{1964-20.1}{B}{Coordinate generators of the ball-algebra ideal}
\GSfield{Kind; source}{question; \GSpage{P000041}}
\GSfield{The notes}{If \(f(0)=0\), can \(f=z_1g_1+z_2g_2\), with all functions holomorphic in the ball and continuous on its closure?}
\GSfield{Statement}{For \(A(B^n)=\mathcal O(B^n)\cap C(\overline B^n)\) and \(a\in B^n\), the ideal \(\{f:f(a)=0\}\) is generated by \(z_j-a_j\). The source asks \(n=2,a=0\).}
\GSfield{Reconstruction}{The general \(n,a\) version records the known extension.}
\GSfield{Fate}{\citet[pp.~595--596, reference Hen]{backlundfallstrom1998gleason} records Leibenzon’s affirmative solution and Henkin’s 1971 publication. \citet[pp.~125--132]{gleason1964ideals} is the contemporary antecedent. Initial circulation of Leibenzon’s proof was not verified.}
\GSfield{Present status}{As of November 2026, the exact ball question is solved; later counterexamples concern other domains.}
\GSfield{Assessment}{literal statement: correct with missing hypotheses; publication match: exact; verification: indirect; development: annotated; confidence: high; problem \hyperref[prob:GP-1964-08]{GP-1964-08}}
\end{GSitem}
\begin{GSitem}{1964-20.2}{E}{Disk and polydisc division}
\GSfield{Kind; source}{construction; \GSpage{P000041}}
\GSfield{The notes}{The one-variable case is immediate and the product domain is easier.}
\GSfield{Statement}{For \(f\in A(\mathbb D)\), \(f(0)=0\), division by \(z\) has a removable singularity and remains boundary-continuous. On \(\mathbb D^n\), successive coordinate telescoping and division produce \(f=\sum z_jg_j\) for \(f(0)=0\).}
\GSfield{Reconstruction}{The telescoping proof is supplied, not a boundary power-series argument.}
\GSfield{Fate}{The one-variable assertion is also recalled by \citet[p.~596]{backlundfallstrom1998gleason}; the telescoping explanation is elementary reconstruction.}
\GSfield{Present status}{As of November 2026, the division holds without assuming uniform boundary convergence of Taylor series.}
\GSfield{Assessment}{literal statement: correct with missing hypotheses; publication match: partial; verification: read the result; development: annotated; confidence: high}
\end{GSitem}
\begin{GSitem}{1964-20.3}{E}{Division without boundary restrictions}
\GSfield{Kind; source}{remark; \GSpage{P000041}}
\GSfield{The notes}{The difficulty disappears when boundary continuity is dropped.}
\GSfield{Statement}{For Stein \(\Omega\subset\mathbb C^n\) and \(a\in\Omega\), the coherent ideal sheaf of \(a\) is generated by \(z_j-a_j\). Cartan B applied to its relation sheaf yields \(f=\sum(z_j-a_j)g_j\), \(g_j\in\mathcal O(\Omega)\), when \(f(a)=0\).}
\GSfield{Reconstruction}{The sheaf proof and Stein hypothesis are editorial.}
\GSfield{Fate}{The holomorphic-versus-Banach-algebra distinction is explicit in \citet[pp.~595--596]{backlundfallstrom1998gleason}.}
\GSfield{Present status}{As of November 2026, the holomorphic statement gives no boundary continuity estimate.}
\GSfield{Assessment}{literal statement: correct with missing hypotheses; publication match: partial; verification: indirect; development: annotated; confidence: high}
\end{GSitem}

\GStalk{21}{1964-11-23}{V. I. Arnold}{P000040-P000046}{WKB geometry, caustics and Morse index}
Action and Jacobian transport are developed through caustics, then related to optical phase and the Morse index. The notebook has inconsistent Schrödinger normalizations; modern formulas are explicitly reconstructed.

\begin{GSitem}{1964-21.1}{E}{Hamilton--Jacobi equation}
\GSfield{Kind; source}{method; \GSpage{P000042-P000043}}
\GSfield{The notes}{An exponential substitution relates Schrödinger evolution to classical trajectories.}
\GSfield{Statement}{For \(ih\psi_t=(-h^2\Delta/2+V)\psi\), \(\psi=e^{iS/h}(a_0+ha_1+\cdots)\) yields \(S_t+H(x,\nabla S)=0\), \(H=|p|^2/2+V\), formally wherever \(S\) is smooth.}
\GSfield{Reconstruction}{Signs, kinetic factor and the formal interpretation are normalized.}
\GSfield{Fate}{This local background is used in \citet[Introduction]{arnold1967characteristic}.}
\GSfield{Present status}{As of November 2026, the formal relation is standard; the notebook coefficients cannot all hold simultaneously.}
\GSfield{Assessment}{literal statement: false as written; publication match: partial; verification: read the result; development: annotated; confidence: high}
\end{GSitem}
\begin{GSitem}{1964-21.2}{E}{Action and half-density}
\GSfield{Kind; source}{construction; \GSpage{P000042-P000043}}
\GSfield{The notes}{The amplitude contains a square-root trajectory Jacobian and the phase \(\int p\,dx\).}
\GSfield{Statement}{For nonsingular \(x=X_t(y)\), \(a_0(t,X_t(y))=a_0(0,y)|\det D_yX_t|^{-1/2}\). The action evolves by \(S(t,X_t(y))=S(0,y)+\int_0^t(p\cdot\dot x-H)\,ds\).}
\GSfield{Reconstruction}{Absolute Jacobian and \(-Hdt\) are supplied.}
\GSfield{Fate}{The geometric setting is \citet[Introduction]{arnold1967characteristic}; these consistent formulas are an editorial reconstruction.}
\GSfield{Present status}{As of November 2026, the formula requires a nonsingular projection; abbreviated spatial action needs a fixed-energy or time convention.}
\GSfield{Assessment}{literal statement: correct with missing hypotheses; publication match: partial; verification: read the result; development: annotated; confidence: high}
\end{GSitem}
\begin{GSitem}{1964-21.3}{A}{Maslov correction at caustics}
\GSfield{Kind; source}{claim; \GSpage{P000044}}
\GSfield{The notes}{A branch sum is corrected by a phase involving \(\pi n_j/2\), associated with Maslov.}
\GSfield{Statement}{For regular Lagrangian propagation, oscillatory-integral chart transitions yield Maslov phase factors; the integral Maslov class governs loop phase, conventionally \(e^{-i\pi\mu/2}\).}
\GSfield{Reconstruction}{The phase sign is conventional; the damaged formula does not determine it.}
\GSfield{Fate}{\citet[Sections~1--3]{arnold1967characteristic} gives the topological formulation, received September 1966 and published 1967, after these lectures.}
\GSfield{Present status}{As of November 2026, the phase class is established; uniform local approximation needs a caustic normal form.}
\GSfield{Assessment}{literal statement: correct with missing hypotheses; publication match: exact; verification: read the result; development: annotated; confidence: high}
\end{GSitem}
\begin{GSitem}{1964-21.4}{E}{Optical quarter-cycle phase}
\GSfield{Kind; source}{remark; \GSpage{P000044}}
\GSfield{The notes}{A ray tangent to an envelope acquires a phase jump \(\pi/2\).}
\GSfield{Statement}{At a generic fold caustic, an Airy model matches oscillatory branches with a relative quarter-cycle phase under standard conventions. Degenerate caustics require different models.}
\GSfield{Reconstruction}{Genericity and phase convention supplied.}
\GSfield{Fate}{This optical motivation precedes the seminar and informs \citet[Introduction]{arnold1967characteristic}.}
\GSfield{Present status}{As of November 2026, the fold qualification is necessary.}
\GSfield{Assessment}{literal statement: correct with missing hypotheses; publication match: partial; verification: read the result; development: annotated; confidence: high}
\end{GSitem}
\begin{GSitem}{1964-21.5}{E}{Morse index and conjugate points}
\GSfield{Kind; source}{theorem; \GSpage{P000044}}
\GSfield{The notes}{The phase index is related to negative squares of the second variation and conjugate points.}
\GSfield{Statement}{For a regular fixed-endpoint mechanical extremal with positive kinetic form and nonconjugate endpoints, the Morse index equals the count of interior conjugate points with multiplicities.}
\GSfield{Reconstruction}{Regularity, endpoints and multiplicity conventions supplied.}
\GSfield{Fate}{The classical variational theorem is background to \citet[Sections~1--3]{arnold1967characteristic}.}
\GSfield{Present status}{As of November 2026, the relation holds in that variational setting; endpoint nullity requires an extra convention.}
\GSfield{Assessment}{literal statement: correct with missing hypotheses; publication match: partial; verification: indirect; development: annotated; confidence: high}
\end{GSitem}

\GStalk{22}{1964-11-23}{I. M. Gelfand}{P000044-P000047}{Phase volume and the fermionic proof of Sturm oscillation}
Berezin’s eigenvalue-counting remark precedes Gelfand’s one-dimensional zeros and determinant argument. The latter uses the fermionic ground state; a full implementation appeared in 2018/2020. The notebook’s sign and indexing require correction.

\begin{GSitem}{1964-22.1}{E}{Phase-volume eigenvalue asymptotics}
\GSfield{Kind; source}{claim; \GSpage{P000045}; F. A. Berezin}
\GSfield{The notes}{Berezin relates a region under the classical Hamiltonian to quantum eigenvalue counting; the written threshold mixes energy and index.}
\GSfield{Statement}{Let \(d\ge1\), \(V\in C^\infty(\mathbb R^d;\mathbb R)\), and \(V(x)\to+\infty\) as \(|x|\to\infty\). Let \(P_h=-h^2\Delta/2+V\) be its self-adjoint Schr\"odinger realization on \(L^2(\mathbb R^d)\), and let \(N_h(E)\) count eigenvalues below \(E\), with multiplicity. For \(E>\min V\) a regular value of \(|p|^2/2+V(x)\), \[ N_h(E)\sim(2\pi h)^{-d} \operatorname{vol}\{(x,p):|p|^2/2+V(x)<E\}\quad(h\downarrow0). \] The phase volume is finite and positive.}
\GSfield{Reconstruction}{Limit, normalization, dimension and confinement are supplied.}
\GSfield{Fate}{The historical eigenvalue-counting antecedent is \citet[pp.~110--117]{weyl1911eigenvalues}. The stated semiclassical Schrödinger version is a modern reconstruction; neither the 1911 paper’s original Laplace problem nor the notebook’s compressed wording is identical to this full formulation.}
\GSfield{Present status}{As of November 2026, the reconstructed Weyl law is established under those hypotheses.}
\GSfield{Assessment}{literal statement: false as written; publication match: partial; verification: publisher metadata checked; development: annotated; confidence: high}
\end{GSitem}
\begin{GSitem}{1964-22.2}{E}{Sturm zeros and interlacing}
\GSfield{Kind; source}{theorem; \GSpage{P000045}}
\GSfield{The notes}{Dirichlet eigenfunctions on \([0,\pi]\) are discussed; the free eigenvalues have a negative sign in the note.}
\GSfield{Statement}{For real continuous \(q\), the Dirichlet operator \(L=-d^2/dx^2+q\) has simple ordered eigenvalues. Its \(j\)-th eigenfunction has \(j-1\) interior zeros interlacing the next. For \(q=0\), \(\lambda_j=j^2\).}
\GSfield{Reconstruction}{Indexing starts at 1 throughout.}
\GSfield{Fate}{\citet[historical discussion and Sturm theorem]{berardhelffer2020sturm} revisits the 1836 result; preprint 2017, publication 2020.}
\GSfield{Present status}{As of November 2026, this is classical; the notebook sign is erroneous for its displayed operator.}
\GSfield{Assessment}{literal statement: false as written; publication match: exact; verification: read the result; development: annotated; confidence: high}
\end{GSitem}
\begin{GSitem}{1964-22.3}{E}{Zeros of the first eigenfunction span}
\GSfield{Kind; source}{theorem; \GSpage{P000045}}
\GSfield{The notes}{A linear combination of the first eigenfunctions has a bounded number of zeros; indexing changes within the note.}
\GSfield{Statement}{For a regular real Dirichlet Sturm--Liouville problem with smooth potential, a nonzero combination of the first \(m\) eigenfunctions has at most \(m-1\) interior zeros, counted with multiplicity.}
\GSfield{Reconstruction}{The indexing and interior multiplicity convention resolve the ambiguity.}
\GSfield{Fate}{\citet[main theorem and historical discussion]{berardhelffer2020sturm} identifies Sturm’s 1836 theorem; the preprint circulated June 2017 before 2020 publication.}
\GSfield{Present status}{As of November 2026, the bound is established; the seminar’s distinctive feature is its proof strategy.}
\GSfield{Assessment}{literal statement: correct with missing hypotheses; publication match: exact; verification: read the result; development: annotated; confidence: high}
\end{GSitem}
\begin{GSitem}{1964-22.4}{B}{Gelfand’s Slater-determinant proof strategy}
\GSfield{Kind; source}{method; \GSpage{P000045}}
\GSfield{The notes}{A determinant of eigenfunction data is related to a many-particle problem; the displayed entries mix values and derivatives.}
\GSfield{Statement}{\(\det[u_j(x_i)]_{i,j=1}^m\) is an antisymmetric eigenfunction of \(\sum_i L_{x_i}\). On \(x_1<\cdots<x_m\), it is a ground state of fixed sign, so the evaluation determinant is nonzero for distinct interior points. Multiplicity of zeros requires an additional argument.}
\GSfield{Reconstruction}{Ordered chamber, ground-state argument and the multiplicity distinction are supplied.}
\GSfield{Fate}{\citet[Sections~3--6]{berardhelffer2020gelfand} explicitly implements Gelfand’s strategy: first preprint 11 July 2018, journal 2020.}
\GSfield{Present status}{As of November 2026, the strategy has a rigorous published realization; the notebook determinant alone is incomplete.}
\GSfield{Assessment}{literal statement: correct with missing hypotheses; publication match: exact; verification: read the result; development: annotated; confidence: high}
\end{GSitem}

\GStalk{23}{}{}{P000045-P000047}{February heading and programme}
A heading/plan with names and future topics, adjoining 8 February 1965. No separately identifiable mathematical assertion occurs. The contradictory catalogue date label is preserved as an undated plan.


\GStalk{24}{1965-02-08}{I. I. Pyateski-Shapiro}{P000045-P000048}{Congruence subgroups and elementary matrices}
The record contrasts the congruence-subgroup property in ranks at least two with the modular group. The main integer-matrix result was already available in 1964. The blanket assertion about rank-one power quotients needs a low-level exception and an orbifold qualification.

\begin{GSitem}{1964-24.1}{E}{Congruence property for integer special linear groups}
\GSfield{Kind; source}{theorem; \GSpage{P000046-P000047}}
\GSfield{The notes}{For \(n\geq3\), every finite-index subgroup of \(SL(n,\mathbb Z)\) contains a congruence subgroup; \(n=2\) differs.}
\GSfield{Statement}{For \(n\geq3\) and finite-index \(\Gamma\leq SL_n(\mathbb Z)\), some \(m\geq1\) satisfies \(\ker(SL_n(\mathbb Z)\to SL_n(\mathbb Z/m\mathbb Z))\subset\Gamma\). For \(n=2\), finite-index noncongruence subgroups exist.}
\GSfield{Reconstruction}{The subgroup need not be normal; passage to its finite-index core is allowed.}
\GSfield{Fate}{\citet[Sections~1.2--1.3, Theorem~1 and historical remark]{serre1968congruence} attributes the rational-integer case to Bass--Lazard--Serre and Mennicke in 1964; the Bourbaki lecture was June 1967, printed 1968.}
\GSfield{Present status}{As of November 2026, the stated integer-matrix theorem is settled, not a February 1965 open question.}
\GSfield{Assessment}{literal statement: correct with missing hypotheses; publication match: exact; verification: read the result; development: annotated; confidence: high}
\end{GSitem}
\begin{GSitem}{1964-24.2}{E}{Reduction to powers of elementary transvections}
\GSfield{Kind; source}{method; \GSpage{P000047}}
\GSfield{The notes}{The normal subgroup generated by powers of elementary matrices is used to reduce the finite-index question.}
\GSfield{Statement}{Given a finite-index subgroup, its normal core contains \(e_{ij}(m)\) for some common positive \(m\), since each elementary generator has finite-order image in the finite quotient. For \(n\geq3\), the congruence theorem then controls the normal elementary subgroup.}
\GSfield{Reconstruction}{Normal core and common multiple of finite orders are supplied.}
\GSfield{Fate}{This is the elementary-matrix route in the congruence programme summarized by \citet[Sections~2--3]{serre1968congruence}.}
\GSfield{Present status}{As of November 2026, the finite-quotient reduction is valid; identifying elementary and principal congruence subgroups at the same level is not automatic.}
\GSfield{Assessment}{literal statement: correct with missing hypotheses; publication match: partial; verification: read the result; development: annotated; confidence: high}
\end{GSitem}
\begin{GSitem}{1964-24.3}{E}{Rank-one power quotients and triangle groups}
\GSfield{Kind; source}{claim; \GSpage{P000047}}
\GSfield{The notes}{For \(n=2\), powers of generators are said to give infinite index, with a quotient interpreted as a surface fundamental group.}
\GSfield{Statement}{In \(PSL_2(\mathbb Z)=\langle s,t:s^2=(st)^3=1\rangle\), adding \(t^m=1\) gives the triangle group \(\Delta(2,3,m)\). For \(2\leq m\leq5\) it is finite, for \(m=6\) Euclidean, and for \(m\geq7\) hyperbolic. It is an orbifold group, not generally a torsion-free surface group.}
\GSfield{Reconstruction}{Passage from SL to PSL and the triangle-group presentation are essential.}
\GSfield{Fate}{The modular-group contrast is classical; \citet[Section~1.2(b)]{serre1968congruence} confirms the failure of the rank-one congruence property. The explicit presentation and low-level check here are editorial algebra.}
\GSfield{Present status}{As of November 2026, the unqualified infinite-index and surface-group assertions are false; a torsion-free finite-index subgroup restores the surface interpretation in the hyperbolic case.}
\GSfield{Assessment}{literal statement: false as written; publication match: partial; verification: read the result; development: annotated; confidence: high}
\end{GSitem}
\begin{GSitem}{1964-24.4}{F}{Unspecified elementary congruence transformation}
\GSfield{Kind; source}{remark; \GSpage{P000047}}
\GSfield{The notes}{An \(a\equiv b\pmod N\) statement is linked to elementary transformations.}
\GSfield{Statement}{No unique mathematical statement can be extracted: the entries, matrix operation, modulus conventions and desired normal form are not defined.}
\GSfield{Fate}{Congruence-subgroup sources explain the surrounding method but do not identify this isolated transformation.}
\GSfield{Present status}{As of November 2026, the fragment remains unidentifiable from the supplied evidence; this does not establish that a mathematical problem is open.}
\GSfield{Assessment}{literal statement: not determinable; publication match: none found; verification: unresolved lead; development: annotated; confidence: low}
\end{GSitem}

\GStalk{25}{1965-02-08}{I. M. Gelfand}{P000046-P000049}{Spin-flavour groups and relativistic symmetry}
An SU(6) discussion is juxtaposed with the Hermitian-matrix model of Minkowski space and proposed larger symmetry groups. The matrix display needs a semidirect-product correction. Coleman--Mandula later answered a restricted precise version, not every unification programme.

\begin{GSitem}{1964-25.1}{E}{SU(6) spin-flavour symmetry}
\GSfield{Kind; source}{construction; \GSpage{P000047}}
\GSfield{The notes}{Spin SU(2) and internal SU(3) are combined in SU(6).}
\GSfield{Statement}{The tensor product on \(\mathbb C^2\otimes\mathbb C^3\) embeds \(SU(2)\times SU(3)\) into \(SU(6)\). This organizes a nonrelativistic approximate spin-flavour model; it is not by itself a Lorentz-covariant exact scattering symmetry.}
\GSfield{Reconstruction}{The approximation status and tensor-product interpretation are supplied.}
\GSfield{Fate}{\citet[pp.~173--175]{gurseyradicati1964spin} appeared 3 August 1964; \citet[pp.~B1756--B1760]{sakita1964supermultiplets} was received 7 August and published 21 December 1964. Both precede this February record.}
\GSfield{Present status}{As of November 2026, the historical SU(6) construction is identified; exact relativistic symmetry needs additional assumptions.}
\GSfield{Assessment}{literal statement: correct with missing hypotheses; publication match: exact; verification: publisher metadata checked; development: annotated; confidence: high}
\end{GSitem}
\begin{GSitem}{1964-25.2}{E}{Hermitian matrices and the Poincaré group}
\GSfield{Kind; source}{construction; \GSpage{P000047-P000048}}
\GSfield{The notes}{Lorentz transformations act on Hermitian matrices; a block matrix with diagonal \(A,(A^*)^{-1}\) and Hermitian upper-right block is written.}
\GSfield{Statement}{For \(A\in SL_2(\mathbb C)\), \(X=X^*\), the map \(X\mapsto AXA^*\) preserves \(\det X=x_0^2-|\mathbf x|^2\). The covering Poincaré group has law \((A,X)(B,Y)=(AB,X+AYA^*)\), represented by \(\begin{pmatrix}A&X(A^*)^{-1}\\0&(A^*)^{-1}\end{pmatrix}\).}
\GSfield{Reconstruction}{The semidirect product and corrected block are editorial. Setting translations to zero gives the Lorentz cover; spatial rotations require \(A\in SU(2)\).}
\GSfield{Fate}{This is the standard spacetime group underlying \citet[p.~1251]{colemanmandula1967symmetries}; the block-law repair is direct matrix algebra.}
\GSfield{Present status}{As of November 2026, the corrected matrices are closed under multiplication. An arbitrary Hermitian upper-right block in the notebook parametrization is not closed.}
\GSfield{Assessment}{literal statement: false as written; publication match: thematic only; verification: read the result; development: annotated; confidence: high}
\end{GSitem}
\begin{GSitem}{1964-25.3}{B}{Combining spacetime and internal symmetries}
\GSfield{Kind; source}{question; \GSpage{P000048}}
\GSfield{The notes}{A larger symmetry containing relativistic and internal transformations is sought.}
\GSfield{Statement}{For a nontrivial analytic scattering matrix in four-dimensional relativistic theory, suppose a connected ordinary Lie symmetry contains the Poincaré group, acts tensorially on multiparticle states, has only finitely many particle types below each mass bound, and satisfies the regularity and nontrivial-scattering hypotheses of Coleman--Mandula. Then it is locally a direct product of Poincaré and internal symmetries.}
\GSfield{Reconstruction}{The precise scattering question is a later formulation; no blanket impossibility theorem is attributed to the notebook.}
\GSfield{Fate}{\citet[p.~1251, theorem and hypotheses~1--5]{colemanmandula1967symmetries}, received March and published July 1967, gives this restricted obstruction. The original broad programme is not equivalent to the theorem’s hypotheses.}
\GSfield{Present status}{As of November 2026, the ordinary-Lie scattering version is settled; graded symmetries and violations of the hypotheses lie outside that conclusion.}
\GSfield{Assessment}{literal statement: correct with missing hypotheses; publication match: partial; verification: read the result; development: annotated; confidence: high}
\end{GSitem}
\begin{GSitem}{1964-25.4}{F}{108- and 144-parameter proposals}
\GSfield{Kind; source}{remark; \GSpage{P000048}}
\GSfield{The notes}{Groups with 108 and 144 parameters are proposed in the discussion.}
\GSfield{Statement}{No unique mathematical statement can be extracted: no Lie brackets, real form, faithful action or particle interpretation identifies the groups uniquely.}
\GSfield{Fate}{SU(6), relativistic supermultiplet and Coleman--Mandula searches contextualize the numbers but do not establish a unique match.}
\GSfield{Present status}{As of November 2026, the fragment remains unidentifiable from the supplied evidence; this does not establish that a mathematical problem is open.}
\GSfield{Assessment}{literal statement: not determinable; publication match: none found; verification: unresolved lead; development: annotated; confidence: low}
\end{GSitem}

\GStalk{26}{1965-02-15}{D. A. Kazhdan}{P000047-P000052}{Primary ideals and differential conditions}
The algebraic discussion extracts first-order conditions from a square-zero thickening, then considers derivations, localization and formal flow. The notebook itself warns that I and J were interchanged. The correct prime/primary inclusions are indispensable.

\begin{GSitem}{1964-26.1}{B}{Primary ideals through differential conditions}
\GSfield{Kind; source}{method; \GSpage{P000048-P000050}}
\GSfield{The notes}{A primary ideal and its prime radical are compared, with a square inclusion and derivations vanishing on the smaller ideal; I and J are later explicitly interchanged.}
\GSfield{Statement}{Let \(A=k[x_1,\ldots,x_n]\), \(\operatorname{char}k=0\), \(P\) prime and \(Q\) \(P\)-primary with \(P^2\subset Q\subset P\). Over \(K=\operatorname{Frac}(A/P)\), linear functionals on \(P A_P/P^2A_P\) annihilating \(QA_P/P^2A_P\), extended as first-order operators, detect membership in \(Q\) together with evaluation modulo \(P\). General primary ideals require higher-order operators.}
\GSfield{Reconstruction}{Prime/primary letters, generic localization, polynomial ring, characteristic zero and operator values are supplied.}
\GSfield{Fate}{The established Noetherian-operator theory is made algebraically and algorithmically explicit by \citet[Theorem~2.1 and Sections~4--6]{cidruizhomssturmfels2021primary}, preprint January 2020, journal 2021. Its introduction traces Palamodov’s correction of constant-coefficient claims.}
\GSfield{Present status}{As of November 2026, the method is established; constant-coefficient operators do not suffice for every primary ideal.}
\GSfield{Assessment}{literal statement: correct with missing hypotheses; publication match: partial; verification: read the result; development: annotated; confidence: high}
\end{GSitem}
\begin{GSitem}{1964-26.2}{E}{Conormal module and derivations}
\GSfield{Kind; source}{construction; \GSpage{P000049}}
\GSfield{The notes}{Derivations are used as linear functionals on an ideal modulo its square.}
\GSfield{Statement}{For a \(k\)-algebra \(A\), ideal \(J\), and \(A/J\)-module \(M\), any \(k\)-derivation \(D:A\to M\) kills \(J^2\) and restricts to an \(A/J\)-linear map \(J/J^2\to M\). Extension of every such functional requires additional smoothness or splitting assumptions.}
\GSfield{Reconstruction}{The module structure and the direction of the assertion are made explicit.}
\GSfield{Fate}{This is standard conormal/differential algebra, used in the modern algebraic framework of \citet[Section~5]{cidruizhomssturmfels2021primary}.}
\GSfield{Present status}{As of November 2026, the restriction statement is valid; surjectivity of the extension map is not automatic for an arbitrary quotient.}
\GSfield{Assessment}{literal statement: correct with missing hypotheses; publication match: partial; verification: read the result; development: annotated; confidence: high}
\end{GSitem}
\begin{GSitem}{1964-26.3}{E}{Lifting across a square-zero ideal}
\GSfield{Kind; source}{construction; \GSpage{P000049-P000050}}
\GSfield{The notes}{A quotient-field embedding is lifted; a square-root equation is corrected by an element of the ideal.}
\GSfield{Statement}{If \(J^2=0\) in a commutative ring and \(x_0^2-y\in J\), with \(2x_0\) a unit, then \(x=x_0-(2x_0)^{-1}(x_0^2-y)\) satisfies \(x^2=y\). Lifting an entire residue field uses separability/formal smoothness, not merely that the quotient is a field.}
\GSfield{Reconstruction}{Square-zero, invertibility and separability conditions are reconstructed.}
\GSfield{Fate}{The displayed square-zero correction is elementary Hensel-style algebra; the general coefficient-field statement needs the extra hypotheses made explicit here.}
\GSfield{Present status}{As of November 2026, the corrected lifting formula holds; a completely arbitrary quotient need not split.}
\GSfield{Assessment}{literal statement: correct with missing hypotheses; publication match: partial; verification: indirect; development: annotated; confidence: high}
\end{GSitem}
\begin{GSitem}{1964-26.4}{E}{Polynomial derivations, localization and formal flow}
\GSfield{Kind; source}{method; \GSpage{P000050}}
\GSfield{The notes}{Operators are written as sums \(P_i\partial_i\), extended to fractions and associated with a formal motion.}
\GSfield{Statement}{A \(k\)-derivation of \(k[x_1,\ldots,x_n]\) is \(\sum D(x_i)\partial_i\); it extends uniquely to a localization by \(D(a/s)=(sD(a)-aD(s))/s^2\). In characteristic zero, \(\exp(tD)\) is an algebra map to \(A[[t]]\); analytic convergence is a separate issue.}
\GSfield{Reconstruction}{Derivation rather than arbitrary differential operator, and formal rather than convergent flow, are supplied.}
\GSfield{Fate}{This is classical differential algebra; the distinction between derivations and higher-order differential operators is explicit in \citet[Section~5]{cidruizhomssturmfels2021primary}.}
\GSfield{Present status}{As of November 2026, these formal identities hold; the notebook formula does not describe every differential operator.}
\GSfield{Assessment}{literal statement: correct with missing hypotheses; publication match: partial; verification: read the result; development: annotated; confidence: high}
\end{GSitem}
\begin{GSitem}{1964-26.5}{F}{The solution should be an ideal}
\GSfield{Kind; source}{remark; \GSpage{P000050}}
\GSfield{The notes}{The discussion concludes that a solution should form an ideal.}
\GSfield{Statement}{No unique mathematical statement can be extracted: the space of solutions, ring action and equations are not specified. Kernels of arbitrary differential operators are not ideals.}
\GSfield{Fate}{Primary-ideal/Noetherian-operator sources clarify a possible intended compatibility condition but do not identify this isolated sentence.}
\GSfield{Present status}{As of November 2026, the fragment remains unidentifiable from the supplied evidence; this does not establish that a mathematical problem is open.}
\GSfield{Assessment}{literal statement: not determinable; publication match: none found; verification: unresolved lead; development: annotated; confidence: low}
\end{GSitem}

\GStalk{27}{1965-02-15}{V. P. Palamodov}{P000050-P000053}{Holomorphic quotients and differential jets}
Palamodov represents a primary-ideal quotient by differential data on its zero variety, with locally realizable compatibility. An appended growth estimate is too incomplete to recover its weighted function space.

\begin{GSitem}{1964-27.1}{A}{Primary-ideal quotient by compatible differential jets}
\GSfield{Kind; source}{theorem; \GSpage{P000051}}
\GSfield{The notes}{For a polynomial primary ideal \(J\), the quotient \(\mathcal H/J\mathcal H\) is represented by restrictions \(d(z,D)F|_N\), with local realizability.}
\GSfield{Statement}{For a coherent ideal \(\mathcal J\) on a Stein domain \(\Omega\), a faithful differential-jet realization of \(\mathcal O_\Omega/\mathcal J\), when chosen, has kernel \(\mathcal J\); compatible locally realizable data globalize by Cartan B. For polynomial primary ideals, suitable Noetherian operators provide such finite-order descriptions with the necessary localization and compatibility conditions.}
\GSfield{Reconstruction}{Steinness, coherent quotient, operator choice and local compatibility are supplied; the 1961 smooth theorem is deliberately kept separate.}
\GSfield{Fate}{\citet[Part~I]{palamodov1970operators} is the speaker’s subsequent systematic treatment (English 1970). \citet[Theorems~1--2]{palamodov1961ideals} is a prior smooth-function analogue, not the identical holomorphic theorem. The exact section and first publication of this formulation were not verified.}
\GSfield{Present status}{As of November 2026, the Noetherian-operator and coherent-quotient framework is established, but arbitrary unrestricted vector data on \(N\) are not the target.}
\GSfield{Assessment}{literal statement: correct with missing hypotheses; publication match: partial; verification: indirect; development: annotated; confidence: medium}
\end{GSitem}
\begin{GSitem}{1964-27.2}{F}{Unspecified growth-controlled division estimate}
\GSfield{Kind; source}{claim; \GSpage{P000051}}
\GSfield{The notes}{An inequality involving the size or growth of the representing function is mentioned after the quotient theorem.}
\GSfield{Statement}{No unique mathematical statement can be extracted: the norm, growth weight, domain geometry and constants are absent.}
\GSfield{Fate}{Palamodov’s 1966 An analytic problem arising in statistics (Doklady 168, 1251--1253; MathNet dan32367 opened) proves a weighted holomorphic division theorem on polydiscs. This is a related outcome, not a verified match to the missing estimate. A later nearby result is \citet[p.~1251, Theorem~1]{palamodov1966statistics}, received 2 August 1965 and published 1966: it treats holomorphic functions of polynomial boundary growth on a polydisc and polynomial-matrix multiplication. Its explicit topology does not identify the unspecified space in the note.}
\GSfield{Present status}{As of November 2026, the fragment remains unidentifiable from the supplied evidence; this does not establish that a mathematical problem is open.}
\GSfield{Assessment}{literal statement: not determinable; publication match: thematic only; verification: unresolved lead; development: annotated; confidence: low}
\end{GSitem}

\GStalk{28}{1965-02-15}{A. S. Dynin}{P000051-P000056}{Elliptic systems, solution sheaves and finiteness}
Local differential operators and symbols lead to solution sheaves and finite cohomology, with fine resolutions and compact functional-analytic comparison. The strongest finiteness assertion needs compatibility hypotheses omitted from the note.

\begin{GSitem}{1964-28.1}{E}{Local linear maps are differential operators}
\GSfield{Kind; source}{theorem; \GSpage{P000052}}
\GSfield{The notes}{A morphism of sheaves of smooth bundle sections is said to be differential.}
\GSfield{Statement}{A linear morphism of sheaves of smooth sections of finite-rank bundles on a smooth manifold is locally a finite-order differential operator (Peetre’s theorem). The order need not be globally bounded on a noncompact base.}
\GSfield{Reconstruction}{Linearity, finite-rank smooth bundles and local order are explicit.}
\GSfield{Fate}{\citet[pp.~211--218]{peetre1959differential}, with \citet[pp.~116--120]{peetre1960correction}, predates the seminar.}
\GSfield{Present status}{As of November 2026, the local finite-order assertion is established; linearity and the sheaf locality condition are necessary.}
\GSfield{Assessment}{literal statement: correct with missing hypotheses; publication match: exact; verification: publisher metadata checked; development: annotated; confidence: high}
\end{GSitem}
\begin{GSitem}{1964-28.2}{E}{Principal symbol and overdetermined ellipticity}
\GSfield{Kind; source}{definition; \GSpage{P000052-P000053}}
\GSfield{The notes}{The principal matrix symbol is placed over a tangent bundle and called elliptic when a monomorphism.}
\GSfield{Statement}{For an order-\(m\) operator \(A:C^\infty(E)\to C^\infty(F)\), \(\sigma_m(A)(x,\xi):E_x\to F_x\) is homogeneous in \(\xi\in T_x^*X\). Injectivity for real \(\xi\ne0\) is overdetermined ellipticity; an isomorphism is determined ellipticity.}
\GSfield{Reconstruction}{Cotangent bundle and real nonzero covectors supplied.}
\GSfield{Fate}{\citet[p.~423]{atiyahsinger1963index} explicitly places the symbol over the cotangent sphere and defines the isomorphism case.}
\GSfield{Present status}{As of November 2026, the modern distinction is standard; tangent and cotangent require a chosen metric to identify.}
\GSfield{Assessment}{literal statement: false as written; publication match: partial; verification: read the result; development: annotated; confidence: high}
\end{GSitem}
\begin{GSitem}{1964-28.3}{F}{Finiteness of every solution-sheaf cohomology group}
\GSfield{Kind; source}{claim; \GSpage{P000053}}
\GSfield{The notes}{For \(\ker A\) on a compact manifold, all cohomology groups are stated to be finite-dimensional.}
\GSfield{Statement}{No unique mathematical statement can be extracted: the stated injective-symbol condition alone does not specify the exact compatibility resolution used in the proof. A sufficient precise theorem assumes a finite elliptic differential complex resolving the solution sheaf; then compactness yields finite-dimensional cohomology.}
\GSfield{Fate}{Elliptic-complex and Cartan finiteness sources identify the intended method but not a theorem in the full unstated generality of this sentence.}
\GSfield{Present status}{As of November 2026, the fragment remains unidentifiable from the supplied evidence; this does not establish that a mathematical problem is open.}
\GSfield{Assessment}{literal statement: not determinable; publication match: none found; verification: unresolved lead; development: annotated; confidence: low}
\end{GSitem}
\begin{GSitem}{1964-28.4}{E}{Two-term fine resolution in the locally solvable case}
\GSfield{Kind; source}{method; \GSpage{P000054}}
\GSfield{The notes}{\(0\to\ker A\to\mathcal E(E_0)\to\mathcal E(E_1)\) is treated as a fine resolution.}
\GSfield{Statement}{If \(A\) is locally surjective as a map of smooth-section sheaves, then \(0\to\ker A\to\mathcal E(E_0)\xrightarrow A\mathcal E(E_1)\to0\) is exact. Therefore \(H^0=\ker A(X)\), \(H^1=\operatorname{coker}A(X)\), and \(H^q=0\) for \(q\ge2\). For a determined elliptic operator on compact boundaryless \(X\), the two displayed global spaces are finite-dimensional.}
\GSfield{Reconstruction}{The terminal zero, local surjectivity and no-boundary condition are supplied.}
\GSfield{Fate}{The resolution principle is \citet[Section~1, Proposition~1]{cartan1954finitude}; elliptic kernel/cokernel finiteness appears in \citet[p.~423]{atiyahsinger1963index}.}
\GSfield{Present status}{As of November 2026, the conclusion is valid with local exactness; omitting it makes the resolution argument incomplete.}
\GSfield{Assessment}{literal statement: correct with missing hypotheses; publication match: exact; verification: read the result; development: annotated; confidence: high}
\end{GSitem}
\begin{GSitem}{1964-28.5}{E}{De Rham and coherent-sheaf finiteness}
\GSfield{Kind; source}{theorem; \GSpage{P000053-P000054}}
\GSfield{The notes}{De Rham forms, holomorphic functions, Kodaira and Cartan are examples of the sheaf approach.}
\GSfield{Statement}{The de Rham complex is a fine resolution of the constant real sheaf. On a compact complex manifold every coherent analytic sheaf has finite-dimensional cohomology. These are separate statements; the holomorphic structure sheaf is not fine.}
\GSfield{Reconstruction}{Coherence and the distinct resolutions are explicit.}
\GSfield{Fate}{\citet[p.~17-1 and Section~1]{cartan1954finitude}, lecture 3 May 1954, states coherent finiteness and the de Rham example; its reference~1 is Cartan--Serre’s 1953 announcement.}
\GSfield{Present status}{As of November 2026, both examples are classical; they require their respective smooth or analytic hypotheses.}
\GSfield{Assessment}{literal statement: correct with missing hypotheses; publication match: exact; verification: read the result; development: annotated; confidence: high}
\end{GSitem}
\begin{GSitem}{1964-28.6}{E}{Compact comparison on resolutions}
\GSfield{Kind; source}{method; \GSpage{P000054-P000055}}
\GSfield{The notes}{Banach/Fréchet sheaves, different resolutions and a compact comparison inducing identity on cohomology are proposed.}
\GSfield{Statement}{If two resolutions have continuous chain maps inducing inverse maps on cohomology, and the composite induces a compact identity on a Hausdorff Banach cohomology space, that space is finite-dimensional. Closed ranges and quotient topologies are necessary in this formulation.}
\GSfield{Reconstruction}{Hausdorff Banach cohomology, continuity and closed-range requirements are supplied.}
\GSfield{Fate}{The analytic finiteness method is exemplified by \citet[proof of the finiteness theorem]{cartan1954finitude}; the source invokes a Schwartz compactness lemma without its precise functional-analytic hypotheses.}
\GSfield{Present status}{As of November 2026, the compact-identity argument works in the reconstructed setting; an arbitrary non-Hausdorff quotient does not qualify.}
\GSfield{Assessment}{literal statement: correct with missing hypotheses; publication match: partial; verification: read the result; development: annotated; confidence: high}
\end{GSitem}

\GStalk{29}{1965-02-15}{I. M. Gelfand}{P000054-P000056}{Fredholm sheaf questions and broader ellipticity}
Gelfand reduces one example to a point and asks about compact perturbations of sheaf maps, infinite-order ellipticity and low-regularity hypoellipticity. Apart from the point case, the operator categories are not specified sufficiently to create definite open problems.

\begin{GSitem}{1964-29.1}{E}{Identity plus compact operator at a point}
\GSfield{Kind; source}{theorem; \GSpage{P000055}}
\GSfield{The notes}{For a one-point space the sheaf question becomes identity plus a compact operator.}
\GSfield{Statement}{For a Banach space \(X\) and compact linear \(K:X\to X\), \(I+K\) is Fredholm with index zero.}
\GSfield{Reconstruction}{Banach topology and linear compactness supplied.}
\GSfield{Fate}{This is classical Riesz--Fredholm theory, contained in the framework for perturbations by compact operators recalled by \citet[p.~861]{smale1965sard}.}
\GSfield{Present status}{As of November 2026, the point case is established, including closed range and finite-dimensional kernel/cokernel.}
\GSfield{Assessment}{literal statement: correct with missing hypotheses; publication match: exact; verification: read the result; development: annotated; confidence: high}
\end{GSitem}
\begin{GSitem}{1964-29.2}{F}{Finite index for a topological sheaf endomorphism}
\GSfield{Kind; source}{question; \GSpage{P000055}}
\GSfield{The notes}{An identity-plus-compact sheaf analogue is asked for.}
\GSfield{Statement}{No unique mathematical statement can be extracted: topologies on sections, meaning of compactness, the cohomology topology and the required index are absent.}
\GSfield{Fate}{Cartan finiteness and Schwartz compact-perturbation sources provide possible frameworks, but none uniquely matches the sentence.}
\GSfield{Present status}{As of November 2026, the fragment remains unidentifiable from the supplied evidence; this does not establish that a mathematical problem is open.}
\GSfield{Assessment}{literal statement: not determinable; publication match: none found; verification: unresolved lead; development: annotated; confidence: low}
\end{GSitem}
\begin{GSitem}{1964-29.3}{F}{Infinite-order elliptic operators}
\GSfield{Kind; source}{question; \GSpage{P000055}}
\GSfield{The notes}{Ellipticity is suggested for infinite-order equations.}
\GSfield{Statement}{No unique mathematical statement can be extracted: no symbol class, growth, function space or notion of ellipticity is given.}
\GSfield{Fate}{Palamodov and pseudodifferential/infinite-order ellipticity searches return many incompatible categories. No precise unaddressed problem follows.}
\GSfield{Present status}{As of November 2026, the fragment remains unidentifiable from the supplied evidence; this does not establish that a mathematical problem is open.}
\GSfield{Assessment}{literal statement: not determinable; publication match: none found; verification: unresolved lead; development: annotated; confidence: low}
\end{GSitem}
\begin{GSitem}{1964-29.4}{F}{Hypoellipticity with continuous coefficients}
\GSfield{Kind; source}{question; \GSpage{P000055}}
\GSfield{The notes}{A question mentions hypoellipticity and continuous coefficients.}
\GSfield{Statement}{No unique mathematical statement can be extracted: operator order, domains, weak-solution meaning and target regularity are not fixed; ordinary \(C^\infty\)-hypoellipticity is not defined by the phrase alone.}
\GSfield{Fate}{Low-regularity elliptic/hypoelliptic searches did not identify a specific conjecture corresponding to this fragment.}
\GSfield{Present status}{As of November 2026, the fragment remains unidentifiable from the supplied evidence; this does not establish that a mathematical problem is open.}
\GSfield{Assessment}{literal statement: not determinable; publication match: none found; verification: unresolved lead; development: annotated; confidence: low}
\end{GSitem}

\GStalk{30}{1965-02-22}{S. P. Novikov}{P000054-P000060}{Adams’s vector-fields theorem and real K-theory}
Novikov explains the vector-fields-on-spheres theorem through elementary constructions, Bott periodicity, real projective K-theory and stunted projective spaces. The Adams theorem was published in 1962. A factor 6 in the periodicity sentence and an order statement are faithful notebook errors.

\begin{GSitem}{1964-30.1}{E}{Elementary sphere obstructions and examples}
\GSfield{Kind; source}{theorem; \GSpage{P000056}}
\GSfield{The notes}{Even spheres have no nonzero field; odd spheres have one; \(S^3,S^7\) are parallelizable.}
\GSfield{Statement}{For even \(n>0\), \(\chi(S^n)=2\) obstructs a nowhere-zero tangent field. On \(S^{2m-1}\subset\mathbb C^m\), \(x\mapsto ix\) is tangent and nonzero. Quaternionic and octonionic multiplication give frames on \(S^3,S^7\).}
\GSfield{Reconstruction}{The explicit complex multiplication and \(n>0\) convention are supplied.}
\GSfield{Fate}{These classical examples precede the upper-bound theorem of \citet[Section~1]{adams1962vector}.}
\GSfield{Present status}{As of November 2026, the assertions hold; \(S^1\) is also parallelizable.}
\GSfield{Assessment}{literal statement: correct with missing hypotheses; publication match: exact; verification: read the result; development: annotated; confidence: high}
\end{GSitem}
\begin{GSitem}{1964-30.2}{E}{Block extension of tangent fields}
\GSfield{Kind; source}{construction; \GSpage{P000056-P000057}}
\GSfield{The notes}{An existing field is applied separately to coordinate blocks, multiplied by their norms.}
\GSfield{Statement}{From continuous independent tangent fields \(v_1,\ldots,v_k\) on \(S^{m-1}\), define on \(S^{pm-1}\) the block field \(w_i(y_1,\ldots,y_p)=(|y_j|v_i(y_j/|y_j|))_j\), with zero at a zero block. These are continuous, tangent and independent; smooth approximation preserves independence.}
\GSfield{Reconstruction}{Zero-block convention and the continuity/smoothing distinction are supplied.}
\GSfield{Fate}{This elementary construction supplies lower-bound intuition in the setting of \citet[Section~1]{adams1962vector}. Its formula is directly checkable from the notebook.}
\GSfield{Present status}{As of November 2026, the extension works; differentiability at a zero block is not automatic for the raw formula.}
\GSfield{Assessment}{literal statement: correct with missing hypotheses; publication match: partial; verification: read the result; development: annotated; confidence: high}
\end{GSitem}
\begin{GSitem}{1964-30.3}{E}{Hurwitz--Radon--Adams number}
\GSfield{Kind; source}{theorem; \GSpage{P000057}}
\GSfield{The notes}{Write \(m=2^b(2a+1)\), \(b=c+4d\), \(0\leq c\leq3\); the maximum number is \(2^c+8d-1\).}
\GSfield{Statement}{If \(m=(2a+1)2^{4d+c}\), then \(S^{m-1}\) has exactly \(\rho(m)-1\) independent continuous tangent fields, where \(\rho(m)=8d+2^c\).}
\GSfield{Reconstruction}{The main boxed formula is correct; a later prose periodicity sentence is corrected separately.}
\GSfield{Fate}{\citet[Theorem~1.1]{adams1962vector}, received November 1961 and published 1962, proves the matching upper bound to the classical lower construction.}
\GSfield{Present status}{As of November 2026, the number is exact in every dimension. Multiplication of \(m\) by 16, not 6, increases it by 8.}
\GSfield{Assessment}{literal statement: correct; publication match: exact; verification: read the result; development: annotated; confidence: high}
\end{GSitem}
\begin{GSitem}{1964-30.4}{E}{Stable vector bundles and KO}
\GSfield{Kind; source}{definition; \GSpage{P000057-P000058}}
\GSfield{The notes}{Stable equivalence, direct sums, tensor products and the real K-functor are described.}
\GSfield{Statement}{For compact \(X\), \(KO^0(X)\) is the Grothendieck group of real vector bundles, with tensor-product multiplication. For connected based \(X\), \(\widetilde{KO}^0(X)\) is the rank-zero subgroup. Stable equivalence of actual bundles alone forgets rank and must be distinguished from full \(KO^0\).}
\GSfield{Reconstruction}{Group completion and reduced group are distinguished.}
\GSfield{Fate}{This is the K-theoretic language of \citet[Sections~1--2]{adams1962vector}; the index-theory use appears already in \citet[p.~423]{atiyahsinger1963index} for complex K-theory.}
\GSfield{Present status}{As of November 2026, the unreduced/reduced distinction resolves the notebook’s missing rank summand.}
\GSfield{Assessment}{literal statement: correct with missing hypotheses; publication match: exact; verification: read the result; development: annotated; confidence: high}
\end{GSitem}
\begin{GSitem}{1964-30.5}{E}{Bott periodicity and stable groups}
\GSfield{Kind; source}{theorem; \GSpage{P000057-P000058}}
\GSfield{The notes}{The stable orthogonal groups recur with period eight and nontrivial entries at residues 0,1,2,4.}
\GSfield{Statement}{\(KO^{-j}(\mathrm{pt})\), for \(j=0,\ldots,7\), is respectively \(\mathbb Z,\mathbb Z/2,\mathbb Z/2,0,\mathbb Z,0,0,0\), repeating with period eight. For \(i>0\), \(\widetilde{KO}^0(S^i)=\pi_i(BO)\).}
\GSfield{Reconstruction}{The table is in KO-degree conventions, not unshifted \(\pi_i(O)\).}
\GSfield{Fate}{\citet[Introduction, stable classical-group periodicity]{bott1959stable}, received November 1958 and published September 1959, is the primary source.}
\GSfield{Present status}{As of November 2026, periodicity is established; the sphere formula requires reduced K-theory.}
\GSfield{Assessment}{literal statement: correct with missing hypotheses; publication match: exact; verification: read the result; development: annotated; confidence: high}
\end{GSitem}
\begin{GSitem}{1964-30.6}{E}{Real K-theory of projective space}
\GSfield{Kind; source}{theorem; \GSpage{P000057-P000059}}
\GSfield{The notes}{The reduced group is cyclic; the tautological class satisfies \(v^2=1\), \(\lambda=v-1\), \(\lambda^2=-2\lambda\). The order is once described as \(f(k)\).}
\GSfield{Statement}{For \(k\ge1\), let \(\phi(k)\) count \(1\le j\le k\) congruent to \(0,1,2,4\pmod8\). Then \(KO^0(\mathbb{RP}^k)=\mathbb Z\oplus\mathbb Z/2^{\phi(k)}\), with reduced generator \(\lambda=[\gamma]-1\), satisfying \(\lambda^2=-2\lambda\).}
\GSfield{Reconstruction}{The tautological bundle is reduced by subtracting the trivial line; the exponent is defined explicitly.}
\GSfield{Fate}{The real-projective computation is part of \citet[projective-space K-theory calculation]{adams1962vector}.}
\GSfield{Present status}{As of November 2026, the correct order is \(2^{\phi(k)}\), in agreement with another line of the same notebook.}
\GSfield{Assessment}{literal statement: false as written; publication match: exact; verification: read the result; development: annotated; confidence: high}
\end{GSitem}
\begin{GSitem}{1964-30.7}{E}{Reflection map and stunted projective space}
\GSfield{Kind; source}{construction; \GSpage{P000059}}
\GSfield{The notes}{An axis determines an orthogonal reflection; a stunted projective quotient maps into a Stiefel quotient.}
\GSfield{Statement}{For integers \(n\ge2\) and \(0\le k\le n-2\), the map \([v]\mapsto I-2vv^t\) from \(\mathbb{RP}^{n-1}\) to \(O(n)\), for unit \(v\), induces \(\mathbb{RP}^{n-1}/\mathbb{RP}^{n-k-2}\to O(n)/O(n-k-1)\) by taking the last \(k+1\) columns. The collapsed subspace is precisely where those coordinates of \(v\) vanish.}
\GSfield{Reconstruction}{A specific block convention for the Stiefel quotient is supplied.}
\GSfield{Fate}{Stunted projective spaces are central to \citet[Theorem~1.2 and ensuing reduction]{adams1962vector}. The explicit reflection formula is editorial clarification of the drawing.}
\GSfield{Present status}{As of November 2026, the construction is well-defined; indexing must match the chosen Stiefel frame size.}
\GSfield{Assessment}{literal statement: correct with missing hypotheses; publication match: partial; verification: read the result; development: annotated; confidence: high}
\end{GSitem}

\GStalk{31}{1965-03-01}{V. P. Palamodov}{P000058-P000061}{Primary-ideal quotient revisited}
P000060 restates the compatible differential-jet realization of a primary polynomial ideal quotient and says Kazhdan obtains it algebraically. It is a recurrence of record 27, not a separate fully specified growth theorem.

\begin{GSitem}{1964-31.1}{A}{Noetherian differential jets and the algebraic proof}
\GSfield{Kind; source}{theorem; \GSpage{P000060}}
\GSfield{The notes}{\(\mathcal H/J\mathcal H\cong\mathcal H\{J\}\), with \(\mathcal H\{J\}\) locally realized by \(d(z,D)F|_N\); Kazhdan is said to obtain this algebraically.}
\GSfield{Statement}{The quotient is represented by compatible finite differential jets on the primary support, once suitable Noetherian operators and the locally realizable target sheaf are fixed. Algebraically, a primary ideal in characteristic-zero polynomial rings is the simultaneous vanishing modulo its radical of finitely many polynomial-coefficient differential operators.}
\GSfield{Reconstruction}{The operator construction and target compatibility are supplied; no attribution of the modern algorithm to Kazhdan is made.}
\GSfield{Fate}{The speaker’s systematic treatment is \citet[Part~I]{palamodov1970operators}; the algebraic operator statement is explicitly formulated in \citet[Theorem~2.1 and equation~(17)]{cidruizhomssturmfels2021primary}. The latter circulated January 2020, published 2021.}
\GSfield{Present status}{As of November 2026, the framework is established; the note does not identify a separate Kazhdan publication or supply all analytic globalization hypotheses.}
\GSfield{Assessment}{literal statement: correct with missing hypotheses; publication match: partial; verification: indirect; development: annotated; confidence: medium}
\end{GSitem}

\GStalk{32}{1965-03-01}{I. M. Gelfand}{P000059-P000062}{Nelson’s commuting differential operators}
P000060--P000061 distinguishes essential self-adjointness from formal symmetry and strong commutation from commutation on a core. The punctured genus-four surface is Nelson’s example, already published in 1959.

\begin{GSitem}{1964-32.1}{E}{Essential self-adjointness and boundary conditions}
\GSfield{Kind; source}{theorem; \GSpage{P000060}}
\GSfield{The notes}{The range of \(A-\lambda\) is required to be dense for every complex \(\lambda\); differentiation on an interval is compared with periodic boundary conditions.}
\GSfield{Statement}{For a densely defined symmetric \(A\), essential self-adjointness is equivalent to density of both \(\operatorname{ran}(A-i)\) and \(\operatorname{ran}(A+i)\). On \(L^2(0,1)\), \(i\,d/dx\) on \(H^1_0(0,1)\) has deficiency indices \((1,1)\); the periodic \(H^1\) domain gives a self-adjoint extension.}
\GSfield{Reconstruction}{Dense domain, the nonreal parameters, and the precise Sobolev domains are supplied.}
\GSfield{Fate}{These are the operator-theoretic preliminaries to \citet[Sections~9--10]{nelson1959analytic}. The range criterion needs a pair of nonreal parameters, not all complex numbers.}
\GSfield{Present status}{As of November 2026, the distinction between a symmetric operator, its closure and a chosen self-adjoint extension is essential.}
\GSfield{Assessment}{literal statement: false as written; publication match: partial; verification: read the result; development: annotated; confidence: high}
\end{GSitem}
\begin{GSitem}{1964-32.2}{E}{Commuting on a domain does not imply strong commutation}
\GSfield{Kind; source}{construction; \GSpage{P000060-P000061}}
\GSfield{The notes}{\(A,B\) commute on an invariant dense domain, and all real linear combinations are essentially self-adjoint, yet their exponentials need not commute.}
\GSfield{Statement}{There exist symmetric \(A,B\) on a common invariant dense domain \(D\), with \(ABf=BAf\) for \(f\in D\), such that every \(aA+bB\), \(a,b\in\mathbb R\), is essentially self-adjoint on \(D\), while the spectral measures of their closures do not commute.}
\GSfield{Reconstruction}{The distinction between algebraic commutation and commuting spectral measures is explicit.}
\GSfield{Fate}{\citet[Section~10]{nelson1959analytic}, published 1959, constructs precisely this obstruction. Essential self-adjointness of a sum of squares is a stronger sufficient integrability condition.}
\GSfield{Present status}{As of November 2026, the counterexample remains valid; separate essential self-adjointness does not license joint functional calculus.}
\GSfield{Assessment}{literal statement: correct; publication match: exact; verification: read the result; development: annotated; confidence: high}
\end{GSitem}
\begin{GSitem}{1964-32.3}{E}{The punctured genus-four translation surface}
\GSfield{Kind; source}{construction; \GSpage{P000061}}
\GSfield{The notes}{A planar region with an inner square removed is glued by translations; after identification of corner points, it is a genus-four surface minus one point.}
\GSfield{Statement}{On Nelson’s translation surface, \(X=\partial_x,Y=\partial_y\) commute locally on \(C_c^\infty\). Almost-everywhere complete straight-line flows generate unitary groups on area \(L^2\); \(iX,iY\), and every real linear combination, are essentially self-adjoint on this domain. The two global flows fail to commute across the missing point.}
\GSfield{Reconstruction}{The invariant area measure, compactly supported smooth domain and factors of \(i\) are supplied.}
\GSfield{Fate}{\citet[Section~10; Russian translation pp.~121--123]{nelson1959analytic} explicitly gives genus four. The notebook drawing agrees; replacing four by two would introduce an error.}
\GSfield{Present status}{As of November 2026, the geometric example is established and the source genus is retained.}
\GSfield{Assessment}{literal statement: correct; publication match: exact; verification: read the result; development: annotated; confidence: high}
\end{GSitem}
\begin{GSitem}{1964-32.4}{E}{The Laplacian obstruction to integrating the Lie algebra}
\GSfield{Kind; source}{theorem; \GSpage{P000061}}
\GSfield{The notes}{If the Laplacian were self-adjoint, all fields would generate the expected group; therefore the example’s Laplacian is not essentially self-adjoint.}
\GSfield{Statement}{For a representation of a finite-dimensional real Lie algebra by skew-symmetric operators \(X_j\) on a common invariant dense domain, essential self-adjointness of the Nelson Laplacian \(-\sum_jX_j^2\) is a sufficient condition for integration to a unitary representation of the simply connected group. Thus this sufficient condition fails for Nelson’s noncommuting translation example.}
\GSfield{Reconstruction}{A common invariant domain, the simply connected integrating group and the sign convention are supplied.}
\GSfield{Fate}{\citet[Section~9, integrability theorem, and Section~10]{nelson1959analytic} proves the criterion and applies it to the example.}
\GSfield{Present status}{As of November 2026, the conclusion concerns essential self-adjointness on the specified domain; it does not deny the existence of self-adjoint extensions.}
\GSfield{Assessment}{literal statement: correct with missing hypotheses; publication match: exact; verification: read the result; development: annotated; confidence: high}
\end{GSitem}
\begin{GSitem}{1964-32.5}{E}{Spectral construction of the unitary evolution}
\GSfield{Kind; source}{method; \GSpage{P000061}}
\GSfield{The notes}{The exponent of a self-adjoint operator is written using an unnormalized resolvent integral, and the differential equation \(u_t=iAu\) is invoked.}
\GSfield{Statement}{For self-adjoint \(A\), \(e^{itA}=\int_{\mathbb R}e^{it\lambda}\,dE_A(\lambda)\). This is a strongly continuous unitary group on all \(H\); strong differentiation solves \(u_t=iAu\) for initial vectors in \(D(A)\). Resolvent boundary values recover \(E_A\) through Stone’s jump formula, not through the unqualified integral in the note.}
\GSfield{Reconstruction}{Spectral measure, normalization and differentiability domain are supplied.}
\GSfield{Fate}{The spectral and unitary-group framework is used throughout \citet[Sections~9--10]{nelson1959analytic}; the displayed correction follows directly from spectral calculus.}
\GSfield{Present status}{As of November 2026, the unitary evolution is established, with a domain condition for strong differentiability.}
\GSfield{Assessment}{literal statement: false as written; publication match: partial; verification: read the result; development: annotated; confidence: high}
\end{GSitem}

\GStalk{33}{1965-03-01}{Ya. G. Sinai}{P000060-P000064}{Simple eigenvalues and countable Lebesgue spectrum}
P000062 and the top of P000063 distinguish simple eigenspaces from the continuous spectral multiplicity of toral automorphisms and Bernoulli shifts. The remark that all known Lebesgue examples have countable multiplicity is historical; no explicit Banach problem is posed.

\begin{GSitem}{1964-33.1}{E}{Simplicity of eigenvalues for an ergodic system}
\GSfield{Kind; source}{theorem; \GSpage{P000062}}
\GSfield{The notes}{For an ergodic system every eigenvalue is simple; invariant functions are constant.}
\GSfield{Statement}{For an ergodic probability-preserving \(\mathbb Z\)-action or \(\mathbb R\)-flow, each Koopman eigenspace is one-dimensional. A nonzero eigenfunction has constant modulus, and the ratio of two eigenfunctions for the same character is invariant. This does not assert simple continuous spectrum.}
\GSfield{Reconstruction}{Probability preservation and the distinction between eigenspaces and continuous multiplicity are supplied.}
\GSfield{Fate}{The classical pure-point classification is restated precisely in \citet[Introduction and Theorem~5.5]{hermlekreidler2023halmos}; its antecedent is Halmos--von Neumann (1942).}
\GSfield{Present status}{As of November 2026, the statement is correct under ergodicity; the source heading concerns discrete spectrum.}
\GSfield{Assessment}{literal statement: correct with missing hypotheses; publication match: partial; verification: read the result; development: annotated; confidence: high}
\end{GSitem}
\begin{GSitem}{1964-33.2}{E}{Absolute continuity versus homogeneous Lebesgue spectrum}
\GSfield{Kind; source}{definition; \GSpage{P000062}}
\GSfield{The notes}{Correlation functions are Fourier transforms of integrable densities, and a constant spectral multiplicity is called homogeneous.}
\GSfield{Statement}{For the centered Koopman space \(L^2_0(X,\mu)\), absolute continuity means every scalar spectral measure is absolutely continuous with respect to Lebesgue measure on \(\mathbb R\) for flows or Haar measure on \(\mathbb T\) for transformations. Lebesgue maximal type additionally means equivalence of measure classes; homogeneous multiplicity means the multiplicity function is almost everywhere constant.}
\GSfield{Reconstruction}{Centered space, measure-class equivalence and the time-group distinction are explicit.}
\GSfield{Fate}{\citet[Introduction, pp.~899--900]{sinai1961lebesgue} distinguishes the homogeneous Lebesgue conclusion and proves countable multiplicity in its class of flows.}
\GSfield{Present status}{As of November 2026, absolute continuity alone does not guarantee full Lebesgue maximal type or homogeneous multiplicity; constants must be removed.}
\GSfield{Assessment}{literal statement: correct with missing hypotheses; publication match: partial; verification: read the result; development: annotated; confidence: high}
\end{GSitem}
\begin{GSitem}{1964-33.3}{F}{Historical multiplicity remark and an editorial Banach comparison}
\GSfield{Kind; source}{remark; \GSpage{P000062}}
\GSfield{The notes}{All known examples of systems with Lebesgue spectrum are said to have countably infinite multiplicity. There is no explicit existence question, no Banach attribution, and no fixed choice between transformations and flows.}
\GSfield{Statement}{The exact source claim is historical. An editorial comparison asks whether an ergodic probability-preserving transformation or flow can have simple Lebesgue spectrum on \(L^2_0\); these are different questions. The notebook does not uniquely specify either as a proposed problem.}
\GSfield{Reconstruction}{The modern comparison is editorial, the flow/transformation distinction is supplied, and the 2026 claim is separated from verified established results.}
\GSfield{Fate}{For flows, \citet[main result]{prikhodko2020flows} gives a positive construction; the paper was received in June 2012, revised October 2019, and published 2020. For transformations, \citet[Introduction, main theorem]{openai2026lebesgue}, dated 23 September 2026, claims a smooth volume-preserving three-torus diffeomorphism with a bilateral Koopman orbit forming an orthonormal basis of \(L^2_0\). Its theorem statement was inspected, but its proof was not independently verified here.}
\GSfield{Present status}{As of November 2026, simple Lebesgue flows have a published construction; a new transformation claim exists, and this audit does not establish an accepted resolution of the transformation problem. Neither modern question is assigned to the notebook as an explicit conjecture.}
\GSfield{Assessment}{literal statement: not determinable; publication match: thematic only; verification: unresolved lead; development: annotated; confidence: medium}
\end{GSitem}
\begin{GSitem}{1964-33.4}{E}{Hyperbolic toral automorphisms have countable Lebesgue spectrum}
\GSfield{Kind; source}{remark; \GSpage{P000062}}
\GSfield{The notes}{The integer-matrix automorphism of the two-torus is diagonalized by Fourier modes whose integer labels run in infinite orbits.}
\GSfield{Statement}{For \(A\in SL(2,\mathbb Z)\) with \(|\operatorname{tr}A|>2\), the automorphism \(x\mapsto Ax\) of the Haar probability torus has on \(L^2_0\) homogeneous countably infinite Lebesgue spectrum: each orbit of \(A^t\) on \(\mathbb Z^2\setminus\{0\}\) gives a bilateral shift.}
\GSfield{Reconstruction}{Haar probability, removal of constants and hyperbolicity are supplied.}
\GSfield{Fate}{\citet[p.~899]{sinai1961lebesgue} identifies compact abelian-group automorphisms among the already known algebraic examples. The Fourier-orbit argument is explicit in the notebook.}
\GSfield{Present status}{As of November 2026, the result holds under hyperbolicity; an arbitrary matrix in \(SL(2,\mathbb Z)\) need not satisfy it.}
\GSfield{Assessment}{literal statement: correct with missing hypotheses; publication match: partial; verification: read the result; development: annotated; confidence: high}
\end{GSitem}
\begin{GSitem}{1964-33.5}{E}{Bernoulli shifts have countable Lebesgue spectrum}
\GSfield{Kind; source}{remark; \GSpage{P000062-P000063}}
\GSfield{The notes}{A tensor product of independent coordinates is organized into translation orbits.}
\GSfield{Statement}{A nontrivial two-sided Bernoulli shift on a standard probability product has homogeneous countably infinite Lebesgue spectrum on \(L^2_0\). Finite tensor words in centered coordinate basis vectors form an orthonormal basis; their shift orbits are free and countably many.}
\GSfield{Reconstruction}{Nontrivial base, two-sided shift, separability and centered space supplied.}
\GSfield{Fate}{This is a classical example within the discrete-time Kolmogorov setting recalled in \citet[Introduction, pp.~899--900]{sinai1961lebesgue}. The tensor-basis proof is directly reconstructed from P000062--P000063.}
\GSfield{Present status}{As of November 2026, the statement holds for the standard probability and separability setting; the constant tensor is excluded.}
\GSfield{Assessment}{literal statement: correct with missing hypotheses; publication match: partial; verification: read the result; development: annotated; confidence: high}
\end{GSitem}
\begin{GSitem}{1964-33.6}{E}{Spectral measures of tensor products are convolutions}
\GSfield{Kind; source}{method; \GSpage{P000063}}
\GSfield{The notes}{Product eigenfunctions produce sums of frequencies, and the associated measure is called a product measure.}
\GSfield{Statement}{For \(\mathbb R\)-representations, the scalar spectral measure of \(v\otimes w\) under \(U_t\otimes V_t\) is the pushforward of \(\mu_v\otimes\nu_w\) by \((\lambda,\eta)\mapsto\lambda+\eta\), hence convolution. For \(\mathbb Z\), use multiplication on \(\mathbb T\).}
\GSfield{Reconstruction}{The acting group is the diagonal subgroup and the spectral pushforward is explicit.}
\GSfield{Fate}{The product-representation setting already appears in \citet[Section~3]{bruhat1954induced}, a December 1952 lecture printed in 1954. The pushforward follows by multiplying the two correlation Fourier transforms.}
\GSfield{Present status}{As of November 2026, the tensor-product formula is established; the measure on the one-dimensional dual is a pushforward, not the raw Cartesian product.}
\GSfield{Assessment}{literal statement: correct with missing hypotheses; publication match: partial; verification: read the result; development: annotated; confidence: high}
\end{GSitem}

\GStalk{34}{1965-03-22}{A. A. Kirillov}{P000062-P000064}{Group actions, spectral reconstruction and cocycles}
The actual record runs from P000063 through P000065, beyond the catalogue envelope. It compares compact homogeneous actions, noncompact representations, cocycles and reconstruction. The finite-group reconstruction question has negative examples; several cohomological remarks remain too indefinite to identify.

\begin{GSitem}{1964-34.1}{E}{Koopman spectra for general group actions}
\GSfield{Kind; source}{definition; \GSpage{P000063}}
\GSfield{The notes}{Dynamical systems are extended to arbitrary groups, with irreducible representations replacing scalar eigenvalues.}
\GSfield{Statement}{A measurable probability-preserving action of a locally compact second-countable group \(G\) on a standard probability space gives a unitary Koopman representation. For a type-I group the usual essentially unique irreducible direct-integral decomposition permits a spectral measure-class and multiplicity description. Arbitrary groups require more care; such data are not automatically a complete classification.}
\GSfield{Reconstruction}{Standard measure spaces, continuity and the type-I restriction on the familiar spectral parametrization are supplied.}
\GSfield{Fate}{The distinction between a representation and an action is central to \citet[Introduction and Section~6]{hermlekreidler2023halmos}. \citet[Section~2]{gelfandfomin1952geodesic} already uses irreducible decomposition for the relevant semisimple groups.}
\GSfield{Present status}{As of November 2026, the representation-theoretic construction is established; the unrestricted classification language needs the type-I and reconstruction qualifications.}
\GSfield{Assessment}{literal statement: correct with missing hypotheses; publication match: partial; verification: read the result; development: annotated; confidence: high}
\end{GSitem}
\begin{GSitem}{1964-34.2}{E}{Ergodic compact-group actions are essentially homogeneous}
\GSfield{Kind; source}{theorem; \GSpage{P000063}}
\GSfield{The notes}{For a compact group, an ergodic action is described by a homogeneous space.}
\GSfield{Statement}{An ergodic probability-preserving action of a second-countable compact group on a standard probability space is, modulo null sets, its action on \(G/H\) for a closed subgroup \(H\), with invariant quotient Haar probability. Equivariant isomorphism corresponds to conjugacy of stabilizers.}
\GSfield{Reconstruction}{Closed stabilizer, quotient Haar measure and the modulo-null-sets conclusion are supplied.}
\GSfield{Fate}{The homogeneous model is the compact instance of \citet[Remark~6.44(iii)]{hermlekreidler2023halmos}; the stabilizer classification is explicit in \citet[Sections~2--3]{mackey1966virtual}, reviewing the earlier elementary transitive theory.}
\GSfield{Present status}{As of November 2026, the theorem holds in the measurable standard-space formulation; pointwise transitivity of an arbitrary representative is stronger than required.}
\GSfield{Assessment}{literal statement: correct with missing hypotheses; publication match: exact; verification: read the result; development: annotated; confidence: high}
\end{GSitem}
\begin{GSitem}{1964-34.3}{F}{A finite-support binary group and an uncertain attribution}
\GSfield{Kind; source}{remark; \GSpage{P000063}}
\GSfield{The notes}{The notation is a product of copies of \(\mathbb Z/2\), but the elements are said to have finitely many nonzero coordinates; names resembling Gordin and Whitehead appear.}
\GSfield{Statement}{No unique mathematical statement can be extracted: the group described is \(\bigoplus_{\mathbb N}\mathbb Z/2\), while its topology, action and the attributed result are not specified.}
\GSfield{Fate}{Searches combining the names, binary groups, spectra and ergodic actions did not identify an exact theorem; the finite-support/product discrepancy is recorded separately.}
\GSfield{Present status}{As of November 2026, the fragment remains unidentifiable from the supplied evidence; this does not establish that a mathematical problem is open.}
\GSfield{Assessment}{literal statement: not determinable; publication match: none found; verification: unresolved lead; development: annotated; confidence: low}
\end{GSitem}
\begin{GSitem}{1964-34.4}{E}{Dyadic roots and their compact dual}
\GSfield{Kind; source}{remark; \GSpage{P000063-P000064}}
\GSfield{The notes}{The dyadic roots of unity are a countable group whose characters form a compact group.}
\GSfield{Statement}{The discrete Prüfer group \(C_{2^\infty}=\mathbb Z[1/2]/\mathbb Z\) is the direct limit of \(C_{2^n}\). Its Pontryagin dual is the inverse limit \(\varprojlim\mathbb Z/2^n\mathbb Z=\mathbb Z_2\), with the compact inverse-limit topology.}
\GSfield{Reconstruction}{Direct-limit and inverse-limit maps are supplied.}
\GSfield{Fate}{The identification is classical and follows directly by restricting a character to each cyclic subgroup: the resulting residues are compatible modulo successive powers of two. No new spectral theorem is asserted.}
\GSfield{Present status}{As of November 2026, the group and dual are unambiguous once the discrete topology is specified.}
\GSfield{Assessment}{literal statement: correct with missing hypotheses; publication match: none found; verification: indirect; development: annotated; confidence: medium}
\end{GSitem}
\begin{GSitem}{1964-34.5}{F}{Unspecified dyadic spectral decomposition}
\GSfield{Kind; source}{remark; \GSpage{P000064}}
\GSfield{The notes}{A remark about point and continuous spectra and pieces of functions follows the dyadic group example.}
\GSfield{Statement}{No unique mathematical statement can be extracted: no action, invariant measure or definite decomposition assertion is supplied.}
\GSfield{Fate}{The dyadic-group identification does not determine the spectrum of an unspecified action; no unique publication match was found.}
\GSfield{Present status}{As of November 2026, the fragment remains unidentifiable from the supplied evidence; this does not establish that a mathematical problem is open.}
\GSfield{Assessment}{literal statement: not determinable; publication match: none found; verification: unresolved lead; development: annotated; confidence: low}
\end{GSitem}
\begin{GSitem}{1964-34.6}{E}{The full real affine group has one infinite-dimensional irreducible}
\GSfield{Kind; source}{theorem; \GSpage{P000064}}
\GSfield{The notes}{The group \(x\mapsto ax+b\), \(a\ne0\), is said to have one infinite-dimensional irreducible representation, besides characters.}
\GSfield{Statement}{For \(G=\mathbb R\rtimes\mathbb R^\times\), the one-dimensional unitary characters are \((b,a)\mapsto |a|^{it}\operatorname{sgn}(a)^\varepsilon\), \(t\in\mathbb R,\varepsilon\in\{0,1\}\), and there is exactly one infinite-dimensional irreducible unitary representation up to equivalence. Restricting the group to \(a>0\) gives two such irreducibles instead.}
\GSfield{Reconstruction}{The full group, continuous unitary representations and both one-dimensional character parameters are explicit.}
\GSfield{Fate}{\citet[Sections~2.2--2.3 and Section~5]{choighandehari2021affine} explicitly records this classical Mackey classification and distinguishes the connected subgroup.}
\GSfield{Present status}{As of November 2026, the notebook’s word one is correct for the full group; replacing it by two would silently change the group.}
\GSfield{Assessment}{literal statement: correct; publication match: exact; verification: read the result; development: annotated; confidence: high}
\end{GSitem}
\begin{GSitem}{1964-34.7}{E}{Discrete full-group spectrum on a compact hyperbolic quotient}
\GSfield{Kind; source}{claim; \GSpage{P000064}}
\GSfield{The notes}{For a compact surface of negative curvature the relevant Lorentz-group representation is said to have discrete spectrum.}
\GSfield{Statement}{Let \(G=PSL(2,\mathbb R)\) and \(\Gamma\) be cocompact discrete. The right representation of \(G\) on \(L^2(\Gamma\backslash G)\) is a discrete Hilbert sum of irreducibles, each with finite multiplicity. Its restriction to the geodesic one-parameter subgroup has Lebesgue spectrum on \(L^2_0\), so the two uses of spectrum must not be identified.}
\GSfield{Reconstruction}{The acting group, cocompactness and the distinction between full-group decomposition and flow spectrum are supplied.}
\GSfield{Fate}{\citet[Sections~2--3, especially p.~132]{gelfandfomin1952geodesic} establishes the continuous-time Lebesgue conclusion. For the full-group discreteness assertion, compactness makes integrated compactly supported continuous kernels compact; the usual compact-operator decomposition supplies the discrete sum.}
\GSfield{Present status}{As of November 2026, the full-group and one-parameter assertions coexist; compactness alone does not give pure-point geodesic-flow spectrum.}
\GSfield{Assessment}{literal statement: correct with missing hypotheses; publication match: partial; verification: read the result; development: annotated; confidence: high}
\end{GSitem}
\begin{GSitem}{1964-34.8}{E}{Measurable cocycles and coboundaries}
\GSfield{Kind; source}{definition; \GSpage{P000064}}
\GSfield{The notes}{Operators have form \(T_gf(x)=f(xg)\chi(x,g)\); equivalent multipliers differ by a ratio of a function at \(xg\) and \(x\).}
\GSfield{Statement}{For a right probability-preserving action, a measurable \(\mathbb T\)-valued cocycle satisfies \(\alpha(x,gh)=\alpha(x,g)\alpha(xg,h)\), almost everywhere with the chosen measurable formulation. It defines unitary \(T_gf(x)=\alpha(x,g)f(xg)\). Replacing \(\alpha\) by \(b(x)^{-1}\alpha(x,g)b(xg)\) conjugates the representation by multiplication by \(b\).}
\GSfield{Reconstruction}{Unit-modulus values, preservation of measure and an explicit right-action convention are supplied.}
\GSfield{Fate}{\citet[Section~5; Russian translation pp.~127--128]{mackey1966virtual} gives the cocycle and cohomology formulas and explains their 1963 announcement.}
\GSfield{Present status}{As of November 2026, the construction is established; measurable equivalence and changes on null sets require the standard cocycle framework.}
\GSfield{Assessment}{literal statement: correct with missing hypotheses; publication match: exact; verification: read the result; development: annotated; confidence: high}
\end{GSitem}
\begin{GSitem}{1964-34.9}{E}{Transitive cocycles and stabilizer representations}
\GSfield{Kind; source}{theorem; \GSpage{P000064}}
\GSfield{The notes}{For a transitive action the cocycle classes are identified with representations of the subgroup.}
\GSfield{Statement}{For the homogeneous right action on \(H\backslash G\), strict Borel cocycles with values in a compact group \(K\), up to Borel cohomology, correspond to continuous homomorphisms \(H\to K\) up to conjugacy. With \(K=\mathbb T\) these are characters of \(H\); with \(K=U(n)\) they are \(n\)-dimensional unitary representations.}
\GSfield{Reconstruction}{Strict Borel representatives, closed stabilizer and the coefficient group are specified.}
\GSfield{Fate}{\citet[Section~5, final transitive-case paragraph]{mackey1966virtual} gives precisely the homomorphism correspondence; \citet[Section~2]{bruhat1954induced} supplies the earlier induced-representation framework.}
\GSfield{Present status}{As of November 2026, the scalar formula yields characters, while higher-dimensional representations require vector-valued cocycles.}
\GSfield{Assessment}{literal statement: correct with missing hypotheses; publication match: exact; verification: read the result; development: annotated; confidence: high}
\end{GSitem}
\begin{GSitem}{1964-34.10}{E}{Virtual subgroups and ergodic groupoids}
\GSfield{Kind; source}{remark; \GSpage{P000064}}
\GSfield{The notes}{Mackey’s generalized subgroups and the analogy between transitive and ergodic actions are discussed.}
\GSfield{Statement}{The identifiable programme replaces a homogeneous action’s stabilizer by a virtual subgroup described through an ergodic measured groupoid. It is a framework for actions, cocycles and induction, not a single unsolved conjecture.}
\GSfield{Reconstruction}{The later groupoid language identifies the programme without assigning every later result to the seminar.}
\GSfield{Fate}{\citet{mackey1963ergodic} announced the programme before this seminar; \citet[Introduction and Sections~9--11]{mackey1966virtual} supplies the systematic explanation and the groupoid formulation.}
\GSfield{Present status}{As of November 2026, the programme became an established part of measured groupoid and ergodic-action theory; the source contains no uniquely delimited remaining question.}
\GSfield{Assessment}{literal statement: correct; publication match: partial; verification: read the result; development: annotated; confidence: high}
\end{GSitem}
\begin{GSitem}{1964-34.11}{F}{Does cohomology determine the dynamical system?}
\GSfield{Kind; source}{question; \GSpage{P000064-P000065}}
\GSfield{The notes}{A broad reconstruction question asks whether cohomology determines a system.}
\GSfield{Statement}{No unique mathematical statement can be extracted: the cohomology degree, coefficient object, additional structure to retain and equivalence of systems are unspecified.}
\GSfield{Fate}{Mackey’s cocycle programme and later reconstruction results delimit several inequivalent readings. They do not identify this fragment with a particular modern open problem.}
\GSfield{Present status}{As of November 2026, the fragment remains unidentifiable from the supplied evidence; this does not establish that a mathematical problem is open.}
\GSfield{Assessment}{literal statement: not determinable; publication match: none found; verification: unresolved lead; development: annotated; confidence: low}
\end{GSitem}
\begin{GSitem}{1964-34.12}{E}{Products of eigenfunctions and a symmetric two-cocycle}
\GSfield{Kind; source}{construction; \GSpage{P000065}}
\GSfield{The notes}{Products of normalized eigenfunctions differ from the eigenfunction of the product eigenvalue by constants, satisfying a cocycle law.}
\GSfield{Statement}{For an ergodic probability-preserving abelian action, choose unit-modulus eigenfunctions \(f_\lambda\). Then \(f_\lambda f_\mu=c(\lambda,\mu)f_{\lambda\mu}\), with \(c\in\mathbb T\); associativity gives the two-cocycle identity and commutativity gives symmetry. Changing phases changes \(c\) by a coboundary. With pure-point spectrum the eigenvalue subgroup classifies the action; otherwise this describes its Kronecker factor only.}
\GSfield{Reconstruction}{Abelian action, normalized modulus, phase equivalence and the pure-point restriction are supplied.}
\GSfield{Fate}{\citet[Theorem~5.5 and the preceding eigenfunction construction]{hermlekreidler2023halmos} gives a modern precise version of the 1942 Halmos--von Neumann classification. The two-cocycle equations follow directly by associating a triple product.}
\GSfield{Present status}{As of November 2026, the source construction is correct; full reconstruction requires pure-point spectrum and the standard probability-space setting.}
\GSfield{Assessment}{literal statement: correct with missing hypotheses; publication match: partial; verification: read the result; development: annotated; confidence: high}
\end{GSitem}
\begin{GSitem}{1964-34.13}{E}{Unitary spectrum does not determine even a finite ergodic action}
\GSfield{Kind; source}{question; \GSpage{P000065}}
\GSfield{The notes}{Whether the spectrum determines the system is said to be unclear even for finite groups; Gelfand answers that it does not.}
\GSfield{Statement}{The answer is negative even for transitive finite actions. Let \(G=GL(3,\mathbb F_2)\) act on the seven nonzero vectors of \(\mathbb F_2^3\), and on the seven nonzero covectors by the dual action. The permutation characters coincide, but the \(G\)-sets are not isomorphic. With uniform probability these are ergodic systems with unitarily equivalent Koopman representations.}
\GSfield{Reconstruction}{The seven-point example is supplied as an elementary check, not attributed to the notebook or asserted to be Todd’s exact presentation. The full chapter’s independent finite-action computation checks all \(168\) matrices and the stabilizer obstruction.}
\GSfield{Fate}{\citet[Remark~6.44(iii)]{hermlekreidler2023halmos} records the obstruction already noted by Mackey (1964, Section~3, Remark~2), citing Todd’s 1950 example. The explicit seven-point verification here is editorial: fixed-vector and fixed-covector counts agree by rank equality, while the stabilizer of a vector fixes no nonzero covector.}
\GSfield{Present status}{As of November 2026, the finite-group question has a negative answer; this is not an open reconstruction problem.}
\GSfield{Assessment}{literal statement: correct with missing hypotheses; publication match: exact; verification: indirect; development: annotated; confidence: high; problem \hyperref[prob:GP-1964-09]{GP-1964-09}}
\end{GSitem}
\begin{GSitem}{1964-34.14}{E}{The icosahedral rotation group needs a central quotient}
\GSfield{Kind; source}{remark; \GSpage{P000065}}
\GSfield{The notes}{The icosahedral rotation group is identified with the unimodular group over five elements.}
\GSfield{Statement}{The orientation-preserving icosahedral group is \(A_5\cong PSL(2,\mathbb F_5)=SL(2,\mathbb F_5)/\{\pm I\}\), of order \(60\). The group \(SL(2,\mathbb F_5)\) itself has order \(120\).}
\GSfield{Reconstruction}{The quotient by the center is supplied.}
\GSfield{Fate}{This is the classical finite-group identification. The immediate order check already forces the central quotient: \(|SL(2,\mathbb F_q)|=q(q^2-1)\), and the center at \(q=5\) has two elements.}
\GSfield{Present status}{As of November 2026, the missing projective quotient is a faithful mathematical error in the note.}
\GSfield{Assessment}{literal statement: false as written; publication match: none found; verification: indirect; development: annotated; confidence: medium}
\end{GSitem}
\begin{GSitem}{1964-34.15}{E}{Multiplicative observables recover the action}
\GSfield{Kind; source}{remark; \GSpage{P000065}}
\GSfield{The notes}{The spectral representation forgets multiplication; a ring of functions with its group action restores the system.}
\GSfield{Statement}{For standard probability spaces, a state-preserving normal unital \(*\)-isomorphism between their commutative \(L^\infty\) algebras that intertwines the group actions gives an isomorphism of measure algebras, hence a spatial conjugacy modulo null sets. Retaining just the \(L^2\) representation discards this multiplicative information.}
\GSfield{Reconstruction}{The algebraic, order/normality and probability-state structure is explicit.}
\GSfield{Fate}{\citet[Section~1.2 and Theorem~4.8]{hermlekreidler2023halmos} makes the reconstruction category explicit through Markov lattice homomorphisms and compact models.}
\GSfield{Present status}{As of November 2026, the structured observable algebra is sufficient in the stated category; an abstract spectrum alone is not.}
\GSfield{Assessment}{literal statement: correct with missing hypotheses; publication match: partial; verification: read the result; development: annotated; confidence: high}
\end{GSitem}

\GStalk{35}{1965-03-22}{I. M. Gelfand; A. Ya. Povzner; F. A. Berezin}{P000062-P000064}{Proposed future directions}
The short programme on P000063 precedes Kirillov physically, though it is record35 in the supplied order. Four directions are named without definite objects or conjectures; they must not be retroactively identified with later named theories.

\begin{GSitem}{1964-35.1}{F}{Noncommutative algebraic geometry}
\GSfield{Kind; source}{remark; \GSpage{P000063}; I. M. Gelfand}
\GSfield{The notes}{Gelfand proposes a version of algebraic geometry for noncommutative rings, with an ideal/matrix remark.}
\GSfield{Statement}{No unique mathematical statement can be extracted: the class of algebras, geometric objects, equivalence and required result are not specified.}
\GSfield{Fate}{The phrase has extensive later literature, but no exact publication or residual problem is identifiable from this proposal.}
\GSfield{Present status}{As of November 2026, the fragment remains unidentifiable from the supplied evidence; this does not establish that a mathematical problem is open.}
\GSfield{Assessment}{literal statement: not determinable; publication match: none found; verification: unresolved lead; development: annotated; confidence: low}
\end{GSitem}
\begin{GSitem}{1964-35.2}{F}{An algebraic theory of differential equations}
\GSfield{Kind; source}{remark; \GSpage{P000063}; I. M. Gelfand}
\GSfield{The notes}{Gelfand proposes an algebraic theory of differential equations.}
\GSfield{Statement}{No unique mathematical statement can be extracted: no differential ring, module category, equation class or theorem is specified.}
\GSfield{Fate}{Several later algebraic and module-theoretic frameworks are thematic matches only. Assigning a particular later theory would exceed the evidence.}
\GSfield{Present status}{As of November 2026, the fragment remains unidentifiable from the supplied evidence; this does not establish that a mathematical problem is open.}
\GSfield{Assessment}{literal statement: not determinable; publication match: none found; verification: unresolved lead; development: annotated; confidence: low}
\end{GSitem}
\begin{GSitem}{1964-35.3}{F}{Small-parameter methods}
\GSfield{Kind; source}{remark; \GSpage{P000063}; A. Ya. Povzner}
\GSfield{The notes}{Povzner names problems involving a small parameter.}
\GSfield{Statement}{No unique mathematical statement can be extracted: the equations, limiting process and desired estimates are absent.}
\GSfield{Fate}{Small-parameter and singular-perturbation searches identify a field, not a unique problem or publication.}
\GSfield{Present status}{As of November 2026, the fragment remains unidentifiable from the supplied evidence; this does not establish that a mathematical problem is open.}
\GSfield{Assessment}{literal statement: not determinable; publication match: none found; verification: unresolved lead; development: annotated; confidence: low}
\end{GSitem}
\begin{GSitem}{1964-35.4}{F}{Special functions of several variables}
\GSfield{Kind; source}{remark; \GSpage{P000063}; F. A. Berezin}
\GSfield{The notes}{Berezin proposes multivariable functions connected to representations and statistical physics; a Povzner addendum follows.}
\GSfield{Statement}{No unique mathematical statement can be extracted: no specific family, representation, equation or classification target is supplied.}
\GSfield{Fate}{Berezin, multivariable special-function and representation searches yield thematic literature, but no exact item can be selected.}
\GSfield{Present status}{As of November 2026, the fragment remains unidentifiable from the supplied evidence; this does not establish that a mathematical problem is open.}
\GSfield{Assessment}{literal statement: not determinable; publication match: none found; verification: unresolved lead; development: annotated; confidence: low}
\end{GSitem}

\GStalk{36}{1965-03-08}{Ya. G. Sinai}{P000062-P000066}{Moduli of ergodic eigenfunctions}
The dated Sinai fragment on P000063 states that a Koopman eigenfunction of an ergodic system has constant modulus. The surrounding pages continue earlier dynamical-system material and begin other speakers; the catalog envelope is retained without attributing all its contents to this fragment.

\begin{GSitem}{1964-36.1}{E}{Ergodic eigenfunctions have constant modulus}
\GSfield{Kind; source}{theorem; \GSpage{P000063}}
\GSfield{The notes}{The notes give \(U_t f_\lambda=e^{i\lambda t}f_\lambda\) and \ru{пост. по mod (у эргод. системы)}.}
\GSfield{Statement}{Let \((X,\mu)\) be a probability space and \((T_t)\) an ergodic measure-preserving measurable flow. If \(0\ne f\in L^2(\mu)\) and \(f\circ T_t=e^{it\lambda}f\) for every \(t\), as equalities in \(L^2\), then \(|f|=c>0\) almost everywhere. More generally the same statement holds for an ergodic action and a unitary character.}
\GSfield{Reconstruction}{The probability-space convention and almost-everywhere qualification make the shorthand precise.}
\GSfield{Fate}{This is an immediate consequence of ergodicity: \(|f|\) is invariant. It is used explicitly, in the normalization \(|f|=1\), by \citet[p.~775]{stepin1966spectra}. That 1966 use is corroboration, not a claim that the elementary observation originated then.}
\GSfield{Present status}{As of November 2026: established with the hypotheses in the modern statement.}
\GSfield{Assessment}{literal statement: correct; publication match: exact; verification: read the result; development: annotated; confidence: high}
\end{GSitem}

\GStalk{37}{1965-03-29}{V. N. Latyshev}{P000064-P000067}{Enveloping algebras, Lie ideals and bounded nilpotence}
The speaker, named only Viktor Nikolaevich in the handwriting, discusses primitive universal enveloping algebras, the passage from associative to Lie ideals, and a Burnside-type nilpotence assertion attributed to Kostrikin. The catalog identification as Latyshev is retained. Two assertions require correction or cannot be identified with a theorem under the recorded hypotheses.

\begin{GSitem}{1964-37.1}{B}{When is a universal enveloping algebra primitive?}
\GSfield{Kind; source}{problem; \GSpage{P000065}}
\GSfield{The notes}{The notes ask when \(U(L)\) is primitive and explain this as admitting a faithful irreducible module; Harish-Chandra and Jacobson are named.}
\GSfield{Statement}{For a specified Lie algebra \(\mathfrak g\) over a field \(k\), determine whether \(U(\mathfrak g)\) has a faithful simple left module. For finite-dimensional complex \(\mathfrak g\), the Dixmier--Moeglin equivalence gives a precise criterion: the zero ideal is primitive if and only if the center of the quotient division ring of \(U(\mathfrak g)\) is \(\mathbb C\). The restriction on the field and dimension is essential to the stated scope.}
\GSfield{Reconstruction}{Finite dimension, characteristic zero and the base field are not recorded; the complex finite-dimensional case is a reconstruction. No claim is made that all \(U(\mathfrak g)\) are primitive.}
\GSfield{Fate}{\citet[p.~2020]{bell2017poisson} explicitly records the primitive--rational--locally-closed equivalence for enveloping algebras of finite-dimensional complex Lie algebras, crediting Dixmier (1977) and Moeglin (1980). Applied to the zero prime ideal, it yields the criterion stated here. This later primary account was read; the original proofs were not. \citet{ooms1974primitive} is a separate, inaccessible lead on solvable Lie algebras and is not needed to certify the criterion.}
\GSfield{Present status}{As of November 2026: the finite-dimensional complex problem belongs to an established primitive-ideal theory and admits the stated criterion. The notebook omits the field and dimension, so an unrestricted field-independent classification is not asserted.}
\GSfield{Assessment}{literal statement: correct with missing hypotheses; publication match: partial; verification: read the result; development: annotated; confidence: high; problem \hyperref[prob:GP-1964-10]{GP-1964-10}}
\end{GSitem}
\begin{GSitem}{1964-37.2}{E}{Two two-dimensional examples of primitive envelopes}
\GSfield{Kind; source}{construction; \GSpage{P000065}}
\GSfield{The notes}{The nonabelian relation is \([e_1,e_2]=e_1\); its envelope is called primitive. The abelian two-generator envelope is said not to be primitive because a commutative primitive ring is a field.}
\GSfield{Statement}{Over \(\mathbb C\), the two-dimensional nonabelian Lie algebra \([x,h]=x\) has primitive enveloping algebra. In contrast \(U(\mathbb C^2)=\mathbb C[x,h]\) is not primitive: every simple module of a commutative ring is a residue field, and faithfulness forces the ring itself to be a field.}
\GSfield{Reconstruction}{The faithful simple module is an editorial verification, not a reconstruction of an omitted notebook proof. The base field is fixed to \(\mathbb C\).}
\GSfield{Fate}{The nonabelian case can be checked directly on \(\mathbb C[t]\): let \(h\) multiply by \(t\) and \(x f(t)=f(t+1)\), so \([x,h]=x\). A submodule is an ideal stable under translation, hence is zero or the whole ring; distinct finite translations are linearly independent over \(\mathbb C[t]\), giving faithfulness. This elementary verification agrees with the later primitive-ideal framework \citep[p.~2020]{bell2017poisson}.}
\GSfield{Present status}{As of November 2026: established with the hypotheses in the modern statement.}
\GSfield{Assessment}{literal statement: correct with missing hypotheses; publication match: exact; verification: read the result; development: annotated; confidence: high}
\end{GSitem}
\begin{GSitem}{1964-37.3}{F}{Associative simplicity does not eliminate Lie ideals}
\GSfield{Kind; source}{claim; \GSpage{P000065}}
\GSfield{The notes}{Under the name Baxter the manuscript says \ru{Если A --- простая, то Ли-ид. тоже нет}, and repeats that absence of ordinary ideals implies absence of Lie ideals.}
\GSfield{Statement}{The literal assertion is false. In the simple associative algebra \(M_2(\mathbb C)\), the center \(\mathbb C I\) is a proper nonzero Lie ideal; so is \(\mathfrak{sl}_2(\mathbb C)\). Correct simplicity theorems concern specified commutator Lie algebras, quotients by their centers, or skew elements of an algebra with involution, with characteristic and small-dimension exceptions.}
\GSfield{Reconstruction}{The corrected theorem is intentionally not guessed: the notebook omits the quotient or involution that would determine which result was intended.}
\GSfield{Fate}{Baxter\textquotesingle s paper \citet[title and theorem scope]{baxter1956lie} concerns a special class of associative rings; its topic cannot validate the unrestricted sentence. The direct matrix counterexamples already settle the literal reading. No new open problem is inferred from the error.}
\GSfield{Present status}{As of November 2026: false literally; the intended Baxter theorem remains incompletely identified.}
\GSfield{Assessment}{literal statement: false as written; publication match: thematic only; verification: indirect; development: annotated; confidence: high}
\end{GSitem}
\begin{GSitem}{1964-37.4}{F}{Uniform nilpotence on a generating Lie algebra}
\GSfield{Kind; source}{claim; \GSpage{P000066}}
\GSfield{The notes}{The Burnside/Kostrikin paragraph has finite generators \(X\), their Lie algebra \([X]\), and a corrected quantifier \(\exists n\) followed by \(e^n=0\). It concludes that the associative algebra is \ru{нильпот.}, glossed as every element being nilpotent.}
\GSfield{Statement}{The strongest recoverable question is: if a finitely generated associative algebra \(A\) over an unspecified field is generated by a finite set \(X\), and every element of the Lie subalgebra generated by \(X\) has a common associative nilpotence bound, what follows about \(A\)? Being nil (each element nilpotent) and being nilpotent (\(A^m=0\) for one \(m\)) are different conclusions.}
\GSfield{Reconstruction}{The membership sign \(e\in[X]\) and universal quantification over \(e\) are interpretive expansions. The field, associativity and distinction between nil and nilpotent are not supplied by the notes.}
\GSfield{Fate}{Searches on Latyshev, Kostrikin, the restricted Burnside problem, Lie-generated associative algebras and bounded nilpotence did not identify a theorem with precisely these assumptions and conclusion. A citation to a restricted-Burnside theorem alone would not supply the omitted characteristic and identity hypotheses.}
\GSfield{Present status}{As of November 2026: insufficiently specified for a reliable fate assignment; not classified as an undeveloped conjecture.}
\GSfield{Assessment}{literal statement: not determinable; publication match: none found; verification: unresolved lead; development: annotated; confidence: low}
\end{GSitem}

\GStalk{38}{1965-03-29}{I. M. Gelfand}{P000065-P000067}{Weyl algebras with different numbers of canonical pairs}
A short intervention attributes to Dixmier the question whether the differential-operator algebras in one and in two variables are isomorphic. This is a comparison of different ranks. Gelfand and Kirillov answered such isomorphism questions in 1966 by an invariant of growth.

\begin{GSitem}{1964-38.1}{A}{Weyl algebras of different ranks}
\GSfield{Kind; source}{question; \GSpage{P000066}}
\GSfield{The notes}{The two rings are written \(P(x,d/dx)\) and \(P(x,d/dx,y,d/dy)\), followed by \ru{Изоморфны ли эти кольца?}.}
\GSfield{Statement}{Over a field \(k\) of characteristic zero, let \(A_n(k)\) be the unital associative \(k\)-algebra generated by \(x_1,\ldots,x_n,\partial_1,\ldots,\partial_n\), with relations \[ [x_i,x_j]=[\partial_i,\partial_j]=0,\qquad [\partial_i,x_j]=\delta_{ij}1\quad(1\le i,j\le n). \] Are \(A_1(k)\) and \(A_2(k)\) isomorphic as \(k\)-algebras? They are not: more generally, \(A_m(k)\cong A_n(k)\) implies \(m=n\).}
\GSfield{Reconstruction}{The base field and preservation of it are made explicit. The scan decisively shows one canonical pair versus two.}
\GSfield{Fate}{\citet[\S4, Lemma~5 and Theorem~2, pp.~11--14]{gelfandkirillov1966corps} distinguish the algebras \(R_{n,k}\), including Weyl algebras, by growth. The paper was received 19 January 1966 and published that year. In modern terms \(\operatorname{GKdim}A_n=2n\). This is unrelated in quantifiers to asking whether every endomorphism of one fixed \(A_n\) is an automorphism.}
\GSfield{Present status}{As of November 2026: solved negatively for isomorphism of distinct ranks; the same-year first index must not label this an open fixed-rank endomorphism conjecture.}
\GSfield{Assessment}{literal statement: correct; publication match: exact; verification: read the result; development: annotated; confidence: high; problem \hyperref[prob:GP-1964-11]{GP-1964-11}}
\end{GSitem}

\GStalk{39}{1965-03-29}{S. I. Gelfand}{P000065-P000070}{Dyson circular ensembles and the logarithmic gas}
The talk explains the circular orthogonal and unitary ensembles, their eigenangle densities and a Coulomb-gas interpretation of an empty arc. It then moves from an exact finite-dimensional integral to a continuum free-energy ansatz. The manuscript fails to keep that approximation separate from an exact gap formula.

\begin{GSitem}{1964-39.1}{E}{The circular orthogonal ensemble as U(N)/O(N)}
\GSfield{Kind; source}{construction; \GSpage{P000066}}
\GSfield{The notes}{Symmetric unitary matrices satisfy \(U^T=U\), \(UU^*=I\), with action \(U\mapsto VUV^T\); the stabilizer at the identity is orthogonal. The eigenangle density has Vandermonde power one.}
\GSfield{Statement}{The pushforward of Haar measure under \(V\mapsto VV^T\) on \(U(N)\) is the invariant probability measure on \(U(N)/O(N)\). Its unordered eigenangles have joint density \(Z_{N,1}^{-1}\prod_{i<j}|e^{i\theta_i}-e^{i\theta_j}|\) relative to Lebesgue measure on \([0,2\pi)^N\), with the ordering convention absorbed in \(Z_{N,1}\).}
\GSfield{Fate}{This is Dyson\textquotesingle s circular orthogonal ensemble, published in 1962 after receipt in September 1961; see \citet[the circular ensembles and eigenvalue measures]{dyson1962statisticalI}.}
\GSfield{Present status}{As of November 2026: established with the hypotheses in the modern statement.}
\GSfield{Assessment}{literal statement: correct; publication match: exact; verification: read the result; development: annotated; confidence: high}
\end{GSitem}
\begin{GSitem}{1964-39.2}{E}{Haar unitary eigenangles and quadratic repulsion}
\GSfield{Kind; source}{construction; \GSpage{P000066}}
\GSfield{The notes}{For all unitary matrices the notes display Vandermonde power two and discuss integration over conjugacy classes. In the action formula the final \(V^*\) is crossed or omitted.}
\GSfield{Statement}{For Haar-distributed \(U\in U(N)\), the eigenangle density is \(Z_{N,2}^{-1}\prod_{i<j}|e^{i\theta_i}-e^{i\theta_j}|^2\). Haar invariance is under left and right multiplication; conjugation \(U\mapsto VUV^*\) is the action relevant to eigenvalue class integration.}
\GSfield{Reconstruction}{Left invariance and conjugacy invariance are separated instead of treating the incomplete action as a new ensemble.}
\GSfield{Fate}{The circular unitary ensemble and its eigenangle law are already in \citet[unitary ensemble]{dyson1962statisticalI}. Coinciding angles have vanishing density, the precise content of the recorded repulsion remark.}
\GSfield{Present status}{As of November 2026: established with the hypotheses in the modern statement.}
\GSfield{Assessment}{literal statement: correct with missing hypotheses; publication match: exact; verification: read the result; development: annotated; confidence: high}
\end{GSitem}
\begin{GSitem}{1964-39.3}{E}{Circular beta gap integrals and logarithmic interaction}
\GSfield{Kind; source}{construction; \GSpage{P000066-P000068}}
\GSfield{The notes}{The forbidden arc is \((-\alpha,\alpha)\); the normalized probability is eventually written \(R_\beta=\Psi_\beta(\alpha)/\Psi_\beta(0)\). The first energy sign conflicts with the later Boltzmann factor.}
\GSfield{Statement}{For \(N\ge1,\beta>0,0\le\alpha<\pi\), define \(\Psi_{N,\beta}(\alpha)=\int_{[\alpha,2\pi-\alpha]^N}\prod_{i<j}|e^{i\theta_i}-e^{i\theta_j}|^\beta\,d\theta\). Then \(R_{N,\beta}(\alpha)=\Psi_{N,\beta}(\alpha)/\Psi_{N,\beta}(0)\) is the no-eigenangle probability. With energy \(W=-\sum_{i<j}\log|e^{i\theta_i}-e^{i\theta_j}|\), the density is proportional to \(e^{-\beta W}\).}
\GSfield{Reconstruction}{The sign convention, product over pairs and normalization are supplied explicitly; \(\beta>0\) is a sufficient integrability range, not the maximal possible range.}
\GSfield{Fate}{Dyson\textquotesingle s log-gas construction is exact at the finite integral level; see \citet[statistical-mechanical interpretation]{dyson1962statisticalI}. The source\textquotesingle s missing product and normalization are repaired by its own next page.}
\GSfield{Present status}{As of November 2026: established with the hypotheses in the modern statement.}
\GSfield{Assessment}{literal statement: false as written; publication match: exact; verification: read the result; development: annotated; confidence: high}
\end{GSitem}
\begin{GSitem}{1964-39.4}{E}{Continuum density and local free-energy approximation}
\GSfield{Kind; source}{method; \GSpage{P000067-P000068}}
\GSfield{The notes}{The notes replace the point density by \(\sigma\), impose \(\int\sigma=N\), and write \(F=V+F_1\) with a logarithmic double integral and a local term \(\int\sigma f_\beta(\sigma)\). They label the construction a physical assumption.}
\GSfield{Statement}{The empirical measure is \(\sum_{j=1}^N\delta_{\theta_j}\), not an ordinary density. A continuum ansatz replaces it by \(\sigma(\theta)\,d\theta\), minimizes a logarithmic-energy functional under the mass and support constraints, and adds a local entropy correction. This does not give an exact finite-\(N\) theorem. If \(F=-\beta^{-1}\log\Psi\), then \(R=e^{-\beta(F_\alpha-F_0)}\).}
\GSfield{Reconstruction}{The distinction between an empirical measure and a density, and between an ansatz and a theorem, is editorial. The source also permits its positive-sign exponential to define an opposite convention for \(F\); it cannot simultaneously be the usual minimized free energy without a sign change.}
\GSfield{Fate}{The method is the 1962 Dyson asymptotic program, not an independently established exact identity. Later determinant asymptotics such as \citet[Introduction and the large-gap formulas]{ehrhardt2007dyson} separate leading energy, logarithmic terms and constants.}
\GSfield{Present status}{As of November 2026: established with the hypotheses in the modern statement.}
\GSfield{Assessment}{literal statement: correct with missing hypotheses; publication match: partial; verification: read the result; development: annotated; confidence: high}
\end{GSitem}
\begin{GSitem}{1964-39.5}{B}{The cosine-power expression is not an exact CUE gap probability}
\GSfield{Kind; source}{claim; \GSpage{P000067}}
\GSfield{The notes}{The manuscript writes \(R_2(\alpha)=\exp(N^2\log\cos(\alpha/2))\) as an equality.}
\GSfield{Statement}{The equality \(R_{N,2}(\alpha)=\cos(\alpha/2)^{N^2}\) is false for the stated circular unitary ensemble. For \(N=1\), the eigenangle is uniform and \(R_{1,2}(\alpha)=1-\alpha/\pi\), which is not \(\cos(\alpha/2)\). The cosine-power captures a leading energy contribution in an appropriate large-\(N\) regime, not the complete finite probability.}
\GSfield{Reconstruction}{The one-eigenvalue check is an editorial diagnosis. No unrecorded special ensemble is inserted to save the equality.}
\GSfield{Fate}{The later asymptotic theory includes essential subleading powers and constants. \citet[Introduction]{ehrhardt2007dyson} verifies the distinction in the sine-kernel limit; it is not a proof of the notebook\textquotesingle s finite equality. The same expression returns in item 1964-42.1.}
\GSfield{Present status}{As of November 2026: the finite identity is disproved by \(N=1\); precise large-gap asymptotics are established.}
\GSfield{Assessment}{literal statement: false as written; publication match: partial; verification: read the result; development: annotated; confidence: high}
\end{GSitem}

\GStalk{40}{1965-04-05}{I. N. Bernstein}{P000068-P000073}{Grothendieck topologies and arithmetic cohomology}
Beginning with finite-field solutions as Frobenius fixed points, Bernstein introduces slice categories, products, coverings and sheaf descent. The title mentions monoidal transformations, but no actual blow-up construction is developed in this fragment. The category of sheaves over a field anticipates Galois cohomology.

\begin{GSitem}{1964-40.1}{B}{Frobenius fixed points and the cohomological counting program}
\GSfield{Kind; source}{method; \GSpage{P000069}}
\GSfield{The notes}{Reduction modulo \(p\) is denoted \(Z_p\). Over its algebraic closure, \(x\mapsto x^p\) fixes exactly the base field; a Lefschetz alternating trace is proposed to count rational points.}
\GSfield{Statement}{For a finite-type scheme \(X/\mathbb F_q\), \(X(\mathbb F_q)\) is the fixed-point set of the geometric Frobenius action on geometric points. The precise cohomological formula is \(\#X(\mathbb F_q)=\sum_i(-1)^i\operatorname{Tr}(\operatorname{Frob}_q\mid H^i_c(X_{\overline{\mathbb F}_q},\mathbb Q_\ell))\), with \(\ell\ne\operatorname{char}\mathbb F_q\) and compatible Frobenius convention.}
\GSfield{Reconstruction}{Finite type, compact supports, \(\ell\)-adic coefficients and the Frobenius convention are later precise reconstruction, not statements supplied in the notes.}
\GSfield{Fate}{The notebook describes the Weil--Grothendieck program rather than a usable topological Lefschetz theorem in the Zariski topology. The étale-site framework was circulating in the 1963--64 SGA seminar; the published source is \citet[Exposés I--IV, site and topos foundations]{sga4I1972}. The full trace formula requires additional étale-cohomological finiteness and trace theorems.}
\GSfield{Present status}{As of November 2026: the arithmetic trace formula is established; the recorded ordinary-topological analogy alone is insufficient.}
\GSfield{Assessment}{literal statement: correct with missing hypotheses; publication match: partial; verification: publisher metadata checked; development: annotated; confidence: high}
\end{GSitem}
\begin{GSitem}{1964-40.2}{E}{Slice categories and relative products}
\GSfield{Kind; source}{definition; \GSpage{P000069-P000070}}
\GSfield{The notes}{The objects over \(X\) are arrows \(U\to X\); a morphism is a commuting triangle. The universal property of products is explained, with open-set intersection and pairs having the same image as examples.}
\GSfield{Statement}{For a category \(\mathcal C\), the slice \(\mathcal C/X\) has objects \(U\to X\). If \(U\times_X V\) exists, it is the product of \(U\to X\) and \(V\to X\) in this slice. Its universal property gives uniqueness up to a unique isomorphism. For the poset of opens of a space, this is intersection.}
\GSfield{Fate}{This is the categorical language used in SGA~4; see \citet[Exposé I]{sga4I1972} and the explicit modern site conventions of \citet[Tag 00VH]{stacks2026sites}. Publication in 1972 must not be mistaken for first circulation: the seminar dates to 1963--64.}
\GSfield{Present status}{As of November 2026: established with the hypotheses in the modern statement.}
\GSfield{Assessment}{literal statement: correct; publication match: exact; verification: read the result; development: annotated; confidence: high}
\end{GSitem}
\begin{GSitem}{1964-40.3}{E}{Covering families and Grothendieck topologies}
\GSfield{Kind; source}{definition; \GSpage{P000070-P000071}}
\GSfield{The notes}{The axioms require identity coverings, composition of covering families and preservation under pullback. Existence of every pairwise product in the slice is recorded separately.}
\GSfield{Statement}{A pretopology on a category with the requisite fiber products is a collection of covering families containing isomorphisms and stable under base change and composition. It generates a Grothendieck topology of covering sieves. A general Grothendieck topology does not require all fiber products in the underlying category; the notes use a convenient stronger presentation.}
\GSfield{Reconstruction}{The distinction between a pretopology of families and a topology of sieves is modern clarification.}
\GSfield{Fate}{This is established pre-lecture Grothendieck-site theory, circulated in SGA~4 (1963--64) and published in \citet[Exposé II]{sga4I1972}. \citet[Tag 00VH and covering axioms]{stacks2026sites} supplies a checked current formulation.}
\GSfield{Present status}{As of November 2026: established with the hypotheses in the modern statement.}
\GSfield{Assessment}{literal statement: correct with missing hypotheses; publication match: exact; verification: read the result; development: annotated; confidence: high}
\end{GSitem}
\begin{GSitem}{1964-40.4}{E}{Sheaf descent and the point represented by a field}
\GSfield{Kind; source}{definition; \GSpage{P000071}}
\GSfield{The notes}{A presheaf is a contravariant abelian-group-valued functor. Compatible sections over a covering glue uniquely. The notes add that a point can mean a field and then give Galois cohomology.}
\GSfield{Statement}{A sheaf \(F\) on a site with fiber products satisfies the equalizer condition \(F(U)\to\prod_iF(U_i)\rightrightarrows\prod_{i,j}F(U_i\times_U U_j)\). On the small étale site of \(\operatorname{Spec}k\), abelian sheaves correspond to discrete continuous modules for \(G_k=\operatorname{Gal}(k^{\mathrm{sep}}/k)\); sheaf cohomology identifies with continuous Galois cohomology.}
\GSfield{Reconstruction}{The small étale topology and continuity of the Galois action are missing but decisive hypotheses.}
\GSfield{Fate}{See \citet[Exposé II]{sga4I1972} for the foundations and \citet[Tags 00VL and 03QR]{stacks2026sites} for descent and the field case. The Zariski site of a one-point space would not produce this Galois information.}
\GSfield{Present status}{As of November 2026: established with the hypotheses in the modern statement.}
\GSfield{Assessment}{literal statement: correct with missing hypotheses; publication match: exact; verification: read the result; development: annotated; confidence: high}
\end{GSitem}

\GStalk{41}{1965-04-12}{I. M. Gelfand}{P000071-P000073}{A curriculum and a symmetric-space research direction}
This is chiefly a planning intervention: subjects needed by seminar participants are listed, alongside a suggestion to extend Dyson\textquotesingle s calculations to symmetric spaces. The item index preserves both the organizational list and the mathematical direction.

\begin{GSitem}{1964-41.1}{F}{Subjects proposed for study}
\GSfield{Kind; source}{remark; \GSpage{P000072}}
\GSfield{The notes}{The list includes Lie groups and representations, sheaves and homological algebra, bundles, algebraic geometry, quantum mechanics, several-variable Cauchy theory and nuclear spaces with measures.}
\GSfield{Statement}{This is a curriculum proposal, not a theorem or a mathematically quantified research problem.}
\GSfield{Fate}{No publication fate can be assigned to a list of study topics. Its value is evidence of the seminar\textquotesingle s program and links to records 39, 40, 48 and 49.}
\GSfield{Present status}{As of November 2026: historical planning evidence; no open mathematical status is applicable.}
\GSfield{Assessment}{literal statement: not determinable; publication match: none found; verification: indirect; development: none; confidence: high}
\end{GSitem}
\begin{GSitem}{1964-41.2}{B}{Dyson-type calculations on symmetric spaces}
\GSfield{Kind; source}{problem; \GSpage{P000072}}
\GSfield{The notes}{Gelfand proposes investigating Dyson\textquotesingle s results for general symmetric spaces.}
\GSfield{Statement}{One precise direction is to push invariant measure on a specified compact symmetric space to a maximal flat and compute the Jacobian in restricted-root coordinates, then study eigenvalue or radial gap observables. A space, normalization, root multiplicities and observable must be specified before an exact theorem can be stated.}
\GSfield{Reconstruction}{The restricted-root formulation and choice of radial observables are an explanatory modern realization, not an assertion that the notebook posed a complete classification problem.}
\GSfield{Fate}{Dyson\textquotesingle s \(U(N)/O(N)\) example is already a compact symmetric space \citep{dyson1962statisticalI}. Later random-matrix symmetry classifications systematically use symmetric spaces \citep[the symmetry-class correspondence]{altlandzirnbauer1997}. This matches the direction, not a single universal gap formula.}
\GSfield{Present status}{As of November 2026: extensively developed as a framework; the short suggestion is too broad to carry an unresolved problem label.}
\GSfield{Assessment}{literal statement: correct with missing hypotheses; publication match: partial; verification: read the result; development: annotated; confidence: medium}
\end{GSitem}

\GStalk{42}{1965-04-12}{S. I. Gelfand}{P000071-P000073}{Circular gap probabilities and an even sine-kernel determinant}
The discussion revisits the cosine-power claim, relates a nearest-neighbor law to derivatives of an empty-interval probability and displays a cosine determinant with a large-gap prediction. Its written exponent and its scaling parameter require explicit critical annotation.

\begin{GSitem}{1964-42.1}{F}{Repeated cosine-power claim and an unnamed replacement model}
\GSfield{Kind; source}{claim; \GSpage{P000072}}
\GSfield{The notes}{The speaker says the cosine-power result has not been proved for the original model but another similar model gives it exactly. That replacement model is not specified.}
\GSfield{Statement}{For the original finite circular unitary ensemble, the asserted exact identity is false as in item 1964-39.5. No exact alternative statement can be recovered without the replacement ensemble and its probability measure.}
\GSfield{Fate}{Dyson\textquotesingle s asymptotic program and subsequent sine-kernel work explain the leading exponential, but do not identify the unrecorded substitute ensemble \citep[Introduction]{ehrhardt2007dyson}. The two claims must be separated: a false literal finite-CUE equality and an unidentifiable possible exact identity for another model.}
\GSfield{Present status}{As of November 2026: the original equality is false; the unnamed model remains unidentified. This is not evidence of an open exact-CUE problem.}
\GSfield{Assessment}{literal statement: not determinable; publication match: thematic only; verification: read the result; development: annotated; confidence: low}
\end{GSitem}
\begin{GSitem}{1964-42.2}{E}{Nearest-neighbor spacing from the empty-interval probability}
\GSfield{Kind; source}{method; \GSpage{P000072}}
\GSfield{The notes}{A translation-invariant two-endpoint expression is written as \(S(\theta_2-\theta_1)\), and the spacing density is related to a second derivative of the empty-interval probability.}
\GSfield{Statement}{For a stationary simple point process on \(\mathbb R\) of intensity \(\rho>0\), let \(E(s)=\Pr\{N(0,s)=0\}\). Under the usual differentiability assumptions, the density \(p(s)\) of the forward nearest-neighbor distance under the Palm law is \(p(s)=E^{\prime\prime}(s)/\rho\). At unit intensity this is \(p=E^{\prime\prime}\). If the gap is parameterized by a half-width, chain-rule factors must also be included.}
\GSfield{Reconstruction}{The Palm nearest-neighbor law, intensity and differentiability hypotheses identify precisely which spacing distribution is meant.}
\GSfield{Fate}{This is the standard gap-to-spacing mechanism in Dyson\textquotesingle s statistical theory \citep[spacing distributions]{dyson1962statisticalIII}. The notebook omits the density normalization; a two-point correlation function is not in general a nearest-neighbor density.}
\GSfield{Present status}{As of November 2026: established with the hypotheses in the modern statement.}
\GSfield{Assessment}{literal statement: correct with missing hypotheses; publication match: exact; verification: indirect; development: annotated; confidence: high}
\end{GSitem}
\begin{GSitem}{1964-42.3}{B}{Even sine-kernel determinant and the large-gap exponent}
\GSfield{Kind; source}{claim; \GSpage{P000073}}
\GSfield{The notes}{The finite matrix is \(\delta_{pq}-\pi^{-1}\int_{-\alpha}^{\alpha}\cos(p\theta)\cos(q\theta)\,d\theta\), with half-integer \(p,q\). The scan shows an asymptotic factor \(t^{17/8}e^{-t^2/4-t/2}\) and \(t=\alpha n\).}
\GSfield{Statement}{For \(K_+(x,y)=\sin(x-y)/(\pi(x-y))+\sin(x+y)/(\pi(x+y))\) on \(L^2(0,s)\), the established asymptotic is \(\log\det(I-K_+)=-s^2/4-s/2-(1/8)\log s+c_++o(1)\) as \(s\to\infty\). Thus the power in this standard normalization is \(-1/8\), not \(17/8\).}
\GSfield{Reconstruction}{The continuum kernel is identified separately from the finite matrix. The power is a faithful manuscript error, while the exact \(m,n,t\) normalization remains an ambiguity; no silent rescaling is made.}
\GSfield{Fate}{\citet[Introduction, formulas for \(K_\pm\)]{ehrhardt2007dyson} proves the large-gap constant as well as the displayed terms (preprint April 2006; publication 2007). The finite cosine matrix is a precursor of this even kernel. The manuscript relation between \(2m\) and \(n\) is itself uncertain (equality or inequality); the natural continuum scaling of the matrix uses \(m\alpha\). The source\textquotesingle s \(t=\alpha n\) is therefore not silently identified with the standard \(s\).}
\GSfield{Present status}{As of November 2026: the standard even-kernel asymptotic is established; the manuscript power and unexplained scaling cannot be used literally.}
\GSfield{Assessment}{literal statement: false as written; publication match: partial; verification: read the result; development: annotated; confidence: high}
\end{GSitem}

\GStalk{43}{1965-04-12}{A. M. Stepin}{P000071-P000074}{A torsion-group action with continuous spectrum}
Stepin answers the preceding dynamical-system discussion with an action of the dyadic torsion group. Finite-order transformations can jointly have a continuous spectral component. His 1966 paper supplies an exact continuation and stronger spectral analysis.

\begin{GSitem}{1964-43.1}{A}{Continuous spectral component for a dyadic torsion-group action}
\GSfield{Kind; source}{construction; \GSpage{P000072}}
\GSfield{The notes}{The group is the dyadic rationals modulo one; the generators satisfy \(T_{1/2^{k+1}}^2=T_{1/2^k}\). The notes describe rearrangements of interval pieces and claim a continuous spectral component.}
\GSfield{Statement}{There is an ergodic probability-preserving action of the discrete Prüfer group \(\mathbb Z[1/2]/\mathbb Z\) whose Koopman representation has a nonzero continuous spectral component, although each individual group element has finite order. Continuous means non-atomic spectral measure on the compact dual group, not continuous time and not necessarily Lebesgue measure.}
\GSfield{Reconstruction}{The discrete topology on the group, joint spectral measure and distinction between continuous and absolutely continuous spectrum are made explicit.}
\GSfield{Fate}{\citet[\S2, pp.~774--776]{stepin1966spectra} constructs precisely this group action on a circle times a two-point space. Its even subspace has discrete spectrum; its odd subspace has simple singular continuous spectrum. The paper was presented 23 November 1965 and published in 1966.}
\GSfield{Present status}{As of November 2026: the announced construction and its singular-continuous refinement are established.}
\GSfield{Assessment}{literal statement: correct; publication match: exact; verification: read the result; development: annotated; confidence: high}
\end{GSitem}

\GStalk{44}{1965-04-12}{A. A. Kirillov}{P000072-P000074}{Recovering convex bodies from central sections and shadows}
Kirillov explains how central section areas of a symmetric body reduce to integration of an even function over great spheres. He then asks the analogous question for projection areas. The symmetry assumption is essential for both uniqueness statements.

\begin{GSitem}{1964-44.1}{E}{Central section volumes determine an origin-symmetric body}
\GSfield{Kind; source}{theorem; \GSpage{P000074}}
\GSfield{The notes}{The body is explicitly centrally symmetric. In three dimensions the notes replace it by \(\rho^2/2\) on \(S^2\) and invoke spherical harmonics to invert integration over great circles.}
\GSfield{Statement}{If \(K,L\subset\mathbb R^n\), \(n\ge2\), are origin-symmetric star bodies with positive continuous radial functions and \(\operatorname{vol}_{n-1}(K\cap u^\perp)=\operatorname{vol}_{n-1}(L\cap u^\perp)\) for all \(u\in S^{n-1}\), then \(K=L\). The section function is the spherical Radon transform of \(\rho_K^{n-1}/(n-1)\); the transform is injective on even functions. In \(n=3\), its spherical-harmonic multiplier at degree \(2j\) is \(2\pi P_{2j}(0)\ne0\), while odd degrees are killed.}
\GSfield{Reconstruction}{Origin symmetry, continuous star bodies and the \((n-1)\)-dimensional measure generalize the explicit three-dimensional argument. Convexity is sufficient but not needed.}
\GSfield{Fate}{The great-circle transform is classical Funk theory, originating before the seminar \citep{funk1913flachen}. The modern geometric formulation is documented by \citet[section functions and geometric tomography]{gardner2026tomography}. The original Funk publication metadata was checked; the general-dimensional formulation is a modern reconstruction.}
\GSfield{Present status}{As of November 2026: established with the hypotheses in the modern statement.}
\GSfield{Assessment}{literal statement: correct with missing hypotheses; publication match: partial; verification: publisher metadata checked; development: annotated; confidence: high}
\end{GSitem}
\begin{GSitem}{1964-44.2}{E}{Do all projection areas determine a convex body?}
\GSfield{Kind; source}{question; \GSpage{P000074}}
\GSfield{The notes}{The final question asks whether the areas of all shadows determine the convex body, followed only by \ru{Сводится}.}
\GSfield{Statement}{For \(n\ge2\), the brightness function \(b_K(u)=\operatorname{vol}_{n-1}(K\mid u^\perp)\) determines an origin-symmetric convex body \(K\subset\mathbb R^n\) uniquely among such bodies. Without a fixed center, translation is invisible; without central symmetry, uniqueness fails, already because \(K\) and \(-K\) have the same brightness.}
\GSfield{Reconstruction}{The inherited central-symmetry assumption is made explicit; if the last line is read independently for arbitrary convex bodies, its affirmative answer is false.}
\GSfield{Fate}{Aleksandrov\textquotesingle s projection theorem dates to 1937; the precise symmetric uniqueness statement and date are given by \citet[Reconstruction from brightness functions; reference to Theorem~3.3.1]{gardner2026tomography}. This is an old theorem under the symmetry convention of the preceding paragraph, not an unresolved projection problem.}
\GSfield{Present status}{As of November 2026: affirmative with origin symmetry, negative as an unrestricted uniqueness claim.}
\GSfield{Assessment}{literal statement: correct with missing hypotheses; publication match: exact; verification: read the result; development: annotated; confidence: high}
\end{GSitem}

\GStalk{45}{1965-04-12}{D. A. Kazhdan}{P000072-P000074}{Dvoretzky sections and an incompletely recorded Gromov method}
Kazhdan reports a result of Gromov in the setting of Dvoretzky\textquotesingle s theorem. Sphere bundles over Grassmannians and disappearance of low spherical harmonics are proposed as a method. The exact additional hypothesis and representation-theoretic quantifiers are not preserved well enough to equate this with Gromov\textquotesingle s later Banach-conjecture paper.

\begin{GSitem}{1964-45.1}{E}{Dvoretzky almost-Euclidean sections}
\GSfield{Kind; source}{theorem; \GSpage{P000073}}
\GSfield{The notes}{For every \(\varepsilon,k\), a sufficiently high-dimensional centrally symmetric convex body has a \(k\)-dimensional section that is nearly a ball.}
\GSfield{Statement}{For every \(k\ge1\) and \(\varepsilon>0\), there is \(N(k,\varepsilon)\) such that every real normed space of dimension at least \(N\) has a \(k\)-dimensional subspace \(E\), a Euclidean norm \(|\cdot|\) on \(E\), and \(a>0\) with \(a|x|\le\|x\|\le(1+\varepsilon)a|x|\) for all \(x\in E\). The corresponding ellipsoid comparison is affine, not a demand for a round section in an arbitrarily fixed ambient metric.}
\GSfield{Reconstruction}{The norm comparison fixes the quantitative meaning of nearly a ball and allows a suitable Euclidean structure.}
\GSfield{Fate}{Dvoretzky\textquotesingle s result appeared in the proceedings of the 1960 Jerusalem symposium, published in 1961. \citet[p.~297, stated theorem]{figiel1976dvoretzky} gives a later short proof and identifies that original publication.}
\GSfield{Present status}{As of November 2026: established with the hypotheses in the modern statement.}
\GSfield{Assessment}{literal statement: correct with missing hypotheses; publication match: exact; verification: read the result; development: annotated; confidence: high}
\end{GSitem}
\begin{GSitem}{1964-45.2}{F}{A sphere-bundle and low-harmonic vanishing proposal}
\GSfield{Kind; source}{method; \GSpage{P000073}}
\GSfield{The notes}{The fiber consists of a unit vector in a \(k\)-plane. The notes claim that, for large \(n\), some fiber kills the first \(\ell\) nonconstant spherical components, then state a broader unitary-representation version. A condition on outer-to-inner radii is left unfinished; weaker and stronger appear in adjacent comments.}
\GSfield{Statement}{The incidence space \(M_{n,k}=\{(E,v):E\in G_k(\mathbb R^n),v\in E,|v|=1\}\) has fiber \(S^{k-1}\). A precise possible question concerns simultaneous vanishing of specified nonconstant harmonic projections of the restriction of a norm to one fiber. The notebook does not specify the representation, parity, regularity, dimension bounds or extra convex-body condition needed to turn its universal sentence into a theorem.}
\GSfield{Reconstruction}{The corrected fiber dimension is certain. A conjectural completion of the harmonic assertion would be editorial invention and is not supplied.}
\GSfield{Fate}{\citet[Theorem~1]{gromov1967banach}, received September 1966 and published 1967, proves results on normed spaces whose subspaces are mutually isometric. Its theorem is not the exact harmonic-vanishing assertion recorded here. Searches of the paper and the Dvoretzky/Grassmannian literature did not establish an exact match.}
\GSfield{Present status}{As of November 2026: the program is recognizable, but this particular theorem announcement remains unidentified; no claim of an orphan open problem is made.}
\GSfield{Assessment}{literal statement: not determinable; publication match: thematic only; verification: read the result; development: annotated; confidence: low}
\end{GSitem}

\GStalk{46}{1965-04-12}{R. A. Minlos}{P000072-P000075}{Gaussian Hermitian matrices, Brownian motion and angles}
The actual Minlos fragment is very short: Gaussian entries in a Hermitian matrix, Brownian motion, and a question about angles between subspaces. The modern interpretations are supplied explicitly as reconstructions, without attributing detailed stochastic equations to the handwriting.

\begin{GSitem}{1964-46.1}{E}{Hermitian matrix Brownian motion and eigenvalue diffusion}
\GSfield{Kind; source}{method; \GSpage{P000073}}
\GSfield{The notes}{The notes say only that the Hermitian matrix entries are Gaussian random variables and then \ru{Брауновское дв-ние}.}
\GSfield{Statement}{A standard realization takes independent real Brownian diagonal entries of variance \(dt\) and independent complex off-diagonal entries with real and imaginary variances \(dt/2\), reflected by Hermitian symmetry. Starting with distinct eigenvalues, its eigenvalues satisfy \(d\lambda_i=dB_i+\sum_{j\ne i}(\lambda_i-\lambda_j)^{-1}dt\) up to, and in fact without, collision. Other entry normalizations rescale time.}
\GSfield{Reconstruction}{The independent increments, variances and stochastic equation are the modern reconstruction. A single Gaussian matrix does not determine a time process.}
\GSfield{Fate}{Dyson\textquotesingle s Brownian-motion model was published in 1962 \citep{dyson1962brownian}. The title and metadata match this interpretation, but the short fragment itself does not distinguish free matrix Brownian motion from a Gaussian stationary Ornstein--Uhlenbeck variant.}
\GSfield{Present status}{As of November 2026: established with the hypotheses in the modern statement.}
\GSfield{Assessment}{literal statement: correct with missing hypotheses; publication match: partial; verification: indirect; development: annotated; confidence: medium}
\end{GSitem}
\begin{GSitem}{1964-46.2}{E}{Angles between two k-dimensional subspaces}
\GSfield{Kind; source}{question; \GSpage{P000073}}
\GSfield{The notes}{The question is simply what an angle between two \(k\)-dimensional subspaces in \(n\) dimensions means.}
\GSfield{Statement}{For subspaces \(E,F\) of a Euclidean or complex Hermitian space, with orthonormal column matrices \(Q_E,Q_F\), the singular values of \(Q_E^*Q_F\) are \(\cos\theta_1,\ldots,\cos\theta_r\), where \(r=\min(\dim E,\dim F)\) and \(0\le\theta_1\le\cdots\le\theta_r\le\pi/2\). These principal angles, including forced zeros from intersections, describe the relative position up to ambient isometry.}
\GSfield{Reconstruction}{The notes ask for an angle; the answer is generally a list of principal angles. The complex version and matrix expression are added for precision.}
\GSfield{Fate}{The theory goes back to \citet[\S\S IV--V; introduction, pp.~104--105]{jordan1875geometrie}, presented in May 1875. The singular-value expression is modern notation for this established collection of invariants.}
\GSfield{Present status}{As of November 2026: established with the hypotheses in the modern statement.}
\GSfield{Assessment}{literal statement: correct; publication match: exact; verification: read the result; development: annotated; confidence: high}
\end{GSitem}

\GStalk{47}{1965-04-19}{L. D. Faddeev}{P000073-P000077}{Characteristic determinants, Jost functions and regularized traces}
Faddeev moves from a finite-interval Sturm--Liouville characteristic function to half-line and whole-line Jost determinants, then to contour trace identities. The notes mix exact identities with formal diagonal-kernel reasoning. A final Gelfand question asks when traces of commutators vanish.

\begin{GSitem}{1964-47.1}{E}{Finite-interval characteristic function and resolvent trace}
\GSfield{Kind; source}{construction; \GSpage{P000074}}
\GSfield{The notes}{For \(L=-d^2/dx^2+q\) on \([0,\pi]\) with Neumann ends, \(\psi(0,\lambda)=1\), \(\psi_x(0,\lambda)=0\), and \(w(\lambda)=\psi_x(\pi,\lambda)\). The finite-matrix identity \(d\log w/d\lambda=-\operatorname{Tr}(L-\lambda)^{-1}\) is proposed.}
\GSfield{Statement}{For real \(q\in L^1(0,\pi)\) and the self-adjoint Neumann realization, \(w\) is entire, with zeros exactly the simple eigenvalues. Outside the spectrum, \(w^{\prime}(\lambda)/w(\lambda)=-\operatorname{Tr}(L-\lambda)^{-1}\). The derivative on \(w\) is with respect to \(\lambda\); the derivative on \(\psi\) is spatial. Multiplying \(w\) by a nonzero constant leaves the identity unchanged.}
\GSfield{Reconstruction}{Integrability, real potential, operator domain and the two different derivatives are made explicit.}
\GSfield{Fate}{This is classical one-dimensional spectral theory. The determinant-to-solution method belongs to the pre-seminar scattering and inverse-spectral framework of \citet[\S\S1--2]{faddeev1959inverse}; the finite-interval formula itself follows from the Green kernel and Lagrange identity. It is not a newly announced 1965 result.}
\GSfield{Present status}{As of November 2026: established with the hypotheses in the modern statement.}
\GSfield{Assessment}{literal statement: correct with missing hypotheses; publication match: partial; verification: read the result; development: annotated; confidence: high}
\end{GSitem}
\begin{GSitem}{1964-47.2}{E}{Half-line Jost function as a relative determinant}
\GSfield{Kind; source}{construction; \GSpage{P000074-P000075}}
\GSfield{The notes}{On the half-line, with \(y(0)=0\) and rapidly decaying \(q\), the resolvent has no ordinary trace. A free part is subtracted; the outgoing solution satisfies \(f(x,k)\sim e^{ikx}\), and \(W(\lambda)=f(0,\sqrt\lambda)\).}
\GSfield{Statement}{For real \(q\) with \(\int_0^\infty(1+x)|q(x)|\,dx<\infty\), take the Dirichlet Schrödinger operator \(H\), free operator \(H_0\), and \(k^2=z\) with \(\operatorname{Im}k>0\). The normalized perturbation determinant is the Jost value \(D(z)=f(0,k)\); off both spectra, \(-d\log D/dz=\operatorname{Tr}((H-z)^{-1}-(H_0-z)^{-1})\). The difference, not the individual resolvent, is trace class.}
\GSfield{Reconstruction}{The relative normalization \(D=1\) for \(q=0\) and a sufficient short-range hypothesis are supplied. The scribbled diagonal asymptotic alone does not justify trace class.}
\GSfield{Fate}{\citet[Theorem~1, p.~13]{buslaevfaddeev1960traces} gives the identity between the relative resolvent trace and the logarithmic derivative of the Jost function before the seminar; the note was received 17 December 1959 and published in 1960. Jost solutions are developed in \citet[\S1 and \S6]{faddeev1959inverse}. The precise Fredholm-determinant equality under weaker \(L^1\) conditions is also verified in \citet[Theorem~1.1]{gesztesy2007jost}.}
\GSfield{Present status}{As of November 2026: established with the hypotheses in the modern statement.}
\GSfield{Assessment}{literal statement: correct with missing hypotheses; publication match: exact; verification: read the result; development: annotated; confidence: high}
\end{GSitem}
\begin{GSitem}{1964-47.3}{E}{Resolvent sheets and free Green kernels}
\GSfield{Kind; source}{remark; \GSpage{P000075}}
\GSfield{The notes}{The notes discuss square-root branching in one dimension and logarithmic branching in two, and say one must leave the Hilbert-space setting to see a multivalued resolvent.}
\GSfield{Statement}{The bounded-operator resolvent \((H-z)^{-1}\) is single-valued and analytic on its resolvent set. For the free one-dimensional operator, its outgoing kernel is \(i e^{ik|x-y|}/(2k)\), \(k^2=z,\operatorname{Im}k>0\), so the kernel has a continuation in the square-root parameter. In two dimensions the outgoing Hankel kernel has logarithmic branching. Continuation across continuous spectrum requires specified weighted spaces or compactly supported test functions.}
\GSfield{Reconstruction}{The distinction between the physical resolvent and continued kernels is essential; no general meromorphic continuation is asserted for every decaying potential.}
\GSfield{Fate}{This interpretation agrees with the analytic scattering framework in \citet[\S\S1--2]{faddeev1959inverse}. It corrects a possible literal reading that an abstract Hilbert-space resolvent is multivalued on its ordinary domain.}
\GSfield{Present status}{As of November 2026: established with the hypotheses in the modern statement.}
\GSfield{Assessment}{literal statement: correct with missing hypotheses; publication match: partial; verification: read the result; development: annotated; confidence: high}
\end{GSitem}
\begin{GSitem}{1964-47.4}{E}{Whole-line Jost Wronskian}
\GSfield{Kind; source}{construction; \GSpage{P000075}}
\GSfield{The notes}{Two solutions behave like \(e^{ikx}\) at plus infinity and \(e^{-ikx}\) at minus infinity. Their Wronskian \(f_1^{\prime} f_2-f_1 f_2^{\prime}\) is called the determinant.}
\GSfield{Statement}{For real short-range \(q\) on \(\mathbb R\), let \(f_\pm(x,k)\sim e^{\pm ikx}\) at \(\pm\infty\), with \(\operatorname{Im}k>0\). The Wronskian \(W(k)=f_+^{\prime} f_- - f_+f_-^{\prime}\) is independent of \(x\); its zeros in the upper half-plane give bound states, with the usual threshold caveat at \(k=0\). The free Wronskian is \(2ik\), so a relative determinant uses \(W(k)/(2ik)\).}
\GSfield{Reconstruction}{The short-range assumption may be taken as \(\int_{\mathbb R}(1+|x|)|q|<\infty\). The normalization and threshold exception are editorial additions.}
\GSfield{Fate}{The construction belongs to the pre-1965 inverse-scattering theory; \citet[whole-line extensions in the introduction and supplementary remarks]{faddeev1959inverse} distinguishes these extensions from the main half-line problem. The unnormalized Wronskian is a characteristic function, not automatically a normalized Fredholm determinant.}
\GSfield{Present status}{As of November 2026: established with the hypotheses in the modern statement.}
\GSfield{Assessment}{literal statement: correct with missing hypotheses; publication match: partial; verification: read the result; development: annotated; confidence: high}
\end{GSitem}
\begin{GSitem}{1964-47.5}{E}{Argument principle and analytically regularized spectral sums}
\GSfield{Kind; source}{method; \GSpage{P000075-P000076}}
\GSfield{The notes}{The notes integrate \(h(\lambda)W^{\prime}(\lambda)/W(\lambda)\), deform the contour, and use \(h=\lambda^s\). The real-axis expression is followed by the annotation that the integral is a jump.}
\GSfield{Statement}{If \(W\) is meromorphic and nonzero on a closed contour and \(h\) is analytic inside it, \((2\pi i)^{-1}\oint hW^{\prime}/W\) equals the zero sum minus the pole sum, with multiplicities. A deformation onto a cut must retain both boundary values, its orientation and the factor \((2\pi i)^{-1}\). For spectral zeta or trace identities, convergence, subtraction of asymptotic terms and continuation must be proved separately.}
\GSfield{Reconstruction}{The argument-principle statement is exact; its application to an unspecified spectral determinant and branch of \(\lambda^s\) remains conditional.}
\GSfield{Fate}{\citet[\S2, pp.~14--15]{buslaevfaddeev1960traces} carries out contour deformation, asymptotic subtraction and Mellin continuation in the singular Sturm--Liouville setting. Its finite-interval discussion on p.~16 uses Dirichlet ends, so it is not silently identified with the notebook\textquotesingle s Neumann expression. This is pre-seminar trace theory, not a universal ordinary-integral identity.}
\GSfield{Present status}{As of November 2026: established with the hypotheses in the modern statement.}
\GSfield{Assessment}{literal statement: correct with missing hypotheses; publication match: partial; verification: read the result; development: annotated; confidence: high}
\end{GSitem}
\begin{GSitem}{1964-47.6}{F}{An unidentified boundary term in a Neumann regularized trace}
\GSfield{Kind; source}{claim; \GSpage{P000076}}
\GSfield{The notes}{The scan reads \(\sum_n(\lambda_n-n^2-c_0)=\tfrac14(y^{\prime}(0)-y^{\prime}(\pi))\). No function \(y\) on that right side is defined.}
\GSfield{Statement}{A standard comparison, for a sufficiently smooth real potential and Neumann eigenvalues numbered \(\lambda_0,\lambda_1,\ldots\), is \(\sum_{n=0}^\infty(\lambda_n-n^2-\bar q)=(q(0)+q(\pi))/4-\bar q/2\), where \(\bar q=\pi^{-1}\int_0^\pi q\). This is a different, specified statement; it cannot identify the written \(y^{\prime}\)-difference without an additional definition.}
\GSfield{Reconstruction}{The comparison formula is supplied to expose the ambiguity; it is not silently substituted into the transcription.}
\GSfield{Fate}{The Gelfand--Levitan trace identity dates to 1953, followed by Gelfand\textquotesingle s 1956 treatment \citep{gelfand1956identities}. The recorded boundary expression is therefore classified as an unidentified or corrupted instance of established trace theory, not a new spectral formula.}
\GSfield{Present status}{As of November 2026: the regularized-trace theory is established, but this exact boundary expression cannot be recovered from the notes.}
\GSfield{Assessment}{literal statement: not determinable; publication match: thematic only; verification: publisher metadata checked; development: annotated; confidence: low}
\end{GSitem}
\begin{GSitem}{1964-47.7}{E}{Formal diagonal kernels and boundary contributions}
\GSfield{Kind; source}{method; \GSpage{P000076}}
\GSfield{The notes}{The distributional kernel \(-\delta^{\prime\prime}(x-y)+q(x)\delta(x-y)\) is restricted formally to the diagonal. The notes write \(\delta(0)=0\), then discuss a boundary coefficient \(1/4\).}
\GSfield{Statement}{The distribution \(\delta(x-y)\) has no pointwise diagonal value \(\delta(0)\), and multiplication by a nonzero potential is not trace class on \(L^2\). A valid regularized trace must first specify a spectral, heat-kernel or resolvent subtraction and account for its endpoint terms. The displayed manipulations are heuristic motivation, not legitimate distributional equalities.}
\GSfield{Reconstruction}{The source\textquotesingle s distinction between periodic and endpoint problems is retained as motivation, while the regularization is not invented.}
\GSfield{Fate}{The endpoint effects belong to established regularized-trace formulas \citep{gelfand1956identities,buslaevfaddeev1960traces}. Their rigorous existence does not validate the intermediate assertion \(\delta(0)=0\).}
\GSfield{Present status}{As of November 2026: established with the hypotheses in the modern statement.}
\GSfield{Assessment}{literal statement: false as written; publication match: partial; verification: publisher metadata checked; development: annotated; confidence: high}
\end{GSitem}
\begin{GSitem}{1964-47.8}{B}{When does a commutator have trace zero?}
\GSfield{Kind; source}{problem; \GSpage{P000076}}
\GSfield{The notes}{Gelfand writes \(\operatorname{Sp}(AB-BA)=0\) and labels it a problem, without operator-ideal assumptions.}
\GSfield{Statement}{For bounded operators on a Hilbert space, \(\operatorname{Tr}(AB)=\operatorname{Tr}(BA)\) when both products are trace class (in particular if one factor is trace class). Merely requiring \([A,B]\) to be trace class is insufficient: for the unilateral shift \(S\) on \(\ell^2(\mathbb N_0)\), \([S^*,S]=P_{e_0}\) has trace one.}
\GSfield{Reconstruction}{The two distinct hypotheses are separated. The rank-one shift example is an explicit diagnosis of the unrestricted formulation.}
\GSfield{Fate}{The question depends completely on the trace domain. Nonzero traces of trace-class commutators became a substantial theory; \citet[Introduction]{heltonhowe1975traces} develops trace formulas for almost commuting operators and relates them to symbols and index theory. This is a later development, not evidence for the unrestricted zero-trace assertion.}
\GSfield{Present status}{As of November 2026: vanishing is established in the standard trace-cyclicity domain and false under the weaker trace-class-commutator assumption.}
\GSfield{Assessment}{literal statement: correct with missing hypotheses; publication match: partial; verification: read the result; development: annotated; confidence: high; problem \hyperref[prob:GP-1964-12]{GP-1964-12}}
\end{GSitem}

\GStalk{48}{1965-04-19}{S. G. Gindikin}{P000075-P000080}{Cauchy--Weil formulas and holomorphic Radon transforms}
Gindikin begins with the product Cauchy formula and analytic polyhedra, introduces Hefer decompositions and skeleton integrals, then proposes a continuous family of defining functions. The final linear case connects directly to his 1965 holomorphic Radon-transform note, announced to the Moscow Mathematical Society in February.

\begin{GSitem}{1964-48.1}{E}{Product Cauchy formula and analytic polyhedra}
\GSfield{Kind; source}{construction; \GSpage{P000076-P000077}}
\GSfield{The notes}{A bicylinder \(D_z\times D_w\) has the iterated Cauchy formula. Weil domains are cut out by holomorphic \(\chi_i(z,w)\in D_i\).}
\GSfield{Statement}{If \(D_1,D_2\subset\mathbb C\) have piecewise smooth positively oriented boundaries and \(f\) is holomorphic near \(\overline{D_1\times D_2}\), then \(f(z,w)=(2\pi i)^{-2}\int_{\partial D_1\times\partial D_2}f(\zeta,\eta)\,d\zeta\wedge d\eta/((\zeta-z)(\eta-w))\). More general analytic polyhedra are relatively compact sets defined by finitely many holomorphic inequalities. Two defining functions are genuine global coordinates only with an injective biholomorphic change of variables; otherwise multiplicities matter.}
\GSfield{Reconstruction}{Boundary regularity, holomorphy near closure and the coordinate-map caveat are supplied.}
\GSfield{Fate}{The product formula is classical, and the polyhedral extension is explained in the pre-seminar work \citet[\S17 and Corollary~18.2]{norguet1961formes}. The notes are expository at this stage.}
\GSfield{Present status}{As of November 2026: established with the hypotheses in the modern statement.}
\GSfield{Assessment}{literal statement: correct with missing hypotheses; publication match: exact; verification: read the result; development: annotated; confidence: high}
\end{GSitem}
\begin{GSitem}{1964-48.2}{E}{Hefer holomorphic difference decomposition}
\GSfield{Kind; source}{theorem; \GSpage{P000077}}
\GSfield{The notes}{The displayed lemma expresses \(f(z,w)-f(\zeta,\eta)\) as \((z-\zeta)P+(w-\eta)Q\), with holomorphic coefficients.}
\GSfield{Statement}{If \(D\subset\mathbb C^n\) is a domain of holomorphy and \(f\in\mathcal O(D)\), there exist \(H_j\in\mathcal O(D\times D)\) such that \(f(z)-f(\zeta)=\sum_j(z_j-\zeta_j)H_j(z,\zeta)\). On a convex domain one may take \(H_j=\int_0^1\partial_j f(\zeta+t(z-\zeta))\,dt\); the global Stein-domain statement uses holomorphic division along the diagonal.}
\GSfield{Reconstruction}{The domain-of-holomorphy hypothesis is essential to the global assertion. The convex-domain integral is a transparent special-case realization.}
\GSfield{Fate}{\citet{hefer1950zerlegung} published the decomposition in 1950. The use and existence of these coefficients for domains of holomorphy are also stated in the opened \citet[p.~78, following Corollary~18.2]{norguet1961formes}.}
\GSfield{Present status}{As of November 2026: established with the hypotheses in the modern statement.}
\GSfield{Assessment}{literal statement: correct with missing hypotheses; publication match: exact; verification: read the result; development: annotated; confidence: high}
\end{GSitem}
\begin{GSitem}{1964-48.3}{E}{Cauchy--Weil representation on the distinguished boundary}
\GSfield{Kind; source}{theorem; \GSpage{P000077-P000078}}
\GSfield{The notes}{The notes sum over pairs of boundary faces, with a determinant of Hefer coefficients divided by two defining-function differences. The indices in the written kernel are incomplete.}
\GSfield{Statement}{For a relatively compact analytic polyhedron \(P=\{|\chi_j|<1\}\subset D\subset\mathbb C^2\), with regular transverse boundary faces and \(D\) a domain of holomorphy, choose \(\chi_i(\zeta)-\chi_i(z)=\sum_{a=1}^2(\zeta_a-z_a)H_{ia}(\zeta,z)\). With the standard induced orientations on \(\Gamma_{ij}=\{|\chi_i|=|\chi_j|=1\}\cap\overline P\), a function holomorphic near \(\overline P\) has the Cauchy--Weil representation \((2\pi i)^{-2}\sum_{i<j}\int_{\Gamma_{ij}} f(\zeta)\det(H_{ia},H_{ja})\,d\zeta_1\wedge d\zeta_2/((\chi_i(\zeta)-\chi_i(z))(\chi_j(\zeta)-\chi_j(z)))\).}
\GSfield{Reconstruction}{The determinant notation means the determinant of the two rows \((H_{i1},H_{i2})\), \((H_{j1},H_{j2})\); regularity and orientation are supplied. This is not a literal reproduction of the notebook\textquotesingle s mixed indices.}
\GSfield{Fate}{The exact architecture of this kernel is already in \citet[\S17; Corollary~18.2]{norguet1961formes}, following Weil and Hefer. The construction is known before the 1965 seminar.}
\GSfield{Present status}{As of November 2026: established with the hypotheses in the modern statement.}
\GSfield{Assessment}{literal statement: correct with missing hypotheses; publication match: exact; verification: read the result; development: annotated; confidence: high}
\end{GSitem}
\begin{GSitem}{1964-48.4}{E}{Polyhedral approximation and decomposition by face integrals}
\GSfield{Kind; source}{method; \GSpage{P000078}}
\GSfield{The notes}{The notes say that domains of holomorphy can be approximated by Weil domains and that the formula decomposes a holomorphic function into terms extending to larger regions.}
\GSfield{Statement}{In a Stein domain, holomorphically convex compact sets admit relatively compact analytic-polyhedron neighborhoods. In a Cauchy--Weil representation, each face integral is holomorphic in any parameter region on which its coefficient functions are holomorphic and its denominator stays nonzero uniformly on the integration cycle. This gives controlled local extension/decomposition; it does not imply that every holomorphic function extends beyond its domain of holomorphy.}
\GSfield{Reconstruction}{The relevant compact-set and nonvanishing conditions are the precise reconstruction of the broad approximation sentence.}
\GSfield{Fate}{The polyhedral and ideal-theoretic foundations precede the seminar \citep{cartan1950ideaux}; \citet[Corollary~18.2]{norguet1961formes} provides the representation used here. The separate extensions of summands must not be confused with extension of their total across every boundary point.}
\GSfield{Present status}{As of November 2026: established with the hypotheses in the modern statement.}
\GSfield{Assessment}{literal statement: correct with missing hypotheses; publication match: partial; verification: read the result; development: annotated; confidence: high}
\end{GSitem}
\begin{GSitem}{1964-48.5}{F}{A continuous family of Weil defining functions}
\GSfield{Kind; source}{method; \GSpage{P000078}}
\GSfield{The notes}{The finite \(\chi_i\) are replaced by \(\chi(\alpha,z,w)\); a kernel containing \(P_1\partial_\alpha P_2-P_2\partial_\alpha P_1\) is sketched.}
\GSfield{Statement}{The fragment proposes a continuous-parameter analogue of the Cauchy--Weil sum. It does not fix the parameter cycle, differentiability in \(\alpha\), the integration manifold, denominator order or orientations needed for an integral representation theorem.}
\GSfield{Reconstruction}{No missing denominator power or integration cycle is guessed.}
\GSfield{Fate}{Integral formulas with continuously varying complex hyperplanes are part of Gindikin\textquotesingle s work, including \citet[formula~(2)]{gindikin1965radon}. That paper identifies the next, linear example; it does not justify every symbol in this more general unfinished kernel. Searches did not recover an exact formula with all the manuscript\textquotesingle s conventions.}
\GSfield{Present status}{As of November 2026: a recognizable method with an incomplete exact statement; no independent open problem is inferred.}
\GSfield{Assessment}{literal statement: not determinable; publication match: thematic only; verification: read the result; development: annotated; confidence: low}
\end{GSitem}
\begin{GSitem}{1964-48.6}{A}{Linear defining functions and the holomorphic Radon transform}
\GSfield{Kind; source}{construction; \GSpage{P000078}}
\GSfield{The notes}{For linear \(\chi\), the notes connect the formulas with Radon transforms, particularly in a tube over a convex domain.}
\GSfield{Statement}{Let \(T_\Omega=\mathbb R^n+i\Omega\), where \(\Omega\) is open convex and contains no affine line. Under the smooth strictly convex assumptions of the initial construction, take a holomorphic \(f\) with sufficient rapid decay in the real directions on each horizontal slice. Its integrals over real \((n-1)\)-planes in suitable complex hyperplanes form a holomorphic Radon transform, with inversion by the specified oriented boundary integral in Gindikin\textquotesingle s formula~(2). The result requires these geometric and decay conditions.}
\GSfield{Reconstruction}{The precise tube, rapid-decay hypothesis and initial boundary assumptions come from the paper, not from the shorthand notes.}
\GSfield{Fate}{This is an exact match to \citet[pp.~237--239, formula~(2)]{gindikin1965radon}. The printed footnote records an announcement to the Moscow Mathematical Society on 23 February 1965, before this April seminar; the publication appeared in 1965. The article also discusses how its initial boundary regularity can be relaxed.}
\GSfield{Present status}{As of November 2026: established published work, with pre-seminar circulation distinguished from publication.}
\GSfield{Assessment}{literal statement: correct with missing hypotheses; publication match: exact; verification: read the result; development: annotated; confidence: high}
\end{GSitem}

\GStalk{49}{1965-04-26}{I. N. Bernstein}{P000078-P000085}{Čech complexes and derived sheaf cohomology}
Bernstein reviews descent, constructs a Čech complex and maps under refinement, and explains why global sections need derived functors. The example of functions vanishing at two points of the projective line is useful and precise. The attribution of a failure of exactness to taking a direct limit is mathematically wrong and is corrected separately from the underlying cohomological construction.

\begin{GSitem}{1964-49.1}{E}{Review of sites, fiber products and sheaf descent}
\GSfield{Kind; source}{definition; \GSpage{P000079-P000080}}
\GSfield{The notes}{The opening diagrams repeat products over an object, pullback of a cover, and compatibility of sections on double overlaps.}
\GSfield{Statement}{For a site presented by covering families and the required fiber products, a sheaf of abelian groups is a contravariant functor satisfying the equalizer condition for every covering. Products in the slice represent overlaps; the universal property fixes all restriction maps used in the Čech complex.}
\GSfield{Fate}{This continues record~40 and follows the pre-seminar site formalism of SGA~4; see \citet[Exposé II]{sga4I1972} and \citet[Tags 00VH, 00VL]{stacks2026sites}.}
\GSfield{Present status}{As of November 2026: established with the hypotheses in the modern statement.}
\GSfield{Assessment}{literal statement: correct; publication match: exact; verification: read the result; development: annotated; confidence: high}
\end{GSitem}
\begin{GSitem}{1964-49.2}{E}{Čech cochains and the alternating differential}
\GSfield{Kind; source}{construction; \GSpage{P000080-P000081}}
\GSfield{The notes}{The terms are products of sections on multiple intersections, with alternating restriction maps; the notes show the first degrees and \(d^2=0\).}
\GSfield{Statement}{For a cover \(\mathcal U=(U_i\to X)\), set \(\check C^q(\mathcal U,F)=\prod_{i_0,\ldots,i_q}F(U_{i_0}\times_X\cdots\times_XU_{i_q})\). The differential is the alternating sum of pullbacks along omission maps. Simplicial identities imply \(d^2=0\), and \(\check H^0(\mathcal U,F)=F(X)\) for a sheaf.}
\GSfield{Fate}{This is classical Čech theory in the language of sites. The explicit construction and extension to complexes are given in \citet[Tags 01FP and 01GZ]{stacks2026cohomology}; the site version was part of the 1963--64 SGA program \citep{sga4I1972}.}
\GSfield{Present status}{As of November 2026: established with the hypotheses in the modern statement.}
\GSfield{Assessment}{literal statement: correct; publication match: exact; verification: read the result; development: annotated; confidence: high}
\end{GSitem}
\begin{GSitem}{1964-49.3}{E}{Refinement maps and direct-limit Čech cohomology}
\GSfield{Kind; source}{construction; \GSpage{P000081}}
\GSfield{The notes}{A refinement is assigned a map of cochain complexes; cohomology over all covers is then considered.}
\GSfield{Statement}{A chosen refinement map induces a cochain map in the opposite direction on Čech complexes. Two choices subordinate to the same refinement induce chain-homotopic maps, hence the same cohomology map. For ordinary open covers, direct-limit Čech cohomology is \(\check H^q(X,F)=\varinjlim_\mathcal U\check H^q(\mathcal U,F)\). On a site, one must use a specified refinement category and smallness convention.}
\GSfield{Reconstruction}{Independence of the chosen refinement and the size convention are supplied; the notes do not prove them.}
\GSfield{Fate}{The refinement construction is explicitly checked in \citet[Tag 01FP, Refinements]{stacks2026cohomology}. Passing to hypercovers gives derived sheaf cohomology on sites \citep[Tag 01GZ, Theorem~25.10.1]{stacks2026cohomology}.}
\GSfield{Present status}{As of November 2026: established with the hypotheses in the modern statement.}
\GSfield{Assessment}{literal statement: correct with missing hypotheses; publication match: exact; verification: read the result; development: annotated; confidence: high}
\end{GSitem}
\begin{GSitem}{1964-49.4}{E}{The direct-limit explanation of Čech exactness is false}
\GSfield{Kind; source}{claim; \GSpage{P000081-P000082}}
\GSfield{The notes}{The source suggests that a long exact sequence exists for each cover but can fail after the passage over refinements.}
\GSfield{Statement}{Filtered colimits of abelian groups are exact, so a compatible filtered system of long exact sequences remains exact. The actual issue is earlier: an epimorphism of sheaves need not be surjective on sections of an overlap, so a short exact sequence of sheaves need not give a short exact sequence of fixed-cover Čech complexes. In higher degree, ordinary Čech cohomology on an arbitrary site need not equal derived sheaf cohomology.}
\GSfield{Reconstruction}{This separates the faithful source error from the correct motivation for derived functors.}
\GSfield{Fate}{The derived-functor remedy was available before this talk in \citet{grothendieck1957tohoku}. The precise hypercover comparison is \citet[Tag 01GZ, Theorem~25.10.1]{stacks2026cohomology}. Neither result supports blaming a filtered direct limit for loss of exactness.}
\GSfield{Present status}{As of November 2026: the exactness issue is settled under the stated category conventions; the manuscript explanation is false.}
\GSfield{Assessment}{literal statement: false as written; publication match: partial; verification: read the result; development: annotated; confidence: high}
\end{GSitem}
\begin{GSitem}{1964-49.5}{E}{Abelian sheaves, presheaves and exactness}
\GSfield{Kind; source}{definition; \GSpage{P000082}}
\GSfield{The notes}{The notes contrast a category of sheaves with one of presheaves and discuss failure of exactness for an inclusion or section functor. A subscript resembling \(p\) is not securely identified as a stalk.}
\GSfield{Statement}{Abelian presheaves form an abelian category with kernels and cokernels computed objectwise. Abelian sheaves also form an abelian category; their kernels are computed as presheaf kernels, while cokernels are obtained by sheafifying presheaf cokernels. The inclusion into presheaves is left exact and generally not right exact. On a topological space, exactness of sheaves is tested on stalks.}
\GSfield{Reconstruction}{The ambiguous subscript is left unresolved; the three functors---inclusion, sections and stalk---are distinguished.}
\GSfield{Fate}{This is part of Grothendieck\textquotesingle s abelian-category foundation \citep{grothendieck1957tohoku}; site formulations appear in SGA~4 \citep[Exposé II]{sga4I1972}. The notebook must not be read as denying stalkwise exactness.}
\GSfield{Present status}{As of November 2026: established with the hypotheses in the modern statement.}
\GSfield{Assessment}{literal statement: correct with missing hypotheses; publication match: exact; verification: publisher metadata checked; development: annotated; confidence: high}
\end{GSitem}
\begin{GSitem}{1964-49.6}{E}{The ideal sheaf of two points on the projective line}
\GSfield{Kind; source}{construction; \GSpage{P000082}}
\GSfield{The notes}{The example takes two points of \(\mathbb P^1\), functions vanishing there, and an exact sequence whose global sections are not right exact.}
\GSfield{Statement}{For distinct \(p,q\in\mathbb P^1_\mathbb C\), let \(J=\mathcal I_{\{p,q\}}\cong\mathcal O(-2)\). The exact sequence \(0\to J\to\mathcal O\to\mathbb C_p\oplus\mathbb C_q\to0\) gives \(H^0(J)=0\), \(H^0(\mathcal O)=\mathbb C\), and the diagonal map \(\mathbb C\to\mathbb C^2\), which is not onto. Its cokernel is \(H^1(J)\cong\mathbb C\), since \(H^1(\mathcal O)=0\).}
\GSfield{Reconstruction}{The points are assumed distinct and the ground field complex, matching the holomorphic formulation. The cohomology calculation completes the source\textquotesingle s example.}
\GSfield{Fate}{The example concretely illustrates the derived-functor theory already established in \citet{grothendieck1957tohoku}. The line-bundle notation identifies the exact obstruction; it is not an open extension problem.}
\GSfield{Present status}{As of November 2026: established with the hypotheses in the modern statement.}
\GSfield{Assessment}{literal statement: correct with missing hypotheses; publication match: exact; verification: publisher metadata checked; development: annotated; confidence: high}
\end{GSitem}
\begin{GSitem}{1964-49.7}{F}{An incomplete differential-form exact sequence}
\GSfield{Kind; source}{construction; \GSpage{P000082-P000083}}
\GSfield{The notes}{A sequence beginning \(0\to C\to\Omega^0\to\Omega\to0\) is sketched without a secure definition of the final sheaf or differential.}
\GSfield{Statement}{On a smooth manifold, \(0\to\underline{\mathbb C}\to\mathcal A^0\xrightarrow{d}\mathcal Z^1\to0\) is locally exact, where \(\mathcal Z^1\) is the sheaf of closed smooth one-forms. The complete de Rham resolution instead continues through \(\mathcal A^1,\mathcal A^2,\ldots\). Replacing \(\mathcal Z^1\) by all one-forms and ending with zero is false in dimension greater than one.}
\GSfield{Reconstruction}{The smooth de Rham example is a possible precise repair, explicitly not a decipherment of the unspecified final symbol.}
\GSfield{Fate}{The Poincaré-lemma/de Rham interpretation is standard pre-seminar material, but the notebook does not decide whether \(\Omega\) means closed forms, an image sheaf or something else. The exact sequence is therefore not identified beyond this possible reconstruction.}
\GSfield{Present status}{As of November 2026: the proposed reconstruction is classical; the literal abbreviated sequence remains underdetermined.}
\GSfield{Assessment}{literal statement: not determinable; publication match: none found; verification: indirect; development: annotated; confidence: low}
\end{GSitem}
\begin{GSitem}{1964-49.8}{E}{Derived functors of sections and cohomology presheaves}
\GSfield{Kind; source}{method; \GSpage{P000083}}
\GSfield{The notes}{The final paragraph proposes right-derived functors of a left-exact functor from sheaves to presheaves, using injective resolutions to recover the long exact sequence.}
\GSfield{Statement}{For abelian sheaves on a site, choose an injective resolution \(F\to I^\bullet\). Then \(H^q(U,F)=H^q(I^\bullet(U))\), and \(U\mapsto H^q(U,F)\) is a cohomology presheaf. Global cohomology is the right-derived functor of \(\Gamma(X,-)\). A short exact sequence of sheaves yields the usual long exact cohomology sequence. The inclusion of sheaves into presheaves is a different left-exact functor, whose derived values encode these local section cohomologies.}
\GSfield{Reconstruction}{Enough injectives, the precise functor and the distinction between a cohomology presheaf and a cohomology sheaf are stated explicitly.}
\GSfield{Fate}{Grothendieck\textquotesingle s 1957 paper established this derived-functor organization before the talk \citep{grothendieck1957tohoku}. \citet[Tags 01FP and 01GZ]{stacks2026cohomology} makes clear how Čech and hypercover constructions compare with it.}
\GSfield{Present status}{As of November 2026: established with the hypotheses in the modern statement.}
\GSfield{Assessment}{literal statement: correct with missing hypotheses; publication match: exact; verification: read the result; development: annotated; confidence: high}
\end{GSitem}

\GStalk{50}{1965-05-10}{Paul R. Halmos; I. M. Gelfand; unidentified participant}{P000083-P000090}{Invariant subspaces, polynomial compactness, unicellular operators and Bishop operators}
Halmos surveys invariant-subspace results and gives the elementary proof of the polynomially compact case. The ensuing questions concern invariant-subspace lattices, weighted shifts, a measure version of Volterra integration and Bishop operators. A final intervention sketches the Sz.-Nagy--Foiaş analytic model. The catalog envelope is preserved, but P000090 begins Kohn on 20 September 1965 and is outside this academic year and record.

\begin{GSitem}{1964-50.1}{C}{The invariant-subspace problem on complex Hilbert space}
\GSfield{Kind; source}{problem; \GSpage{P000084; P000088}}
\GSfield{The notes}{The notes explicitly require a closed invariant subspace different from zero and the whole complex Hilbert space. On P000088 this is reformulated as whether the lattice can consist only of \(\{0,H\}\).}
\GSfield{Statement}{For every bounded linear operator \(T\) on a separable infinite-dimensional complex Hilbert space \(H\), must there be a closed subspace \(M\) with \(0\ne M\ne H\) and \(TM\subseteq M\)? Equivalently, can every nonzero vector be cyclic for one such \(T\)?}
\GSfield{Fate}{The 1964 de Branges--Rovnyak announcement was retracted in 1965, as explicitly identified by Rovnyak\textquotesingle s author bibliography \citep{debrangesrovnyak1964existence,debrangesrovnyak1965correction}. The primary preprint \citet[abstract and Introduction, pp.~1--2]{marcoux2026generators}, dated 21 August 2026, still calls the precisely scoped Hilbert problem unsolved. Targeted searches through 12 November found claimed solutions and partial results; no generally accepted general solution was verified. A claim in a preprint is not treated as an established theorem.}
\GSfield{Present status}{As of 12 November 2026: no generally accepted solution was verified in the literature checked. This qualified status concerns bounded operators on complex infinite-dimensional Hilbert space, not arbitrary Banach spaces or special classes.}
\GSfield{Assessment}{literal statement: correct; publication match: exact; verification: read the result; development: annotated; confidence: high; problem \hyperref[prob:GP-1964-07]{GP-1964-07}}
\GSfield{Worth developing}{Yes, as the central recognized problem, but an unrestricted attack is far beyond a bounded chapter development. The tractable work here is confined to the precisely specified Bishop family and finite-orbit certificates.}
\end{GSitem}
\begin{GSitem}{1964-50.2}{E}{Normal and self-adjoint operators have invariant subspaces}
\GSfield{Kind; source}{theorem; \GSpage{P000084}}
\GSfield{The notes}{The normal and Hermitian cases are described as known.}
\GSfield{Statement}{Every normal bounded operator on a complex Hilbert space of dimension greater than one has a nonzero proper closed invariant subspace. A nontrivial spectral projection gives a reducing subspace; when the operator is scalar, any one-dimensional subspace suffices. Self-adjoint operators are included.}
\GSfield{Reconstruction}{The dimension restriction and the scalar exceptional argument prevent the shorthand from being applied to a one-dimensional space.}
\GSfield{Fate}{This is classical spectral theory, already treated as established background in the contemporary polynomially compact discussion \citep[p.~433]{halmos1966polynomial}. It does not require compactness or an eigenvalue in the nonscalar continuous-spectrum case.}
\GSfield{Present status}{As of November 2026: established with the hypotheses in the modern statement.}
\GSfield{Assessment}{literal statement: correct with missing hypotheses; publication match: partial; verification: read the result; development: annotated; confidence: high}
\end{GSitem}
\begin{GSitem}{1964-50.3}{E}{Compact operators and the Aronszajn--Smith theorem}
\GSfield{Kind; source}{theorem; \GSpage{P000084}}
\GSfield{The notes}{Completely continuous operators are attributed to von Neumann, followed by Aronszajn and Smith.}
\GSfield{Statement}{Every compact bounded linear operator on an infinite-dimensional complex Banach space has a nonzero proper closed invariant subspace. In particular this holds on complex Hilbert space; for finite dimension greater than one, ordinary eigenvector theory supplies the result.}
\GSfield{Fate}{\citet{aronszajnsmith1954} published the Banach-space theorem in 1954. The later primary account \citet[Introduction]{chamizo2020bishop} explicitly records the earlier unpublished von Neumann Hilbert-space result. The accessible Aronszajn--Smith record verified publication metadata; the complete original proof was not read.}
\GSfield{Present status}{As of November 2026: established with the hypotheses in the modern statement.}
\GSfield{Assessment}{literal statement: correct with missing hypotheses; publication match: exact; verification: publisher metadata checked; development: annotated; confidence: high}
\end{GSitem}
\begin{GSitem}{1964-50.4}{E}{Subnormality and the unilateral shift}
\GSfield{Kind; source}{definition; \GSpage{P000084}}
\GSfield{The notes}{Subnormal means having a normal extension on a larger Hilbert space. The unilateral shift becomes normal after negative-index basis vectors are added; finite-dimensional subnormal operators are normal.}
\GSfield{Statement}{An operator \(T\in B(H)\) is subnormal if \(H\) is an invariant subspace of a Hilbert space \(K\) for a normal \(N\in B(K)\), and \(T=N|_H\). The unilateral shift on \(\ell^2(\mathbb N_0)\) is the restriction of the bilateral unitary shift on \(\ell^2(\mathbb Z)\). Finite-dimensional subnormal operators are normal. Taking adjoints need not preserve subnormality.}
\GSfield{Reconstruction}{Invariance of the embedded space is essential: a compression of a normal operator to an arbitrary subspace would not be the stated notion.}
\GSfield{Fate}{This is the established Halmos--Bram theory. The definition and intrinsic positivity criteria are recalled and extended by \citet[p.~99]{lambert1984subnormal}; Bram\textquotesingle s original 1955 paper is identified in its bibliography \citep{bram1955subnormal}.}
\GSfield{Present status}{As of November 2026: established with the hypotheses in the modern statement.}
\GSfield{Assessment}{literal statement: correct with missing hypotheses; publication match: partial; verification: read the result; development: annotated; confidence: high}
\end{GSitem}
\begin{GSitem}{1964-50.5}{B}{Invariant subspaces for subnormal operators}
\GSfield{Kind; source}{problem; \GSpage{P000084}}
\GSfield{The notes}{After defining subnormality, the notes say that its invariant-subspace question is still unsolved.}
\GSfield{Statement}{Does every subnormal bounded operator on a complex Hilbert space of dimension greater than one possess a nonzero proper closed invariant subspace? The answer is affirmative.}
\GSfield{Fate}{\citet[abstract and pp.~310--333]{brown1978subnormal} proved the affirmative result, published in September 1978 after receipt on 21 July 1978. The publisher abstract was opened and explicitly states this result; the full subscription article was not read. The 1965 unsolved label is historically appropriate but cannot be retained as a current status.}
\GSfield{Present status}{As of November 2026: solved affirmatively by Brown in 1978; this special-class problem is distinct from item 1964-50.1.}
\GSfield{Assessment}{literal statement: correct; publication match: exact; verification: publisher metadata checked; development: annotated; confidence: high; problem \hyperref[prob:GP-1964-13]{GP-1964-13}}
\end{GSitem}
\begin{GSitem}{1964-50.6}{E}{Hyponormality is necessary but not sufficient for subnormality}
\GSfield{Kind; source}{claim; \GSpage{P000084}}
\GSfield{The notes}{The overwritten qualification after \(T^*T-TT^*\ge0\) reads necessary and sufficient; the scan leaves the horizontal sign slightly ambiguous.}
\GSfield{Statement}{Subnormality implies hyponormality \(T^*T-TT^*\ge0\), but the converse is false. The bounded unilateral shift with weights \(1,1,2,2,\ldots\) is hyponormal since the weights are nondecreasing. If it were subnormal, its moment measure would have moments \(\gamma_0=\gamma_1=\gamma_2=1\), forcing concentration at \(1\), contrary to \(\gamma_3=4\).}
\GSfield{Reconstruction}{The minus sign is the mathematically coherent reading. The weighted-shift counterexample and moment argument are editorial verification; they are not attributed to the speaker.}
\GSfield{Fate}{The full Bram--Halmos positivity conditions involve all finite arrays of vectors and powers, not just the first commutator; see \citet[Introduction and criterion]{lambert1984subnormal}. The recorded biconditional is a faithful mathematical error. The explicit shift supplies a short diagnostic counterexample, independently of the uncertain overwritten punctuation.}
\GSfield{Present status}{As of November 2026: the single positivity condition is insufficient; complete intrinsic subnormality criteria are established.}
\GSfield{Assessment}{literal statement: false as written; publication match: partial; verification: read the result; development: annotated; confidence: high}
\end{GSitem}
\begin{GSitem}{1964-50.7}{A}{Polynomially compact operators have invariant subspaces}
\GSfield{Kind; source}{theorem; \GSpage{P000084-P000087}}
\GSfield{The notes}{The case \(T^2\) compact is attributed to Robinson and Bernstein using predicate calculus. Halmos then proves the general case \(p(T)\) compact for a nonzero polynomial.}
\GSfield{Statement}{If \(H\) is an infinite-dimensional complex Hilbert space, \(T\in B(H)\), and \(p\in\mathbb C[z]\setminus\{0\}\) has \(p(T)\) compact, then \(T\) has a nonzero proper closed invariant subspace. The case \(p(T)=0\) is included.}
\GSfield{Reconstruction}{Infinite dimension and a nonzero polynomial are explicit. In finite dimension greater than one the conclusion is elementary; in dimension one it fails.}
\GSfield{Fate}{\citet[Theorem~1.1]{bernsteinrobinson1966} published the nonstandard-analysis proof in 1966; their paper was received 5 July 1964 and revised 10 December 1964. \citet[pp.~433--437]{halmos1966polynomial} gave the elementary proof, received 10 October 1964. Both manuscripts had been submitted before this May 1965 lecture, although publication followed it; the receipt dates alone do not establish public circulation. The coauthor is Allen R. Bernstein, not seminar participant I. N. Bernstein.}
\GSfield{Present status}{As of November 2026: established; the notebook is an early oral record of the subsequently published elementary proof.}
\GSfield{Assessment}{literal statement: correct with missing hypotheses; publication match: exact; verification: read the result; development: annotated; confidence: high}
\end{GSitem}
\begin{GSitem}{1964-50.8}{A}{Cyclic Hessenberg matrices and small subdiagonal products}
\GSfield{Kind; source}{method; \GSpage{P000084-P000085}}
\GSfield{The notes}{Gram--Schmidt on \(e,Te,T^2e,\ldots\) gives a matrix with one nonzero subdiagonal. Products along its lowest diagonal tend to zero by compactness of \(p(T)\).}
\GSfield{Statement}{If one nonzero cyclic vector has linearly independent iterates, Gram--Schmidt gives \(a_{mn}=\langle Te_n,e_m\rangle=0\) for \(m>n+1\). For \(k\ge1\), \((T^k)_{n+k,n}=\prod_{j=1}^k a_{n+j,n+j-1}\). If \(p\) has degree \(k\), nonzero leading coefficient, and \(p(T)\) is compact, these products tend to zero, so some increasing subsequence of subdiagonal entries tends to zero.}
\GSfield{Reconstruction}{The leading coefficient, consistent row-column convention and exceptional finite-dimensional orbit are made explicit.}
\GSfield{Fate}{This is the reduction on \citet[pp.~433--434]{halmos1966polynomial}. A noncyclic orbit closure or a finite dependence already yields an invariant subspace. P000085 reverses the matrix indices in one definition; maintaining the convention from P000084 repairs the argument.}
\GSfield{Present status}{As of November 2026: established with the hypotheses in the modern statement.}
\GSfield{Assessment}{literal statement: false as written; publication match: exact; verification: read the result; development: annotated; confidence: high}
\end{GSitem}
\begin{GSitem}{1964-50.9}{A}{Almost invariant finite compressions and polynomials}
\GSfield{Kind; source}{method; \GSpage{P000085-P000086}}
\GSfield{The notes}{The small subdiagonal entry is used to compress \(T\) to initial finite-dimensional spans. Bounded vector sequences are compared modulo norm convergence to zero, with a special warning about constant terms.}
\GSfield{Statement}{Let \(P_n\) be finite-rank orthogonal projections with \(P_n\to I\) strongly and \(\|(I-P_n)TP_n\|\to0\). For every polynomial \(p\) and bounded \(x_n\) with \(\|(I-P_n)x_n\|\to0\), one has \(\|p(T)x_n-p(P_nTP_n)x_n\|\to0\). On \(P_nH\), the compression admits a complete invariant flag.}
\GSfield{Reconstruction}{The norm of the escaping corner, boundedness of the vector sequence, and asymptotic membership in \(P_nH\) are the exact sufficient hypotheses.}
\GSfield{Fate}{The approximation lemma is in \citet[p.~435]{halmos1966polynomial}. The source\textquotesingle s two differently placed \(P_n\) factors explain its warning at exponent zero; dropping them silently would create a false formula for arbitrary bounded sequences.}
\GSfield{Present status}{As of November 2026: established with the hypotheses in the modern statement.}
\GSfield{Assessment}{literal statement: correct with missing hypotheses; publication match: exact; verification: read the result; development: annotated; confidence: high}
\end{GSitem}
\begin{GSitem}{1964-50.10}{A}{Adjacent flags and weak limits in the compactness argument}
\GSfield{Kind; source}{method; \GSpage{P000086-P000087}}
\GSfield{The notes}{Two adjacent projections from each finite-dimensional invariant flag are chosen at a norm threshold. Their weak limits define spaces on which they converge strongly to the identity; compactness provides a nonzero limiting vector.}
\GSfield{Statement}{For adjacent projections \(Q_n^-\le Q_n^+\), with rank-one difference and asymptotically invariant ranges, choose weak-operator limits \(Q^\pm\). The subspaces \(M^\pm=\ker(I-Q^\pm)=\{x:Q_n^\pm x\to x\}\) are closed and invariant. The compact polynomial and a correct threshold selection give \(M^-\ne H\), \(M^+\ne0\); rank-one adjacency excludes \(M^-=0\) and \(M^+=H\) simultaneously.}
\GSfield{Reconstruction}{The weak limits are positive contractions. The mathematical mechanism is reconstructed from the exact publication; no unrecorded monotonicity of the displayed norms is assumed.}
\GSfield{Fate}{This is the closing argument of \citet[pp.~435--437]{halmos1966polynomial}. Weak limits of projections need not be projections. The notebook\textquotesingle s missing terminal vector and conflicting threshold directions are not reproduced as a proof; the published selection supplies their precise repair.}
\GSfield{Present status}{As of November 2026: established with the hypotheses in the modern statement.}
\GSfield{Assessment}{literal statement: false as written; publication match: exact; verification: read the result; development: annotated; confidence: high}
\end{GSitem}
\begin{GSitem}{1964-50.11}{A}{An entire function giving a compact operator reduces to a polynomial}
\GSfield{Kind; source}{theorem; \GSpage{P000088}}
\GSfield{The notes}{The notes state that if \(p\) is entire and \(p(T)\) compact, the polynomial result immediately gives an invariant subspace.}
\GSfield{Statement}{Let \(f\) be a nonzero entire function and \(T\in B(H)\) on an infinite-dimensional complex Hilbert space. If \(f(T)\) is compact, factor the finitely many zeros in a compact neighborhood of \(\sigma(T)\): \(f=p g\), where \(p\) is polynomial and \(g(T)\) is invertible. Then \(p(T)=g(T)^{-1}f(T)\) is compact, so the polynomial theorem applies.}
\GSfield{Reconstruction}{The finite-zero factorization takes place near the compact spectrum, not over all of \(\mathbb C\). The analytic functional calculus supplies invertibility of the zero-free factor.}
\GSfield{Fate}{This functional-calculus consequence uses the theorem of \citet{bernsteinrobinson1966,halmos1966polynomial}. The zero entire function must be excluded: allowing it would make the assumption hold for every bounded operator and simply restate the unsolved general problem.}
\GSfield{Present status}{As of November 2026: established with the hypotheses in the modern statement.}
\GSfield{Assessment}{literal statement: correct with missing hypotheses; publication match: partial; verification: read the result; development: annotated; confidence: high}
\end{GSitem}
\begin{GSitem}{1964-50.12}{F}{The de Branges--Rovnyak compact-defect assertion}
\GSfield{Kind; source}{claim; \GSpage{P000088}}
\GSfield{The notes}{After noting the error in their general announcement, the manuscript claims a consequence under \(\|T\|<1\) or \(\|T\|\le1\) and compact \(I-T^*T\). The scan does not securely distinguish the two norm signs.}
\GSfield{Statement}{If \(\|T\|<1\), then \(I-T^*T\ge(1-\|T\|^2)I\); on infinite-dimensional \(H\), this operator cannot be compact. The strict reading is vacuous. The weak reading is a materially different assertion about a contraction with compact defect and requires a separate theorem; it is not certified by the retracted general proof.}
\GSfield{Reconstruction}{The coercivity argument diagnoses the strict reading without choosing an uncertain glyph. No current open problem is manufactured from the weak reading.}
\GSfield{Fate}{Rovnyak\textquotesingle s author bibliography identifies the 1964 announcement and its 1965 retraction \citep{debrangesrovnyak1964existence,debrangesrovnyak1965correction}. It does not verify the exact compact-defect corollary. The general retraction and the unresolved norm sign are kept separate.}
\GSfield{Present status}{As of November 2026: the strict hypothesis is impossible in the intended infinite-dimensional setting; an exact source for the non-strict compact-defect assertion was not verified.}
\GSfield{Assessment}{literal statement: not determinable; publication match: thematic only; verification: publisher metadata checked; development: annotated; confidence: low}
\end{GSitem}
\begin{GSitem}{1964-50.13}{E}{Trace-class imaginary parts and triangular models}
\GSfield{Kind; source}{remark; \GSpage{P000088}}
\GSfield{The notes}{Sakhnovich, Livshits and Brodskii are named for \(A+iB\), with \(A\) self-adjoint and \(B\) small in some sense, for example nuclear.}
\GSfield{Statement}{One precise realization is the following. Let \(T=A+iB\) be bounded on a separable complex Hilbert space, with \(A=A^*\), \(B=B^*\), and \(B\) trace class. The restriction of \(T\) to the orthogonal complement of its maximal reducing self-adjoint part is unitarily equivalent to the corresponding non-self-adjoint part of a Liv\v{s}ic triangular model. The additional reducing component is self-adjoint. Here nuclear means trace class; small operator norm is not a substitute for this hypothesis.}
\GSfield{Reconstruction}{The bounded separable Hilbert-space setting and the canonical self-adjoint real and imaginary parts are supplied explicitly. The nuclear example in the source identifies this established theorem, but its unspecified additional classes are not reconstructed. The stronger Hilbert--Schmidt criterion is reported by a later survey; its attribution was not independently confirmed in the cited 1959 original.}
\GSfield{Fate}{\citet[Introduction, p.~145; \S11, Theorem~4, pp.~180--181]{livsic1954spectral} gives the triangular-model theorem for trace-class imaginary part. Its introduction dates the preliminary announcements to two 1952 notes. \citet[Introduction and theorem, pp.~180--181]{sakhnovich1959triangular} recalls this nuclear case and constructs a triangular representation under an explicit invariant-subspace interpolation condition. The latter paper was received 6 December 1957 and published in 1959. The broader survey discussion \citep[\S7, pp.~57--59]{dymkatsnelson2007schur} is kept distinct from the original results actually read.}
\GSfield{Present status}{As of November 2026: the nuclear example is an established pre-seminar result. The match is partial because the notebook does not identify its further small-imaginary-part classes or a particular theorem of each named author.}
\GSfield{Assessment}{literal statement: correct with missing hypotheses; publication match: partial; verification: read the result; development: annotated; confidence: high}
\end{GSitem}
\begin{GSitem}{1964-50.14}{E}{Contractions with neither forward nor adjoint powers strongly stable}
\GSfield{Kind; source}{theorem; \GSpage{P000088}}
\GSfield{The notes}{Both arrows in \(T^n\not\to0\), \((T^*)^n\not\to0\) are crossed; the norm condition is \(\|T\|\le1\). The result is attributed to Nagy and Foias.}
\GSfield{Statement}{If \(T\) is a contraction on a complex Hilbert space with \(\dim H>1\), and neither \(T^n\) nor \((T^*)^n\) converges strongly to zero, then \(T\) has a nonzero proper closed invariant subspace.}
\GSfield{Reconstruction}{Strong nonconvergence means failure for at least one vector for each sequence; it does not require every vector to have a nonzero limiting norm.}
\GSfield{Fate}{This is exactly \citet[\S6, Theorem~3 and Remark~(ii), p.~314]{sznagyfoias1964ix}; footnote~29 handles the unitary case and dimension restriction. The paper was received 1 May 1964 and published in 1964. Its 1965 correction explicitly says that the error in Proposition~5.4 does not affect this theorem or its proof \citep[p.~193]{sznagyfoias1965corrections}.}
\GSfield{Present status}{As of November 2026: established with the hypotheses in the modern statement.}
\GSfield{Assessment}{literal statement: correct with missing hypotheses; publication match: exact; verification: read the result; development: annotated; confidence: high}
\end{GSitem}
\begin{GSitem}{1964-50.15}{E}{Donoghue lattice of finite initial spans}
\GSfield{Kind; source}{construction; \GSpage{P000088}}
\GSfield{The notes}{The lattice \(\{0,1,2,\ldots,\infty\}\) is attributed to Donoghue (1957). The upper-diagonal weights begin \(1/2,1/4\), followed by an uncertain \(1/8\) or \(1/6\).}
\GSfield{Statement}{On \(\ell^2(\mathbb N_0)\), the backward weighted shift \(Be_0=0\), \(Be_n=2^{-n}e_{n-1}\) for \(n\ge1\) is compact and quasinilpotent, and its closed invariant subspaces are exactly \(0\), \(\operatorname{span}\{e_0,\ldots,e_{m-1}\}\) for \(m\ge1\), and the whole space.}
\GSfield{Reconstruction}{The geometric weight sequence is a literature-supported interpretation, not a confident decipherment of its overwritten third entry. The endpoints are the zero and whole subspaces.}
\GSfield{Fate}{\citet[Examples~1--2, pp.~1031--1033]{donoghue1957lattice} proves the forward-shift tail lattice and its adjoint version. A scalar multiple yields the notebook\textquotesingle s normalization. The 1957 date is correct; the paper was received 25 June 1956.}
\GSfield{Present status}{As of November 2026: established with the hypotheses in the modern statement.}
\GSfield{Assessment}{literal statement: correct with missing hypotheses; publication match: exact; verification: read the result; development: annotated; confidence: high}
\end{GSitem}
\begin{GSitem}{1964-50.16}{B}{Weight conditions for the Donoghue invariant-subspace lattice}
\GSfield{Kind; source}{problem; \GSpage{P000088}}
\GSfield{The notes}{The exact conditions on the superdiagonal numbers that give precisely the Donoghue chain are said to be unknown.}
\GSfield{Statement}{For a bounded sequence of nonzero weights \((w_n)_{n\ge1}\), determine criteria expressed in the weights for the backward shift \(Be_n=w_ne_{n-1}\), \(Be_0=0\), to have only finite initial spans and the whole space as closed invariant subspaces. Equivalently its forward-shift adjoint must be unicellular, with only coordinate tails.}
\GSfield{Reconstruction}{Nonzero bounded weights and the Hilbert adjoint equivalence supply a precise classification question. The broad request for conditions is not silently replaced by an unverified demand for a novel universal closed-form criterion.}
\GSfield{Fate}{\citet[Theorem~1 and subsequent counterexamples]{nikolskii1967weighted} gives weight conditions and examples delimiting them. \citet[\S1; Lemmas~2.1--2.2]{harrison1971unicellularity} develops further criteria. The opened abstract of \citet{yakubovich1985weighted} proves that positive moduli monotonically decreasing to zero suffice on \(\ell^p\), \(1\le p<\infty\); Hilbert adjoints give the backward form. \citet[p.~232]{latapokhriyalsingh2023powers} also credits Nikol\textquotesingle skii\textquotesingle s 1967 and 1968 work with additional necessary and sufficient conditions. The 1968 original was verified only at metadata level \citep{nikolskii1968basicity}; its exact scope is not guessed.}
\GSfield{Present status}{As of November 2026: developed and answered in substantial verified weight classes. The generality of a complete explicit characterization for arbitrary bounded nonmonotone weights was not ascertained in this audit. That verification gap is not classified as an established open problem.}
\GSfield{Assessment}{literal statement: correct; publication match: partial; verification: read the result; development: annotated; confidence: medium; problem \hyperref[prob:GP-1964-14]{GP-1964-14}}
\end{GSitem}
\begin{GSitem}{1964-50.17}{B}{An invariant-subspace chain indexed by all integers}
\GSfield{Kind; source}{question; \GSpage{P000088}}
\GSfield{The notes}{The question whether there is an operator with the structure of all integers is followed by unknown.}
\GSfield{Statement}{There is a bounded operator on \(\ell^2(\mathbb Z)\) with full invariant-subspace lattice \(\{0,H\}\cup\{M_m:m\in\mathbb Z\}\), where \(M_m=\overline{\operatorname{span}}\{e_j:j\ge -m\}\). In particular take \(Se_j=\exp(-2|2j+1|)e_{j+1}\). The \(M_m\) increase strictly, have intersection zero and dense union. The literal ordered set \(\mathbb Z\) needs the endpoints \(0,H\).}
\GSfield{Reconstruction}{The endpoints and reversal of indexing are editorial. The concrete choice of Domar\textquotesingle s weights and the unitary conversion were checked independently in this audit; no novelty is claimed.}
\GSfield{Fate}{\citet[Theorem~2, pp.~134--135; definitions p.~133]{domar1981translation} supplies an exact construction through weighted translation. Set \(w_n=\exp(-2n|n|)\). Then \(\log w\) is convex for \(n\le0\), concave for \(n\ge0\), \(\sum_{n\ne0}w_n^{1/n}=2\sum_{n\ge1}e^{-2n}<\infty\), and its quadratic limits are \(2>\log3\) at minus infinity and \(-2<-\log3\) at plus infinity. The unitary map \(c\mapsto(w_nc_n)\) turns right translation into the displayed shift. Domar proves that all invariant subspaces are the standard tails. Received 9 June 1980; published 1981.}
\GSfield{Present status}{As of November 2026: the completed integer-chain existence question is solved affirmatively by the established bilateral weighted-shift construction. The literal omission of least and greatest lattice elements is only a shorthand defect.}
\GSfield{Assessment}{literal statement: correct with missing hypotheses; publication match: exact; verification: read the result; development: annotated; confidence: high; problem \hyperref[prob:GP-1964-15]{GP-1964-15}}
\end{GSitem}
\begin{GSitem}{1964-50.18}{E}{Volterra integration and a continuous chain}
\GSfield{Kind; source}{construction; \GSpage{P000088}}
\GSfield{The notes}{The interval \([0,1]\) is represented by \((Vf)(x)=\int_0^x f(y)\,dy\), with the result dated to the 1950s.}
\GSfield{Statement}{On \(L^2(0,1)\), the nonzero proper closed invariant subspaces of the Volterra operator are precisely \(M_a=\{f:f=0\ \text{a.e. on }(0,a)\}\), \(0<a<1\). With \(M_0=H\), \(M_1=0\), these form a continuous totally ordered lattice, reverse-ordered by \(a\), hence order-isomorphic to \([0,1]\).}
\GSfield{Reconstruction}{The direction of inclusion and the endpoint subspaces are supplied; the Hilbert case is \(p=2\).}
\GSfield{Fate}{\citet[Example~3, pp.~1033--1034]{donoghue1957lattice} gives the theorem and proof, including \(L^p\), \(1\le p<\infty\), and an appropriate continuous-function version. The existence assertion was established well before the seminar.}
\GSfield{Present status}{As of November 2026: established with the hypotheses in the modern statement.}
\GSfield{Assessment}{literal statement: correct; publication match: exact; verification: read the result; development: annotated; confidence: high}
\end{GSitem}
\begin{GSitem}{1964-50.19}{E}{A measure-Volterra lattice is not determined by the support as written}
\GSfield{Kind; source}{conjecture; \GSpage{P000088}}
\GSfield{The notes}{A measure on the whole real line is drawn, \((Af)(x)=\int_0^x f(y)\,d\mu(y)\), and the lattice is conjectured to be isomorphic to \(\operatorname{supp}\mu\). The normalization resembles total mass one but is not secure.}
\GSfield{Statement}{The literal conjecture is false even for an atomless probability measure: take \(d\mu=\tfrac12\mathbf1_{[-1,1]}dx\), interpreting the integral from zero in the oriented sense for negative \(x\). The subspaces supported on \([-1,0]\) and \([0,1]\) are closed, nonzero, proper and invariant, since every integration interval stays on its side of zero. They are incomparable, whereas \(\operatorname{supp}\mu=[-1,1]\) is linearly ordered.}
\GSfield{Reconstruction}{The counterexample has no atoms, so endpoint conventions cannot explain away the failure. Its total mass is one. The boundedness of this particular integral operator follows directly from Cauchy--Schwarz.}
\GSfield{Fate}{A repaired one-sided statement is classical: a positive continuous probability density on \([a,b]\), with lower integration limit \(a\), is transformed by its cumulative distribution into ordinary Volterra integration. The latter lattice is given by \citet[Example~3]{donoghue1957lattice}. This does not validate the two-sided source formula, and no general unresolved measure classification is inferred from that false conjecture.}
\GSfield{Present status}{As of November 2026: false under the literal two-sided formulation; the indicated regular one-sided repair is established.}
\GSfield{Assessment}{literal statement: false as written; publication match: partial; verification: read the result; development: annotated; confidence: high}
\end{GSitem}
\begin{GSitem}{1964-50.20}{C}{Bishop operators as possible invariant-subspace counterexamples}
\GSfield{Kind; source}{conjecture; \GSpage{P000089}}
\GSfield{The notes}{The weighted composition operator is specialized to \(Tf(x)=x f(x+\alpha\bmod1)\), irrational \(\alpha\). It is proposed as having no invariant subspace, equivalently every vector cyclic.}
\GSfield{Statement}{For irrational \(\alpha\in(0,1)\), let \(T_\alpha f(x)=x f(\{x+\alpha\})\) on complex \(L^2(0,1)\). Does every such operator have a nonzero proper closed invariant subspace? The source\textquotesingle s contrary suggestion is false for many parameters, but the universal affirmative question is a distinct remaining problem. The equivalence with cyclicity requires every nonzero vector, not the zero vector.}
\GSfield{Reconstruction}{The universal affirmative invariant-subspace question is a modern reformulation. The source suggests a counterexample without precisely delimiting the parameter set; its statement that every vector is cyclic must in any case exclude the zero vector.}
\GSfield{Fate}{Davie\textquotesingle s 1974 theorem covers non-Liouville parameters; \citet[Theorem~3.7]{chamizo2020bishop} extends this to convergent denominators satisfying \(\log q_{j+1}=O(q_j/(\log q_j)^3)\), and hence includes some Liouville numbers. The same paper proves biquasitriangularity and norm approximation by nilpotent operators, neither of which alone produces an invariant subspace. \citet[Introduction]{behani2023bishop} explicitly retains the all-irrational question.}
\GSfield{Present status}{As of November 2026: the checked established results cover the stated arithmetic class, not every irrational parameter. No established all-parameter resolution was verified; the original no-subspace suggestion is already disproved for non-Liouville parameters.}
\GSfield{Assessment}{literal statement: false as written; publication match: partial; verification: read the result; development: annotated; confidence: high; problem \hyperref[prob:GP-1964-16]{GP-1964-16}}
\GSfield{Worth developing}{Yes. Exact finite-rotation products and quantitative rational approximation give concrete test cases. An invariant-subspace proof for the exceptional arithmetic regime still requires an argument beyond the finite-orbit continuity developed here.}
\end{GSitem}
\begin{GSitem}{1964-50.21}{C}{Cyclic vectors for irrational Bishop operators}
\GSfield{Kind; source}{problem; \GSpage{P000089}}
\GSfield{The notes}{After conjecturing that every vector is cyclic, the speaker says that not even one cyclic vector is known.}
\GSfield{Statement}{For each irrational \(\alpha\in(0,1)\), does there exist \(f\in L^2(0,1)\) with \(\overline{\operatorname{span}}\{T_\alpha^n f:n\ge0\}=L^2(0,1)\)? A stronger concrete question asks when the constant function \(1\) is cyclic. Existence of one cyclic vector is far weaker than every nonzero vector being cyclic.}
\GSfield{Fate}{\citet[Theorems~4.4--4.5, 4.9 and 4.16]{behani2023bishop} establishes rational cyclicity, a co-meager set of cyclic irrational parameters, and sufficient continued-fraction conditions. The preprint circulated in July 2022 and the paper appeared in 2023. Its Question~5.4 asks for an explicit approximation threshold. The chapter gives a bounded rational-fibre certificate and continuity estimate, not an all-irrational proof. A 2026 preprint lead was inaccessible and is not used as an established result.}
\GSfield{Present status}{As of November 2026: cyclicity is established for a co-meager parameter set and explicit sufficient classes in the checked work; universal irrational cyclicity was not verified. The historical absence of any known vector is obsolete.}
\GSfield{Assessment}{literal statement: correct; publication match: partial; verification: read the result; development: developed in this chapter; confidence: high; problem \hyperref[prob:GP-1964-17]{GP-1964-17}}
\GSfield{Worth developing}{Yes. Rational cyclicity, exact finite products and strong continuity reduce selected questions to finite polynomial approximations with controlled error. The chapter develops such certificates and states their remaining arithmetic and conditioning obstacles explicitly.}
\end{GSitem}
\begin{GSitem}{1964-50.22}{F}{The unspecified Ornstein cyclicity example}
\GSfield{Kind; source}{claim; \GSpage{P000089}}
\GSfield{The notes}{An operator on \(L^2(0,1)\) is attributed to Ornstein, with every vector in \(L^{2+\varepsilon}\) cyclic. A nonsingular transformation without invariant measure is mentioned, but no operator is written.}
\GSfield{Statement}{At minimum the cyclicity assertion must concern nonzero vectors and specify whether \(\varepsilon>0\) is fixed or varies. The operator, weight, transformation, and intended class of invariant measures are absent. A nonsingular transformation alone does not specify a bounded weighted-composition operator with the asserted cyclicity.}
\GSfield{Fate}{Targeted searches on Ornstein, cyclic vectors, \(L^{2+\varepsilon}\), nonsingular transformations and invariant measures did not produce an opened primary theorem with these exact hypotheses. This anecdote is not identified with unrelated Ornstein results in ergodic theory, and no claimed counterexample to the Hilbert invariant-subspace problem is extracted from it.}
\GSfield{Present status}{As of November 2026: insufficiently specified and not securely identified; retained with explicit search limits.}
\GSfield{Assessment}{literal statement: not determinable; publication match: none found; verification: unresolved lead; development: annotated; confidence: low}
\end{GSitem}
\begin{GSitem}{1964-50.23}{F}{Why invariant subspaces, and what could replace them?}
\GSfield{Kind; source}{question; \GSpage{P000089}}
\GSfield{The notes}{Gelfand asks why invariant subspaces are needed and what could replace them.}
\GSfield{Statement}{This is a methodological question about structural invariants and operator models, not a quantified mathematical conjecture. Spectral measures, characteristic functions, similarity invariants and invariant-subspace lattices answer different versions of it.}
\GSfield{Fate}{The next intervention gives one concrete contemporary response through analytic characteristic functions. No solved/open status is meaningful for the unrestricted methodological question itself.}
\GSfield{Present status}{As of November 2026: historical methodological discussion, not an independently unresolved theorem.}
\GSfield{Assessment}{literal statement: not determinable; publication match: none found; verification: indirect; development: none; confidence: high}
\end{GSitem}
\begin{GSitem}{1964-50.24}{E}{Analytic characteristic functions and the contraction model}
\GSfield{Kind; source}{construction; \GSpage{P000089}; Unidentified participant (possibly C. Foiaş)}
\GSfield{The notes}{An English-speaking visitor with an uncertain F-name describes an analytic contraction from \(E\) to \(E_*\), Hardy boundary functions and analogues of Jordan blocks. The final finite-dimensional affine-map sentence does not specify the equivalence relation.}
\GSfield{Statement}{For a completely nonunitary contraction \(T\), its characteristic function \(\Theta_T(z):\mathcal D_T\to\mathcal D_{T^*}\) is analytic and contractive for \(|z|<1\). It determines \(T\) up to unitary equivalence, modulo constant unitary identifications of the defect spaces. The standard model uses \([H^2(\mathcal D_{T^*})\oplus\overline{\Delta L^2(\mathcal D_T)}]\ominus\{\Theta_T f\oplus\Delta f:f\in H^2(\mathcal D_T)\}\), with \(\Delta=(I-\Theta_T^*\Theta_T)^{1/2}\).}
\GSfield{Reconstruction}{The missing upper Fourier-integral limit is \(2\pi\); the positive exponent \(e^{int}\), \(n\ge1\), correctly kills negative Fourier coefficients. Complete nonunitarity and the defect-space constraints are essential. The separate finite-dimensional affine-map/Jordan analogy remains underdetermined.}
\GSfield{Fate}{\citet[Theorem~1, p.~50; Theorem~3, p.~58]{sznagyfoias1964viii} gives the model and complete unitary invariant; it was received April 1963, revised December 1963 and published in 1964. \citet[Theorem~1 and formulas (3.16)--(3.17), pp.~294--295]{sznagyfoias1964ix} gives the factorization/invariant-subspace correspondence. This is an exact mathematical match, but not independent proof of the visitor\textquotesingle s identity.}
\GSfield{Present status}{As of November 2026: established with the hypotheses in the modern statement.}
\GSfield{Assessment}{literal statement: correct with missing hypotheses; publication match: exact; verification: read the result; development: annotated; confidence: medium}
\end{GSitem}

\section{Developments}\label{dev:1964}
\subsection{Bishop operators: exact finite certificates for the cyclicity question}
\label{dev:bishop}

\paragraph{The question supported by the notebook.}
On P000089 the notes give, on complex \(L^2(0,1)\),
\[
 (B_\alpha f)(x)=x f(\{x+\alpha\}),\qquad
 \alpha\in(0,1)\setminus\mathbb Q,
\]
where \(\{\cdot\}\) denotes fractional part.  The margin says
\ru{Изучить подробно этот оператор}.  The text then juxtaposes the suggestion
\ru{нет инв. подпр.}\ with the admission
\ru{неизв. ни одного циклич. вектора}.  These are different quantifiers.
For a bounded operator \(B\), absence of a nonzero proper closed invariant
subspace is equivalent to \emph{every nonzero} vector being cyclic; existence
of \emph{one} cyclic vector is much weaker.  Indeed, the closed linear span
of the orbit of a nonzero vector is an invariant subspace, and conversely a
vector inside a proper invariant subspace cannot be cyclic.  The zero vector
must be excluded from the notebook's shorthand \(\forall\) assertion.

The precise questions are therefore: for which irrational \(\alpha\) does
\(B_\alpha\) have a nonzero proper closed invariant subspace, and for which
does it possess a cyclic vector?  The stronger, plausible cyclicity question
considered here asks whether the \emph{specific} vector \(1\) is cyclic for
every irrational \(\alpha\).  This last universal formulation is an editorial
research question, not a statement attributed to Halmos or Gelfand.
No universal answer is claimed below.

\paragraph{Verified literature and the scope of the development.}
\citet[Theorem~3.7]{chamizo2020bishop}, first circulated on 19 December 2018
and published in 2020, prove the existence of nontrivial hyperinvariant
subspaces if the continued-fraction denominators satisfy
\[
 \log q_{n+1}=O\!\left(\frac{q_n}{(\log q_n)^3}\right).
\]
Their Theorem~4.1 supplies an obstruction to the particular Beurling-algebra
method for sufficiently extreme approximability; it does not prove absence
of invariant subspaces.  \citet[Theorems~4.4, 4.5, 4.9, 4.16]{behani2023bishop},
first posted on 23 July 2022 and published in 2023, give the rational
fibre-determinant criterion, rational cyclicity of holomorphic vectors
nonzero at zero, generic irrational cyclicity, and a sufficient large-gap
condition.  Locators here refer to the opened arXiv version.

A search dated 12 November 2026 also found the abstract of
\citet{ma2026rational}, which announces quantitative determinant and cyclicity
results and specifically mentions \((r,q)=(9,16)\).  Its primary full-text
and abstract endpoints did not open in this run.  It is marked
\emph{unverified}; none of its announced theorems is used as an assumption.
Its existence prevents a novelty claim for the same finite example or for
the general programme of explicit gap estimates.  No all-irrational
cyclicity or invariant-subspace theorem was verified in the sources checked
through 12 November 2026.

The bounded contribution is an exact-arithmetic certificate format, an
independent proof of the transfer estimate, and a fully specified terminating
procedure for a sufficient gap function.  The rational and generic
cyclicity conclusions are established literature, re-proved here for the
certificate's foundations.  The label \ours{} means that this companion
supplies and checks a proof; it does not assert historical priority.

\paragraph{Rational fibres and the endpoint determinant.}
Fix relatively prime integers \(1\le r<q\), \(q\ge2\), and write
\(a_i=(ir\bmod q)/q\), with indices read modulo \(q\).
Splitting \([0,1)\) into the intervals \(a_i+[0,1/q)\) identifies
\(L^2(0,1)\) unitarily with
\(L^2((0,1/q);\mathbb C^q)\) by
\[
 (Uf)(t)_i=f(t+a_i),\qquad 0<t<1/q.
\]
Lebesgue measure is used on the base and the ordinary Euclidean norm on the
fibres, so no normalization factor is missing.  In these coordinates,
\[
 (\mathcal B(t)v)_i=(t+a_i)v_{i+1},\qquad
 \mathcal B(t)^q=m_q(t)I,
 \quad m_q(t)=\prod_{j=0}^{q-1}\left(t+\frac jq\right).
\]
The continuous function \(m_q\) is strictly increasing on \([0,1/q]\).
Let \(K_{r,q}(t)\) be the matrix whose column \(k\), \(0\le k<q\), is
\(\mathcal B(t)^k\mathbf1\), and let \(D_{r,q}(t)=\det K_{r,q}(t)\).

\begin{lemma}[A nonzero endpoint determinant]\ours{}
\label{lem:bishop-endpoint}
For the coprime pair above,
\[
 K_{r,q}(t)_{i,k}=\prod_{h=0}^{k-1}(t+a_{i+h}),\qquad
 D_{r,q}(0)=(-1)^{(q-1)(q-2)/2}\prod_{j=1}^{q-1}a_j^{\,j}\ne0.
\]
An empty product is one.  Consequently \(D_{r,q}\) is a nonzero polynomial
and has only finitely many zeros in the fundamental interval.
\end{lemma}
\begin{proof}
The product formula follows by iterating the displayed weighted cyclic
matrix.  Since \(a_0=0\) and every other \(a_i\) is positive, row zero at
\(t=0\) is \((1,0,\ldots,0)\).  In row \(i\ge1\), the last possibly
nonzero entry is column \(q-i\), equal to
\(\prod_{j=i}^{q-1}a_j\).  Delete the first row and column and reverse the
remaining \(q-1\) rows.  The resulting triangular matrix has these entries
on its diagonal.  Reversal has sign
\((-1)^{(q-1)(q-2)/2}\); multiplying the diagonal entries counts \(a_j\)
exactly \(j\) times.  This proves the formula and all its consequences.
\end{proof}

\begin{proposition}[Rational cyclicity of the constant vector]\ours{}
\label{prop:bishop-rational}
The vector \(1\) is cyclic for \(B_{r/q}\) on complex \(L^2(0,1)\).
\end{proposition}
\begin{proof}
Let \(\mathcal C\) be the closed span of \(B_{r/q}^n1\), \(n\ge0\),
in the fibre representation.  For each \(0\le k<q\) it contains
\(p(m_q(t))\mathcal B(t)^k\mathbf1\) for every polynomial \(p\), since
\(B_{r/q}^{q\ell+k}1\) has precisely that form for \(p(z)=z^\ell\).
The map \(m_q\) is a homeomorphism onto a compact real interval.
Polynomial approximation therefore gives every continuous scalar
coefficient \(h(t)\) times this bounded vector field.  Approximation of
bounded measurable coefficients by continuous functions in the appropriate
finite weighted \(L^2\) measure gives the same conclusion for bounded
measurable \(h\).

For arbitrary \(g\in L^2((0,1/q);\mathbb C^q)\), define
\[
 E_N=\{t:\|K_{r,q}(t)^{-1}\|\le N,\ \|g(t)\|\le N\},
 \qquad g_N=\mathbf1_{E_N}g.
\]
The finitely many singular fibres can be discarded.  By the lemma,
\(E_N\) exhausts the base up to a null set; hence \(g_N\to g\) in
\(L^2\).  The components of
\(\mathbf1_{E_N}K_{r,q}^{-1}g\) are bounded measurable functions.
Expressing \(g_N\) in the columns of \(K_{r,q}\) puts \(g_N\) in
\(\mathcal C\).  Closedness now gives \(g\in\mathcal C\).
\end{proof}

\paragraph{Why a uniformly invertible fibre matrix is the wrong requirement.}
The proof needs almost-everywhere invertibility and truncation, not a
positive lower bound for \(|D_{r,q}(t)|\) on the whole interval.  This
distinction is testable without estimating any irrational orbit.

\begin{example}[A singular interior fibre]\oursComputed{}
\label{ex:bishop-singular}
For \((r,q)=(9,16)\), exact arithmetic gives
\(D_{9,16}(0)<0<D_{9,16}(1/16)\).  Thus at least one
\(t\in(0,1/16)\) has \(D_{9,16}(t)=0\), although the constant vector
is cyclic by Proposition~\ref{prop:bishop-rational}.
\end{example}
For reproducibility, use the integer matrix
\[
 C(y)_{i,k}=\prod_{h=0}^{k-1}\bigl(y+(9(i+h)\bmod16)\bigr),
 \qquad \det C(y)=16^{120}D_{9,16}(y/16).
\]
The two exact determinants, including their integer signs, are in
\nolinkurl{evidence/independent_bishop_check.json}.  The first computation
used rational symbolic Gaussian elimination; the second used a separately
implemented fraction-free integer Bareiss algorithm, with every division
checked for exactness.  Both agree.  Continuity and the intermediate value
theorem, rather than a plotted sample, establish the existence of the zero.
No assertion is made about its uniqueness, multiplicity, or the location
of zeros for all other coprime pairs.

For small denominators the same program obtains, for example,
\[
 D_{1,2}(t)=\frac12,\quad
 D_{1,3}(t)=-\frac{9t+4}{27},\quad
 D_{2,3}(t)=-\frac{9t+2}{27}.
\]
Checks of the endpoint formula and matrix-power identity for all 31 coprime
pairs with \(2\le q\le10\) passed.  Those checks are finite checks;
Lemma~\ref{lem:bishop-endpoint} is the independent all-\(q\) proof.

\paragraph{A quantitative transfer bound that accounts for wrap points.}
For \(k\ge0\), put \(w_{\alpha,k}=B_\alpha^k1\).  Thus
\(w_{\alpha,0}=1\) and
\[
 w_{\alpha,k}(x)=\prod_{j=0}^{k-1}\{x+j\alpha\}\quad(k\ge1).
\]
All factors lie in \([0,1]\).  Treating them as globally Lipschitz
functions of \(\alpha\) would overlook their jumps at wrap points.

\begin{lemma}[Finite-orbit continuity]\ours{}
\label{lem:bishop-continuity}
For all real \(\alpha,\beta\) and integers \(k\ge0\),
\[
 \|w_{\alpha,k}-w_{\beta,k}\|_2
 \le |\alpha-\beta|^{1/2}\sum_{j=1}^{k-1}\sqrt j
 \le \frac{k(k-1)}2|\alpha-\beta|^{1/2},
\]
where the sum and final bound are zero for \(k=0,1\).
Consequently, if \(Q(z)=\sum_{k=0}^N c_kz^k\), then
\[
 \|Q(B_\alpha)1-Q(B_\beta)1\|_2
 \le R(Q)|\alpha-\beta|^{1/2},\qquad
 R(Q)=\sum_{k=0}^N|c_k|\frac{k(k-1)}2.
\]
\end{lemma}
\begin{proof}
For \(0\le d\le1\), a sawtooth translated by \(d\) changes by
\(d\) on an interval of length \(1-d\) and by \(d-1\) on the remaining
interval.  Its squared \(L^2\) difference is
\((1-d)d^2+d(1-d)^2=d(1-d)\).  The same identity applies to a
negative shift by replacing it by its absolute value modulo one.
Translation invariance then gives
\[
 \|\{x+j\alpha\}-\{x+j\beta\}\|_2^2
 \le j|\alpha-\beta|.
\]
For \(j|\alpha-\beta|\le1\), use the identity; for a larger difference,
the identity is at most \(1/4\le j|\alpha-\beta|\).
Expand the difference of the two products by telescoping.  Every other
factor has absolute value at most one, so the triangle inequality in
\(L^2\) gives the first bound.  For integers \(j\ge1\),
\(\sqrt j\le j\), giving the second.  Another triangle inequality,
now over the coefficients of \(Q\), proves the final claim.
\end{proof}

\begin{example}[A complete finite approximation certificate]\oursComputed{}
\label{ex:bishop-certificate}
Let
\[
\begin{split}
 Q(z)=\frac{1}{1032182056}\bigl(&-24558043+609112200z
       -604809944z^2\\
       &+2030693280z^3-987436800z^4\bigr).
\end{split}
\]
Direct piecewise polynomial integration gives
\[
 \|Q(B_{1/2})1-x^2\|_2^2
 =\frac{2117561}{24772369344}<\frac1{10000},
 \qquad R(Q)=\frac{1577688823}{129022757}<13.
\]
Therefore, for every real \(\alpha\),
\[
 |\alpha-\tfrac12|<\frac1{1690000}
 \quad\Longrightarrow\quad \|Q(B_\alpha)1-x^2\|_2<\frac1{50}.
\]
This is a statement about this target and this accuracy, valid for rational
and irrational parameters in the indicated neighbourhood.
\end{example}
Indeed, Lemma~\ref{lem:bishop-continuity} bounds the additional error by
\(13/1300=1/100\), and the rational approximation error is strictly less
than \(1/100\).  It does not follow from a single such certificate that
\(1\) is cyclic for any fixed irrational parameter.

The coefficients were obtained by exact least squares in the span of the
first five orbit vectors.  One program used the normal equations and then
integrated the squared residual independently.  A second program used only
Python's rational-number arithmetic, rebuilt the orbit polynomials, and
reintegrated both cells.  It obtained the identical rational number.
The same two checks agree on all five certificates in the output:
\begin{center}
\begin{tabular}{@{}rrrcc@{}}
\toprule
\(r\)&\(q\)&Degree&\(L^2\) error for \(x^2\)&\(R(Q)\)\\
\midrule
1&2&2&0.07448495&0.02173913\\
1&2&4&0.009245581&12.22799\\
1&2&8&0.0007486373&332.0757\\
1&3&6&0.01199983&91.78700\\
2&3&6&0.008914569&58.38182\\
\bottomrule
\end{tabular}
\end{center}
These rounded values describe exact fractions saved with the code.  The
tradeoff is visible: a smaller rational error can require a much larger
coefficient sum, reducing the neighbourhood to which the certificate
transfers.  Small residuals alone do not control irrational parameters.

\paragraph{A terminating sufficient criterion, with all quantifiers stated.}
The finite certificates can be organized into a theorem.  This produces an
algorithmically computable function; no modest closed-form growth bound for
it is asserted.

\begin{theorem}[A computable continued-fraction gap rule]\ours{}
\label{thm:bishop-gap}
There is a total computable function \(\Psi:\{2,3,\ldots\}\to\mathbb N\)
with the following property.  If an irrational \(\alpha\in(0,1)\), with
continued-fraction convergents \(r_n/q_n\), satisfies
\(q_{n+1}>\Psi(q_n)\) for infinitely many \(n\) with \(q_n\ge2\), then
the vector \(1\) is cyclic for \(B_\alpha\) on complex \(L^2(0,1)\).
\end{theorem}
\begin{proof}
For each fixed \(q\), each \(1\le r<q\) coprime to \(q\), and each
\(1\le j\le q\), enumerate polynomials \(Q\in\mathbb Q[z]\) in a fixed
effective order until
\[
 \|Q(B_{r/q})1-x^{j-1}\|_2^2<\frac1{4q^2}.
 \tag{*}
\]
This is a decidable strict rational inequality: split \([0,1)\) at the
points \(i/q\), express the orbit products as rational polynomials on
each cell, and integrate their squared residual.  To specify the enumeration,
at stage \(h=1,2,\ldots\) list degrees \(0,\ldots,h\) and all coefficient
tuples \(c_k=a_k/b_k\), \(|a_k|\le h\), \(1\le b_k\le h\), in
lexicographic order of the integer tuples; repetitions are harmless.
Every rational polynomial eventually occurs.  The search terminates
by Proposition~\ref{prop:bishop-rational}.  Since the target and the orbit
vectors are real, take real parts of an approximating complex polynomial
and then approximate its finitely many coefficients by rationals; the
strict target error remains attainable.

Call the first polynomial so found \(Q_{r,q,j}\), and define the positive
rational number
\[
 A_q=\max_{\substack{1\le r<q,\ (r,q)=1\\1\le j\le q}}
       \bigl(1+R(Q_{r,q,j})\bigr),
 \qquad \Psi(q)=\lceil4q A_q^2\rceil.
\]
Only finitely many terminating searches occur for each input \(q\), so
\(\Psi\) is total and computable.  The usual convergent inequality is
\(|\alpha-r_n/q_n|<1/(q_nq_{n+1})\).  On the stipulated subsequence it
therefore gives
\[
 |\alpha-r_n/q_n|<\frac1{4q_n^2 A_{q_n}^2}.
\]
Combining (*) with the continuity lemma yields, for every \(j\le q_n\),
\[
 \|Q_{r_n,q_n,j}(B_\alpha)1-x^{j-1}\|_2<\frac1{q_n}.
\]
The denominators on any infinite subsequence tend to infinity.  For each
fixed \(j\), the monomial \(x^{j-1}\) thus belongs to the closed cyclic
subspace.  Polynomials are dense in \(L^2(0,1)\), proving cyclicity.
\end{proof}

\begin{corollary}[The fixed-vector generic conclusion]\ours{}
The set of irrational \(\alpha\in(0,1)\) for which \(1\) is cyclic for
\(B_\alpha\) contains a dense \(G_\delta\) set.
\end{corollary}
\begin{proof}
For each positive pair \(j,m\), take the union, over rational-coefficient
polynomials \(Q\), of the open sets
\(\{\alpha:\|Q(B_\alpha)1-x^{j-1}\|_2<1/m\}\).
Openness follows from Lemma~\ref{lem:bishop-continuity}; each set in the
countable family of unions contains every rational \(\alpha\in(0,1)\)
by Proposition~\ref{prop:bishop-rational}, so is dense.  The Baire theorem
and intersection with the irrational parameters give the conclusion.
All monomials are approximable at the parameters in this intersection.
\end{proof}

\paragraph{Independent proof check and the remaining obstruction.}
The fibre proof was re-derived using the direct cyclic matrix rather than
the orbit-product recurrence; the factors and row-reversal sign agree.
The continuity bound was checked separately by splitting at the sawtooth
jump and obtaining \(d(1-d)\); a pointwise Lipschitz bound was rejected
because it fails at that jump.  The gap proof was re-checked in the order
fixed target, infinite convergent subsequence, denominator tending to
infinity, and finally polynomial density.  The small cases \(k=0,1\) and
\(q=2\) are covered explicitly.  The two independent exact programs agree
on the determinant signs and every reported integral, including both
numerator orders at denominator three.  No finite-denominator observation
has been promoted into an all-denominator assertion.

The certificate method lacks a quantitative bound on the degree and
coefficient size needed as \(q\) grows.  Even a practical bound would only
give a sufficient approximation class unless the final transfer argument
were strengthened.  It also cannot turn existence of one cyclic vector
into cyclicity of every nonzero vector or settle the invariant-subspace
question.  Determining the full scope of arbitrary-weight classification
criteria and addressing the unrestricted Hilbert invariant-subspace
problem lie outside this bounded development; the ambiguous walk question
does not yet specify a unique probability law.

One further research run should first obtain and check the 2026 determinant
preprint, compare its quantitative bounds with these exact certificates,
and then optimize the rational approximation step with proved coefficient
bounds.  It could supply a certified comparison on additional finite
denominators and possibly a sharper sufficient arithmetic condition.
Neither those computations nor this chapter justify a forecast of a
universal irrational theorem.

\section{Problem index}\label{problems:1964}
Earlier-year occurrence identifiers are retained as text; current-year identifiers link to their entries.

\subsection{GP-1964-01: Growth and displacement of planar self-avoiding walks: a conditional identification}\label{prob:GP-1964-01}
For the uniform probability on nearest-neighbor self-avoiding walks \(\gamma\) of fixed length \(n\) from the origin in \(\mathbb Z^2\), determine the asymptotic order of \(\mathbb E\|\gamma_n\|\) and the square-lattice connective constant \(\mu=\lim c_n^{1/n}\). This is a definite modern problem associated with the ambiguous notebook rule, which might instead mean nonbacktracking or kinetic growth.
\par\noindent\textit{Occurrences:} \hyperref[item:1964-3.11]{1964-3.11}, \hyperref[item:1964-3.12]{1964-3.12}.
\par\noindent\textit{Status (2026-11):} As of November 2026: the uniform-SAW interpretation is sub-ballistic by \citet[Theorem 1.1 and Corollary 1.2]{duminil2013subballistic}, first posted May 2012 and published October 2013. The conjectural square-lattice endpoint exponent \(3/4\) is not a theorem in the reviewed sources. The proposed count base \(e\) is already excluded by the rigorous 2000 bound \(\mu\le2.679192495\) of \citet{ponitz2000bounds}; an exact square-lattice constant is not obtained. Nonbacktracking walks have the different count \(4\cdot3^{n-1}\). The later published bound is \(\mu\le2.662343\), from \citet{couronne2025bound} (preprint first posted 29 November 2022; journal publication 31 October 2025). An author-posted preprint of 21 October 2026, \citet{tomar2026connective}, claims \(\mu\le2.6586513\); its certificate was not independently rerun here. Neither value is an exact determination of the square-lattice constant.

\subsection{GP-1964-02: Invariant solutions for invariant data}\label{prob:GP-1964-02}
Given a group acting on distributions and an equivariant differential operator \(P\), does the solvability of \(Pu=f\) for invariant \(f\) imply solvability by an invariant \(u\)?
\par\noindent\textit{Occurrences:} \hyperref[item:1964-12.5]{1964-12.5}.
\par\noindent\textit{Status (2026-11):} As of November 2026: true for compact groups by averaging; false without further hypotheses for noncompact groups. Translations of \(\mathbb R\), \(P=d/dx\), \(f=1\) give an elementary counterexample. The Lorentz action discussed in \citet{methee1954lorentz} is a structured special case, whose full surjectivity theorem was not read in this run.

\subsection{GP-1964-03: Fundamental solutions of invariant operators on groups}\label{prob:GP-1964-03}
For a specified left-invariant differential operator \(P\) on a Lie group \(G\), determine whether a distribution \(E\) exists with \(PE=\delta_e\). On a compact group, distinguish this equation from its reduced version with the constant or null-space projection removed.
\par\noindent\textit{Occurrences:} \hyperref[item:1964-12.6]{1964-12.6}, \hyperref[item:1964-12.7]{1964-12.7}.
\par\noindent\textit{Status (2026-11):} As of November 2026: a universal positive answer is false. Complex invariant Lewy-type operators can be locally nonsolvable on noncompact groups; compact groups have zero-mode and representation-theoretic obstructions. \citet[Theorem 1]{cerezo1969fundamental} gives a compact-group criterion. For \(\partial_x-\alpha\partial_y\) on \(\mathbb T^2\), the reduced distributional problem is governed by the polynomial small-divisor condition in \citet{greenfield1972global}; the unreduced delta equation fails for every slope.

\subsection{GP-1964-04: Extreme tracial characters of SL2 over the rationals}\label{prob:GP-1964-04}
Classify extreme normalized positive-definite class functions on the discrete group \(\mathrm{SL}_2(\mathbb Q)\). The source initially says only two-by-two rational matrices; the later determinant-one formula motivates, but does not prove, this precise interpretation.
\par\noindent\textit{Occurrences:} \hyperref[item:1964-13.9]{1964-13.9}.
\par\noindent\textit{Status (2026-11):} As of November 2026: solved by \citet[Theorem 2.4]{petersonthom2016characters}, preprint March 2013, online publication March 2014, print July 2016. The extreme characters are the trivial character and the two central characters extended by zero off \(\{\pm I\}\). The result does not claim a classification of the entire unitary dual or an unqualified \(\mathrm{GL}_2(\mathbb Q)\) problem.

\subsection{GP-1964-05: Index of nonlinear Fredholm maps}\label{prob:GP-1964-05}
A \(C^1\) map \(F:M\to N\) of Banach manifolds is Fredholm if every \(DF_x\) is Fredholm. Its derivative index is locally constant, hence constant on connected \(M\). This does not define vector-space kernels or cokernels for \(F\) itself.
\par\noindent\textit{Occurrences:} \hyperref[item:1964-15.3]{1964-15.3}.
\par\noindent\textit{Status (2026-11):} As of November 2026, Fredholm-map theory supplies the precise answer in this framework.

\subsection{GP-1964-06: Stable index of a continuous Fredholm family}\label{prob:GP-1964-06}
For norm-continuous \(T:X\to\operatorname{Fred}(H)\), with compact Hausdorff \(X\), choose finite-dimensional \(V\) with \(T_xH+V=H\). Then \(E_x=T_x^{-1}V\) is a bundle and \([E]-[X\times V]\in K^0(X)\) is independent of \(V\).
\par\noindent\textit{Occurrences:} \hyperref[item:1964-15.4]{1964-15.4}.
\par\noindent\textit{Status (2026-11):} As of November 2026, the index bundle answers the stable question; its rank is the numerical index.

\subsection{GP-1964-07: Invariant subspaces of bounded Hilbert-space operators}\label{prob:GP-1964-07}
Does every bounded linear operator on a separable infinite-dimensional complex Hilbert space \(H\) admit a closed invariant subspace \(M\) with \(0\ne M\ne H\)?
\par\noindent\textit{Occurrences:} \hyperref[item:1964-17.1]{1964-17.1}, \hyperref[item:1964-50.1]{1964-50.1}.
\par\noindent\textit{Status (2026-11):} As of 12 November 2026, no generally accepted solution was verified in the literature checked. \citet[Introduction]{marcoux2026generators}, first posted 21 August 2026, explicitly states this Hilbert-space question as unsolved. The 1964 announcement \citep{debrangesrovnyak1964existence} was retracted in 1965 \citep{debrangesrovnyak1965correction}; these records are bibliographically verified, not treated as valid proofs. Polynomially compact and other special classes in record 50 do not settle the unrestricted question.

\subsection{GP-1964-08: Coordinate generators for the ball-algebra maximal ideal}\label{prob:GP-1964-08}
For \(A(B^n)=\mathcal O(B^n)\cap C(\overline B^n)\) and \(a\in B^n\), the ideal \(\{f:f(a)=0\}\) is generated by \(z_j-a_j\). The source asks \(n=2,a=0\).
\par\noindent\textit{Occurrences:} \hyperref[item:1964-20.1]{1964-20.1}.
\par\noindent\textit{Status (2026-11):} As of November 2026, the exact ball question is solved; later counterexamples concern other domains.

\subsection{GP-1964-09: Does Koopman spectrum determine a finite ergodic group action?}\label{prob:GP-1964-09}
The answer is negative even for transitive finite actions. Let \(G=GL(3,\mathbb F_2)\) act on the seven nonzero vectors of \(\mathbb F_2^3\), and on the seven nonzero covectors by the dual action. The permutation characters coincide, but the \(G\)-sets are not isomorphic. With uniform probability these are ergodic systems with unitarily equivalent Koopman representations.
\par\noindent\textit{Occurrences:} \hyperref[item:1964-34.13]{1964-34.13}.
\par\noindent\textit{Status (2026-11):} As of November 2026, the finite-group question has a negative answer; this is not an open reconstruction problem.

\subsection{GP-1964-10: Primitivity of universal enveloping algebras}\label{prob:GP-1964-10}
For a Lie algebra \(\mathfrak g\), when does \(U(\mathfrak g)\) have a faithful irreducible module? For the finite-dimensional complex reconstruction, this asks whether its zero prime ideal is primitive.
\par\noindent\textit{Occurrences:} \hyperref[item:1964-37.1]{1964-37.1}.
\par\noindent\textit{Status (2026-11):} As of November 2026: established primitive-ideal theory gives the Dixmier--Moeglin criterion in the finite-dimensional complex case. The source omits field and dimension, so the audit does not assert an unrestricted field-independent answer.

\subsection{GP-1964-11: Isomorphism of Weyl algebras of different ranks}\label{prob:GP-1964-11}
For characteristic-zero \(k\), can \(A_m(k)\cong A_n(k)\) as \(k\)-algebras with \(m\ne n\)? The notebook asks \(m=1,n=2\).
\par\noindent\textit{Occurrences:} \hyperref[item:1964-38.1]{1964-38.1}.
\par\noindent\textit{Status (2026-11):} Solved: no. Growth dimension \(2n\) distinguishes the rank; Gelfand--Kirillov, received January 1966, published 1966. Status checked November 2026.

\subsection{GP-1964-12: Trace of a commutator}\label{prob:GP-1964-12}
Under what hypotheses on bounded \(A,B\) and the trace is \(\operatorname{Tr}[A,B]=0\)?
\par\noindent\textit{Occurrences:} \hyperref[item:1964-47.8]{1964-47.8}.
\par\noindent\textit{Status (2026-11):} As of November 2026: both products trace class suffice; trace class of the commutator alone does not, as \([S^*,S]\) shows. Later trace formulas treat the nonvanishing case.

\subsection{GP-1964-13: Invariant subspaces for subnormal operators}\label{prob:GP-1964-13}
Does every bounded subnormal operator on a complex Hilbert space of dimension greater than one admit a nonzero proper closed invariant subspace?
\par\noindent\textit{Occurrences:} \hyperref[item:1964-50.5]{1964-50.5}.
\par\noindent\textit{Status (2026-11):} As of November 2026: solved affirmatively by \citet{brown1978subnormal}, published in September 1978. The publisher's abstract explicitly states the affirmative theorem; the full subscription article was not read in this audit. The notebook's then-open special-class question is distinct from the unrestricted Hilbert-space problem.

\subsection{GP-1964-14: Weight criteria for a unicellular backward shift}\label{prob:GP-1964-14}
Characterize bounded nonzero weight sequences for which the backward shift on \(\ell^2\) has exactly the finite initial-span invariant subspaces and the whole space.
\par\noindent\textit{Occurrences:} \hyperref[item:1964-50.16]{1964-50.16}.
\par\noindent\textit{Status (2026-11):} As of November 2026: developed and partly answered in verified weight classes, including monotone moduli tending to zero. The scope of a universal explicit arbitrary-weight criterion was not ascertained in this audit; it is not declared solved or open by default.

\subsection{GP-1964-15: An invariant-subspace lattice with a bilateral integer chain}\label{prob:GP-1964-15}
Does a bounded Hilbert-space operator have full invariant-subspace lattice \(\{0,H\}\cup\{M_m:m\in\mathbb Z\}\), with strict successive inclusion, intersection zero and dense union?
\par\noindent\textit{Occurrences:} \hyperref[item:1964-50.17]{1964-50.17}.
\par\noindent\textit{Status (2026-11):} As of November 2026: yes. Domar\textquotesingle s Theorem~2 (1981) gives an exact example; \(Se_j=\exp(-2|2j+1|)e_{j+1}\) on \(\ell^2(\mathbb Z)\) has precisely the coordinate-tail nest and the two endpoints.

\subsection{GP-1964-16: Invariant subspaces of every irrational Bishop operator}\label{prob:GP-1964-16}
For every irrational \(\alpha\), must \(T_\alpha f(x)=x f(\{x+\alpha\})\) on \(L^2(0,1)\) have a nonzero proper closed invariant subspace?
\par\noindent\textit{Occurrences:} \hyperref[item:1964-50.20]{1964-50.20}.
\par\noindent\textit{Status (2026-11):} Affirmative for non-Liouville parameters and a larger arithmetic class, including the continued-fraction condition in Chamizo et al., Theorem 3.7. No established all-parameter resolution was verified through November 2026.

\subsection{GP-1964-17: Cyclicity of irrational Bishop operators}\label{prob:GP-1964-17}
Does every irrational Bishop operator admit a cyclic vector, and for which parameters is the constant function cyclic?
\par\noindent\textit{Occurrences:} \hyperref[item:1964-50.21]{1964-50.21}.
\par\noindent\textit{Status (2026-11):} Béhani proves a co-meager cyclic parameter set and sufficient approximation conditions. The chapter develops a bounded rational-orbit certificate and continuity estimates. It does not prove universal irrational cyclicity; an inaccessible 2026 lead is not relied on.

\section{Corrections to the transcription and translation}\label{corr:1964}
The evidence distinguishes a transcription correction from a mathematical emendation of wording faithfully preserved in the notebook.
\subsection*{P000001 (both; faithful error): resolution paragraph}
\noindent\textit{Transcription reads.} A smooth hypersurface replaces the inverse image of the singular divisor\par
\noindent\textit{Critical reading.} Distinguish the smooth strict transform from the normal-crossing total transform.\par
\noindent\textit{Evidence.} Scan preserves the compressed claim; Hironaka characteristic-zero theorem, item 1964-2.2.\par
\subsection*{P000001-P000002 (both; faithful error): universal blowup algebra}
\noindent\textit{Transcription reads.} \(S_A(J)=\operatorname{Sym}_A(J)\) used without qualification\par
\noindent\textit{Critical reading.} For general \(J\) use the Rees algebra \(\bigoplus J^n\); \(\operatorname{Sym}(J)\) works when \(J\) is of linear type.\par
\noindent\textit{Evidence.} Scan and both transcripts; Stacks tags 01OF, 0806; item 1964-2.3.\par
\subsection*{P000002 (both; ambiguity): homogeneous prime selection}
\noindent\textit{Transcription reads.} \ru{Выбр.} ideals containing all positive-degree polynomials\par
\noindent\textit{Critical reading.} Retain the abbreviation; the intended Proj construction excludes such ideals.\par
\noindent\textit{Evidence.} Scan does not distinguish choose/discard expansion; item 1964-2.4.\par
\subsection*{P000002 (both; ambiguity): local fractions}
\noindent\textit{Transcription reads.} Fractions of degree 1: \(J_U\)\par
\noindent\textit{Critical reading.} Degree-zero fractions give the structure sheaf; degree-one fractions give the twist in the appropriate convention.\par
\noindent\textit{Evidence.} The source distinguishes \(O_U\) and \(J_U\) but the surrounding localization wording is compressed; item 1964-2.4.\par
\subsection*{P000004-P000006 (both; faithful error): thermodynamic potential}
\noindent\textit{Transcription reads.} \(f=N^{-1}\log Z\) called convex\par
\noindent\textit{Critical reading.} With the positive-log convention \(f\) is concave; minus the potential is convex.\par
\noindent\textit{Evidence.} Both transcripts and scan; perspective \(f(v)=v\,s(1/v)\), item 1964-3.2. Russian convexity vocabulary can be older/ambiguous, so this is a mathematical convention note.\par
\subsection*{P000006-P000008 (both; faithful error): pressure plateau}
\noindent\textit{Transcription reads.} \(\partial f/\partial v=0\)\par
\noindent\textit{Critical reading.} \(\partial f/\partial v\) is constant on an affine coexistence segment; it need not vanish.\par
\noindent\textit{Evidence.} Visible zero retained in the transcription. Pressure \(p=T\partial_v f\); item 1964-3.3.\par
\subsection*{P000006 (both; faithful error): partition-function asymptotic}
\noindent\textit{Transcription reads.} \(Z(V,N)\) approximately \(\exp(-Nf(v,T))\)\par
\noindent\textit{Critical reading.} With the earlier definition \(f=\lim N^{-1}\log Z\), the exponent is \(+Nf\), up to \(o(N)\).\par
\noindent\textit{Evidence.} Minus signs are visible; this is a change of convention or faithful algebraic error, not an OCR correction. Item 1964-3.2.\par
\subsection*{P000007 (both; faithful error): large-contour scale}
\noindent\textit{Transcription reads.} \(D\ell^{(\nu-1)/\nu}\)\par
\noindent\textit{Critical reading.} \(D\ell^{\nu-1}\), since \(\ell\) is side length; alternatively \(DV^{(\nu-1)/\nu}\).\par
\noindent\textit{Evidence.} The fractional exponent is visible on P7; P10 writes the surface scale \(\ell^{\nu-1}\). Item 1964-3.8.\par
\subsection*{P000007 (both; faithful error): small-contour index range}
\noindent\textit{Transcription reads.} all remaining small contours followed by \(\sum_{i=1}^{M}\)\par
\noindent\textit{Critical reading.} For the remaining contours use \(i>M\), as on P10.\par
\noindent\textit{Evidence.} The visible bounds conflict with the accompanying description; P10 supplies the consistent statement. Item 1964-3.8.\par
\subsection*{P000008 (both; ambiguity): Gibbs exponent}
\noindent\textit{Transcription reads.} \(e^{U/T}\)\par
\noindent\textit{Critical reading.} Reconcile \(U\) with the earlier negative energy before choosing the sign; no silent replacement.\par
\noindent\textit{Evidence.} Earlier P3 uses negative occupied-bond energy with \(e^{-U/T}\); scan P8 changes convention or sign.\par
\subsection*{P000009 (both; ambiguity): \(Q(T)\) series}
\noindent\textit{Transcription reads.} \(S_\nu e(1-e^{-1/(2T)})\) and \(3\exp\{-1/(2T)\}^k\)\par
\noindent\textit{Critical reading.} Retain the uncertain factor and power placement; the source is insufficient for a numerical threshold.\par
\noindent\textit{Evidence.} Scan P9 checked repeatedly. Published Dobrushin equations (1.6)--(1.7) located but not visually recovered confidently; item 1964-3.10.\par
\subsection*{P000009 (both; ambiguity): isoperimetric constant}
\noindent\textit{Transcription reads.} \(S_3=1/(C\sqrt{\pi})\)\par
\noindent\textit{Critical reading.} The geometrically intended constant is \(1/(6\sqrt{\pi})\); C/c/6 remains a glyph ambiguity.\par
\noindent\textit{Evidence.} Volume of a ball of surface area one; scan permits 6, but mathematical identification is not a definitive handwriting reading.\par
\subsection*{P000009 (both; faithful error): subcube deviation probability}
\noindent\textit{Transcription reads.} A closing brace precedes \(>\varepsilon\)\par
\noindent\textit{Critical reading.} Place the comparison with \(\varepsilon\) inside the event braces in the modern statement.\par
\noindent\textit{Evidence.} The extra brace is visible in the scan; item 1964-3.6.\par
\subsection*{P000010 (both; faithful error): final subcube fraction}
\noindent\textit{Transcription reads.} \((G\cap S_\ell)/V_\ell\)\par
\noindent\textit{Critical reading.} For a density in \(S_\ell\) divide its volume by \(|S_\ell|\); distinguish this from a global-volume fraction.\par
\noindent\textit{Evidence.} Scan visibly has \(V_\ell\); the claimed local-density deduction needs the subcube normalization.\par
\subsection*{P000011 (both; ambiguity): walk probability law}
\noindent\textit{Transcription reads.} Turn with probability \(1/3\), but returning is forbidden\par
\noindent\textit{Critical reading.} Keep unresolved whether return means immediate reversal or visiting any earlier site, and how choices are weighted.\par
\noindent\textit{Evidence.} Scan and both languages; uniform fixed-length SAW and kinetic/no-backtracking laws differ. Items 1964-3.11--3.12.\par
\subsection*{P000011 (both; faithful error): walk count conjecture}
\noindent\textit{Transcription reads.} \(2.6^n<\text{count}<2.8^n\); conjecture \(e^n\)\par
\noindent\textit{Critical reading.} Under the square-lattice SAW reading the exponential base is below \(e\); under nonbacktracking the exact count is \(4\) times \(3^{n-1}\).\par
\noindent\textit{Evidence.} Pönitz--Tittmann 2000 bound 2.679192495; model assignment remains conditional. The repeated phrase allowing returns is not silently reversed.\par
\subsection*{P000013 (both; ambiguity): equal-time mixed commutator}
\noindent\textit{Transcription reads.} Unclear second-field mark and abbreviated delta\par
\noindent\textit{Critical reading.} The mathematical CCR is \([\phi(f),\pi(g)]=i\hbar(\int fg)I\), after smearing and domain specification.\par
\noindent\textit{Evidence.} Scan checked; no unique handwritten dot/prime reading asserted. Item 1964-5.4.\par
\subsection*{P000012 (both; faithful error): wave energy norm}
\noindent\textit{Transcription reads.} \(\lVert[\phi,\omega]\rVert=\lVert B\phi\rVert^2+\lVert\omega\rVert^2\)\par
\noindent\textit{Critical reading.} The left side must be the squared norm.\par
\noindent\textit{Evidence.} No square is visible on the left in the scan; both right terms are squares. Item 1964-4.3.\par
\subsection*{P000012 (both; faithful error): wave generator adjoint}
\noindent\textit{Transcription reads.} self-adjoint in this norm\par
\noindent\textit{Critical reading.} Skew-adjoint in the energy inner product, consistent with \(A=-A^*\) on P11.\par
\noindent\textit{Evidence.} The block generator \(\bigl(\begin{smallmatrix}0&I\\-B^2&0\end{smallmatrix}\bigr)\) is skew-adjoint; scan says self-adjoint twice. Item 1964-4.3.\par
\subsection*{P000012 (both; faithful error): variational equation}
\noindent\textit{Transcription reads.} \(v'=Av+K_{u(\cdot)}\)\par
\noindent\textit{Critical reading.} \(v'=Av+DK(u(\cdot))v\)\par
\noindent\textit{Evidence.} The factor \(v\) after the derivative of \(K\) is absent in the scan. Item 1964-4.5.\par
\subsection*{P000014 (both; ambiguity): historical motivation}
\noindent\textit{Transcription reads.} Radon measures exist only on locally compact spaces\par
\noindent\textit{Critical reading.} This describes the initial functional-definition restriction, not the scope of the later compact-regular definition.\par
\noindent\textit{Evidence.} P15 explicitly announces removal of local compactness; item 1964-6.1.\par
\subsection*{P000014 (both; faithful error): pushforward set location}
\noindent\textit{Transcription reads.} arbitrary subset \(A\subset X\) before \(f(\mu)(A)\)\par
\noindent\textit{Critical reading.} The set tested for image-measure measurability must be \(A\subset Y\).\par
\noindent\textit{Evidence.} The source alternates \(A\) and \(B\) and visibly says \(X\); definition of pushforward, item 1964-6.2.\par
\subsection*{P000015 (both; faithful error): counterexample}
\noindent\textit{Transcription reads.} Take a non-Borel subset of \(Y\setminus X\); it is nonmeasurable\par
\noindent\textit{Critical reading.} Choose a non-Lebesgue-measurable subset; non-Borel alone is insufficient.\par
\noindent\textit{Evidence.} Visible parenthesis reads non-Borel. The complement \(Y\setminus X\) itself works under the preceding construction; item 1964-6.3.\par
\subsection*{P000015 (both; ambiguity): compactification}
\noindent\textit{Transcription reads.} maximal ideals of the algebra of continuous functions\par
\noindent\textit{Critical reading.} For the Stone--Cech compactification use bounded continuous functions; the handwritten abbreviation is not a secure extra boundedness word.\par
\noindent\textit{Evidence.} Scan reviewed; Courrege1963 explicitly uses bounded functions; item 1964-6.4.\par
\subsection*{P000015 (both; faithful error): lower regularity display}
\noindent\textit{Transcription reads.} for every \(B\) compact there exists a compact \(k\)\par
\noindent\textit{Critical reading.} For every Borel \(B\) and \(\varepsilon>0\) require compact \(K\subset B\) with \(\mu(B\setminus K)<\varepsilon\).\par
\noindent\textit{Evidence.} Scan shows compact after \(B\); as written the condition is tautological by \(K=B\). Item 1964-6.4.\par
\subsection*{P000015 (both; ambiguity): concentration on X}
\noindent\textit{Transcription reads.} \(\mu(\widehat X\setminus X)=0\)\par
\noindent\textit{Critical reading.} Use completed measurability or outer measure when \(X\) need not be Borel in the compactification.\par
\noindent\textit{Evidence.} No Borel hypothesis on \(X\) is stated; item 1964-6.4.\par
\subsection*{P000015 (both; faithful error): last sentence}
\noindent\textit{Transcription reads.} one cannot define the support of arbitrary Radon measures\par
\noindent\textit{Critical reading.} For finite Hausdorff Radon measures support is well defined and has full measure, without local compactness.\par
\noindent\textit{Evidence.} Scan plainly retains the sentence. Finite-subcover argument plus inner regularity in item 1964-6.9.\par
\subsection*{P000016 (both; ambiguity): operator-space claim}
\noindent\textit{Transcription reads.} \(\operatorname{Hom}(E,F)\) is of the same kind for any topology between pointwise and compact convergence\par
\noindent\textit{Critical reading.} Specify the locally convex/Suslin hypotheses and topology; do not assert arbitrary \(\operatorname{Hom}(E,F)\) is itself Frechet.\par
\noindent\textit{Evidence.} Broad wording confirmed in both languages; exact general theorem not read. Item 1964-6.6.\par
\subsection*{P000016-P000017 (both; faithful error): continuous-path event}
\noindent\textit{Transcription reads.} Probability of continuity always exists\par
\noindent\textit{Critical reading.} A completed Radon law on the function space suffices; a bare cylinder sigma algebra need not contain the event.\par
\noindent\textit{Evidence.} Scan and measure-space distinction; item 1964-6.7.\par
\subsection*{P000017 (both; faithful error): common modulus}
\noindent\textit{Transcription reads.} With probability one all functions have a common modulus of continuity\par
\noindent\textit{Critical reading.} Use countably many compact layers, or a deterministic modulus multiplied by an almost surely finite random constant.\par
\noindent\textit{Evidence.} Counterexample \(X_t=Zt\) with unbounded \(Z\) excludes a coefficient-one common modulus; item 1964-6.8.\par
\subsection*{P000017 (both; ambiguity): Weil group exception}
\noindent\textit{Transcription reads.} no matrices of order 2 among real components\par
\noindent\textit{Critical reading.} A safe formulation excludes simple factors locally isomorphic to \(\mathrm{SL}_2(\mathbb R)\); the shorthand is not a general matrix-size criterion.\par
\noindent\textit{Evidence.} Weil1964 Section 3; item 1964-8.1.\par
\subsection*{P000017 (both; ambiguity): nondiscrete aside}
\noindent\textit{Transcription reads.} If \(\Gamma\) is not discrete, Weil theorem holds although \(H^1\) is nonzero\par
\noindent\textit{Critical reading.} Retain as unidentified; topology and cohomology convention are unspecified.\par
\noindent\textit{Evidence.} Scan reads not discrete, so it cannot be silently changed to noncompact. Item 1964-8.3.\par
\subsection*{P000018 (both; faithful error): Hilbert inclusion topology}
\noindent\textit{Transcription reads.} The Hilbert topology does not coincide with the subspace topology\par
\noindent\textit{Critical reading.} It need not coincide; it can coincide, for example on the finite-dimensional constant subspace.\par
\noindent\textit{Evidence.} Scan and immediately following constant example; item 1964-7.1.\par
\subsection*{P000018 (both; faithful error): sum-norm display}
\noindent\textit{Transcription reads.} \(\lVert h\rVert=\inf(\lVert h_1\rVert^2+\lVert h_2\rVert^2)\)\par
\noindent\textit{Critical reading.} \(\lVert h\rVert^2=\inf(\lVert h_1\rVert^2+\lVert h_2\rVert^2)\)\par
\noindent\textit{Evidence.} No left square is visible; quotient-Hilbert homogeneity and Schwartz1962 Section 1; item 1964-7.2.\par
\subsection*{P000019-P000020 (both; ambiguity): strict cone wording}
\noindent\textit{Transcription reads.} strictly convex cone\par
\noindent\textit{Critical reading.} In the kernel correspondence it is pointed and cancellative, not strictly convex in the modern geometric sense.\par
\noindent\textit{Evidence.} Scan P20 says strictly convex. Positive semidefinite matrices have nontrivial faces; item 1964-7.3.\par
\subsection*{P000019 (both; faithful error): scalar order equivalence}
\noindent\textit{Transcription reads.} \(\lambda H\ge H\) iff \(\lambda\ge1\)\par
\noindent\textit{Critical reading.} Add \(H\) nonzero; the zero Hilbert subspace satisfies the order for every nonnegative scalar.\par
\noindent\textit{Evidence.} Visible definition includes the zero subspace; item 1964-7.3.\par
\subsection*{P000022 (both; faithful error): Sobolev kernel}
\noindent\textit{Transcription reads.} \(H^1\) corresponds to \(-\Delta+\lambda\)\par
\noindent\textit{Critical reading.} Its positive reproducing kernel as a subspace of \(D'\) is \((-\Delta+\lambda)^{-1}\), for \(\lambda>0\).\par
\noindent\textit{Evidence.} Scan and transcripts agree on missing inverse. Fourier multiplier/energy-duality check; item 1964-7.6.\par
\subsection*{P000023 (both; ambiguity): partition refinement entropy}
\noindent\textit{Transcription reads.} \(\lim H(\text{join})/n=h(\xi)\)\par
\noindent\textit{Critical reading.} Distinguish one-step \(H(\xi)\) from the entropy rate \(h(T,\xi)\).\par
\noindent\textit{Evidence.} Scan reuses \(h\); modern definition and inverse-image refinements in item 1964-10.1.\par
\subsection*{P000023-P000024 (both; provenance): sharp entropy coefficient}
\noindent\textit{Transcription reads.} The bound is attained; \(n\) may be replaced by \(n-1\) for invertible maps\par
\noindent\textit{Critical reading.} The 1965 published theorem is \(n\log_2\lambda\); the \(n-1\) bound is justified here by a later Lyapunov/Ruelle route.\par
\noindent\textit{Evidence.} Kushnirenko1965 Theorem 1 and final examples read. Item 1964-10.3 separates publication from reconstruction.\par
\subsection*{P000024 (both; ambiguity): canonical free energy}
\noindent\textit{Transcription reads.} minus \(\log Z\) per particle called free energy, then differentiated\par
\noindent\textit{Critical reading.} Physical \(F_N=-T\log Z_N\); the displayed minus-log quantity is reduced free energy \(F_N/(NT)\).\par
\noindent\textit{Evidence.} Scan clearly lacks the temperature prefactor; item 1964-9.2.\par
\subsection*{P000025 (both; faithful error): entropy convexity}
\noindent\textit{Transcription reads.} \(S\) is convex in two variables\par
\noindent\textit{Critical reading.} With positive logarithmic entropy and the usual stable short-range limit, \(S\) is concave in extensive densities.\par
\noindent\textit{Evidence.} Same positive-log convention issue as record 3; item 1964-9.3.\par
\subsection*{P000025 (both; provenance): Wigner construction}
\noindent\textit{Transcription reads.} A phase-space expression with position and momentum marginals is called for\par
\noindent\textit{Critical reading.} Do not infer positivity or uniqueness from the two marginals; Wigner formula is supplied from literature/editorial reconstruction.\par
\noindent\textit{Evidence.} The scan does not actually display the bilinear Fourier formula; item 1964-9.4.\par
\subsection*{P000027 (both; ambiguity): Lorentz-invariant reduction}
\noindent\textit{Transcription reads.} Every invariant function is \(T\) of the quadratic form\par
\noindent\textit{Critical reading.} Specify group component and separate time-oriented orbits when necessary; distributions may include cone and origin terms.\par
\noindent\textit{Evidence.} One-variable shorthand is visible, but orbit and distribution qualifications are absent; item 1964-12.3.\par
\subsection*{P000028 (both; faithful error): finite-dimensional fixed point}
\noindent\textit{Transcription reads.} A finite-dimensional representation of a semisimple Lie group has a fixed point\par
\noindent\textit{Critical reading.} The useful extension argument concerns an affine action and vanishing of the finite-dimensional cocycle class; a linear action always fixes zero.\par
\noindent\textit{Evidence.} Scan says representation and fixed point. Item 1964-12.4 restores the affine-cocycle content.\par
\subsection*{P000028 (both; ambiguity): torus elementary solution}
\noindent\textit{Transcription reads.} No solution; irrational \(\alpha\) well approximated\par
\noindent\textit{Critical reading.} \(PE=\delta\) always fails for a derivative on a compact torus because of the constant mode. The reduced equation additionally has small-divisor obstructions.\par
\noindent\textit{Evidence.} Scan shows \(\alpha\) with no independent power; Greenfield--Wallach1972 theorem/proposition; item 1964-12.7.\par
\subsection*{P000029 (both; faithful error): one-dimensional finite affine characters}
\noindent\textit{Transcription reads.} \(e^{2\pi i a/p}\)\par
\noindent\textit{Critical reading.} Use multiplicative characters: \(\chi_j(u^r)=\exp(2\pi i j r/(p-1))\).\par
\noindent\textit{Evidence.} The visible expression uses the additive denominator \(p\); item 1964-13.2.\par
\subsection*{P000030 (both; faithful error): translation eigenvector transport}
\noindent\textit{Transcription reads.} \(S_aT_bS_a^{-1}=T_{ab}\) and \(S_ae_k=e_{ak}\)\par
\noindent\textit{Critical reading.} These conventions give \(S_ae_k\) proportional to \(e_{a^{-1}k}\); alternatively reverse the conjugation convention consistently.\par
\noindent\textit{Evidence.} Both displayed formulas visible. P31 switches the conjugation convention; item 1964-13.2.\par
\subsection*{P000031 (both; faithful error): permutation realization}
\noindent\textit{Transcription reads.} \(T_gf(x)=f(gx)\)\par
\noindent\textit{Critical reading.} For a left representation take \(T_gf(x)=f(g^{-1}x)\), or use a right action, with the antirepresentation convention stated consistently.\par
\noindent\textit{Evidence.} Composition reverses the group order; trace values unchanged. Item 1964-13.3.\par
\subsection*{P000032 (both; faithful error): first row of normalized affine table}
\noindent\textit{Transcription reads.} \(a\ne0\)\par
\noindent\textit{Critical reading.} \(a\ne1\)\par
\noindent\textit{Evidence.} Scan distinctly has zero; first row must exclude the separate translation rows \(a=1\). Item 1964-13.3.\par
\subsection*{P000032 (both; ambiguity): normal-subgroup lemma}
\noindent\textit{Transcription reads.} Extreme points coincide with those for the original group\par
\noindent\textit{Critical reading.} The invariant convex set is unspecified; unrestricted restriction need not preserve extremality.\par
\noindent\textit{Evidence.} \(S_3\) standard character restricts to the mean of two \(A_3\) characters. Scan remains abbreviated; item 1964-13.6.\par
\subsection*{P000032-P000033 (both; faithful error): infinite symmetric-group characters}
\noindent\textit{Transcription reads.} arbitrary real \(a_k\) in product \(\prod a_k^{n_k}\)\par
\noindent\textit{Critical reading.} For extreme characters \(a_k=\sum\alpha_i^k+(-1)^{k-1}\sum\beta_i^k\), with nonnegative decreasing parameters and sum \(\le1\).\par
\noindent\textit{Evidence.} Thoma formula verified in Gohm--Koestler2010 Theorem 0.1. No parameter restrictions are written in the scan; item 1964-13.7.\par
\subsection*{P000033 (both; faithful error): type-I criterion}
\noindent\textit{Transcription reads.} \(G_0\) isomorphic to \(G\) and \(G/G_0\) finite\par
\noindent\textit{Critical reading.} \(G\) must have an abelian subgroup of finite index; it may be chosen normal.\par
\noindent\textit{Evidence.} The scan reads izom. \(G\), not abelian. This is a mathematical emendation, not an OCR correction. Item 1964-13.8.\par
\subsection*{P000033 (both; ambiguity): rational matrix problem}
\noindent\textit{Transcription reads.} characters of two-by-two rational matrices\par
\noindent\textit{Critical reading.} Keep determinant/topology unspecified in the source; the \(\mathrm{SL}_2(\mathbb Q)\) solution is an explicitly conditional identification.\par
\noindent\textit{Evidence.} Later determinant-one matrix supports the interpretation, but not an unqualified \(\mathrm{GL}_2\) claim; item 1964-13.9.\par
\subsection*{P000033 (both; faithful error): conjugacy by trace}
\noindent\textit{Transcription reads.} conjugacy class determined by \(\alpha+\delta\)\par
\noindent\textit{Critical reading.} Valid for regular semisimple \(\mathrm{SL}_2\) over an algebraically closed field; false universally.\par
\noindent\textit{Evidence.} Identity and a nontrivial unipotent both have trace two; rational conjugacy can split further. Item 1964-13.15.\par
\subsection*{P000034 (both; ambiguity): real discrete-series brace and author}
\noindent\textit{Transcription reads.} undefined minimum eigenvalue; missing case condition; uncertain 1947 surname\par
\noindent\textit{Critical reading.} Retain the unresolved formula and surname. A rotation-weight boundary formula may be given only as an explicit modern reconstruction.\par
\noindent\textit{Evidence.} Scan reviewed; no silent substitution of Bargmann or complete case table. Item 1964-13.12.\par
\subsection*{P000034 (en; translation): p-adic eigenvalue cases}
\noindent\textit{Transcription reads.} real / nonreal eigenvalue\par
\noindent\textit{Critical reading.} Read the mathematical case distinction as eigenvalue in the base local field / in a quadratic extension.\par
\noindent\textit{Evidence.} Gelfand--Graev1963 Section 5.1 formula (5); item 1964-13.13.\par
\subsection*{P000034 (both; provenance): finite-field attribution}
\noindent\textit{Transcription reads.} Frobenius 1895\par
\noindent\textit{Critical reading.} Retain as the notebook attribution, not a verified exact citation for the complete \(\mathrm{SL}_2\) finite-field classification.\par
\noindent\textit{Evidence.} Prasad2007 introduction instead distinguishes the Jordan/Schur1907 classifications; item 1964-13.14.\par
\subsection*{P000035 (both; faithful error): adele definition}
\noindent\textit{Transcription reads.} direct sum of all p-adic groups containing diagonal rationals\par
\noindent\textit{Critical reading.} Use the restricted product; for full adeles include the real place.\par
\noindent\textit{Evidence.} Nonzero diagonal rationals do not belong to the algebraic direct sum. Gelfand--Graev--Pyatetskii-Shapiro1964 Section 1; item 1964-13.16.\par
\subsection*{P000036 (both; provenance): speaker initials}
\noindent\textit{Transcription reads.} I. I. / I. M. Pyatetsky-Shapiro\par
\noindent\textit{Critical reading.} I. I. Pyatetskii-Shapiro as identification; preserve the merged glyph uncertainty.\par
\noindent\textit{Evidence.} The same line explicitly expands the name as Ilya Iosifovich; talk 14.\par
\subsection*{P000036 (en; translation): sphere dimension}
\noindent\textit{Transcription reads.} surface of the 3-dimensional sphere\par
\noindent\textit{Critical reading.} The two-dimensional sphere in three-dimensional Euclidean space.\par
\noindent\textit{Evidence.} Equation has three coordinates \(x,y,z\). This removes the English intrinsic-dimension ambiguity.\par
\subsection*{P000036 (both; faithful error): splitting-prime condition}
\noindent\textit{Transcription reads.} \(\xi^2+m=0\bmod q\) is solvable\par
\noindent\textit{Critical reading.} For the cited theorem require fixed \(q>3\) and Legendre symbol \(\bigl(\frac{-m}{q}\bigr)=1\), together with admissible \(m\) modulo 8.\par
\noindent\textit{Evidence.} Linnik1957 Theorem 1, page 258, read. Zero-root divisibility alone is insufficient; item 1964-14.2.\par
\subsection*{P000036 (both; ambiguity): arithmetic dynamics degree}
\noindent\textit{Transcription reads.} each point has its own 48 orthogonal matrices\par
\noindent\textit{Critical reading.} Retain 48 as unresolved shorthand; do not identify it with a universal \(q\)-dependent transition degree.\par
\noindent\textit{Evidence.} Linnik1957 Section 3 uses a reduced two-direction flow. The order 48 of integral orthogonal symmetries is not proof of intended meaning; item 1964-14.3.\par
\subsection*{P000018-P000022 (both; provenance): catalogue span for record 7}
\noindent\textit{Transcription reads.} Record 7 assigned P000016-P000018\par
\noindent\textit{Critical reading.} Keep catalogue span in talk metadata but locate actual Hilbert-kernel items on P000018-P000022.\par
\noindent\textit{Evidence.} Source headings and consecutive mathematics are explicit.\par
\subsection*{P000024-P000026 (both; provenance): catalogue overlap for records 9--12}
\noindent\textit{Transcription reads.} Talk page spans overlap several distinct speakers\par
\noindent\textit{Critical reading.} Preserve catalogue records while assigning actual content pages to each item.\par
\noindent\textit{Evidence.} P24--25 Povzner, P25 Kirillov notice, P25--28 Schwartz; no duplicated mathematical content.\par
\subsection*{P000029-P000036 (both; provenance): catalogue overlap for records 13--14}
\noindent\textit{Transcription reads.} Record 13 spans P000028-P000036; record14 spans P000034-P000036\par
\noindent\textit{Critical reading.} Kirillov mathematics is P29--top P35; the Linnik report is P36.\par
\noindent\textit{Evidence.} Headings and scan boundaries checked. P35 separately announces Linnik for 16 November; do not redatetime the catalogue report solely from that announcement.\par
\subsection*{P000035 (both; faithful error): Fredholm definition}
\noindent\textit{Transcription reads.} \(\operatorname{coker}K\)\par
\noindent\textit{Critical reading.} \(\operatorname{coker}T\)\par
\noindent\textit{Evidence.} The scan introduces no second operator \(K\); this is a notebook notation slip.\par
\subsection*{P000037 (both; provenance): Povzner heading}
\noindent\textit{Transcription reads.} Povzner. Naum Yakovlevich\par
\noindent\textit{Critical reading.} Retain attribution conflict with catalogue A. Ya. Povzner\par
\noindent\textit{Evidence.} The scan contains the unexpected given name; no supported substitution is available.\par
\subsection*{P000042 (both; faithful error): Schrödinger and action formulas}
\noindent\textit{Transcription reads.} \((i/h)\psi_t=\Delta\psi-U\psi/h^2\); \(\int p\,dx\)\par
\noindent\textit{Critical reading.} Consistent modern normalization \(ih\psi_t=-h^2\Delta\psi/2+U\psi\); action integral \(\int(p\,dx-H\,dt)\)\par
\noindent\textit{Evidence.} The coefficients and shorthand are faithfully visible in the scan; the emendation is mathematical.\par
\subsection*{P000045 (both; faithful error): free Dirichlet spectrum}
\noindent\textit{Transcription reads.} \(\lambda=-n^2\)\par
\noindent\textit{Critical reading.} \(\lambda_n=n^2\) for the displayed \(-y''\) operator\par
\noindent\textit{Evidence.} The scan faithfully has both the operator sign and negative eigenvalue.\par
\subsection*{P000045 (both; ambiguity): Morse and zero indexing}
\noindent\textit{Transcription reads.} first eigenfunction has Morse index 1; \(n\) zeros\par
\noindent\textit{Critical reading.} With indexing from 1, index of \(L-\lambda_j\) is \(j-1\) and the first \(m\) span has at most \(m-1\) interior zeros\par
\noindent\textit{Evidence.} The notebook mixes zero-based and one-based labels; modern conventions stated explicitly.\par
\subsection*{P000047 (both; faithful error): rank-one quotient}
\noindent\textit{Transcription reads.} infinite index for all \(N\); fundamental group of a surface\par
\noindent\textit{Critical reading.} Triangle quotient \(\Delta(2,3,N)\); finite for \(N=2,3,4,5\) and generally an orbifold group\par
\noindent\textit{Evidence.} The scan gives the unqualified assertion. The modular presentation exhibits the missing small-level and torsion qualifications.\par
\subsection*{P000048 (both; faithful error): block Poincaré display}
\noindent\textit{Transcription reads.} upper-right \(H\) Hermitian independently of \(A\)\par
\noindent\textit{Critical reading.} upper-right \(X(A^*)^{-1}\), \(X\) Hermitian, or use pairs \((A,X)\)\par
\noindent\textit{Evidence.} The displayed set in the scan is not closed under multiplication; the corrected semidirect representation is checked algebraically.\par
\subsection*{P000048-P000050 (both; ambiguity): prime and primary inclusions}
\noindent\textit{Transcription reads.} \(I\) primary, \(I\) contains \(J\) prime, \(J\) contains \(I^2\)\par
\noindent\textit{Critical reading.} Use \(P\) prime and \(Q\) \(P\)-primary with \(P^2\) contained in \(Q\) contained in \(P\)\par
\noindent\textit{Evidence.} P000050 explicitly says \(I\) and \(J\) were interchanged; scans confirm the inconsistency.\par
\subsection*{P000050 (both; ambiguity): differential-operator formula}
\noindent\textit{Transcription reads.} all operators \(\sum P_i\partial_i\)\par
\noindent\textit{Critical reading.} all derivations \(\sum P_i\partial_i\); general differential operators include higher derivatives\par
\noindent\textit{Evidence.} The surrounding Leibniz rule and localization distinguish derivations from general operators.\par
\subsection*{P000052 (both; faithful error): symbol base}
\noindent\textit{Transcription reads.} \(T(X)\), tangent bundle\par
\noindent\textit{Critical reading.} \(T^*(X)\), cotangent bundle, unless a metric identification is chosen\par
\noindent\textit{Evidence.} The scan says tangent; symbol variables are covectors. Atiyah--Singer 1963 p423 confirms the standard base.\par
\subsection*{P000057 (both; faithful error): periodicity sentence}
\noindent\textit{Transcription reads.} \(m\) increases by factor 6, fields increase by 8\par
\noindent\textit{Critical reading.} \(m\) increases by factor 16, fields increase by 8\par
\noindent\textit{Evidence.} Scan shows 6; the displayed \(m=2^{c+4d}(2a+1)\) and \(d+1\) imply 16.\par
\subsection*{P000058 (both; ambiguity): K-theory of sphere}
\noindent\textit{Transcription reads.} \(K_R(S^i)=\pi_i(BO)\)\par
\noindent\textit{Critical reading.} Reduced \(KO\) of \(S^i\) equals \(\pi_i(BO)\) for \(i>0\)\par
\noindent\textit{Evidence.} Rank summand omitted in the scan; the surrounding stable-bundle discussion supports reduced interpretation.\par
\subsection*{P000059 (both; faithful error): cyclic group order}
\noindent\textit{Transcription reads.} cyclic group of order \(f(k)\)\par
\noindent\textit{Critical reading.} cyclic group of order \(2^{f(k)}\)\par
\noindent\textit{Evidence.} The scan states order \(f(k)\) but later says add the generator \(2^{f(k)}\) times; the latter matches Adams.\par
\subsection*{P000060 (both; faithful error): dense-range criterion}
\noindent\textit{Transcription reads.} the image of \(A-\lambda\) is dense for all complex \(\lambda\)\par
\noindent\textit{Critical reading.} density for \(A-i\) and \(A+i\) (equivalently the appropriate nonreal parameters)\par
\noindent\textit{Evidence.} The scan includes all complex parameters; real eigenvalues of a self-adjoint operator contradict that wording.\par
\subsection*{P000061 (both; faithful error): resolvent formula for evolution}
\noindent\textit{Transcription reads.} \(\int \exp(it\lambda)/(iA-\lambda)\,d\lambda\)\par
\noindent\textit{Critical reading.} spectral integral \(\int\exp(it\lambda)\,dE_A(\lambda)\), or a specified normalized resolvent-jump formula\par
\noindent\textit{Evidence.} The scan omits contour/boundary prescription and normalization; spectral calculus gives the precise replacement.\par
\subsection*{P000062 (both; ambiguity): Lebesgue spectrum convention}
\noindent\textit{Transcription reads.} densities for all \(f\in L^2\)\par
\noindent\textit{Critical reading.} use \(L^2_0\) for a probability-preserving system and distinguish absolute continuity from equivalent Lebesgue maximal type\par
\noindent\textit{Evidence.} The constant function has an atom at the trivial character; the scan does not exclude it.\par
\subsection*{P000063 (both; ambiguity): tensor spectral measure}
\noindent\textit{Transcription reads.} product measure on the spectral line\par
\noindent\textit{Critical reading.} pushforward of the product measure under addition, or multiplication on the circle\par
\noindent\textit{Evidence.} The preceding addition of frequencies in the scan determines convolution.\par
\subsection*{P000063 (both; ambiguity): binary-group notation}
\noindent\textit{Transcription reads.} product of \(\mathbb Z/2\) with finitely many nonzero coordinates\par
\noindent\textit{Critical reading.} direct sum of \(\mathbb Z/2\) for the stated finite-support group\par
\noindent\textit{Evidence.} The finite-support clause is legible in the scan and determines a direct sum.\par
\subsection*{P000064 (both; ambiguity): transitive scalar cocycles}
\noindent\textit{Transcription reads.} representations of the stabilizer\par
\noindent\textit{Critical reading.} characters for scalar unit-modulus cocycles; unitary representations for vector-valued cocycles\par
\noindent\textit{Evidence.} The preceding operator formula is scalar; Mackey Section5 specifies homomorphisms to the chosen coefficient group.\par
\subsection*{P000065 (both; faithful error): icosahedral group}
\noindent\textit{Transcription reads.} unimodular group over five elements\par
\noindent\textit{Critical reading.} projective special linear group \(\mathrm{PSL}(2,\mathbb F_5)\)\par
\noindent\textit{Evidence.} The scan omits the quotient; \(\mathrm{SL}(2,\mathbb F_5)\) has order120 whereas the rotational group has order60.\par
\subsection*{P000065 (both; provenance): record34 catalogue envelope}
\noindent\textit{Transcription reads.} record34 pages P000062-P000064\par
\noindent\textit{Critical reading.} catalogue envelope retained; actual item passages continue on P000065\par
\noindent\textit{Evidence.} The scan continues the Kirillov discussion through cocycles, the finite-group question and the algebra of observables.\par
\subsection*{P000065 (both; faithful error): Baxter paragraph, final two lines}
\noindent\textit{Transcription reads.} If \(A\) is simple, there are no Lie ideals either.\par
\noindent\textit{Critical reading.} False literally: for \(A=M_2(\mathbb C)\), \(\mathbb C I\) and \(\mathfrak{sl}_2(\mathbb C)\) are proper nonzero Lie ideals. The intended quotient/involution theorem is not specified.\par
\noindent\textit{Evidence.} Russian and English agree; scan confirms the two unrestricted sentences. Direct matrix examples refute the literal assertion. Item 1964-37.3.\par
\subsection*{P000066 (both; ambiguity): Kostrikin/Burnside paragraph}
\noindent\textit{Transcription reads.} Uniform \(e^n=0\); the associative algebra is nilpotent, glossed as every element nilpotent.\par
\noindent\textit{Critical reading.} Distinguish nil from nilpotent; retain uncertainty about the field, the quantified Lie subalgebra, and the intended conclusion.\par
\noindent\textit{Evidence.} Russian abbreviation and explanatory clause do not give the same condition. The scan supports the uniform-\(n\) correction but does not supply characteristic or a theorem statement. Item 1964-37.4.\par
\subsection*{P000066 (both; provenance): Dixmier question, two differential-operator rings}
\noindent\textit{Transcription reads.} \(P(x,d/dx)\) and \(P(x,d/dx,y,d/dy)\); are they isomorphic?\par
\noindent\textit{Critical reading.} Interpret as \(A_1\) versus \(A_2\), not the fixed-rank Dixmier endomorphism conjecture.\par
\noindent\textit{Evidence.} Scan shows one canonical pair and then two; Gelfand--Kirillov 1966 Section 4 distinguishes ranks by growth. The error is an index/provenance risk, not a false transcription.\par
\subsection*{P000066-P000068 (both; faithful error): Log-gas energy and Boltzmann factor}
\noindent\textit{Transcription reads.} Positive sum of pair logarithms called \(W\), together with \(\exp(-\beta W)\).\par
\noindent\textit{Critical reading.} With \(\exp(-\beta W)\), set \[W=-\sum_{i<j}\log|\exp(i\theta_i)-\exp(i\theta_j)|.\]\par
\noindent\textit{Evidence.} The eigenangle density requires a positive power of the Vandermonde; the later integral uses \(\exp(+\beta\sum\log)\). Item 1964-39.3.\par
\subsection*{P000067 (both; faithful error): Claimed exact \(\beta=2\) gap expression}
\noindent\textit{Transcription reads.} \(R_2(\alpha)=\exp(N^2\log\cos(\alpha/2))\).\par
\noindent\textit{Critical reading.} Not an exact finite-CUE probability; \(N=1\) gives \(1-\alpha/\pi\). Treat only an appropriately qualified leading-energy approximation.\par
\noindent\textit{Evidence.} Scan and both transcriptions give equality. Elementary \(N=1\) check and the distinction from later determinant asymptotics. Items 1964-39.5 and 42.1.\par
\subsection*{P000068 (both; faithful error): First free-energy relation}
\noindent\textit{Transcription reads.} \(R_\beta(\alpha)=\exp(+\beta(F(\alpha)-F(0)))\) while \(F\) is minimized.\par
\noindent\textit{Critical reading.} For \(F=-\beta^{-1}\log\Psi\), use \(\exp(-\beta(F(\alpha)-F(0)))\); otherwise explicitly define the opposite sign convention for \(F\).\par
\noindent\textit{Evidence.} Scan confirms positive exponent. The normalization \(\Psi(\alpha)/\Psi(0)\) and usual minimized-free-energy convention fix the sign. Item 1964-39.4.\par
\subsection*{P000068 (both; faithful error): Middle \(W\) formula}
\noindent\textit{Transcription reads.} \(W=\sum\log|\exp(i\theta_i-i\theta_j)|\).\par
\noindent\textit{Critical reading.} The interaction depends on \(|\exp(i\theta_i)-\exp(i\theta_j)|\), with its energy sign separately specified.\par
\noindent\textit{Evidence.} Scan shows a single exponential containing the phase difference; its modulus is identically one. The next multiple integral restores the difference of exponentials. Item 1964-39.3.\par
\subsection*{P000069 (both; ambiguity): Notation \(Z_p\) in finite-field counting discussion}
\noindent\textit{Transcription reads.} \(Z_p\) and its algebraic closure.\par
\noindent\textit{Critical reading.} Here \(Z_p\) means the finite field \(\mathbb F_p\), not the \(p\)-adic integers.\par
\noindent\textit{Evidence.} Russian finite-equation and Frobenius context identifies the notation. Clarification of historical notation, not replacement of a wrong glyph. Item 1964-40.1.\par
\subsection*{P000073 (both; faithful error): Even-kernel asymptotic}
\noindent\textit{Transcription reads.} \(R_1(t)\sim At^{17/8}\exp(-t^2/4-t/2)\).\par
\noindent\textit{Critical reading.} For the identified standard even sine kernel, the power is \(-1/8\); retain the manuscript \(17/8\) in the notes field.\par
\noindent\textit{Evidence.} Scan and Russian EvidenceBlock both confirm \(17/8\); Ehrhardt Introduction gives \(-(1/8)\log(s)\). The parameter normalization remains separately uncertain. Item 1964-42.3.\par
\subsection*{P000073 (both; ambiguity): Relation between \(m,n\) and gap parameter}
\noindent\textit{Transcription reads.} \(2m\ge n\) or \(2m=n\); \(t=\alpha n\).\par
\noindent\textit{Critical reading.} Keep the relation ambiguous; the finite cosine matrix naturally scales with \(m\alpha\) and cannot automatically be identified with the standard \(s\).\par
\noindent\textit{Evidence.} The handwritten relation is thick and not securely an equality; \(t=\alpha n\) is underlined twice. No invented rescaling. Item 1964-42.3.\par
\subsection*{P000073 (both; ambiguity): Sphere bundle over the Grassmannian}
\noindent\textit{Transcription reads.} \(M_{n,k}\to G_{n,k}\) with fibre written \(S_k\).\par
\noindent\textit{Critical reading.} The unit sphere in a \(k\)-dimensional real plane is \(S^{k-1}\); retain \(S_k\) only as historical notation.\par
\noindent\textit{Evidence.} The definition says a \(k\)-plane and one point of its unit sphere; the subscript is not a claim of dimension \(k\). Item 1964-45.2.\par
\subsection*{P000074-P000075 (both; ambiguity): Half-line characteristic function and resolvent}
\noindent\textit{Transcription reads.} Whole function/trace terminology carried over from the finite interval.\par
\noindent\textit{Critical reading.} For the half-line use the Jost function on the specified \(\sqrt z\) sheet and a relative resolvent trace, not the ordinary trace of the individual resolvent.\par
\noindent\textit{Evidence.} Russian text distinguishes subtraction of the free term; Buslaev--Faddeev 1960 Theorem 1 gives the exact relative identity. Items 1964-47.1--47.3.\par
\subsection*{P000076 (both; ambiguity): Real-axis trace integral}
\noindent\textit{Transcription reads.} Integral of \(|\lambda|^sW'/W\) equals sum of eigenvalue powers.\par
\noindent\textit{Critical reading.} Read the manuscript parenthesis “integral is a jump”: retain both boundary values, contour orientation, normalization and regularizing subtractions.\par
\noindent\textit{Evidence.} The handwritten parenthesis and Russian text explicitly say jump. Buslaev--Faddeev Section 2 implements the contour/Mellin argument. Item 1964-47.5.\par
\subsection*{P000076 (both; ambiguity): Neumann trace boundary expression}
\noindent\textit{Transcription reads.} \((1/4)(y'(0)-y'(\pi))\).\par
\noindent\textit{Critical reading.} Retain \(y'\) and mark it unidentified; do not silently replace it by \(q(0)+q(\pi)\).\par
\noindent\textit{Evidence.} Scan confirms \(y\) with prime in both terms. No definition of this \(y\) is supplied. Item 1964-47.6.\par
\subsection*{P000076 (both; faithful error): Formal diagonal of the delta kernel}
\noindent\textit{Transcription reads.} \(\delta(0)=0\).\par
\noindent\textit{Critical reading.} \(\delta(0)\) is not a defined pointwise value; state an explicit spectral/resolvent regularization before taking a trace.\par
\noindent\textit{Evidence.} Scan confirms the equality. Distribution theory does not permit this diagonal substitution. Item 1964-47.7.\par
\subsection*{P000076 (both; ambiguity): Gelfand commutator trace question}
\noindent\textit{Transcription reads.} \(\operatorname{Tr}(AB-BA)=0\) without assumptions.\par
\noindent\textit{Critical reading.} Both \(AB\) and \(BA\) trace class suffice; trace class of the commutator alone does not. \([S^*,S]\) has rank one and trace one.\par
\noindent\textit{Evidence.} The source labels this a problem rather than a fully quantified theorem; the distinction is necessary for a precise reconstruction. Item 1964-47.8.\par
\subsection*{P000078 (both; ambiguity): Continuous-parameter Weil kernel}
\noindent\textit{Transcription reads.} Kernel involving \(P_1P'_2-P_2P'_1\) over a single displayed defining-function difference.\par
\noindent\textit{Critical reading.} Retain as an unfinished formula; no denominator order, parameter cycle or orientation is supplied editorially.\par
\noindent\textit{Evidence.} Scan shows the incomplete formula and the subsequent linear/Radon specialization. The latter has a precise 1965 publication, but it does not identify the whole general kernel. Item 1964-48.5.\par
\subsection*{P000081 (both; faithful error): Exactness after refinements}
\noindent\textit{Transcription reads.} Long exact sequences exist for each cover but may lose exactness upon passage to the limit.\par
\noindent\textit{Critical reading.} Filtered colimits of abelian groups preserve exactness. The fixed-cover section complexes may already fail to form a short exact sequence.\par
\noindent\textit{Evidence.} Scan confirms the claim about passage to a limit. The distinction between an epimorphism of sheaves and lifting sections locally gives the actual obstruction. Item 1964-49.4.\par
\subsection*{P000082 (both; ambiguity): Subscript \(p\) in exactness discussion}
\noindent\textit{Transcription reads.} \(F'_p\to F_p\to F''_p\) is not necessarily onto.\par
\noindent\textit{Critical reading.} Do not interpret \(p\) automatically as a stalk. The line labels presheaves; stalks of an exact sequence of sheaves are exact.\par
\noindent\textit{Evidence.} Scan has “\ru{предпучки}” under the first object and distinguishes the inclusion functor below. The handwriting alone does not define \(p\). Item 1964-49.5.\par
\subsection*{P000082-P000083 (both; ambiguity): Truncated differential-form sequence}
\noindent\textit{Transcription reads.} \(0\to C\to\Omega^0\to\Omega\to0\).\par
\noindent\textit{Critical reading.} Identify the final sheaf and differential before claiming exactness; one classical repair uses closed one-forms, not all one-forms in higher dimension.\par
\noindent\textit{Evidence.} Russian and English leave the final \(\Omega\) unspecified; scan confirms the truncated sequence. Item 1964-49.7.\par
\subsection*{P000084 (both; faithful error): Subnormality condition}
\noindent\textit{Transcription reads.} \(T^*T-TT^*\ge0\) is necessary and sufficient.\par
\noindent\textit{Critical reading.} Necessary only; subnormality requires the full positivity conditions. Weights \(1,1,2,2,\ldots\) give a hyponormal non-subnormal unilateral shift.\par
\noindent\textit{Evidence.} Scan shows the rewritten necessary-and-sufficient phrase. The moment counterexample is checked in item 1964-50.6; uncertainty of the horizontal glyph does not save the biconditional.\par
\subsection*{P000084-P000085 (both; faithful error): Matrix coefficients in the polynomial proof}
\noindent\textit{Transcription reads.} P84 \(a_{mn}=(Ae_n,e_m)\), but P85 defines \(a_{mn}^{(p)}=(p(A)e_m,e_n)\).\par
\noindent\textit{Critical reading.} Use \(a_{mn}^{(p)}=(p(A)e_n,e_m)\) consistently with the lower-Hessenberg convention.\par
\noindent\textit{Evidence.} Scans of both pages and Halmos 1966 pp.433--434 agree on the convention required by the displayed subdiagonal product. Item 1964-50.8.\par
\subsection*{P000086 (both; faithful error): Positive norm limit}
\noindent\textit{Transcription reads.} \(\lim\lVert p(P_nAP_n)\rVert=d>0\) after convergence applied to \(e\).\par
\noindent\textit{Critical reading.} The inference supplies \[\lim\lVert p(P_nAP_n)e\rVert=\lVert p(A)e\rVert>0;\] it does not supply operator-norm convergence.\par
\noindent\textit{Evidence.} Scan visibly omits \(e\) in this one line while surrounding formulas contain it; the published Halmos proof uses the vector norm. Item 1964-50.10.\par
\subsection*{P000086-P000087 (both; faithful error): Adjacent-flag threshold and cancelled monotonicity}
\noindent\textit{Transcription reads.} Conflicting half-\(d\) threshold directions; assertion of monotonic decrease is cancelled.\par
\noindent\textit{Critical reading.} Keep the cancellation; select the first adjacent crossing using the endpoint values, as in the published proof, without assuming monotonicity of these norms.\par
\noindent\textit{Evidence.} Both scans reviewed. The terminal norm is zero and initial norm exceeds \(3d/4\), but monotonicity is neither needed nor established. Halmos 1966 pp.435--437. Item 1964-50.10.\par
\subsection*{P000088 (both; ambiguity): Compact-defect norm assumption}
\noindent\textit{Transcription reads.} \(\lVert T\rVert<1\) or \(\lVert T\rVert\le1\); \(I-T^*T\) compact.\par
\noindent\textit{Critical reading.} Keep the sign uncertain. Strict inequality is impossible with compact defect in infinite dimension; non-strict inequality is a separate theorem request.\par
\noindent\textit{Evidence.} The scan does not securely distinguish the signs. Coercivity \(I-T^*T\ge(1-\lVert T\rVert^2)I\) diagnoses the strict reading. Item 1964-50.12.\par
\subsection*{P000088 (both; ambiguity): Sz.-Nagy--Foias powers}
\noindent\textit{Transcription reads.} Both \(T^n\to0\) and \((T^*)^n\to0\) have crossed arrows.\par
\noindent\textit{Critical reading.} Retain failure of strong convergence to zero for both power sequences, with \(\lVert T\rVert\le1\) and \(\dim H>1\).\par
\noindent\textit{Evidence.} Scan confirms crossed arrows; exact primary theorem is Sz.-Nagy--Foias IX, Section 6, Theorem 3 and Remark (ii), p.314, unaffected by the 1965 correction. Item 1964-50.14.\par
\subsection*{P000088 (both; ambiguity): Donoghue superdiagonal weights}
\noindent\textit{Transcription reads.} \(1/2,1/4\), then \(1/8\) or \(1/6\).\par
\noindent\textit{Critical reading.} Retain the third-entry uncertainty; the geometric \(2^{-n}\) choice is a literature-supported reconstruction, not a certain decipherment.\par
\noindent\textit{Evidence.} Scan and diplomatic uncertainty compared with Donoghue Examples 1--2. Item 1964-50.15.\par
\subsection*{P000088 (both; ambiguity): Integer-chain lattice}
\noindent\textit{Transcription reads.} The structure of all integers.\par
\noindent\textit{Critical reading.} Adjoin the obligatory lattice endpoints \(0\) and \(H\); the completed integer-chain question has an affirmative Domar construction.\par
\noindent\textit{Evidence.} Every full invariant-subspace lattice contains least and greatest elements. Domar 1981 Theorem 2 supplies the exact completed order type. Item 1964-50.17.\par
\subsection*{P000088 (both; faithful error): Measure-Volterra conjecture}
\noindent\textit{Transcription reads.} \(Af(x)=\int_0^x f(y)\,d\mu(y)\) on the real line; lattice isomorphic to \(\operatorname{supp}(\mu)\).\par
\noindent\textit{Critical reading.} False for \(\mu=(1/2)\mathbf1_{[-1,1]}\,dx\) with oriented integration: negative-side and positive-side support spaces are incomparable invariant subspaces.\par
\noindent\textit{Evidence.} Scan gives the real line and lower limit zero. The atomless probability counterexample avoids endpoint-convention ambiguity. The one-sided classical repair is stated separately. Item 1964-50.19.\par
\subsection*{P000089 (both; faithful error): Every vector cyclic, including the Ornstein remark}
\noindent\textit{Transcription reads.} Every vector is cyclic.\par
\noindent\textit{Critical reading.} Every nonzero vector is the meaningful formulation; the zero vector cannot be cyclic in a nonzero space.\par
\noindent\textit{Evidence.} Scan and both transcriptions omit nonzero. The Bishop universal claim is also disproved for known arithmetic parameter classes. Items 1964-50.20--50.22.\par
\subsection*{P000089 (both; ambiguity): Hardy Fourier condition}
\noindent\textit{Transcription reads.} \(\int_0 u(t)\exp(int)\,dt=0\), \(n=1,2,\ldots\), with no upper limit.\par
\noindent\textit{Critical reading.} Use upper limit \(2\pi\) in the stated circle normalization. Keep the positive exponent: it correctly eliminates negative Fourier coefficients.\par
\noindent\textit{Evidence.} Scan shows lower zero only and the circle \(2\pi\) immediately above. Item 1964-50.24.\par
\subsection*{P000089 (both; provenance): Final intervention speaker}
\noindent\textit{Transcription reads.} An English-speaking professor with a doubtful F-name, Hungarian from Romania.\par
\noindent\textit{Critical reading.} Possible C. Foias identification remains qualified; exact theorem identification does not establish speaker identity.\par
\noindent\textit{Evidence.} Scan surname is overwritten; analytic-contraction content matches Sz.-Nagy--Foias 1964 but is not independent biographical evidence. Item 1964-50.24.\par
\subsection*{P000090 (both; provenance): Academic-year boundary at page heading}
\noindent\textit{Transcription reads.} Page included in the supplied 1964--65 translation envelope.\par
\noindent\textit{Critical reading.} Opening 20/IX, Kohn, belongs to 20 September 1965 and the next academic year; do not attribute it to record 50.\par
\noindent\textit{Evidence.} Scan and Russian heading explicitly show 20/IX and Kohn. Catalog page envelopes overlap; provenance is preserved in the record-50 summary.\par
